\chapter{XII 葬礼仪式}
% 第18章章内目录，按法文原书第433—436页条目顺序生成。
\par\bigskip\noindent{\Large\headfont 目录}\par\medskip
\begingroup\small\setlength{\parskip}{0pt}
\par\noindent\hspace*{0.00em}\hyperref[ch18:local:h1]{18.1.~葬礼悼词}\nobreak\dotfill\nobreak\hyperref[ch18:local:h1]{\pageref*{ch18:local:h1}}\par
\par\noindent\hspace*{1.35em}\hyperref[ch18:local:h2]{18.1.1.~(1) 恭维}\nobreak\dotfill\nobreak\hyperref[ch18:local:h2]{\pageref*{ch18:local:h2}}\par
\par\noindent\hspace*{4.05em}\hyperref[ch18:local:n1]{注~104}\nobreak\dotfill\nobreak\hyperref[ch18:local:n1]{\pageref*{ch18:local:n1}}\par
\par\noindent\hspace*{1.35em}\hyperref[ch18:local:h3]{18.1.2.~(2) 力量与光环}\nobreak\dotfill\nobreak\hyperref[ch18:local:h3]{\pageref*{ch18:local:h3}}\par
\par\noindent\hspace*{4.05em}\hyperref[ch18:local:n2]{注~105}\nobreak\dotfill\nobreak\hyperref[ch18:local:n2]{\pageref*{ch18:local:n2}}\par
\par\noindent\hspace*{0.00em}\hyperref[ch18:local:h4]{18.2.~阴阳之钥}\nobreak\dotfill\nobreak\hyperref[ch18:local:h4]{\pageref*{ch18:local:h4}}\par
\par\noindent\hspace*{1.35em}\hyperref[ch18:local:h5]{18.2.1.~(1) 肌肉与肚肠(阳葬阴(1))}\nobreak\dotfill\nobreak\hyperref[ch18:local:h5]{\pageref*{ch18:local:h5}}\par
\par\noindent\hspace*{4.05em}\hyperref[ch18:local:n3]{注~106}\nobreak\dotfill\nobreak\hyperref[ch18:local:n3]{\pageref*{ch18:local:n3}}\par
\par\noindent\hspace*{4.05em}\hyperref[ch18:local:n4]{注~106\textsuperscript{1}}\nobreak\dotfill\nobreak\hyperref[ch18:local:n4]{\pageref*{ch18:local:n4}}\par
\par\noindent\hspace*{1.35em}\hyperref[ch18:local:h6]{18.2.2.~(2) 一段生命的故事:一个三乐章的循环}\nobreak\dotfill\nobreak\hyperref[ch18:local:h6]{\pageref*{ch18:local:h6}}\par
\par\noindent\hspace*{2.70em}\hyperref[ch18:local:h7]{18.2.2.1.~(a) 天真(阴与阳的联姻)}\nobreak\dotfill\nobreak\hyperref[ch18:local:h7]{\pageref*{ch18:local:h7}}\par
\par\noindent\hspace*{4.05em}\hyperref[ch18:local:n5]{注~107}\nobreak\dotfill\nobreak\hyperref[ch18:local:n5]{\pageref*{ch18:local:n5}}\par
\par\noindent\hspace*{2.70em}\hyperref[ch18:local:h8]{18.2.2.2.~(b) 超级父亲(阳葬阴(2))}\nobreak\dotfill\nobreak\hyperref[ch18:local:h8]{\pageref*{ch18:local:h8}}\par
\par\noindent\hspace*{4.05em}\hyperref[ch18:local:n6]{注~108}\nobreak\dotfill\nobreak\hyperref[ch18:local:n6]{\pageref*{ch18:local:n6}}\par
\par\noindent\hspace*{4.05em}\hyperref[ch18:local:n7]{注~108\textsuperscript{1}}\nobreak\dotfill\nobreak\hyperref[ch18:local:n7]{\pageref*{ch18:local:n7}}\par
\par\noindent\hspace*{4.05em}\hyperref[ch18:local:n8]{注~108\textsuperscript{2}}\nobreak\dotfill\nobreak\hyperref[ch18:local:n8]{\pageref*{ch18:local:n8}}\par
\par\noindent\hspace*{2.70em}\hyperref[ch18:local:h9]{18.2.2.3.~(c) 重逢(阴的苏醒(1))}\nobreak\dotfill\nobreak\hyperref[ch18:local:h9]{\pageref*{ch18:local:h9}}\par
\par\noindent\hspace*{4.05em}\hyperref[ch18:local:n9]{注~109}\nobreak\dotfill\nobreak\hyperref[ch18:local:n9]{\pageref*{ch18:local:n9}}\par
\par\noindent\hspace*{2.70em}\hyperref[ch18:local:h10]{18.2.2.4.~(d) 接纳(阴的苏醒(2))}\nobreak\dotfill\nobreak\hyperref[ch18:local:h10]{\pageref*{ch18:local:h10}}\par
\par\noindent\hspace*{4.05em}\hyperref[ch18:local:n10]{注~110}\nobreak\dotfill\nobreak\hyperref[ch18:local:n10]{\pageref*{ch18:local:n10}}\par
\par\noindent\hspace*{1.35em}\hyperref[ch18:local:h11]{18.2.3.~(3) 伴侣}\nobreak\dotfill\nobreak\hyperref[ch18:local:h11]{\pageref*{ch18:local:h11}}\par
\par\noindent\hspace*{2.70em}\hyperref[ch18:local:h12]{18.2.3.1.~(a) 事物的动力学(阴阳的和谐)}\nobreak\dotfill\nobreak\hyperref[ch18:local:h12]{\pageref*{ch18:local:h12}}\par
\par\noindent\hspace*{4.05em}\hyperref[ch18:local:n11]{注~111}\nobreak\dotfill\nobreak\hyperref[ch18:local:n11]{\pageref*{ch18:local:n11}}\par
\par\noindent\hspace*{2.70em}\hyperref[ch18:local:h13]{18.2.3.2.~(b) 敌对的配偶(阳葬阴(3))}\nobreak\dotfill\nobreak\hyperref[ch18:local:h13]{\pageref*{ch18:local:h13}}\par
\par\noindent\hspace*{4.05em}\hyperref[ch18:local:n12]{注~111$'$}\nobreak\dotfill\nobreak\hyperref[ch18:local:n12]{\pageref*{ch18:local:n12}}\par
\par\noindent\hspace*{2.70em}\hyperref[ch18:local:h14]{18.2.3.3.~(c) 一半与全——或裂缝}\nobreak\dotfill\nobreak\hyperref[ch18:local:h14]{\pageref*{ch18:local:h14}}\par
\par\noindent\hspace*{4.05em}\hyperref[ch18:local:n13]{注~112}\nobreak\dotfill\nobreak\hyperref[ch18:local:n13]{\pageref*{ch18:local:n13}}\par
\par\noindent\hspace*{2.70em}\hyperref[ch18:local:h15]{18.2.3.4.~(d) 原型认识与制约}\nobreak\dotfill\nobreak\hyperref[ch18:local:h15]{\pageref*{ch18:local:h15}}\par
\par\noindent\hspace*{4.05em}\hyperref[ch18:local:n14]{注~112$'$}\nobreak\dotfill\nobreak\hyperref[ch18:local:n14]{\pageref*{ch18:local:n14}}\par
\par\noindent\hspace*{1.35em}\hyperref[ch18:local:h16]{18.2.4.~(4) 我们的母亲,死亡}\nobreak\dotfill\nobreak\hyperref[ch18:local:h16]{\pageref*{ch18:local:h16}}\par
\par\noindent\hspace*{2.70em}\hyperref[ch18:local:h17]{18.2.4.1.~(a) 行为}\nobreak\dotfill\nobreak\hyperref[ch18:local:h17]{\pageref*{ch18:local:h17}}\par
\par\noindent\hspace*{4.05em}\hyperref[ch18:local:n15]{注~113}\nobreak\dotfill\nobreak\hyperref[ch18:local:n15]{\pageref*{ch18:local:n15}}\par
\par\noindent\hspace*{2.70em}\hyperref[ch18:local:h18]{18.2.4.2.~(b) 挚爱者}\nobreak\dotfill\nobreak\hyperref[ch18:local:h18]{\pageref*{ch18:local:h18}}\par
\par\noindent\hspace*{4.05em}\hyperref[ch18:local:n16]{注~114}\nobreak\dotfill\nobreak\hyperref[ch18:local:n16]{\pageref*{ch18:local:n16}}\par
\par\noindent\hspace*{2.70em}\hyperref[ch18:local:h19]{18.2.4.3.~(c) 信使}\nobreak\dotfill\nobreak\hyperref[ch18:local:h19]{\pageref*{ch18:local:h19}}\par
\par\noindent\hspace*{4.05em}\hyperref[ch18:local:n17]{注~114$'$}\nobreak\dotfill\nobreak\hyperref[ch18:local:n17]{\pageref*{ch18:local:n17}}\par
\par\noindent\hspace*{2.70em}\hyperref[ch18:local:h20]{18.2.4.4.~(d) 安格拉——或告别与再会}\nobreak\dotfill\nobreak\hyperref[ch18:local:h20]{\pageref*{ch18:local:h20}}\par
\par\noindent\hspace*{4.05em}\hyperref[ch18:local:n18]{注~115}\nobreak\dotfill\nobreak\hyperref[ch18:local:n18]{\pageref*{ch18:local:n18}}\par
\par\noindent\hspace*{1.35em}\hyperref[ch18:local:h21]{18.2.5.~(5) 拒受与接纳}\nobreak\dotfill\nobreak\hyperref[ch18:local:h21]{\pageref*{ch18:local:h21}}\par
\par\noindent\hspace*{2.70em}\hyperref[ch18:local:h22]{18.2.5.1.~(a) 失乐园}\nobreak\dotfill\nobreak\hyperref[ch18:local:h22]{\pageref*{ch18:local:h22}}\par
\par\noindent\hspace*{4.05em}\hyperref[ch18:local:n19]{注~116}\nobreak\dotfill\nobreak\hyperref[ch18:local:n19]{\pageref*{ch18:local:n19}}\par
\par\noindent\hspace*{2.70em}\hyperref[ch18:local:h23]{18.2.5.2.~(b) 循环}\nobreak\dotfill\nobreak\hyperref[ch18:local:h23]{\pageref*{ch18:local:h23}}\par
\par\noindent\hspace*{4.05em}\hyperref[ch18:local:n20]{注~116$'$}\nobreak\dotfill\nobreak\hyperref[ch18:local:n20]{\pageref*{ch18:local:n20}}\par
\par\noindent\hspace*{2.70em}\hyperref[ch18:local:h24]{18.2.5.3.~(c) 伴侣——或“恶”之谜}\nobreak\dotfill\nobreak\hyperref[ch18:local:h24]{\pageref*{ch18:local:h24}}\par
\par\noindent\hspace*{4.05em}\hyperref[ch18:local:n21]{注~117}\nobreak\dotfill\nobreak\hyperref[ch18:local:n21]{\pageref*{ch18:local:n21}}\par
\par\noindent\hspace*{2.70em}\hyperref[ch18:local:h25]{18.2.5.4.~(d) 阳扮阴——或主人的角色}\nobreak\dotfill\nobreak\hyperref[ch18:local:h25]{\pageref*{ch18:local:h25}}\par
\par\noindent\hspace*{4.05em}\hyperref[ch18:local:n22]{注~118}\nobreak\dotfill\nobreak\hyperref[ch18:local:n22]{\pageref*{ch18:local:n22}}\par
\par\noindent\hspace*{1.35em}\hyperref[ch18:local:h26]{18.2.6.~数学的阴与阳}\nobreak\dotfill\nobreak\hyperref[ch18:local:h26]{\pageref*{ch18:local:h26}}\par
\par\noindent\hspace*{2.70em}\hyperref[ch18:local:h27]{18.2.6.1.~(a) 最“大男子气”的艺术}\nobreak\dotfill\nobreak\hyperref[ch18:local:h27]{\pageref*{ch18:local:h27}}\par
\par\noindent\hspace*{4.05em}\hyperref[ch18:local:n23]{注~119}\nobreak\dotfill\nobreak\hyperref[ch18:local:n23]{\pageref*{ch18:local:n23}}\par
\par\noindent\hspace*{2.70em}\hyperref[ch18:local:h28]{18.2.6.2.~(b) 陌生的美人}\nobreak\dotfill\nobreak\hyperref[ch18:local:h28]{\pageref*{ch18:local:h28}}\par
\par\noindent\hspace*{4.05em}\hyperref[ch18:local:n24]{注~120}\nobreak\dotfill\nobreak\hyperref[ch18:local:n24]{\pageref*{ch18:local:n24}}\par
\par\noindent\hspace*{2.70em}\hyperref[ch18:local:h29]{18.2.6.3.~(c) 欲与严格}\nobreak\dotfill\nobreak\hyperref[ch18:local:h29]{\pageref*{ch18:local:h29}}\par
\par\noindent\hspace*{4.05em}\hyperref[ch18:local:n25]{注~121}\nobreak\dotfill\nobreak\hyperref[ch18:local:n25]{\pageref*{ch18:local:n25}}\par
\par\noindent\hspace*{2.70em}\hyperref[ch18:local:h30]{18.2.6.4.~(d) 上涨的海……}\nobreak\dotfill\nobreak\hyperref[ch18:local:h30]{\pageref*{ch18:local:h30}}\par
\par\noindent\hspace*{4.05em}\hyperref[ch18:local:n26]{注~122}\nobreak\dotfill\nobreak\hyperref[ch18:local:n26]{\pageref*{ch18:local:n26}}\par
\par\noindent\hspace*{2.70em}\hyperref[ch18:local:h31]{18.2.6.5.~(e) 九个月与上涨的海}\nobreak\dotfill\nobreak\hyperref[ch18:local:h31]{\pageref*{ch18:local:h31}}\par
\par\noindent\hspace*{4.05em}\hyperref[ch18:local:n27]{注~123}\nobreak\dotfill\nobreak\hyperref[ch18:local:n27]{\pageref*{ch18:local:n27}}\par
\par\noindent\hspace*{2.70em}\hyperref[ch18:local:h32]{18.2.6.6.~(f) 阴的葬礼(阳葬阴(4))}\nobreak\dotfill\nobreak\hyperref[ch18:local:h32]{\pageref*{ch18:local:h32}}\par
\par\noindent\hspace*{4.05em}\hyperref[ch18:local:n28]{注~124}\nobreak\dotfill\nobreak\hyperref[ch18:local:n28]{\pageref*{ch18:local:n28}}\par
\par\noindent\hspace*{2.70em}\hyperref[ch18:local:h33]{18.2.6.7.~(g) 超级妈妈,还是超级爸爸}\nobreak\dotfill\nobreak\hyperref[ch18:local:h33]{\pageref*{ch18:local:h33}}\par
\par\noindent\hspace*{4.05em}\hyperref[ch18:local:n29]{注~125}\nobreak\dotfill\nobreak\hyperref[ch18:local:n29]{\pageref*{ch18:local:n29}}\par
\par\noindent\hspace*{1.35em}\hyperref[ch18:local:h34]{18.2.7.~阴阳的颠倒}\nobreak\dotfill\nobreak\hyperref[ch18:local:h34]{\pageref*{ch18:local:h34}}\par
\par\noindent\hspace*{2.70em}\hyperref[ch18:local:h35]{18.2.7.1.~(a) 颠倒(1)——或激烈泼辣的妻子}\nobreak\dotfill\nobreak\hyperref[ch18:local:h35]{\pageref*{ch18:local:h35}}\par
\par\noindent\hspace*{4.05em}\hyperref[ch18:local:n30]{注~126}\nobreak\dotfill\nobreak\hyperref[ch18:local:n30]{\pageref*{ch18:local:n30}}\par
\par\noindent\hspace*{2.70em}\hyperref[ch18:local:h36]{18.2.7.2.~回顾(1)——或一幅画的三联}\nobreak\dotfill\nobreak\hyperref[ch18:local:h36]{\pageref*{ch18:local:h36}}\par
\par\noindent\hspace*{4.05em}\hyperref[ch18:local:n31]{注~127}\nobreak\dotfill\nobreak\hyperref[ch18:local:n31]{\pageref*{ch18:local:n31}}\par
\par\noindent\hspace*{2.70em}\hyperref[ch18:local:h37]{18.2.7.3.~回顾(2)——或死结}\nobreak\dotfill\nobreak\hyperref[ch18:local:h37]{\pageref*{ch18:local:h37}}\par
\par\noindent\hspace*{4.05em}\hyperref[ch18:local:n32]{注~127$'$}\nobreak\dotfill\nobreak\hyperref[ch18:local:n32]{\pageref*{ch18:local:n32}}\par
\par\noindent\hspace*{2.70em}\hyperref[ch18:local:h38]{18.2.7.4.~(d) 父母——或冲突的心脏}\nobreak\dotfill\nobreak\hyperref[ch18:local:h38]{\pageref*{ch18:local:h38}}\par
\par\noindent\hspace*{4.05em}\hyperref[ch18:local:n33]{注~128}\nobreak\dotfill\nobreak\hyperref[ch18:local:n33]{\pageref*{ch18:local:n33}}\par
\par\noindent\hspace*{4.05em}\hyperref[ch18:local:n34]{注~128$^1$}\nobreak\dotfill\nobreak\hyperref[ch18:local:n34]{\pageref*{ch18:local:n34}}\par
\par\noindent\hspace*{2.70em}\hyperref[ch18:local:h39]{18.2.7.5.~(e) 敌对的父亲(3)——或阳葬阳}\nobreak\dotfill\nobreak\hyperref[ch18:local:h39]{\pageref*{ch18:local:h39}}\par
\par\noindent\hspace*{4.05em}\hyperref[ch18:local:n35]{注~129}\nobreak\dotfill\nobreak\hyperref[ch18:local:n35]{\pageref*{ch18:local:n35}}\par
\par\noindent\hspace*{2.70em}\hyperref[ch18:local:h40]{18.2.7.6.~(f) 箭与浪}\nobreak\dotfill\nobreak\hyperref[ch18:local:h40]{\pageref*{ch18:local:h40}}\par
\par\noindent\hspace*{4.05em}\hyperref[ch18:local:n36]{注~130}\nobreak\dotfill\nobreak\hyperref[ch18:local:n36]{\pageref*{ch18:local:n36}}\par
\par\noindent\hspace*{2.70em}\hyperref[ch18:local:h41]{18.2.7.7.~(g) 冲突之谜}\nobreak\dotfill\nobreak\hyperref[ch18:local:h41]{\pageref*{ch18:local:h41}}\par
\par\noindent\hspace*{4.05em}\hyperref[ch18:local:n37]{注~131}\nobreak\dotfill\nobreak\hyperref[ch18:local:n37]{\pageref*{ch18:local:n37}}\par
\par\noindent\hspace*{2.70em}\hyperref[ch18:local:h42]{18.2.7.8.~(h) 颠倒(2)——或暧昧的反叛}\nobreak\dotfill\nobreak\hyperref[ch18:local:h42]{\pageref*{ch18:local:h42}}\par
\par\noindent\hspace*{4.05em}\hyperref[ch18:local:n38]{注~132}\nobreak\dotfill\nobreak\hyperref[ch18:local:n38]{\pageref*{ch18:local:n38}}\par
\par\noindent\hspace*{1.35em}\hyperref[ch18:local:h43]{18.2.8.~主人们与仆人}\nobreak\dotfill\nobreak\hyperref[ch18:local:h43]{\pageref*{ch18:local:h43}}\par
\par\noindent\hspace*{2.70em}\hyperref[ch18:local:h44]{18.2.8.1.~(a) 天鹅绒的脚掌——或微笑}\nobreak\dotfill\nobreak\hyperref[ch18:local:h44]{\pageref*{ch18:local:h44}}\par
\par\noindent\hspace*{4.05em}\hyperref[ch18:local:n39]{注~133}\nobreak\dotfill\nobreak\hyperref[ch18:local:n39]{\pageref*{ch18:local:n39}}\par
\par\noindent\hspace*{2.70em}\hyperref[ch18:local:h45]{18.2.8.2.~(b) 兄弟与配偶——或双重签名}\nobreak\dotfill\nobreak\hyperref[ch18:local:h45]{\pageref*{ch18:local:h45}}\par
\par\noindent\hspace*{4.05em}\hyperref[ch18:local:n40]{注~134}\nobreak\dotfill\nobreak\hyperref[ch18:local:n40]{\pageref*{ch18:local:n40}}\par
\par\noindent\hspace*{4.05em}\hyperref[ch18:local:n41]{注~134$^1$}\nobreak\dotfill\nobreak\hyperref[ch18:local:n41]{\pageref*{ch18:local:n41}}\par
\par\noindent\hspace*{2.70em}\hyperref[ch18:local:h46]{18.2.8.3.~(c) 阴,作为仆人——与新主人}\nobreak\dotfill\nobreak\hyperref[ch18:local:h46]{\pageref*{ch18:local:h46}}\par
\par\noindent\hspace*{4.05em}\hyperref[ch18:local:n42]{注~135}\nobreak\dotfill\nobreak\hyperref[ch18:local:n42]{\pageref*{ch18:local:n42}}\par
\par\noindent\hspace*{2.70em}\hyperref[ch18:local:h47]{18.2.8.4.~(d) 阴,作为仆人(2)——与慷慨}\nobreak\dotfill\nobreak\hyperref[ch18:local:h47]{\pageref*{ch18:local:h47}}\par
\par\noindent\hspace*{4.05em}\hyperref[ch18:local:n43]{注~136}\nobreak\dotfill\nobreak\hyperref[ch18:local:n43]{\pageref*{ch18:local:n43}}\par
\par\noindent\hspace*{4.05em}\hyperref[ch18:local:n44]{注~136$^1$}\nobreak\dotfill\nobreak\hyperref[ch18:local:n44]{\pageref*{ch18:local:n44}}\par
\par\noindent\hspace*{1.35em}\hyperref[ch18:local:h48]{18.2.9.~天鹅绒中的利爪}\nobreak\dotfill\nobreak\hyperref[ch18:local:h48]{\pageref*{ch18:local:h48}}\par
\par\noindent\hspace*{2.70em}\hyperref[ch18:local:h49]{18.2.9.1.~(a) 天鹅绒的脚掌——或微笑}\nobreak\dotfill\nobreak\hyperref[ch18:local:h49]{\pageref*{ch18:local:h49}}\par
\par\noindent\hspace*{4.05em}\hyperref[ch18:local:n45]{注~137}\nobreak\dotfill\nobreak\hyperref[ch18:local:n45]{\pageref*{ch18:local:n45}}\par
\par\noindent\hspace*{2.70em}\hyperref[ch18:local:h50]{18.2.9.2.~(b) 颠倒(4)——或夫妻的马戏}\nobreak\dotfill\nobreak\hyperref[ch18:local:h50]{\pageref*{ch18:local:h50}}\par
\par\noindent\hspace*{4.05em}\hyperref[ch18:local:n46]{注~138}\nobreak\dotfill\nobreak\hyperref[ch18:local:n46]{\pageref*{ch18:local:n46}}\par
\par\noindent\hspace*{2.70em}\hyperref[ch18:local:h51]{18.2.9.3.~(c) 天真的暴力——或传递}\nobreak\dotfill\nobreak\hyperref[ch18:local:h51]{\pageref*{ch18:local:h51}}\par
\par\noindent\hspace*{4.05em}\hyperref[ch18:local:n47]{注~139}\nobreak\dotfill\nobreak\hyperref[ch18:local:n47]{\pageref*{ch18:local:n47}}\par
\par\noindent\hspace*{2.70em}\hyperref[ch18:local:h52]{18.2.9.4.~(d) 奴隶与傀儡——或闸门}\nobreak\dotfill\nobreak\hyperref[ch18:local:h52]{\pageref*{ch18:local:h52}}\par
\par\noindent\hspace*{4.05em}\hyperref[ch18:local:n48]{注~140}\nobreak\dotfill\nobreak\hyperref[ch18:local:n48]{\pageref*{ch18:local:n48}}\par
\par\noindent\hspace*{1.35em}\hyperref[ch18:local:h53]{18.2.10.~暴力——或游戏与刺棒}\nobreak\dotfill\nobreak\hyperref[ch18:local:h53]{\pageref*{ch18:local:h53}}\par
\par\noindent\hspace*{2.70em}\hyperref[ch18:local:h54]{18.2.10.1.~(a) 正义者的暴力}\nobreak\dotfill\nobreak\hyperref[ch18:local:h54]{\pageref*{ch18:local:h54}}\par
\par\noindent\hspace*{4.05em}\hyperref[ch18:local:n49]{注~141}\nobreak\dotfill\nobreak\hyperref[ch18:local:n49]{\pageref*{ch18:local:n49}}\par
\par\noindent\hspace*{2.70em}\hyperref[ch18:local:h55]{18.2.10.2.~(b) 机械与自由}\nobreak\dotfill\nobreak\hyperref[ch18:local:h55]{\pageref*{ch18:local:h55}}\par
\par\noindent\hspace*{4.05em}\hyperref[ch18:local:n50]{注~142}\nobreak\dotfill\nobreak\hyperref[ch18:local:n50]{\pageref*{ch18:local:n50}}\par
\par\noindent\hspace*{2.70em}\hyperref[ch18:local:h56]{18.2.10.3.~(c) 贪婪——或蚀本买卖}\nobreak\dotfill\nobreak\hyperref[ch18:local:h56]{\pageref*{ch18:local:h56}}\par
\par\noindent\hspace*{4.05em}\hyperref[ch18:local:n51]{注~143}\nobreak\dotfill\nobreak\hyperref[ch18:local:n51]{\pageref*{ch18:local:n51}}\par
\par\noindent\hspace*{2.70em}\hyperref[ch18:local:h57]{18.2.10.4.~(d) 两种认识——或对认识的恐惧}\nobreak\dotfill\nobreak\hyperref[ch18:local:h57]{\pageref*{ch18:local:h57}}\par
\par\noindent\hspace*{4.05em}\hyperref[ch18:local:n52]{注~144}\nobreak\dotfill\nobreak\hyperref[ch18:local:n52]{\pageref*{ch18:local:n52}}\par
\par\noindent\hspace*{2.70em}\hyperref[ch18:local:h58]{18.2.10.5.~(e) 秘密的神经}\nobreak\dotfill\nobreak\hyperref[ch18:local:h58]{\pageref*{ch18:local:h58}}\par
\par\noindent\hspace*{4.05em}\hyperref[ch18:local:n53]{注~145}\nobreak\dotfill\nobreak\hyperref[ch18:local:n53]{\pageref*{ch18:local:n53}}\par
\par\noindent\hspace*{2.70em}\hyperref[ch18:local:h59]{18.2.10.6.~(f) 激情与饥渴——或逐步加码}\nobreak\dotfill\nobreak\hyperref[ch18:local:h59]{\pageref*{ch18:local:h59}}\par
\par\noindent\hspace*{4.05em}\hyperref[ch18:local:n54]{注~146}\nobreak\dotfill\nobreak\hyperref[ch18:local:n54]{\pageref*{ch18:local:n54}}\par
\par\noindent\hspace*{2.70em}\hyperref[ch18:local:h60]{18.2.10.7.~(g) 心软的爸爸}\nobreak\dotfill\nobreak\hyperref[ch18:local:h60]{\pageref*{ch18:local:h60}}\par
\par\noindent\hspace*{4.05em}\hyperref[ch18:local:n55]{注~147}\nobreak\dotfill\nobreak\hyperref[ch18:local:n55]{\pageref*{ch18:local:n55}}\par
\par\noindent\hspace*{2.70em}\hyperref[ch18:local:h61]{18.2.10.8.~神经中的神经——或侏儒与巨人}\nobreak\dotfill\nobreak\hyperref[ch18:local:h61]{\pageref*{ch18:local:h61}}\par
\par\noindent\hspace*{4.05em}\hyperref[ch18:local:n56]{注~148}\nobreak\dotfill\nobreak\hyperref[ch18:local:n56]{\pageref*{ch18:local:n56}}\par
\par\noindent\hspace*{1.35em}\hyperref[ch18:local:h62]{18.2.11.~另一个自己}\nobreak\dotfill\nobreak\hyperref[ch18:local:h62]{\pageref*{ch18:local:h62}}\par
\par\noindent\hspace*{2.70em}\hyperref[ch18:local:h63]{18.2.11.1.~(a) 缓期中的宿怨——或事物的回转(2)}\nobreak\dotfill\nobreak\hyperref[ch18:local:h63]{\pageref*{ch18:local:h63}}\par
\par\noindent\hspace*{4.05em}\hyperref[ch18:local:n57]{注~149}\nobreak\dotfill\nobreak\hyperref[ch18:local:n57]{\pageref*{ch18:local:n57}}\par
\par\noindent\hspace*{2.70em}\hyperref[ch18:local:h64]{18.2.11.2.~(b) 天真与冲突——或绊脚石}\nobreak\dotfill\nobreak\hyperref[ch18:local:h64]{\pageref*{ch18:local:h64}}\par
\par\noindent\hspace*{4.05em}\hyperref[ch18:local:n58]{注~150}\nobreak\dotfill\nobreak\hyperref[ch18:local:n58]{\pageref*{ch18:local:n58}}\par
\par\noindent\hspace*{2.70em}\hyperref[ch18:local:h65]{18.2.11.3.~(c) 天赐的情势——或封神}\nobreak\dotfill\nobreak\hyperref[ch18:local:h65]{\pageref*{ch18:local:h65}}\par
\par\noindent\hspace*{4.05em}\hyperref[ch18:local:n59]{注~151}\nobreak\dotfill\nobreak\hyperref[ch18:local:n59]{\pageref*{ch18:local:n59}}\par
\par\noindent\hspace*{2.70em}\hyperref[ch18:local:h66]{18.2.11.4.~(d) 否认(1)——或提醒}\nobreak\dotfill\nobreak\hyperref[ch18:local:h66]{\pageref*{ch18:local:h66}}\par
\par\noindent\hspace*{4.05em}\hyperref[ch18:local:n60]{注~152}\nobreak\dotfill\nobreak\hyperref[ch18:local:n60]{\pageref*{ch18:local:n60}}\par
\par\noindent\hspace*{4.05em}\hyperref[ch18:local:n61]{注~153}\nobreak\dotfill\nobreak\hyperref[ch18:local:n61]{\pageref*{ch18:local:n61}}\par
\par\noindent\hspace*{2.70em}\hyperref[ch18:local:h67]{18.2.11.5.~(f) 布景——或“第二天性”}\nobreak\dotfill\nobreak\hyperref[ch18:local:h67]{\pageref*{ch18:local:h67}}\par
\par\noindent\hspace*{4.05em}\hyperref[ch18:local:n62]{注~154}\nobreak\dotfill\nobreak\hyperref[ch18:local:n62]{\pageref*{ch18:local:n62}}\par
\par\noindent\hspace*{2.70em}\hyperref[ch18:local:h68]{18.2.11.6.~(g) 另一个自己——或认同与冲突}\nobreak\dotfill\nobreak\hyperref[ch18:local:h68]{\pageref*{ch18:local:h68}}\par
\par\noindent\hspace*{4.05em}\hyperref[ch18:local:n63]{注~155}\nobreak\dotfill\nobreak\hyperref[ch18:local:n63]{\pageref*{ch18:local:n63}}\par
\par\noindent\hspace*{2.70em}\hyperref[ch18:local:h69]{18.2.11.7.~(h) 敌对的兄弟——或传递(2)}\nobreak\dotfill\nobreak\hyperref[ch18:local:h69]{\pageref*{ch18:local:h69}}\par
\par\noindent\hspace*{4.05em}\hyperref[ch18:local:n64]{注~156}\nobreak\dotfill\nobreak\hyperref[ch18:local:n64]{\pageref*{ch18:local:n64}}\par
\par\noindent\hspace*{4.05em}\hyperref[ch18:local:n65]{注~156$^1$}\nobreak\dotfill\nobreak\hyperref[ch18:local:n65]{\pageref*{ch18:local:n65}}\par
\par\noindent\hspace*{1.35em}\hyperref[ch18:local:h70]{18.2.12.~冲突与发现——或恶之谜}\nobreak\dotfill\nobreak\hyperref[ch18:local:h70]{\pageref*{ch18:local:h70}}\par
\par\noindent\hspace*{2.70em}\hyperref[ch18:local:h71]{18.2.12.1.~(a) 无恨亦无容赦}\nobreak\dotfill\nobreak\hyperref[ch18:local:h71]{\pageref*{ch18:local:h71}}\par
\par\noindent\hspace*{4.05em}\hyperref[ch18:local:n66]{注~157}\nobreak\dotfill\nobreak\hyperref[ch18:local:n66]{\pageref*{ch18:local:n66}}\par
\par\noindent\hspace*{2.70em}\hyperref[ch18:local:h72]{18.2.12.2.~(b) 领会与更新}\nobreak\dotfill\nobreak\hyperref[ch18:local:h72]{\pageref*{ch18:local:h72}}\par
\par\noindent\hspace*{4.05em}\hyperref[ch18:local:n67]{注~158}\nobreak\dotfill\nobreak\hyperref[ch18:local:n67]{\pageref*{ch18:local:n67}}\par
\par\noindent\hspace*{2.70em}\hyperref[ch18:local:h73]{18.2.12.3.~(c) 无缘之暴力的缘由}\nobreak\dotfill\nobreak\hyperref[ch18:local:h73]{\pageref*{ch18:local:h73}}\par
\par\noindent\hspace*{4.05em}\hyperref[ch18:local:n68]{注~159}\nobreak\dotfill\nobreak\hyperref[ch18:local:n68]{\pageref*{ch18:local:n68}}\par
\par\noindent\hspace*{2.70em}\hyperref[ch18:local:h74]{18.2.12.4.~(d) 藤井日达古鲁吉——或太阳与其行星}\nobreak\dotfill\nobreak\hyperref[ch18:local:h74]{\pageref*{ch18:local:h74}}\par
\par\noindent\hspace*{4.05em}\hyperref[ch18:local:n69]{注~160}\nobreak\dotfill\nobreak\hyperref[ch18:local:n69]{\pageref*{ch18:local:n69}}\par
\par\noindent\hspace*{2.70em}\hyperref[ch18:local:h75]{18.2.12.5.~(e) 祈祷与冲突}\nobreak\dotfill\nobreak\hyperref[ch18:local:h75]{\pageref*{ch18:local:h75}}\par
\par\noindent\hspace*{4.05em}\hyperref[ch18:local:n70]{注~161}\nobreak\dotfill\nobreak\hyperref[ch18:local:n70]{\pageref*{ch18:local:n70}}\par
\par\noindent\hspace*{2.70em}\hyperref[ch18:local:h76]{18.2.12.6.~(f) 信念与认识}\nobreak\dotfill\nobreak\hyperref[ch18:local:h76]{\pageref*{ch18:local:h76}}\par
\par\noindent\hspace*{4.05em}\hyperref[ch18:local:n71]{注~162}\nobreak\dotfill\nobreak\hyperref[ch18:local:n71]{\pageref*{ch18:local:n71}}\par
\par\noindent\hspace*{2.70em}\hyperref[ch18:local:h77]{18.2.12.7.~(g) 最烫的铁——或转折}\nobreak\dotfill\nobreak\hyperref[ch18:local:h77]{\pageref*{ch18:local:h77}}\par
\par\noindent\hspace*{4.05em}\hyperref[ch18:local:n72]{注~162$'$}\nobreak\dotfill\nobreak\hyperref[ch18:local:n72]{\pageref*{ch18:local:n72}}\par
\par\noindent\hspace*{2.70em}\hyperref[ch18:local:h78]{18.2.12.8.~(h) 无穷的链——或传递(3)}\nobreak\dotfill\nobreak\hyperref[ch18:local:h78]{\pageref*{ch18:local:h78}}\par
\par\noindent\hspace*{4.05em}\hyperref[ch18:local:n73]{注~162$''$}\nobreak\dotfill\nobreak\hyperref[ch18:local:n73]{\pageref*{ch18:local:n73}}\par
\endgroup\medskip



\section{葬礼悼词}
\label{ch18:local:h1}

\subsection{(1) 恭维}
\label{ch18:local:h2}

\notehead{104}{}
\phantomsection\label{ch18:local:n1}

\dmark{}(5月12日)\fnote{1}{(*)}{(5月18日)下面这条注,“出自一条(注n\textsuperscript{o}~47的)脚注——它膨胀到了犯规的规模”。我把它插在这里,心想这一次这个次序比编年的次序更自然。自这条注写成的时刻起,我就感到需要再把它发展些许——这将由紧随本注之后的一条注来完成;写下这些文字时,它还没有写成。这两条注的全套,从现在起就顶上了那个顺理成章的名字:“葬礼悼词”!}在这本\origpage{447}\fnote{2}{(**)}{(5月18日)说的是IHES(高等科学研究所)为庆祝建所二十五周年于1983年出版的小册子。注“有益的挣脱”(n\textsuperscript{o}~42)的脚注里已提到它;“拒受遗产——或矛盾的代价”一注(n\textsuperscript{o}~47)的开头再次提到;本注(葬礼悼词(1))即关涉后者(见前一条脚注)。}小册子里就我的著作所写的那篇小传里,“上同调”或“同调”一词竟不曾出现!“概形”一词也没有。文中当然谈到(情势所需——我当时毕竟充当着“IHES的第一枚菲尔兹奖章”)我著作的“巨人般面貌”、出版的卷数、提炼出的本质问题、“最大的自然一般性”(奇怪的法语!),极为讲究的术语,影射“Grothendieck群”(我敢打赌,又是一项“最大的自然一般性”!),甚至topos及其在逻辑中的用处(但尤其不在别处!)……却没有任何一处影射一项\textbf{结果},或一项我本可能发展、且或许派得上用场的\textbf{理论}——看来只好相信:那二十卷巨人般的著作,严格说是空的,或至多是一些(从未解决的)问题与观念的汇编,当然都带着最大的自然一般性:Grothendieck群,判给了(既然我的名字已经粘在后面),呈献为代数K理论的“始祖”(!)(而它当然与拓扑K理论毫无干系——对后者则讳莫如深\fnote{3}{(***)}{我关于Riemann--Roch定理的工作,是代数\textbf{K}理论的第一记强力启动,而全然不是什么“始祖”。拓扑\textbf{K}理论,恰恰诞生于我证明Riemann--Roch--Grothendieck定理的那一年(1957),继我在赫策布鲁赫研讨班上的报告之后。这位被讳莫如深的“后代”的“始祖”,当时还不满一岁!代数K理论(加上巴斯引入的函子$K_i$[原书此处仅印$K$],以补我引入的函子$K^0$)在其后的年月里发展起来,受“始祖”与其第一个“后代”的双重影响。况且,早在六十年代下半叶,我就有了一条通向描述高阶$K^i$(对一个原文称为“monomiale”的范畴,比如加性范畴)的进路,循着Sinh~夫人的论文的路线。它停留在启发的层面,因为它奠基于包裹Picard的$\infty$-\textbf{范畴}的直觉;而彼时(乃至此后)还没有人费心发展(非严格的)$\infty$-范畴的观念——即我如今以topos上的$\infty$-叠(点的topos)之名所称的观念。随着我即将在《追寻栈》里发展的、一套叠的上同调-同伦形式主义的奠基草图(正处在我1955至1965年间所发展观念的直线上),这条通往高阶K-不变量理论的“几何”进路,终将唾手可得。};至于Riemann--Roch定理,想必是“始祖”的后代们去打理的——那些做真正定理、做正经事情的人!

在一个时尚以轻蔑为荣、对一般性(被那句有点可笑的“最大的自然一般性”的措辞不动声色地讥讽着……)不屑一顾的时代,\dmark{}那支费心为我撰写葬礼悼词的匿名手笔,把如今正被交付轻蔑的东西\dmark{}\origpage{448},毫不吝惜地赏给了我\fnote{4}{(*)}{(5月18日)而且这还是我删削过的!我的葬礼悼词的全文引证,见“葬礼悼词(2)”一注。}。我也赏识了(也许我是第一个……)同一支匿名手笔在悼词的这一段里的全部幽默:

“他在IHES创建了一个代数几何学派,围绕他主持的研讨班聚集起来,\textbf{并以他传播自己观念时的那份慷慨为滋养}”(着重号是我加的)。遗憾的是,一如我的“巨人般著作”,这个被我滋养得如此之好的“代数几何学派”,严格说也是空的——一个名字也不曾提起;而且没有人来抱怨自己被遗忘了——至少没有向我抱怨。

然而我似乎记得,年轻的Deligne曾在1965年(他当时该有十九岁)至1969年间,忠实地出没于那场(想必空空如也的)研讨班,并在这一研讨班里、在我们的两两对谈里,既学得了概形的技术,也学得了上同调技术与平展上同调——也就是说,恰恰是那些在其著作(至少在我见过的那些部分里)的每一页上都被使用的工具。在同一本小册子里献给Deligne的那篇小传里,同样没有任何影射,能使读者疑心\dmark{}\origpage{449}他可能从我这里学到了什么。然而,引人注目的是,在这篇(这一回可毫不“葬礼”的)Deligne颂——“IHES的第三枚菲尔兹奖章”——里,我的名字被提到了三次!甚至有一处迂回的措辞,带着照例必须环绕我寒微身量的那份含混,影射我曾“在任意域上的几何中构造了上同调理论”——而且肯定还是“带着最大的自然一般性”:一股Grothendieck学问的浓重气味,劈面而来\fnote{5}{(*)}{(5月18日)在葬礼悼词里,谈到我曾对术语怀有的那份“高度留意”。在“最大的自然一般性”或“在任意域上的几何中的上同调理论”这类离奇的措辞的使用里,我分明察觉到把这份留意变成笑柄的意图。我对给予事物之名字的那份极度的用心,自然出于我对这些事物的敬重——事物的名字本应表达其本质,或至少其某个本质的侧面。从传到我耳中的回声看,这份敬重的态度——它如今俨然合乎时宜的轻蔑所针对的对象——已不止一次令我震惊;那轻蔑,有时借给一些重要观念取些阿布拉卡达布拉式的名字来表达。对此另可参见“反常性”一注(n\textsuperscript{o}~76)。}。把上下文完整引出来是值得的;这是此类体裁的一件小小的杰作:

\begin{quote}
“从那里出发[经典Hodge理论],并从Grothendieck所提示的$\ell$-adic类比出发[人们纳闷:Grothendieck哪里找来的时间,一边写着他那二十卷最大的自然一般性,一边学得起这样正经的东西],他[Deligne]提炼出混合Hodge结构的观念,并使任何复代数簇的上同调带有这一结构。于是在$\ell$-adic上同调里[?],对有限域上的簇,他证明了困难尽人皆知的Weil猜想。这一结果显得愈发令人惊讶[!!],因为Grothendieck在任意域上的几何中构造了上同调理论之后[人们实在纳闷他还到那里去找什么],曾把余下的猜想[???]归约为一系列猜想——它们今天一如既往地无从下手。”
\end{quote}

说白了:这些Grothendieck学问(其名字足以吓跑最顽强的博物学通才)非但对证明这个困难如此“尽人皆知”的惊人结果毫无贡献,而且至多恰到好处,又用\textbf{猜想}(他从来不做别的!)把我们缠住一层,而且(谁又能料到呢)无从下手,一如他突发奇想提出它们的那时,也一如今天。

\dmark{}\origpage{450}然而我记得自己曾着手对付过它们——这些无从下手的猜想;但那想必是因为我消息不灵。那是在我出发——抱歉,我是想说“亡故”——的时分;而我那些比我消息灵通的后裔们,则小心从不把鼻子伸进这类玩意,因为Deligne言之凿凿:那是无从下手的!

我分明认得这套做派:人们尽了本分,大量引用Grothendieck(无论他还是谁,都不能声称人们在这庄严的日子埋葬他);人们甚至来了一个“暂停”式的影射,提到某些“$\ell$-adic类比”,它们曾在混合Hodge理论的启动中起过作用。这想必是十三年前那句大名鼎鼎的碑文式半行之后的第二遭\fnote{6}{(*)}{这句“碑文式的半行”,见于Deligne1970年尼斯国际数学家大会的报告“Hodge理论I”。评语见注n\textsuperscript{o}~78\textsuperscript{2}。};两处影射,与某篇1968年文章里那些“重量方面的考虑”出奇地相似\fnote{7}{(**)}{对此可参见“罐装的权与十二年的秘密”一注(n\textsuperscript{o}~49)的开头,及“驱逐”一注(n\textsuperscript{o}~63)中更周详的检视。}:人们处于“暂停”的豁免状态,同时牵着读者的鼻子走!在这里,托庄严场合之福,那记“暂停”式的引证干得比消灭罪证更妙——这段文字想就那位大名鼎鼎的Grothendieck给出的印象,恰恰是这几年我感到的那股时尚之“风”所承载的;那股风,我\textbf{今天}就有机会感受到了\fnote{8}{(***)}{见同日的“屠杀”一注,n\textsuperscript{o}~87.},不再是在当众致辞、宾客盈场的葬礼悼词的音调里,而是在屠杀的音调里……

我继续引用——它值得:

\begin{quote}
“这条定理(前Weil猜想)有助于使$\ell$-adic上同调成为一件强大的工具,无须提名这件如此强大的工具的发明者——那位光彩夺目而谦虚的人……;它适用于一些表面上远离代数几何的问题,比如Ramanujam猜想。

最近,他研究了阿贝尔簇上的Hodge循环,朝着Grothendieck曾梦见的“Motif”理论迈出了第一步。他还证明了“相交上同调”的代数机制——Mac~Pherson与Goresky的拓扑理论。这使得把它移植到$\ell$-adic理论成为可能;在那里,它显出惊人的用处。”
\end{quote}

就这样,一支匿名的手笔(我猜是同一支),在那部“值得纪念的文集”问世一年之后\fnote{9}{(****)}{即1982年出版的Lecture~Notes~n\textsuperscript{o}~900那一卷;“一个梦的回忆——或Motif的诞生”与尤其“下葬——或新父亲”两注(n\textsuperscript{o}~51、52)谈到了它。正是在那一卷里,Motif被“开棺出土”(在对它们十二年的死一般的沉默之后),顶着一份(心照不宣的)替补父权。},终于弥补了该卷里的一处小“疏漏”。想必毕竟有人提了一个问题;Deligne在这里以他自己的方式尽了他弥补疏漏的手续(毕竟,当终于轮到正经数学的时候,还肯引这位做梦人\dmark{}Grothendieck一把,也算客气!)\origpage{451}而且照旧欺瞒读者:所谓“第一步”,早在1968年就迈出了——随着Deligne启动Hodge--Deligne理论;而那理论,扎的根恰是Motif瑜伽:在先前的四年里,他的确是经由与我的接触“吸收”了它。他那部著作所自出、他始终无法摆脱却偏要否认的这部瑜伽,实际上在第一段引文的迂回措辞里,被以“$\ell$-adic类比”之名打发了。一位既不极灵通也不极留心的读者,断然不会疑心这些“$\ell$-adic类比”(它们对Hodge--Deligne理论扮演过一个起点的角色——但至多如此……\fnote{10}{(*)}{这套Hodge--Deligne理论至今仍停留在童年状态:因为“Hodge--Deligne复形”的观念——架在$\mathbb{C}$上任意有限型概形之上——与这些“系数”的六种运算形式主义,始终不曾得到发展。对这样一套理论的需要,对Deligne一如对我,早在其混合Hodge结构的最初工作之前就已分明;它以分明的方式出自Motif瑜伽。但自我退出数学舞台起,Deligne身上便发展出一道对我引入同调代数的关键观念(导出范畴、六种运算——更不必说topos)的“阻塞”;这使得一套开局如此惊艳的理论的自然发育归于停顿。})与一门我确实梦到过的“Motif理论”(而且是一个精确得吓人的梦!)之间的联系——除非经由这一条:还是那位做梦人Grothendieck,(凭“最大的自然一般性”的本事!)硬要向真正的数学家们提示一些类比,而做正经工作,则是他们的事了。

至于那条大名鼎鼎的“‘相交上同调’的代数机制”,我们这就到了“反常层会议”的正中心\fnote{11}{(**)}{关于这场研讨会,见第VII行列“研讨会——或Mebkhout的层与反常性”。}(尽管“反常”一词并未出口)。对“IHES的四枚菲尔兹奖章”中的一位\dmark{},人们毕竟还戴上手套——考虑到场合的庄严;但对这同一位Grothendieck的遗孤学生,就无须顾忌了。在这个特殊的场合——聚光灯下、部长致辞、等等等等——我自己的下葬,不是经由沉默的下葬,而是经由\textbf{恭维}的下葬:调配得当,按方服用。\origpage{452}但明摆着:在 Mac~Pherson 与 Goresky 被点名之处,对遗孤学生Zoghman~Mebkhout,沉默才是规矩——一如两年前的反常层会议,一如今天依然如此。

\subsection{(2) 力量与光环}
\label{ch18:local:h3}

\notehead{105}{}
\phantomsection\label{ch18:local:n2}

(9月29日)那条“前面”的注——“葬礼悼词(1)——或恭维”(n\textsuperscript{o}~104)——写于5月12日,距今已四个多月。它当初是作为“拒受遗产,或矛盾的代价”一注(注n\textsuperscript{o}~47,3月末)的一条脚注起手的,为的是顺带记下我刚察觉的一桩小“怪事”。可是写着写着,我顺着行与页发觉:我正在评点的这两篇看似平淡的短文,竟是一座真正的“矿”\fnote{12}{(*)}{关于此事的回顾性评语,见9月24日“表面与深处”一注(n\textsuperscript{o}~101)的开头。}——我事先既没有料到,也没有刻意求之。那一天,也正是我刚画完一幅屠杀图景之日(注n\textsuperscript{o}~87);那幅图景,在此前的几个星期里,已从迷雾中渐渐显形。在这里,它骤然落实,单凭一段铺陈的描述之力取得肉身;此刻它正在用力唤我。屠杀,与“恭维”——对可敬的亡者的葬礼悼词——宛如同一幅触目图景的相辅相成的两联,在同一天显现!

这当然足以使我满盈!次日起,“我的手发痒”要乘胜追击,尤其要更深地探测那座我意外按上手的小小矿中珍宝。变得清楚的是:头一件该做的事,是把小册子里那两段文字全文引出——同时也是更好地接触这些文本、更好地浸透其真正讯息——那份“字里行间”的讯息——的最好办法\fnote{13}{(**)}{对此可参见“论破译一份讯息的艺术——或书写的颂歌”一注(n\textsuperscript{o}~102),它紧随前一条脚注所引的那条注之后。}。\origpage{453}甚至\dmark{}还没来得及把那两篇文本誊清,前一天的接触就足以在我身上唤起或唤醒好几个联想;我感到它们多汁。我急于追随它们,虽然还不大知道它们会把我引向何方……

最终,在其后的日子与星期里,我并没有沿着这条势头接下去——尽管整整这段时间里,我一直许诺自己近日就回到它。一桩意外的“健康事故”,使《收获与播种》的一切反思工作、乃至一切智力工作,中断了三个多月\fnote{14}{(*)}{对此可参见“事件——或身体与精神”与“陷阱——或省力与衰竭”两注,n\textsuperscript{o}~98、99.}。那几天刚刚敞开的、有利于沿这个方向继续反思的“热的时刻”,从此过去了。它是否会回来,甚至我是否肯下力气去“吹”(把热的吹回来!)以便不惜一切代价唤回它,都不确定。干脆说:我此刻真正的愿望,是走向那条终注,为整场名为葬礼的反思做一份临时\textbf{盘点}——并画上一道\textbf{最后的线}!至于本注,我至少要兑现我许下的(而且外加许给读者的)全文引证;或许还会就这两篇文本(或许还有把它们白纸黑字重抄这件事本身)在我身上激起的某些联想,给几句简要的指示。

所说的两篇文本(分别载于1983年题名《高等科学研究所》的庆典小册子第13与15页),属于一组“速写肖像”:自1958年建所以来在IHES任过职的“终身教授”与“长期访问者”,按到所的编年为序。那是相当简短的文本,各约半页,各带在所年月与身份(教授或长期访问者)、主要荣誉、主要兴趣领域与最重要的贡献,并(视情形)附某些合作者的名字。然而对我寒微的身量而言,关于这后三个“客观”侧面——兴趣领域、主要贡献、主要合作者或学生——的显著空白,却由那些酒神颂风格的“恭维”填满;其中一些,前一条注里已经检出并引证……\dmark{}\origpage{454}

我有幸开篇的这一系列,由以下数学家与物理学家构成:A.~Grothendieck、L.~Michel、R.~Thom、D.~Ruelle、P.~Deligne、N.H.~Kuiper、D.~Sullivan、P.~Cartier、H.~Epstein、J.~Fröhlich、A.~Connes、K.~Gawedzki、M.~Gromov、O.~Lanford。

我记得Dieudonné与我同时在IHES任教授;而在这份名单上我发现并非如此——他想必只满足于主持《数学出版物》。然而我现在发现,在小册子第3页,IHES的“履历”里,事实并非如此:Dieudonné分明同我一样,自1958年(直到1964年)是“终身教授”——至少名义上如此。一桩有点奇怪的矛盾!我把“履历”的头两“条”抄在这里,1958与1961:

\begin{quote}
\textbf{1958}高等科学研究所协会由Léon~Motchane在巴黎创建,辅以声名世界级的科学顾问团……与一批欧洲工业家。科学活动由两位数学家启动:Jean~Dieudonné(→1964)与Alexandre~Grothendieck(→1970),聘为终身教授。《IHES数学出版物》第1期出版。

\textbf{1961}承认为公益机构。

……
\end{quote}

顺带我注意到:在这份简短的履历里,似乎有必要提到(略带象征意味的)《数学出版物》第1期的出版(它由一篇24页的文章构成,作者 G.E.~Wall与这个新生机构并无特殊瓜葛),却不提代数几何研讨班——它们以家喻户晓的缩写SGA~1与SGA~2为人所知:我正是经由它们,在那个机构几乎还只“纸上存在”的年月里,开始独自支撑起它的科学声誉。况且,直到《数学出版物》大约第24卷为止,这些出版物的主体,是由《代数几何原理》(EGA)的相继各卷(1至4)构成的\fnote{15}{(*)}{我是其作者,与 J.~Dieudonné合作。};其余各卷,每卷围绕五十页上下打转(其科学水准之高,自不待言\dmark{}\origpage{455})。另一方面,在第19页(在“速写肖像”系列——Dieudonné缺席其中,天知道为什么\fnote{16}{(*)}{(9月30日)一个念头来到心头:缘由恐怕是这个:为了不必说,在所论的年月里(1958--1964),Dieudonné的时间,主要用于写《代数几何原理》(在那里,我不合时宜地俨然第一作者)与Bourbaki的成文工作——更不必说钢琴与烹饪(Dieudonné既是行家音乐家又是行家厨师);这些,在那本小册子里,当然无从谈起:它太挑剔了,连一个顺带的微笑都挤不进去……}——之后),人们读到一段排版极“广告招贴”的文字(配一幅诱人的照片:那套显赫出版物全帙堆成的巍峨书山):

\begin{quote}
数学出版物

正是Jean~Dieudonné,独自[!],自1959年起把《数学出版物》带上了世界卓越的峰巅。

自1979年起,它们以每年400页的定期刊物形式出版,由一个编辑部主持,主编是Jacques~Tits。

发行由……承担(等等)
\end{quote}

倘若《数学出版物》被以\textbf{这种}方式——就在那里,在这份显赫机构的庆典呈献里(它的主要天职,从来不是当一份期刊的出版者)——捧上云端,那么毫无疑问,这是为了叫某些人忘掉一桩不太愉快的事实\fnote{17}{(**)}{这可委屈了我的朋友尼科(他当时已是那所庆祝周年的机构的十二年所长);他想必(在这个场合一如在别的场合)对此什么也没看出来……}:倘若不是在某几个关键的年头里,某一位无名氏——他在自己的角落里执拗地推进自己的一些想法(它们有幸咬住了某些人,包括“上流社会”里的人)——顶风冒雨\fnote{18}{(***)}{顶风冒雨:在这四年里,我全然没有被那些持续的警告与流言吓住——说一场(如消息灵通的朋友们所暗示的……)全然不现实的冒险、甚或带点骗局味道的冒险,眼看就要破产!事实是:IHES当时没有任何财务或地产的根基;它的生命,始终悬在几位心境或好或坏的工业家的短期捐赠上。我对此并不太操心,只是信任创办人兼所长Léon~Motchane:他年复一年,凭财务戏法与“公共关系”的奇技,“保住了局面”。毕竟,在那个宽厚的年代,倘若它塌了,我很有希望很快找到一个不那么成问题的落脚点!相反,倘若我赢了我押在IHES上的这一注(有Dieudonné的鼓励——他认识Motchane,而我对他全然信任),我在IHES的职位,便比我所知道的任何别的位置都合我的心意。}\origpage{456}带给它一份担保与信誉——那是天底下最漂亮的协会章程、乃至最漂亮的“声名世界级的科学顾问”(原文如此),\dmark{}都无力给予的——那么,这家机构恐怕早已自生自灭,被人遗忘了很久。

(9月30日)这份庆典小册子(我会越来越熟识它的!)的“炫示”与“四面八方涂脂抹粉”的风格——抱歉,我是想说(极)高规格的“公共关系”——当然不是我朋友皮埃尔的风格,也不是尼科的:他们两位肯定有更要紧的事,犯不着撰写这类应景文字。相反,显而易见:占据着我的这两篇速写肖像——一篇写我,另一篇写Deligne——并非不经后者至少供给关键词就写成;单凭他是IHES唯一在行的人,就足够了。而同样对我分明的:这两篇文本,至少,未经这同一位Deligne先读过并点头,是不会付印的。于是,一开始对我就很清楚:这两篇文本,无论如何,首先映出我朋友的心境与意图——他力图就我的这个人、就他自己的这个人,向他自己、也向数学公众给出的形象。当然,正是以这个身份,这两段文字使我感兴趣。这份兴趣,并不取决于Deligne是否这些泄露心迹的文字的作者,或作者是别人(想必是把整本小册子“构思”成篇的那一位):后者出于这种或那种缘由,大概也认领了这个我的朋友执意要传递的“讯息”。

到底了,这两篇速写肖像如下,摘自题为“终身教授与长期访问教授的活动”的肖像画廊(第13--19页)。

\begin{quote}
Alexandre~Grothendieck,数学家,1958至1970年任IHES教授,菲尔兹奖章得主。

在他留在研究所的12年里,A.~Grothendieck革新了代数几何的基础与方法,并为它开启了新的应用,尤其是算术方面的。他在IHES创建了一个代数几何学派,围绕他主持的研讨班聚集起来,并以他传播自己观念时的那份慷慨为滋养。\dmark{}\origpage{457}其著作的巨人般面貌,映照在其出版物中:与Jean~Dieudonné合作的专著“代数几何原理”(8分册),以及与众多学生合作的“布瓦玛丽代数几何研讨班”12卷。

在代数几何中,他提炼出本质的问题,并给每一个概念以最大的自然一般性。所引入的观念,显出远超代数几何的本质性。它们往往显得如此自然,以致我们难以想象它们曾费过多大力气。倘若它们今天显得理所当然,这想必得益于他对术语的高度留意。

还应当记得:“Grothendieck群”——在代数几何中与相交理论相连,并在拓扑中被使用——是代数K理论的始祖。为把先前在$\mathbb{C}$上经拓扑途径证明的结果移植过来而在一般基底域上的代数几何中引入的topos,如今用于逻辑。

他在1970年离开IHES,那时他对数学的热情正在消退。难道应当认为,他沿着自己划定的道路提出的问题,已经变得过于困难了吗?

\begin{center}$\cdots$\end{center}

Pierre~Deligne,数学家,自1970年起任IHES教授,菲尔兹奖章、Henri~Poincaré金质奖章得主,法国科学院外籍院士。

他的研究以“理解代数簇的上同调”为主轴。若复代数簇$X$是非奇异射影簇,调和积分理论便在$H^*(X)$上给出一个Hodge结构。他由此出发,并借助Grothendieck所提示的$\ell$进类比,提炼出混合Hodge结构的概念,并使每个复代数簇的上同调都具有这种结构。在$\ell$进上同调中,因而对于有限域上的簇,他证明了难度之高众所周知的Weil猜想。这一结果显得更加令人惊讶,因为Grothendieck在构造了任意域上的上同调理论之后,曾把尚余的猜想归结为一系列猜想,而这些猜想今天仍与当时一样难以攻克。

这一定理促使$\ell$进上同调成为一种强有力的工具,可以用于看似远离代数几何的问题,例如Ramanujam猜想。

\origpage{458}最近,他还研究了阿贝尔簇上的Hodge循环,朝Grothendieck梦寐以求的“动机理论”迈出了第一步。他也证明了“交叉上同调”——Mac Pherson与Goresky的拓扑理论——的代数机制。由此得以将它移植到$\ell$进理论中,而它在那里显得出人意料地有用。

他目前对非交换调和分析感兴趣(实或$p$进Lie群——或有限经典群——以及某些齐性空间上的函数理论);这延续了他关于自守形式(Ramanujam猜想),以及他与G.~Lusztig合作开展的关于有限群表示的研究。

他吸收、洞察整个数学领域的速度极快;因此,面对人们向他提出的每个问题,他都能作出富于启发且有建设性的回应。
\end{quote}

这两篇文本,还当补上第三篇:在其中,Deligne与我一口气并列。我在小册子里散页的一张上找到它;标题与肖像画廊所在的那章相同——“IHES的研究取向”,副标题:“关于‘科学活动之前景’的概要”。它实质上是那幅肖像画廊的一个苛刻的“缩编”——这一回只收“终身教授”(在任的与往任的)\fnote{19}{(*)}{(10月1日)为凑“分量”,Connes也被塞了进去(尽管他不过是“访问者”);这总归给收藏家又添一枚“菲尔兹奖章”。相反,我的朋友Nico~Kuiper却被弃置不顾。为这场合隐身让位而出难题的,想必不会是他……},每人两三行。他们依次是(按被提及的顺序):我自己、Deligne、Michel、Thom、Ruelle、Sullivan、Connes、Lanford三世、Gromov。这与那幅更详尽的肖像画廊的次序相同;只除了一点:这一回,Deligne“升了位”,以便有幸与我一口气并列。有趣的细节:在这份文本里,被巡视的显赫人物的名字,全都带着着重号出现;唯独我寒微的身量例外\fnote{20}{(**)}{(10月1日)这一手光辉的排印花招(其意图或许并非有意)取得的排版效果是:将要引出的这一段,俨然是献给Pierre~Deligne的——他的名字在字形上俨然“终身教授”世系之首,把我排除在外;而我在其中俨然一位“合作者”、机构的局外人!编年的次序倒是遵守了,无可挑剔——然而产生的效果(而且肯定是有意求得的),是一场角色的\textbf{颠倒},在我身上唤起一些熟悉的联想(“颠倒”“驱逐”“暂停”等注里谈过,n\textsuperscript{o}~68$'$、63、77)。于是我同时也认出那套据有的\textbf{做派}——“暂停!”的做派——它向我明白无误地指出这份讯息的\textbf{真正}作者。}!以下就是关于我的朋友与我的那一段:

\origpage{459}\begin{quote}
Alexandre~Grothendieck深邃得近乎传奇的理论与Pierre~Deligne的光辉发现(两位都是菲尔兹奖章得主),经由“跨学科”的手段(上同调),把拓扑、代数几何与数论连在了一起。就在最近,这使得联邦德国的 G.~Faltings(他已在IHES工作过)得以证明一条艰深的定理:它在数论中立标,并照亮了那条大名鼎鼎的“Fermat定理”。
\end{quote}

顺带我注意到:“菲尔兹奖章”在这幅微型画廊里配得上一个大写的M;而“跨学科性”,自IHES建所之初,就是其创办人的一大偏爱的主题。或许正是托这份情势的福,在这份摘要里,人们似乎终究暗示:我的这个人,兴许与某一种称为“上同调”的“跨学科手段”有点关系(而它碰巧也是Deligne著作的“主导轴心”——不知是哪门子的偶然)。

\dmark{}可瞧,我又开始从细处抠这份文本了!那句应景的Faltings引证——他一夜之间以其轰动性的结果(这里称为“艰深”,仿佛问题就在这里;不过对我的话题无关紧要……)蹿升到科学时事的头排——也属于文本的“细处”:说到底,是那位抄写者的“落款”,不值得我停留。分明是对Deligne与我的第一句话,才载有该段的“讯息”之本质。

它就我朋友兼前学生身上的某些心境,向我透了不少——首先是一份深层的“Unsicherheit”(不安全感,缺乏底气,缺乏深层的立足)\fnote{21}{(*)}{这个来到我心的德文词“Unsicherheit”,在法语里没有对应词,(我相信)英语里也没有。其直译“不安全”几乎不能用来指一种心灵特质。那个否定的词“缺乏底气”,是另一种将就的近似。应当理解:这里的“底气”是在一个深层的层面上;其缺乏,可以在某些时刻被知觉,而表面上占上风的,却是一种有底气、完美从容的印象;它们如同一副保护性的甲壳,带着往往可观的惯性与“坚固性”,无懈可击……}。这里,一如在他署名的已刊文本里的任何一处\fnote{22}{(**)}{至少在我迄今过目的那些里。}、一如在前两篇速写肖像里,没有任何东西\origpage{460}能使人设想:我的朋友可能在某个时刻从我这里学到过什么。可瞧,他在这里,以清清楚楚的措辞,把自己呈现为一份宏阔的统一视野的\textbf{另一位父亲}——那视野是“取”自他人的\fnote{23}{(*)}{这里还有一层特别的讽刺:这份视野——在这里被当作他自己的“光环”取自他人——实际上,自“师父亡故”以来,恰恰被那位俨然继承人、一边划清界限一边否认遗产的人,交付给轻蔑,并被系统地阻挠。对此可参见“继承人”“共同继承人……”“……与电锯”三注(n\textsuperscript{o}~90、91、92);其他例证,见第X行列(灵车),由四口“棺材”1至4与掘墓人构成(注n\textsuperscript{o}~93至97)。},仿佛被一种亲身的信念所慑服:相信自己深深无力构想并让那些\textbf{他自己的}视野——同样宏大乃至更宏大——在他身上发芽绽放;仿佛为了“是”并“显得”伟大,他从此只剩下那份寒微的资源:把那道光环\textbf{记在自己的账上}——那道光环,他年轻时曾乐于环绕在一位显赫、如今已亡故(或无论如何,被一份天赐共识宣告如此……)的长辈身边。去攫取一道\textbf{光环},而不是让自己身上那些尚且无形无名的、等待出生并被命名的东西发芽绽放——而不是活出栖于他身上的、同样在等待的\textbf{他自己的力量}……

(10月1日)昨夜,我觉得自己又一次触到了那桩冲突的心脏——那桩冲突,早在《收获与播种》的最开头(在“他人之无谬与对己之轻蔑”一节,n\textsuperscript{o}~4),即将八个月前,我已以一般的措辞谈到它;而临近葬礼的开端(在“死结”一注,4月26日的n\textsuperscript{o}~65),我又“在一个极端而格外触目的个案里”重新找到了它。这又是一场不期而遇:在一处引文的转角——我求个心安,终究把它与前两处一道收了进来!几天前,重翻那本大名鼎鼎的小册子时,我就已注意到了这一段;当时它确实打动了我,但我没有停留。可昨天,一经白纸黑字写下,它立刻显得比那两段我刚抄的、本该构成我正在写的这条注之主题的应景文字,意味更重、更加触目。然而,在那两段文字里,并不缺使人“翻然醒悟”的地方;它们本会在四个月前还使我毫不迟疑地\dmark{}紧接着再写十页——倘若不是二十页——的联想。但忽然我觉得:我本可以这样展开的东西,说到底——至多除一处例外——\origpage{461}是已知的东西的再度确认,只是角度或许稍有不同;尤其:它们终究是些附带的侧面——那一类我已在5月“恭维”那条注里、乃至在整个葬礼反思中充分铺陈过的侧面。而第三段引文,却把我带回某种本质的东西——在那场漫长的“侦查”(葬礼上的工作是其中一部分)中,我一直有点失去视线的东西。

当时我一度想就此打住,至少不试图以言词把握那短短四行的碑文式句子对我说的话——它在某个层面上,分明已被“听见”。我终究越过了这份犹豫。词句上升得又慢又迟疑,而那起初弥散的印象,却随着书写渐渐澄清。一经白纸黑字写下、删去看似多余的部分,我便知道:我已把我“听见”的东西,把握到我所能把握的最好程度。

夜已深得犯规,真的该停在那里了。我心满意足地睡下,却还不确定:是否会把刚写下的东西,收进我这份注定要发表的见证。毕竟,我也满可以让读者——倘若他有兴趣穿透一份讯息的表面——自己去打捞他在其中听到的东西!只是到今天,我才知道我会把这段文字收进去:它分明表达着我(或自以为)对某件事的某种知觉或领会;那件事,在我看来是重要的,甚至是这场葬礼的一个深层内驱力的关键。

\section{阴阳之钥}
\label{ch18:local:h4}

\subsection{(1) 肌肉与肚肠(阳葬阴(1))}
\label{ch18:local:h5}

\notehead{106}{}
\phantomsection\label{ch18:local:n3}

(10月2日)我想再至少追一笔由三联式葬礼悼词(昨天我终于给出了全文引证)激起的联想。这一联想,早在5月13日——我刚写完“葬礼悼词(1)——或恭维”一注(n\textsuperscript{o}~104)的次日——就已强加于我。它触及事物的一个常常在不知不觉中错失的侧面;我只是近五六年来,才开始\dmark{}\origpage{462}真正意识到它。

在所检视的文本的字里行间,人们看到某些\textbf{价值}的崇拜在宣示自己。比如,关于Deligne所证明的Weil猜想,被突出渲染的,是它们的“\textbf{困难}”——而非其\textbf{美}、其单纯、其自Weil陈述它们的那一刻起便已敞开的宏阔视野\fnote{24}{(*)}{(10月3日)而且还被称作“尽人皆知的”困难!这几乎毫无意义,除非意在镇唬那些不在局中的人!一条猜想的“困难”,只有在其得到证明之后才可能真正掂量;而它的\textbf{多产},却可以预先感到,并且往往在证明之前,就已经由它所启发的种种工作,客观地显形。那些“大”猜想之有别于其他猜想,不在于其“困难”(那是未知的——姑且假定这个词还有点意义……),而分明在于其\textbf{多产}。顺带指出:这是事物的一个典型的“阴”的、女性的侧面,而“困难”是一个典型的“阳”的、“男性”的价值。}。我也想到这些隐约瞥见的视野在其被证明之前就已结出的果实,想到一旦越过通往其证明的漫长旅程的最后一步,便应时成熟的其他隐约瞥见的果实。正是这些猜想的美、它们内在的惊人自洽、以及它们使人隐约瞥见的此前无人料到的联系,使它们成为两代几何学家与算术学家如此强而多产的灵感之源。我著作的最深部分(无论“彻底做到底”的那部分,还是“Motif之梦”)都直接受启于此(经由Serre的中介——他善于领会并传递其猜想中所表达的那份视野的全部力量)。没有它们,无论$\ell$-进上同调,乃至topos的语言,恐怕都无从问世。或者不如说:我在自己生命中大约十五年间执意发展的那幅关于(代数)几何、拓扑与算术的“宏阔统一视野”,正是在这些“Weil猜想”里,找到了它的一幅最初的、触目惊心的草图。而随着这视野在幅度与成熟上步步增长,恰是这视野本身、连同它使人得以一一把握的那些此前隐藏的事物,一步一步向我低语:该做什么,该从哪个头“下手”眼前伸手可及的东西。Weil猜想证明中的最后一步,不多不少,正是一场漫长而迷人的旅程中的一步——这场旅程始于何时,我说不上;肯定远在我的诞生之前,而且在我的死亡之后,也远不会终结!

\dmark{}\origpage{463}但照所引文本里可以辨出的那股精神,人们会以为“Weil猜想”是一桩举重的事:瞧,要“生拉硬拽”举起的重量!两百公斤,可不是闹着玩的;困难尽人皆知;很多人试过,还没有一个人成功——直到“H日”(“H”如“赫拉克勒斯”)!结果是惊人的(注106\textsuperscript{1}):您掂掂,整整两公担——本来谁也不信有人举得起来……在论Faltings所证“艰深定理”的那句简短评语里,人们知觉到的,也是同一股精神:在那里,同样,在这件我们对事物之认识的新台阶的命名本身之中,被突出渲染的又是\textbf{困难},以博取芸芸众生的叹服——而非那自一座新攀上的峰巅敞开的视野\fnote{25}{(*)}{而最打动我的,从我亲手拿到Faltings证明三条关键猜想——包括此处所说的Mordell猜想——的预印本那一刻起,恰恰是其进路的惊人的\textbf{单纯}:他凭四十来页,证明了这些本被当作“高不可攀”的结果!(可与注n\textsuperscript{o}~3对读。)}。连“Mordell猜想”这个名字(诚然,数学之外的公众不认识它)似乎都不值得提——仿佛对猜想的领会与陈述(在这里,是Mordell的)是件附带事,因为“容易”。取而代之的,是对“Fermat定理”(据说它被“照亮”了)的一记凭空的展望。诚然,后者举世皆知(甚至在数学界之外),作为一块足足三百公斤的重量(它顶住了三个世纪的攻势)。

我头一个想说的要点是:这些文本里(带着合乎场合的审慎)被高扬的价值,是那些可以称作\textbf{肌肉}的价值——在这里,即“脑力肌肉”的价值:那种使人得以凭手腕之力,超越那些尽人皆知的“困难”纪录的价值。

这些价值,不只是这里被捧上云端的这位英雄的价值,也不只是某本庆典小册子的作者的价值(作者匿名;我自信认得出)。它们也是(似乎)越来越支配数学界、进而更普遍地支配科学界的价值。甚至越出这个还相对有限的世界的边界\dmark{}\origpage{464},可以说:它们也是、而且越来越是一种被称作“西方的”特定“文化”的价值\fnote{26}{(*)}{当我在这里谈到今日所见的我们“文化”的“价值”时,我指的当然是那些“官方的”价值——由学校、媒体、家庭所传递,并在各职业圈子里获得普遍共识的价值。这并不意味着这些价值被人人毫无保留地接受,也不意味着它们构成所有人的态度与行为中的基调。而且,正是怀着痛心,正派的人们、媒体与称职的专业文献(出于教育家、社会学家、精神科医生等的笔下),谈到某“一部分青年”:他们分明不“合调”,有损一幅体面的图景!}。久已以来,这种“文化”及其价值已征服了我们星球的表面,消灭了所有其他文化——其优越性无可辩驳的明证。这些价值的行星象征、其英勇的化身,就是密封甲胄里的宇航员:在千百万屏息凝神、瘫坐在屏幕前的电视观众的注视下,第一个踏上某颗无法想象的遥远而荒凉的行星。

这些价值——由于无从更贴切地把握,我姑且以一个带象征意味的粗略字眼“肌肉”称之——并非始于昨日。用人类学的行话说,人们也可以称之为“父权制的”。依我看,最早以不容置辩之力(!)宣示其至尊地位的成文文本之一,是《旧约》(尤其摩西诸书)。然而,只需读一读这份远古时代的迷人文献,便会发觉:“父权”价值的至尊——男人之于女人,或“精神”之于“身体”或“物质”——远未走到否定或轻蔑那些补充价值的地步(那些价值,当时或许还不曾被知觉为“对立的”或“敌对的”\fnote{27}{(**)}{比如,对母亲的崇敬,是犹太文化里一种扎根极深的传统;它想必对“圣典”里所高扬的那些(不妨这样说的)“官方”价值起着补偿的作用。这一传统,以一种改造而更炽烈的形态,重现于天主教传统之中,即对(童贞女!)马利亚的崇敬。})。我不知道这两组互补价值之坎坷浮沉的历史是否已被写出——追索这段历史,穿越一个个世纪与千年,从摩西的时代直到今天,想必是一件迷人的事。它想必也是一部历史:一种“价值”平衡——一边是“父权的”或“男性的”,另一边是“母权的”或“女性的”;一边是“肌肉”与“肚肠”,“精神”与“物质”——的逐渐败坏;这败坏分明朝着“阳刚”(或在东方传统辩证法里的“阳”)价值的方向,以牺牲“阴柔”(或“阴”)价值为代价。

\dmark{}\origpage{465}我觉得,我们这个时代的特点,正是这种文化败坏的变本加厉。这段历史的最近的几幕,是两大对抗的超级大国(浸透着本质上相同的价值)之间的“太空竞赛”与(尤其是核)军备竞赛。作为这类“力量”或“权力”的加码竞赛中一场狂飙演进的终极一幕与可能的结局,人们此刻就可以预告:某场行星规模的核浩劫(或别的什么——可挑选的余地有的是……)。它也许有这份功劳:一举而永远地解决所有问题……

不过,我在这里的主旨,并不是要描画一幅诱人的“世界末日”图景(这事不必等我),更不是要向“肌肉”、或向“大脑”(别名“精神”)宣战。我很清楚:连我的“肚肠”也从中捞不到什么!我珍爱我的肌肉与我的大脑——它们对我有用,自不待言;一如我也珍爱我的“肚肠”,它同样少不了。毋宁说,我觉得有用的,是在这里以几句话(倘能做到)说明:这场由周围文化所传递的、在这两类价值之间的深层冲突,是如何在我自己的身上上演的。用更接地气的话说,这也是我对自己这个人身上两个同样真实、同样具体、按本性不可分而互补、本身绝非敌对的\textbf{侧面}或\textbf{面孔}的态度(接受乃至推崇,或拒斥)的沿革。我可以称它们为我身上的“\textbf{男人}”与“\textbf{女人}”,或者(取两个不那么“有负担”、因而不太容易误导的称呼):“\textbf{阳}”与“\textbf{阴}”。

在大多数人身上,“骰子似乎早在幼年初期就已掷定”:那些将在一生中默默主宰态度与行为、带着一台调校完美的自动机之效率的本质机制,就在那时各就其位。这些机制的核心,是对我们身上这样或那样特质、或这样或那样深层冲动(带着阳的或阴的“署名”)、或这样那样一组给定署名的特质与冲动的“包裹”、乃至对“阳”包裹或“阴”包裹之全体的肯定或拒斥的机制。正是这些机制,在很大程度上,决定着我们“自我”之结构的其余一切选择(肯定或拒斥)机制。

\dmark{}\origpage{466}由于一些至今对我仍是谜的缘由,在我自己的个案里,“自我”(“老板”)与我这个人身上的“男性”与“女性”之间的关系史(无论有意识的还是无意识的;无论在“老板”本身之中,还是在“工人”之中——这两者都系于万物的阴阳双面)——这段历史,比寻常的要动荡得多。我在其中辨出三个时期。最后一个时期,在某种意义上,又接上了第一个时期:它铺展在我童年最初的五年。这第三个时期——我可以称之为\textbf{成熟}的时期——可以被看作向那个童年的某种“回归”,或看作与“\textbf{童年状态}”的逐渐重逢,与我存在之中“阴”与“阳”毫无风波的联姻之和谐的逐渐重逢。这些重逢,始于1976年7月,时年四十八岁——也正是在那一年,我(三个月后)发现了此前一直不为我所知的一种力量:禅观的力量\fnote{28}{(*)}{见“欲与禅观”与“惊叹”两节,n\textsuperscript{o}~36、37.}。

我的父母——母亲与父亲——各自身上占支配地位的价值,都是阳的价值:意志、智力(即:智力上的力量)、自制、对他人的支配、不妥协、“Konsequenz”(德语,意为对——尤其意识形态的——取舍的极端一贯)、“理想主义”——无论在政治层面还是实践层面……在我母亲身上,这份价值化自少女时代起就取得了一种加剧的力量;那是她对“她自己身上的女人”(并由此,对一般的女性)所怀的一份真正的憎恨的背面。她身上的这份憎恨,最终取得了一种愈发放肆、愈发摧毁性的力量——恰恰因为它终生全然被遮蔽。(我自己直到五年前——禅观出现在我生命里三年之后——才发现这些事。)在这样一个双亲的语境里,我竟能在我童年最初的五年里全然绽放——直到那场从父母身边的剥离、直到我的原生家庭(由我父母、比我年长四岁的姐姐与我构成)的毁灭——这是一个谜(然而是一个对我毫不含糊的事实);那毁灭,出于我母亲的意志,并借1933年政治事件的东风(倘可这样说)。

\notehead{106\textsuperscript{1}}{}
\phantomsection\label{ch18:local:n4}

\dmark{}(10月3日)Weil猜想可能不成立——对此,我与Deligne从不曾有丝毫怀疑;我也不记得听过谁表达过这样的怀疑。把那个“结果”(即这些猜想的证明)称作“令人惊讶”,再度暴露了那种蓄意哗众的用心。况且,自“拓扑”与平展上同调引入以来的任何时刻,我从不曾感到这些猜想高不可攀;毋宁说(自1963年起),它们不出近几年就必被证明。1970年我离去之时,我几乎毫不怀疑:Deligne——他是所有人中最有资格的——不会迟迟不去证明它们(他也确实没有迟疑),同时还有“代数循环的标准猜想”——那些更强的猜想(他却转而去败坏它们的名声)。

\origpage{467}Deligne对这后一些猜想的效验有所保留,他做得有理;我同他一样不确信。但一条猜想的分量,不取决于它到头来是真还是假,也不取决于其所谓的“困难”品格——那品格会使它“高不可攀”;那品格,本是纯主观的。它唯独取决于:该猜想把手指按上去的那个\textbf{问题}(在它被提出之前,无人知觉过这个问题)——这个问题,是否触及对我们的认识而言真正本质的东西。而这一点是明摆着的(至少对我而言!):只要这些猜想之效验的问题未获解决,人们就不可能有对代数循环的良好领会,也不可能有对代数簇上同调的那些所谓“算术”性质(或再或者说,“Motif几何”)的良好领会。今天,一如1968年的孟买大会之时,我把这个问题——连同奇点消解问题——视为代数几何里摆在我们面前的最根本的两个问题之一。我分明感到两者的分量!这份潜在的多产,必定会显现出来,只等人们不再满足于跌跌撞撞地绕开一条被判定“过于困难”的猜想,只等有人终于肯挽起袖子、跟它真刀真枪地干一场!

\subsection{(2) 一段生命的故事:一个三乐章的循环}
\label{ch18:local:h6}

\subsubsection*{(a) 天真(阴与阳的联姻)}
\tocsubsub{(a) 天真(阴与阳的联姻)}
\label{ch18:local:h7}

\notehead{107}{}
\phantomsection\label{ch18:local:n5}

\dmark{}(10月4日)我已有机会提到我一生最初五年的一个重要侧面,把它称作一项头奖的“特权”\fnote{29}{(*)}{见“屠杀”一注,n\textsuperscript{o}~87.}:一种与我的父亲之间深的、毫无问题的认同;它从不曾被恐惧或嫉妒触及。这一情势——连同这种对我父亲的认同的存在本身、及其无声的力量——我是仅仅四年前(在那场1979年8月至1980年3月关于我父母的禅观之后、对童年与我此生的禅观之中)才意识到的。这份认同,宛如另一种认同的宁静而有力的心脏:对我父母、我姐姐(她比我年长四岁)与我组成的那个家庭的认同。我对父亲一如对母亲,怀着无涯的敬佩与爱。他们的身量,对我而言,是万物的尺度。

这绝不意味着我对他们的态度是一种自动的赞许、一种景仰的憨态。我当然不知道他们是我万物的尺度;但我十分清楚,他们同我一样会犯错误;我心中没有任何恐惧,会妨碍我察觉一个分歧并把它明白表示出来。在我周围的种种冲突里,我不怕以我自己的方式表态。这丝毫触及不着那份信、那份安然的确信——它们构成我存在深处不可动摇的立足;毋宁说,这恰恰自发地出自那份信、那份确信本身。有过这样的时候:我父亲在无可奈何的愤怒发作之中——而我姐姐(不动声色地)以撩拨他为乐——动手粗暴地打她;而每一次,我都义愤填膺,怀着对我姐姐毫无保留的团结。我想,这些是我与“我的父亲”的关系里仅有的几片过路的浓云(与母亲之间则没有)。这并不是我赞成我姐姐那些有时挺可恶的把戏;我想它们也并不真正困扰我——对我而言,万物的尺度不是\textbf{她}。她的把戏(其缘由,想必同对我父亲一样,向我滑脱;而他回回“中计”,或我母亲则无论事先事后都小心不介入)——这些把戏,在某种意义上,对我并不真正产生后果。她是我的姐姐,她就是她那样,如此而已。可是\textbf{我的父亲}竟放任自己那种盲目的粗暴……

最亲近的三个存在——它们一起构成我最初年月的母体——被冲突撕裂着:冲突使每一位既\dmark{}\origpage{469}反对自己、又反对另外两位:我母亲与我姐姐之间,是一场阴险的、不动声色的冲突;我父亲与另一边的我母亲及我姐姐之间,则是一场激烈爆发的冲突;她们各用自己的方式(而且在我父母生前,谁也不曾装作察觉……)牵着他的鼻子走。神秘而非凡的是:在人生这些最敏感、最关键的年月里,被如此包围在冲突之中,这场冲突竟对我保持\textbf{外在};它在那几年的光阴里,不曾真正“咬进”我的存在,不曾在我身上安营扎寨。

我存在之中的那份分裂——它给我的生命打下的印记,一如给任何别人的生命——并不是在那几年里进驻我的;而是在其后的两三年里,从我六岁到八岁上下。在某个时刻(我自信可以定位到误差几个月;它当落在我的第八个年头),在同父母(他们不大肯给我一点活着的音讯)与姐姐分离了两年多之后,发生了一次\textbf{翻转}。那首先是一场与我\textbf{童年}的决裂——从那一刻起,它被一套有效的遗忘机制“埋葬”(那套机制,差不了多少,一直延续到今天)。在某个深的层面(还不是最深的层面……),我那时把我的父母判成了“外人”,一如我的童年从此被判为“外来”。在某种意义上,我\textbf{退位}了:为了被接纳进从此环绕我的那个世界,我决定变得像“他们”、像在那里定规矩的成年人——去取得并发展在那里逼人敬重的武器,去在一个唯有某种“力量”才被接受、被看重的世界里,以对等的武器作战……

况且,那也正是我父母偏爱的那种力量——环绕我最初年月的,正是他们。而在这里,我回到了那件“神秘的事”(我刚循着另一条由它激起的联想的线索,离开了它):我生命最初年月里\textbf{我身上分裂的缺席}。

也许,谜已不在于这份缺席,而在于这一点:我的父母——父亲与母亲——那时竟各自把我\textbf{整个地}接受、彻底地接受:接受我身上“男性的”、“男人的”部分,也接受“女性的”部分。或者换个说法:我的父母,各自被冲突撕裂着,各自否认着自己存在的一个本质部分——各自无力对自己、对对方、对我姐姐,怀有一份慈爱的敞开……然而他们竟对我——他们的儿子\dmark{}\origpage{470}——找到了这样一份敞开,一份毫无保留的接纳。

再换个说法:在我一生最初的五年里,任何时刻,我都不曾体验过\textbf{羞于我之所是}的感觉——无论在我的身体及其官能上,还是在我的冲动、我的癖好、我的行动上。任何时刻,我都不曾须否认我身上的什么,以换取周围环境的接纳、与之和平共处。

当然,我也做过不“过关”的事:一如所有的孩子,当我来劲时,我肯定也惹人烦、乃至让人受不了;而有时明摆着,我得拨正船头。我不立规矩,也无意立——我无须补偿什么隐秘的残缺。在我父母对我的爱里,本不会有谄媚的位置、对任性耍脾气的纵容——不会有无条件的首肯。但是,倘若我免不了挨我父亲或我母亲的申饬(一如反过来有时也发生),在那几年里,他们谁也不曾使我因某个不中其意的行动或举止而蒙羞。

在一道对父亲毫无歧义的深深认同的底色上,作为孩子的我,今天在我眼前,显出同时打着男性气概与女性气质的印记,两者同样强。

我觉得,在每一个存在、每一件事物之中——在其阴性与阳性品质那些不可分解而流变的联姻之中,那使他成为他的联姻;其精微的平衡,正是活在这存在或事物之中的美、那和谐——在这种阴阳的亲密结合里,往往(或许永远)有一个基调、一个“主音”:或阴,或阳。这个基调,在一个人身上未必容易辨出:因为那些或多或少的有效而彻底的压抑机制,会歪曲牌局,以一副借来的形象替换原初的和谐。比如,我四十年来的“品牌形象”,是一副几乎纯然男性气概的形象——而且从来不曾被质疑、甚至不曾被(无论我自己还是他人,依我看)辨认为如此;直到我四十八岁。然而我倾向于相信:与生俱来的那个基调,终生都在场——至少在某些或许永远无缘公开亮相的深层。在我自己的个案里,奇怪的是,直到今天我都说不好这主音是什么——也就是\dmark{}\origpage{471}那个浸透了我最初的童年、在我出生时就已是“我的”的主音。不止一次,各种征兆使我疑心:这个基调是“阴”的;当我的存在得以自发显现之时——在那些它摆脱了自我童年起积攒在我身上的种种制约的时刻——主宰它的,是“女性”的品质。换个说法:也许,在我身体与精神之中充当创造之力的那个东西——我有时称之为我身上的“孩子”或“工人”(以别于代表自我结构的“老板”,即我身上受制约的部分、我个人累积的制约之总和或结果)——这份力量,或许更“女性”而非“男性”(尽管按其本性与必然,它两者都是)。

这里不是逐一巡视这些“征兆”的地方。况且,要紧的并不是:我身上那个深层的主音是“女性”还是“男性”。要紧的毋宁是:在每一刻都知道\textbf{做我自己}——毫无迟疑地接纳我身上使我成为“女人”的那些特质与冲动,一如接纳使我成为“男人”的那些——并容它们自由地表达自己。

在孩子时代的最初几年,常有不熟的人把我当作女孩——而这从来不在我的心头造成丝毫的不自在、丝毫的不安全感。我想,主要是我的声音造成这种效果:一副非常清亮、尖细的嗓子——更不必说我有一头长发(多半蓬乱着),也许仅仅因为我母亲(她另有操不完的心)不常费心替我剪一剪。此外我壮得像头小牛,略带粗暴或冒险的游戏并不使我生厌;这丝毫不妨碍一份对寂静乃至孤独的偏好,以及一份对抱娃娃游戏的偏好\fnote{30}{(*)}{这一偏好在小男孩身上显得少见,我相信主要是因为它被周围环境系统地劝阻。}。我不记得有人为此嘲笑过我;但这类事想必难免此处彼处地发生过。倘若这类小插曲过去而不留创伤或羞辱的痕迹,那肯定是因为:它们在我身上找不到任何回音或放大器——找不到任何不安全感可资凭依;而“我之所是”被那些唯一真正于我有分量的人所接纳,是超出一切疑问的。嘲笑本伤不到我;它只会反噬那位发难者——在我眼里,那人只会显得十分愚蠢:竟装模作样地对天底下最自然的事情挑三拣四。

\dmark{}\origpage{472}况且我深知:这类有点古怪的蠢事绝非稀罕;单是裸体,就可以成为惊世骇俗的缘由!然而,就我记忆所及,我早已有机会见过我的母亲、父亲与姐姐的裸体;也有的是机会满足我那份正当的好奇:看看他们每一位连同我自己,究竟是怎生造出来的。分明,男人的或女人的身形构造里,没有任何惊世骇俗之处;它们在我看来,就那样,好得很——尤其(我对此毫不讳言)女人的身形。

\subsubsection*{(b) 超级父亲(阳葬阴(2))}
\tocsubsub{(b) 超级父亲(阳葬阴(2))}
\label{ch18:local:h8}

\notehead{108}{}
\phantomsection\label{ch18:local:n6}

(10月5日)1933年——我生命的第六个年头——安放了我一生中第一个关键的转折;它同时也是我母亲、我父亲的生命中,在他们彼此的关系、他们与孩子的关系中,一个关键的转折。那就是我们四人组成的家庭遭受暴力的、决定性毁灭的情节;四十六年后,头一个也是唯一一个对它作出盘点、并追索其来龙去脉的人,是我——经由我父母的通信、以及一两份苍白、费解而顽强的记忆残片,耐心地探掘并破译——那是在我父亲去世、我母亲去世很久之后\fnote{31}{(*)}{我父亲1942年死于奥斯维辛;我母亲卒于1957年。我在这里谈的那项工作,是在1979年8月至1980年10月之间进行的。}

在这里铺陈我在这项漫长的工作中就这一情节的分量与意味所学到的、所懂得的,并非我的题旨。三天前,我已影射过这一转折\fnote{32}{(**)}{见“肌肉与肚肠(阳葬阴(1))”一注(n\textsuperscript{o}~106)的末尾。},把它作为我身上阴阳联姻的历史中三大时期的第一个的粗暴终结。1933年12月,我被人急匆匆地打发到一个陌生的家庭——无论我还是把我从柏林送来的母亲,都不曾见过他们。实际上,她把我送去的这些素不相识的人,不过是最先肯收我做“寄宿生”的人——收费微薄之至,而且没有任何形式的担保说这笔钱哪一天真能付;而母亲则准备尽快去巴黎,与正在那里望眼欲穿等她的父亲会合。\dmark{}\origpage{473}父母之间是说定了的:一切都将会很好——对在布兰肯塞(汉堡附近)的我如此,对我姐姐亦然:几个月前,她终究也被打发到柏林的一家残疾儿童机构(人家倒肯收她,尽管她并不比我或我们的父母更残疾)。

作为奇异六个月的结局——那六个月,满载着隐隐的威胁与焦虑——我一夜之间发现自己置身于一个全然不同的世界:与我的生活里唯一认识过的那个世界——由我父母、我姐姐与我构成的那个——全然不同。我置身其中,像一群寄宿生中的一员:我们与那家人分桌吃饭,在那家的孩子们眼里,俨然二等的孩子;他们自成一个世界,居高临下地看我们。在我待在那里的五年里(直到1939年——大战前夜,迫于时势,我才终于去投奔父母),母亲偶尔寄来一封匆忙而拘谨的信;父亲则从未有过一行亲笔。

收留我的那对夫妇很快喜欢上了我。丈夫是位辞去圣职的前牧师,靠一份微薄的年金与拉丁文、希腊文、数学的私人授课度日;妻子生气勃勃,有时透着狡黠——两位都不是寻常人物,在许多方面都令人依恋。丈夫是一位文化视野开阔的人文学者,在政治上有点迷失,与纳粹政权有过节;后者到头来放过了他。战后我重新与他们联络,并保持着持续的往来,直到两位先后谢世\fnote{33}{(*)}{她两年前去世,享年九十九岁;我还有机会看到她的遗容——在葬礼的前一天,与她单独相对。}。

从他、尤其从她那里,一如从我的父母那里,我领受了最好的,也领受了最坏的。今天,隔着长长的回顾,我为那份“最好”感谢他们(一如我感谢我的父母),也为那份“最坏”。正是这份最先由我父母、继由他们二位给予的最好与最坏,构成了我在童年分得的那个浩繁“包裹”的绝大部分(一如人人都分得自己的一份),那是一个属于我的包裹……——它等着我去开箱、去检视。它们属于我过去的实体、我过去的财富;滋养我的当下,全在我自己。

我的新环境,方方面面都“得体”之至、因循之至;尤其对一切涉及身体、涉及性的东西,都带着惯常的压抑态度。\dmark{}\origpage{474}然而,我想,过了好几年,我才把那些态度内化、据为己有——比如以裸身为耻,连同一种对身体的暧昧关系。这份自幼灌输的羞耻,是一道深层分裂的一个侧面:身体遭一种心照不宣的轻蔑,而所谓的“文化”价值(与记诵之类的智力能力混为一谈)却被捧上云端。我身上的这道分裂,一直无人知晓,直到我四十八岁,它才开始消解。那是我一生的第二个大转折:它标志着我与自己之关系的历史中“第三个时期”的降临——也就是说,我与我的身体、与我身上的“男人”与“女人”之关系的历史。但在此之前,我已有充分的机会把这道分裂传给我的孩子\fnote{34}{(*)}{至少,传给了五个孩子中由我参与抚养的四个。第五个也是最后一个,由他母亲抚养;迄今,连相识的机会都还不曾有过一个。},又眼看他们把它传下去……

昨天我已影射过\fnote{35}{(**)}{见前一条注“力量的绽放——或联姻”(注n\textsuperscript{o}~107)的开头。}我身上终于发生的那次“翻转”。在脱离最初的家庭环境(或不如说,在那个环境被\textbf{毁灭})之后两年多——这次翻转,标志着惯常的压抑机制的到位;我的童年,曾有那份罕见的运气,一直免于此。迄今我侦察到两大具有压抑本性的力量,它们主宰了我的成年生活、我童年的大部分(注108\textsuperscript{1})。我自信可以说:它们的出现并非渐进;在我的个案里,这些机制多少是一夜之间、以其全部力量出现的——作为无意识层面上一个蓄意\textbf{选择}的后果。我先前把这一选择称作“退位”;但与此同时,它也是一个强劲的行动原则:“我要像‘他们’那样”(而不是“像我那样”),也意味着:我要押“脑袋”的宝——我的脑袋反正不比任何人的差——用他们自己的武器作战,并打赢“他们”!

这些机制之一——也是我在这里最关心的——是最寻常不过的机制之一:那就是对我身上“女性”特质的\textbf{压制}(或被时下共识感知为女性特质的那些),以让位于“男性”价值。这枚硬币的正面,当然是倾其所有投注于那些被感知为“男性”的特质与才能,并把它们发展到失度的地步:它们占去了一片大得离谱的天地。

\dmark{}\origpage{475}倘若这里有什么超出常轨,那当然不是这一双重机制的单纯\textbf{在场},也不是(依我看)严格意义的那份“压抑”成分的力量——即对“阴”的特质、态度、冲动的压制之力。在这件事上,与我母亲身上的情形没有共同的尺度:她的一生(以及她亲近者的生活)被她对那使她成为女人的东西的憎恨(终生不为人知)摧残着。我相信,我的为人处事,从来都不曾全然没有某种温和,乃至温情;它们执拗地磨圆我自童年起为自己裁定的那个人物的棱角,常为我赢得同情与好感。出格之处,毋宁在于我投资的失度、在于我倾注于任务的能量的\textbf{失度}——目不斜视,不分心!在正经工作之外,我的精神不断地被抛向完成、抛向工作的这一站或那一站的抵达。这份态度(德语作“Zielgerichtetheit”,英语作“aimdirectedness”),是卓越地一种阳的态度:一种\textbf{绷紧}的态度,对一切不显得与任务直接相关者\textbf{封闭}的态度。这份失度,足以在他人身上唤起某种“超人”或“超级男性”的形象\origpage{476}——诚然可叹地令人叹服!(照时行的价值);却立刻(多半在无意识的层面上)激起本能的防御反应、乃至对抗反应:面对这样一份力量的铺展——它被感受为有威胁乃至有侵略性,或无论如何危险(注108\textsuperscript{2})。尤其,这形象无可抗拒地唤起“\textbf{超级父亲}”的形象,随即发动那环绕着永恒的父亲的冲突而纠结起来的、暧昧的吸引与排斥的万千反应……这正是在那些\textbf{暧昧}关系之中\textbf{我的}贡献;那些关系,在我的生活里如此常见,而我在《收获与播种》里又有那么多次与它们照面。这份暧昧,因我身上阴特质的持存而加深,而非减弱:那些特质喂养着一份同情——倘若只有阳特质增生成一个巨型“超人”,它是无从唤起这份同情的。

于是我再一次确认:在这些永恒的“暧昧关系”里,我仍旧不过是在收割我自己播下的种子——哪怕每一次收成都显得意外(而且不合时宜……)!因为推动我身上的“老板”在作品的积累中不断自我超越的动机(或至少\textbf{一项}动机),难道不恰恰是迫使并不断刷新我的同侪(首先)与非同侪(此外)的敬意?听我最优秀的几位哀叹:他们跟不上我向前奔跑的那个节奏?!是的,我身上确有这份隐秘的欲望:在他人身上(一如在我自己身上)唤起这幅“超乎自然”的、失度的形象——一如它所映照的那个人本身;而它执拗地经由他人回返:或以明白响亮的言词,经由那份预期的赞颂(并被当作应得之物照单全收)——\textbf{而且},经由阴沉的敌意与冲突的幽暗深层之路……\fnote{36}{(*)}{(10月6日)干脆说:我方才把手指按上去的那份“隐秘的欲望”,直到今天仍未烧尽——尽管它终于(不过几年之前……)被察觉,尽管它今天已不如当年那样吞噬一切。}

\notehead{108\textsuperscript{1}}{}
\phantomsection\label{ch18:local:n7}

\dmark{}(10月6日)我是说:在我的生活中起过作用的那些具有压抑本性的力量,似乎主要地(倘若不是排他地)采取这两种特定形态之一:埋葬过去,以及抬高我身上“男性”的特质、贬抑“女性”的特质。我绝无意说,这两股“力量”——它们各自都具有压抑的本性(即:指向一种“压抑”、一份对某一实在的藏掖)——就是“主宰了我的生活”的仅有力量!那会忘掉我存在中全部非自我本位的侧面:那股求知的冲动——它在身体的层面与精神的层面同样表达着自己。(对此尤其可见“我的诸热情”第35节。)

即便在结构着自我的那些力量之中——即“老板”的流出物之中——至少有一份,其本性并非压抑,而且远早于那些压抑的力量、在我的生活中的作用还更要本质:那就是对我父亲的认同;它宛如我自身力量之感“宁静而有力的心脏”。这份认同,丝毫不带有抬高某些价值或品质(比如男性的)、贬抑另一些(“女性的”)的意味。不问父亲所信奉的价值如何,他这个人(直到1933年——那年他身上发生了一次翻转\fnote{37}{(*)}{引人注目的是:我父亲身上这一“翻转”(他时年43岁)的方向,竟是朝着一个\textbf{超级-阴}的状态,朝着一种帕夏式的被动;与我母亲——扮演着超级阳的角色——密切共谋。她转而照管他,撇下自己的孩子们。(他们被弃置不顾、自生自灭,至少直到1939年——那一年,迫于时势,违背本意,她终于把我接到身边……)我父亲的这种依附关系、我父母之间阴阳角色的这场颠倒,一直持续到我父亲1942年消失。})始终带着一份强的阴阳平衡:直觉与自发,在其中占有的份额,不亚于智力与意志。

最后,作为另一股自我本性的重要“力量”——它与那些压抑机制血肉相连(或者不如说,它本身就具有“压抑”的本性)——还应当算上那份永恒的\textbf{虚荣};它的作用,在我的生活里,一如在任何人那里,同样沉重。但这份“力量”的本性是如此普遍——一如它在每个人的生活里(以或粗或精的形态)所扮演的支配角色——以致在一份关于自我结构之力与机制在个人身上所采取的特定形态、并赋予自我其特殊面貌与立足的清单里,几乎不值得专列一条。\origpage{477}

\notehead{108\textsuperscript{2}}{}
\phantomsection\label{ch18:local:n8}

\dmark{}(10月6日)在这份“力量的铺展”里,没有任何现行意义的“侵略”意图——无论有意识还是无意识;只有一份无意识的欲望:要令人印象深刻,要博得敬意。诚然,“博得敬意”这个自发来到我心的措辞,本身已带着一份\textbf{强制}的意味,近于“侵略”。这份无意识的强制意图——它同样在无意识的层面上被知觉——想必常常被经验为一种侵略(即便这份经验始终被遮蔽,一如它所引发的对抗反应)。与此同时,这份经验想必常与一些类似的、上溯到童年的经验混同——在那里,主角是父亲;而他显身为压抑性权威的主要持有者,甚至俨然一个压倒性的对手,被艳羡也被憎恨。

即便没有这样的混同,即便撇开他人对我身上任何“强制”意图的知觉,也想必常常有一份知觉:在这份(至少在其精神与意图上)纯然阳的“力量铺展”里,有一种强的\textbf{失衡}、一种根本的不和谐。这份失度,对当事者本人——即我自己——是有害的;而且说到底,对他的肉体生存都分明“危险”(近年来的几桩健康事故,已充分证明了这一点!)。想必正是这一点,潜在我思想的底层,当我写道:这样一份“力量的铺展”被感受为“无论如何危险”——“按本性”危险,是一个尤其不足为法的例子……!单是这样的感受,肯定就足以引发“防御反应”——即便没有任何侵略或侵略的意图。

诚然,这样的暧昧关系,1976年之后又重现过,与我的一些学生之间,在那些毫无数学投资的时刻——那时我的生活中没有任何明显的“力量铺展”。同样真确的是:那些\textbf{过去}的“铺展”造就了一份\textbf{名声},它至今贴着我的皮肤,尤其在我的职业生活里;而且在某种程度上,它取代了对\textbf{当下}的我的知觉。更不必说:我在操持某些数学主题时习得的从容如此之甚,以致即便在我的数学时期之外、加上名声的推波助澜,这份从容或天然的驾驭,对一些动力不足的学生,也可以已经有“力量铺展”的效果,使他们尽管看见我身上也有一些亲切乃至令人安心的特质,仍把我感受为某种超人(边缘上还有点超级父亲!)\dmark{}\origpage{478}。

况且,作为我所说那份从容的背面,我常常倾向于低估:对某个学生而言,取得某份行装、发展某件工具,会有怎样的困难——这往往把他摆进一个对我的期望的尴尬境地。(对此可参见“教学的失败(1)”一注,n\textsuperscript{o}~23\textsuperscript{iv}。)这样的局面,想必常常是一项与父亲之间的虚伪关系的重要配料之一……

\subsubsection*{(c) 重逢(阴的苏醒(1))}
\tocsubsub{(c) 重逢(阴的苏醒(1))}
\label{ch18:local:h9}

\notehead{109}{}
\phantomsection\label{ch18:local:n9}

(10月9日)四天前,在收束前一条注\fnote{38}{(*)}{见“阳葬阴——或超级父亲”一注,n\textsuperscript{o}~108.}之时,我感到满心欢喜。我不期然地重新接上了一份直觉——它曾于某个星期日,1976年10月17日来到我心头(距今将近八年):即我身上某股“特定的力量”在我的生活中、一如在我母亲的生活中的毁灭性效果的直觉。那是我一生中头一回,为我的平生——尤其我的童年——作一场反思,哪怕极其简略。那也是我发现禅观力量两天之后的日子\fnote{39}{(**)}{见“欲与禅观”一节,n\textsuperscript{o}~39.};而且,自那一刻以来,这是我头一回动用这份久被忽视的力量。那一天,反思径自转向我的童年,并无蓄意,而是出于一股深层的冲动,仿佛被一个非常可靠的本能所驱动。只是带着事后的认识,我才掂量出:当时把我推向的,恰是我真正的力量之源,同时也是我身上冲突与分裂之源;一份深层的知的渴望推着我走向它。将近三年,我没有再回到那里——那些年,我被一些“日程事务”的问题独占了心思,竟不觉得自己一直逗留在我生活中那个冲突的周边,却执拗地远离其心脏:那段淹没在迷雾里的童年;它显得无限遥远……

我刚刚“斜着”重新翻阅了那场我生命中至关重要的禅观的那十八页——密度超常的十八页。正是在那场禅观之后的当夜,或不如说,在那禅观之夜过后的凌晨,我做了一个震撼人心的梦——那也是我一生中头一个被我从头到尾热烈破译其讯息的梦。当时我对自己的去向、对正在发生的事,同两天前“发现禅观”之时一样,毫无所知。\dmark{}\origpage{479}整整四个小时,我沉入那份经验、那个寓言式之梦的意味之中,穿过一层层越来越灼热的含义,才抵达讯息的心脏、其简单而分明的意味。

那不是“智力”之领悟的骤然咔嗒一声,甚至不是暗处或半明处的一线突然之光。那宁可像一股在我体内诞生的深深的浪,蓦地贯穿我全身;在它浩瀚的水面上,给我送来那份此前一直躲避着我的意味:我在这一刻重新找回了一个至亲至贵的存在——一个自我童年起便失落了的存在……

这一刻被经验为一次\textbf{诞生},一场深刻的更新。这份感觉,那一整天、其后的几天里,都十分强烈。带着八年的回顾,这一刻在我今天眼里,依然是我一生中头等的创造时刻,是我的精神历程中一个本质转折的时刻。它当然是由此前几天、几个月里的许多别的“时刻”所预备的。也许头一个先兆,是十多年之前的那次“有益的挣脱”——脱离那个我本打算终老其中的机构\fnote{40}{(*)}{见同名注,n\textsuperscript{o}~42.}。这些在先的时刻,在我眼前宛如那些配料,或不如说,供我调遣的\textbf{手段}:凭它们,我才得以越过眼前那道我未曾察觉的门槛——它安放在比我已经越过的其他门槛更深、更隐蔽的层面。几天或几个小时以来,一切都已齐聚,只待我越过它;而我尽可以越过它,也尽可以不越过——日复一日,一生如此……

而一旦这道门槛真的越过,路便向进一步的跨越敞开,向别的“觉醒”或“苏醒”敞开——其中每一个,按其本性,都是一次更新,而且多少是一次“新生”、一次再生。有过这样的情形:我把其中某一道回避了几个月乃至几年,才终于跨出那一步,一路卸下某种顽强的幻象——它一生拦在我与我生活、与我周围世界的全部滋味之间。肯定,也有一些,直到我写下这些文字的时刻,我仍在回避……

\dmark{}\origpage{480}在这几天的反思的眼光里,正是这个与我的童年重逢的时刻——人们曾以为它失落、死灭了漫长的一生——标志着我精神行程“第二个时期”的终结:那是我个人生活中\textbf{自我本位机制}占支配地位的时期,对着创造的力量、认识与更新的力量;那些力量,曾经历近四十年的近乎完全的停滞。那也是一个“特定的力量”占上风的时代:一股几乎纯然“男性”性格的力量,映照着周围世界里所尊崇的价值;以牺牲我存在中“女性”的深层侧面与力量为代价——它们遭忽视、遭压抑(谢天谢地,从未全然得逞!)。

对这份力量之毁灭本性的头一份直觉——它曾主宰我的生活,一如主宰我母亲的生活、也主宰在我生命中扮演过重要角色的其他女人的生活——这份直觉,在这些强烈成熟的日子里短暂现身;那无疑是借着“阴”的、“女性”的能量在我对事物的意识把握中重新涌出之便。与我匆匆记起的相反,这份现身,不发生在重逢前夜的那场禅观里,而在重逢之后几个小时——在一场对刚刚发生之事之意味的简短禅观里。这份直觉,在这场禅观那几页笔记的末尾诞生并成形。我知觉到这份“力量”(今天我会称之为“超级阳的力量”,即阳占压倒主导的力量)的毁灭本性——先在我母亲身上,继而在其他女人身上;然后接上这几行收尾的文字:

\begin{quote}
“至于我自己身上的‘力量’,正是它,使我成了一个久待的靶子与对象——在我年轻的一生里,等着M.、然后J.、然后S.~的隐秘的恨与怨——那份恨,早在她们认识我的存在之前许久,就已存进她们心里,在那失去慈爱的童年的张皇岁月里。”
\end{quote}

“童年”一词——在见证我一生中一个头等重要的日子的这最后一行里——也在此后近三年间,最后一次出现!至于那份关于我身上超级阳的力量——作为对抗反应乃至恨与怨的挑动者——之本性的直觉,它(依我看)往往沉入遗忘,直到最近这些天。更确切地说:它只在我对生活中某些重要关系的知觉里(尤其与我爱过的女人们的关系)保持在场。相反,它却不曾真正深入那些多少“不速而至”\fnote{41}{(*)}{或被如此对待……}\origpage{481}的冲突局面——尤其与某些学生的冲突;我在《收获与播种》里已有多次机会检视或提及。在整个这段反思里,这一个事实——即经由一种并非自愿的“挑衅”,我自己也给我此处彼处提及或检视的那些冲突局面,贡献了我自己的一份——常常全然被遮蔽;而对方的贡献,却显得一清二楚。这当然是一个最普遍的反射,不妨说是一个普遍的反射!最近这些天的反思,终于把它拆了引信,同时又使我重新在我自己身上侦察到它——使我突然,在道路的转弯处(在一场关于阴与阳的反思里……)与我自己的鼻子撞了个正着——与一个\textbf{特定的}我自己,至少。

四天前那场简短的反思,也只不过勉强触及我这个人身上的诸多侧面——那“人物”(自我童年起我扮演的角色)中阳失衡借以显形的诸多侧面;以及这份失衡有时对他人产生的碾压效果。尤其对那些阳型力量尚未扎根的人——首先对我自己的孩子。我在这里尤其想到某种“武断自信”的做派:在一切(而且为数众多)我或对或错地有自己一套看法、感受或定见的事情上,我都是照着它运转的。诚然,我从未想过把我的看法强加给任何人,对我的孩子更甚于对任何别人——而且,仗着我身上(至少在有意识的层面)没有任何强制意图,我大半生都无力觉察:然而我的这些为人处事之道(它们在我看来是自发的、自然的;我远没有辨明其复杂的本性……)——它们对我的孩子、对别人,竟有着与强制同样的效果;或者不如说,一个更阴险的效果:在对方心中激起或滋养一份\textbf{不安}——对其自身的感受、看法、意见之价值的不安——仿佛它们(面对我那份毫无破绽的自信、乃至我那份痛心的惊讶)竟连\textbf{存在}的资格都没有。

况且我预感到:我身上这一癖好的发展——尤其在与我的孩子们的关系里——可能相当复杂,与我婚姻生活的坎坷血肉交织。这里不是追索其奥秘的地方;也不是对我这个人身上显出这份失衡的其余侧面试作一份或多或少\dmark{}\origpage{482}完整清单的地方——前一条注里,我试着勾勒了其一个格外显眼的侧面:那份“力量的铺展”。

不应当以为:这份被培养了终生的失衡,连同它借以显现的众多心理机制,会像魔杖一挥,一夜之间烟消云散。我本来也不作此想——无论在重逢的那一天,还是在其后的日子与星期里。

(10月10日)那是冰河解冻的日子,由一股新能量的强劲涌入所托举——是内心工作与惊叹的日子:面对这些崭新的世界,它们日复一日在我眼前半启,诞生于日常琐事的寒微织体之中,在一双渴望看见的、强烈的眼睛的作用下铺展。那也是我开始第一次预感到一个未知存在之丰富的日子:它忽然向我发问,而就在前一天我还对它浑然不觉。我从刚刚向我显露的那些“片段”去把握它——它们在重逢的当下、在随之而来的那段无法预料的旅程中现身。我分明感到:我刚经历的这场“诞生”,只是某件全然未知之事的开端;或不如说,是某件曾被中断、被割断或窒息、如今神秘地重新发动之事的重启。说来,这场强烈的“生成”,在此前几个月便已重新运转,只是还处在内省的思考几乎未曾参与的层面……

这场重新活过来的生成、这项重新开始的工作,其深层的一面,是我身上“女人”与“男人”、阴与阳之间原初平衡的逐步恢复——贯穿其后的日日、周周、年年。在某种意义上,我可以说:自重逢之时起,“童年”或童年状态便一直潜在地在场;我对自己的本性、对自己本质上不可摧毁的统一,有一份深刻而不可磨灭的认识,超越了某种“分裂”的效果——它还常常扰动我存在的表面。用“孩子”或“童年”这个词来称呼这种存在的统一,\dmark{}\origpage{483}倒是好几年后才有的事,约在我开始于有意识的思考层面认识万物的阴阳双面之时。也正是在那时,我开始知道(或至少预感到):童年状态、创造状态,乃是阴阳力量与能量完美平衡的状态,是阴阳“联姻”的状态,表现为一种创造的和谐。

我觉得,在某个层面上,这种对我根本统一的认识,每时每刻都在场,也每时每刻都在起作用。诚然,它的作用随时刻而或显或隐、或强或弱;而它绝不意味着比较持久地消除、甚至整块摧毁那些自我本位的力量,即“老板”——甚至不意味着消除那些压抑的力量(它们构成“自我”的很大一部分,倘若还不是其全部……)。这些力量,暗中偷换环绕我的现实与在我身上展开的现实;它们默默而执拗地运作,顶风冒雨地维持种种顽固的幻象——没有它们支撑,那些幻象便会立即因自身的重量而坍塌……其中一些压抑机制被我逐一认出,随后消失。我摆脱了一些沉重压在我身上的幻象,也澄清了几个顽固的疑点:它们曾终生被扔进地下的垃圾窖(受“老板”照料……),从未受过检视。它们的讯息终于被听见之后,这些疑点便消失,留下一份宁静而喜悦的认识。我也认出一些威力极大的压抑机制,深植于自我之中;近年来我意识到:它们在我生活中的分量,今天一如往昔,依然可观。它们朝着阳失衡的方向,朝着某些阴的力量与官能被遮蔽的方向。我不知道这些机制哪一天会被拆除引信——而我知道,这全在我自己。无疑,它们将在我对生活中那场冲突的根源,进入得比我迄今所做的更深、更彻底的那一天——也唯独那一天——烟消云散。

眼下,鉴于我的生活正朝着一项重要的数学投资转向,我尽可以说:这事毫无进展的迹象!

\subsubsection*{(d) 接纳(阴的苏醒(2))}
\tocsubsub{(d) 接纳(阴的苏醒(2))}
\label{ch18:local:h10}

\notehead{110}{}
\phantomsection\label{ch18:local:n10}

\dmark{}(10月11日)一两天以来,我一直想用几句话\origpage{484}盘点一下:八年之后,我身上这场“阴阳平衡的逐步修复”进展到了哪里。

也许最重要的变化,在于对我这个人——如其每时每刻真实所是——的一份比过去\textbf{大得多}的\textbf{接纳}。换个说法:我身上的压抑机制已大为松弛。如我昨天所说,一些机制在被发现并理解之后消失了;另一些——我一生不曾察觉的——则在其日常显现中变得为我所熟悉。我看着它们运作,不是当作必须不惜代价拔除的敌人,而是当作我受制约存在的诸多侧面之一部分,由此,也是当下“给定”之丰富的一部分——那份“给定”,忠实地映照着我的过往:既有我的种种制约、我存在中分裂之根源的那段“古老的”历史,也有我成熟的那段较近的历史——即那项工作:经由它,我终于开箱取出、并“吃下”消化了父母及其继任者传给我的最初包裹。因此,我身上的这份“接纳”,不仅包括“孩子”的那些我长期忽视并压抑的冲动与特质(尤其那些映出我身上女性侧面的),也包括“老板”特有的那些压抑机制——也就是说,恰恰是那些积习已深的“不接纳”机制!接纳后者,与“栽培”它们或强化它们,毫无共同之处。相反,它是一步不可或缺的第一步:经由一份好奇而慈爱的注意的效果,把它们多少解开或拆除引信。这八年的经验使我深信:只要这份注意沉得够深、直达压抑的根,压抑便会消解、消失,并释放出可观的能量——那份直到那时还被冻结起来、用以顶风冒雨维持某一套压抑机制、以及用以维持它们的思维习惯与其他习惯的能量。

不过,这份对我的新接纳,首先并不显现于对我这个人那些本性上“纠结”的侧面。它来时不声不响,早在禅观发现之前——因此也早在紧随其后的那场“重逢”之前。那是1976年7月,在与一位名叫G.~的年轻女子的一段短暂恋情中——她为人处事,或许比我先前爱过的女人们略多几分“男子气”。偶然(? )安排的情势是:环绕这段恋爱的物质条件,竟把我摆进一个典型的“女性”角色。我料理家务、\dmark{}\origpage{485}准备晚饭,等那位伴侣结束一天漫长而劳累的工作归来:她得在山里看管一百五十只山羊,晚上还得挤奶。这个不寻常的持家人妻的角色,竟像手套一样合我。事情看似微不足道——然而当时就“咔嗒”了一声。这一联,在我心中系向了我恋爱生活中的某些冲动与欲望——它们当时并头一次在几首情诗里表达出来;诗中,恋爱的经验毫不含糊地显为“女性”的。我那时懂得了——不经反思、不经“努力”,没有丝毫迟疑或扭捏的念头:在我的身体里一如在我的欲望里,在我的感情里一如在我的精神里,我是女人,同时我也是男人——而我存在中这两个深的实在之间,没有任何冲突。在那几天里,主音是女性的;我怀着感激、在默默的惊讶中,接纳了这件事。想到它时,我心中有一种宁静的、极温柔的喜悦。

这份喜悦自足自立,无须形诸言词——无论对我自己,还是对他人。我不知道是否对那位以我为男情人(或许是女情人……)的女子谈起过。肯定,在某个层面上,她知道,无须我说。

这份喜悦没有走味;它活着,直到今天。它出自一份活的认识,如香气伴随着花。在我生活的某些时刻或时期,这份认识——以及作为其标志的这份喜悦——比别的时刻更在场,作用更强。但我不相信它有一天会弃我而去。

在其后的星期与年月里,当我此处彼处有机会谈到这一经验与这份认识时,每一次都像是把一件贵重的珍宝传给别人——在我感到他敞开着、哪怕只消几分钟肯领受我身上这份喜悦的些许之时。我从不曾感到一份使我难于启齿的扭捏——仿佛这是一桩多少有点有伤风化的事。(倘若我身上“男人”的实在与力量不是高悬于一切怀疑之上,或许有时倒会有这样的扭捏!)我还记得有一回,我分明是在招摇过市,极力炫耀自己一箭双雕——只差像人人一样来例假、并且立马生个娃了。

我这份新的女性身份,叠加在我的男性身份之上,给我的恋爱生活带来一种立竿见影的更新。它在我此后做过情人的那些女人身上,激起一份\dmark{}\origpage{486}极强的回响:唤醒了恋人身上的男性冲动——那些冲动终生被小心压抑,此前只得以“偷偷摸摸”的方式表达,如某种毛边,不配进入有意识的恋爱经验。

无意识的恋爱经验富含原型冲动,其中最强有力的一种,是对母亲的回归、对原始怀抱的回归。这样的原型,在恋爱经验的深层,男人身上有,女人身上也有。在女人身上,对在两人的恋爱经验中满足这样一份冲动的抵抗,比男人更强:在男人那里,它撞上一道关键的禁忌;在她那里,撞上两道。无论对他还是对她,这些冲动在共同经验中的满足,往往多半停留在象征层面,尤其,避着意识。当这样的原型与这份经验从深层浮上意识之光的表层、进入有意识之眼的视野,这份经验旋即转变,获得一个新的维度。同时,可观能量的解放随之而来——它们先前被压抑机制压缩着,或被压抑的任务束缚着。其效果,是\textbf{性爱}冲动的一份即刻的\textbf{解放}:表现为焕然一新的强度,与恋爱经验中一种新的丰满。

由以上所述,想必已经显明:对我自身的这份新接纳,是与对他人的接纳并行的。两者密不可分。当然,这里说的是完整意义上的“接纳”:它绝不意味着对这样那样被感受为(唉)不可避免的“毛病”或“缺点”的一份(往往是酸溜溜甜丝丝的)\textbf{容忍}——对那种人们只好“凑合着过”的不幸。在那样一种态度里,我首先感到的是一份认命,甚至一份退位;肯定不是喜悦之源,也不是一种认知的冲动——认知那值得认识的东西:在人们肯于容忍的那些“缺点”或“毛病”的平板表面背后,那份隐约感到的、未知的深处……

这里说的是一份喜悦的、创造性的接纳,但这绝不意味着这份接纳是全然的——昨天我已确认,事情并非如此。留心的读者,在《收获与播种》里想必已不止一次亲自确认过这一点;一如我此处彼处顺带地有所察觉:当我发现自己又一次面对我身上那个永恒的\textbf{拒斥}机制——对着任何以不讨人喜欢的外观现身者,无论在他人身上还是在我自己身上\dmark{}\origpage{487}。(不过当事关自己时,这机制最常有的效果,是压根不去知觉那件不讨人喜欢的事……)

我所说的接纳,扎根于一份对被“接纳”之物的\textbf{兴趣}——无论在自己身上还是在他人身上。接纳本身,是一种典型“阴”性的内在心性;而它在我身上带的这层“兴趣”的意味,则是“阳”性的——在阴阳无限交织的那份精微的中国辩证法里,这是“阴中之阳”……我差点顺势说:接纳(真正的接纳!)与这份兴趣、这份好奇,是纯粹的同一。然而稍微端详一下这件事,我发现确也还有另一种接纳方式,其本性比我所惯常的更全然是“阴”的。那像是对被接纳之物的\textbf{迎纳},而非一种向它扑去加以探测的冲动。(这份迎纳的细微差别,顿时间在我眼前显为“\textbf{阴}中之阴”——瞧,来了!)兴趣的冲动与迎纳的态度,都可以构成对他人或对自己的接纳的基调。两者共有的是\textbf{同情}。这也是爱的一种形态。倘若这里有什么深层的东西可供把握,那就是这份确认:\textbf{接纳是爱的一种形态}。对自己的爱,对他人的爱;两者密不可分……

除寥寥的瞬间之外,当事关我自己的时候,我的兴趣比事关他人时更为强烈地投入。正是这份对我自己的炽热兴趣,鼓舞着这八年来的那些长程禅观。诚然,自我认识正是认识他人与世界的心脏,而不是相反;我感到,我的新热情——禅观——从前到现在都把我引向事物的心脏、最本质之处。这些年,我对他人的兴趣,连同由此而来的接纳,显现得更零碎,也更迟疑。它的一种具体表现是:与人相处时,我不再那么爱说话,而愿意倾听。我一生的大部分时间几乎全无这种倾听的能力。即使在重逢的那个大转折之后,我也屡屡发现自己因为没有倾听、缺乏辨别,\dmark{}\origpage{488}在不合时宜的时候说了话;过了好一阵,这积习才开始退去。

后来它不再那样泛滥,甚至几乎消失,丝毫不是我给自己定下了什么规矩的结果(譬如:“只有在……时你才许开口”)。只是因为,在我明知说了也无益、对别人对自己都带不来什么——至少带不来我眼里有价值的东西——的时刻,我失去了开口的兴致。倘若如今我常能感觉到这些,想必是因为我变得更留心了。这也不是纪律的成果(“你可得竖起耳朵仔细听……”);它究竟怎样来的,我说不上。无论如何,我因此感觉更好,生活也更有趣(尤其更安静了!)。别人也因此感觉更好……

我想,我开始真正少说话,是在我身上那股总想纠正别人身上所谓“错误”(不论我判断得对错)的力量几乎消失之时——仿佛察觉并更正我自己的错误还不够似的!这也是那份推着我(有时至今仍推着我)想要不惜一切代价说服他人相信这个或那个的力量,而不是简单看一看:为什么某人执拗地宁肯信这个而不信那个(那个在我看来是“对的”,我很想说服他!)——或者,为什么我如此看重他信那个而不是这个。这份我们身上近乎普遍的力量——它不断推着我们到他人的赞许中(哪怕只是一个人的……)寻求我们所信为真者之正当性的确认——这份深深扎根于自我的力量,我相信,终于在我身上松手了。那是一大解脱,是一场惊人的能量涣散的终结。正是两年前,当我终于意识到这份力量在我生活中的分量、其本性、以及它所代表的那份惊人的能量涣散时,它被拆除了引信——我顿时如释“千钧重负”。毫无迟疑地面对他人向我们的人格回掷的回声,而不被一份对赞许或确认的欲望或“需要”(无论藏得多深)所缚——这才是真正地“摆脱了他”。正是这样一种需要或欲望,构成了那个真正的“挂钩”——不起眼却坚固无比——冲突由此得以在我们身上“挂住”;我们由此(无论我们愿意与否、承认与否)处于他人的依附之下、处于其\dmark{}\origpage{489}\textbf{随心所欲}之下——总之,由此他“拿住”我们,并且(不动声色地)随意摆布我们……

照正常的道理,对他人的接纳,理当也蕴含着对其看待事物方式的接纳——无论这些方式在我们看来对错,甚至当事关他看待我们这个宝贵自身的方式时(包括我们自己的看法……)。然而恰恰在这里最疼——那正是接纳他人的命门,而绝不在于接纳那些不直接关涉我们自身的、多少恼人的寻常“缺点”。况且常常,我们拒斥他人的某些“缺点”,主要是因为经由它们,我们直接感到自己被问责——单是面对一些(又是对错难辨)似乎与我们相反的为人之道。换言之,是我们身上的一份\textbf{不安}——它以虚荣的(或多或少显性或隐性的)反应显形——才是妨碍我们接纳他人的大障碍。而这份扎根极深、由虚荣的动作所补偿的不安,在我看来,与对我们自己的不接纳密不可分;它宛如其形影不离的影子。

于是,对自己的完整接纳,在此显为把我们向对他人的接纳敞开的那把钥匙。而这道刚刚向我显现的联系,又接上另一道我久已认识、或许始终认识的深刻联系:对自己的爱,是对他人之爱的宁静而有力的心脏。

\subsection{(3) 伴侣}
\label{ch18:local:h11}

\subsubsection*{(a) 事物的动力学(阴阳的和谐)}
\tocsubsub{(a) 事物的动力学(阴阳的和谐)}
\label{ch18:local:h12}

\notehead{111}{}
\phantomsection\label{ch18:local:n11}

(10月13日)昨天我没有继续写注。取而代之,我兴致勃勃地把一批“阴阳成对”的东西巡视了一遍。从脑子里随手冒出来的那些开始,我渐渐被咬住了钩,最后以一场力所能及的“普查”收场:凡是我能抓到手的,统统入册。我起头,是因为我盘算过:我近来写下的东西,有相当一部分,对一个对事物的阴阳双面尚不多少熟悉的读者来说,恐怕会整个“从头顶掠过”。费心给出若干醒目的这类成对的例子,以外加最近这些天里顺便溜进来的那些,或许不无用处。\dmark{}\origpage{490}接着,经不住我身上那位系统学的小魔鬼(或小天使,说不好……)的撺掇,我最终翻出了五年前在这个主题上的旧反思。当时我曾饶有兴味地“搜集”了一两百对十分有启发性的对子;它们按亲和性聚成约二十个组。由于那场反思是在我正在写的那部大名鼎鼎的“诗作”的边缘进行的,我忍不住把这些组按亲和与意义的承袭,鱼贯排开。昨晚,带着回顾重拾这场反思,脖子上不再套着诗的轭具,我找到了十八个组(而不是二十个),分组或许更严格了些。况且我疑心,应当还有许多别的组,甚至无数个组,对应于一些我在那项工作里(或许迄今为止都)不曾想过的把握实在的方式。至于我确实辑得的这十八个组,我尽力把它们装配成一幅图(或“图示”),依着联结它们的主要亲和纽带。其中有些纽带,还是在一遍遍画图草稿的过程中,才强加于我的注意。这份工作,与那十分熟悉的数学工作实在相近:当我们力求以尽可能醒目的方式,把一组或多或少复杂的关系(比如由“映射”给出的关系,以箭头表示)——存在于若干“集合”或“范畴”之间,它们充当所构“图示”的“顶点”——按图形方式把握起来。在那里,同样是一些本质上属于审美本性的要求——尤其对称与结构透明的要求——常常引导人们引入(因而在必要时发现乃至发明)一些起初不曾想到的“箭头”或纽带,有时甚至一些新的“顶点”。总之,经过五六次相继的草稿,我最终得到一幅图,大致呈圣诞树的形状;它暂时令我满意——更何况夜确实已深得犯规!

我心满意足地睡下;我感到自己没有虚掷光阴,即便那些注一个字也没有推进\fnote{42}{(*)}{作为补偿,我满可以为一项发明申请专利:一种新的诗体,即所谓“非线性”或“图表式”的诗。}。但我重新与一些分明多汁的事物接上了头\dmark{}\origpage{491}——这些组中的每一个都富于分量与神秘;而构成它们的(不如说,全体一道\textbf{指涉}着它,却丝毫无意穷尽它的)每一对阴阳,都有些精微而重要的东西要对我说:关于我所栖居的这个世界的本性,常常也关于我自己的本性。我以新的力量,重新找回五年前便已在场的那份感觉:阴阳——万物之中“女性”与“男性”——的精微游戏,是一条无可比拟的引线,通往对世界与自我的理解。它把我们径直引向那些本质的问题。常常,阴阳的“瑜伽”本身——我是说,单单留意事物与事件的那一个以阴阳平衡与失衡来表达的侧面——就为更好地理解这些问题、走向一个回答,供给了第一把钥匙。

我该向某些读者道个歉:过去一两页,我给人的印象也许是在谈“天使的性别”;他们可能连我说的那些阴阳“对子”究竟是什么都还不清楚,更不用说由其中若干对子组成的“组”,以及这些组最终又被装配成一幅“图”(数学到底还是有用的!)。这里我至少该举出一组。我想选昨天自发地从它开始的那一组:在反思中,它最终也显为“原始”组;其余各组似乎都通过一层层“衍生”从它生出(在我的那幅图上延续了八“代”……)。以下便是我列出的构成这一原始组的对子(可用其中第一对,称其为“行动——不行动”组):

\begin{itemize}
\item 行动——不行动
\item 活动——被动
\item 清醒——睡眠
\item 主体——客体
\item 生育——受孕\fnote{43}{(*)}{(11月6日)实际上,还有一个更原始的组,可以称为“\textbf{父——母}”组。关于这一“疏漏”,见“我们的母亲死亡——或行动与禁忌”一注(n\textsuperscript{o}~113)。“生育——受孕”与“执行——构想”这两对——我把它们列在了上面那个(号称“原始的”)“行动——不行动”组里——分明以更自然的方式,插入那个围绕“父——母”对子形成的“母组”之中。}
\item 执行——构想\textsuperscript{43(*)}
\item \dmark{}动力——平衡
\item 冲劲——立足\origpage{492}
\item 热切——坚忍
\item 猛进——耐心
\item 激情——宁静
\item 执着——超脱。
\end{itemize}

还有几组对子,是昨天反思的势头在今天早上又送来的十来个“迟到者”里的两个,我愿意把它们添在下面:

\begin{itemize}
\item 知——识
\item 解释——领会。
\end{itemize}

无须说破:在这些对子里,放在前头的,是“阳”的或“男性”的那个词——照我们父权社会的惯例,由男人给这对子起名。相反,中国传统社会比我们的社会父权得多,可当人们循中国的惯例谈论阴阳的关系时,却总把“阴”(“女性”)放在前头——比如说“阴阳平衡”(而非“阳阴平衡”)。这一惯例的意味,肯定在于那份原型直觉:阳生于阴;阴是两者中“最原始”的本原,而不是相反……

这里不是对这一对或那一对子加以评说的地方。对那看了它们却“毫无感觉”的读者,反正也是白费劲;而对那个感到被它们唤起的人——他感到(哪怕模糊地)它们每一个都有些话要对他说,关于世界、关于他自己,关于平衡与失衡、关于存在与事物的内在动力学——这个人无须周详的评注,尽可以把这份唤起当作他自己反思的起点。

\subsubsection*{(b) 敌对的配偶(阳葬阴(3))}
\tocsubsub{(b) 敌对的配偶(阳葬阴(3))}
\label{ch18:local:h13}

\notehead{111$'$}{}
\phantomsection\label{ch18:local:n12}

这里只有一点我想强调——它毫无例外地为一切“阴阳”对子所共有。我觉得,这也是为了理解阴与阳之间关系的本性、并由此理解这两个本原(或能量,或侧面,或力量……)在宇宙中之本性,最为要紧的一点。这就是:这些对子中任何一个的两端——比如行动——不行动,在\textbf{缺了另一端}的情况下\dmark{}\fnote{44}{(*)}{(10月16日)实际上,这份“缺席”从来不是全然的,依我看——在任何事物里,阴或阳都不曾以纯粹状态在场,而无其补充者的同时在场,无论多么微弱。所以我说的“失衡”,其特征不是两个补充端之一的全然缺席(那是从未实现的事),而是那一端的\textbf{过弱}状态。另一类失衡、或病态,则出现在\textbf{两}端都“缺席”之时——或更确切地说,都在场却都极弱。比如在“行动——不行动”这对子的情形,一种\textbf{躁动}状态:它并不严格地“行动”(除非为了自我延续、维持一场混乱),却四散着能量——想必可以算作这样一种“亏空型”(阴与阳两亏)的失衡。}\origpage{493},便构成一种严重的失衡状态;而(当所说的“缺席”近乎完全且持续时)在极限处,是一种导向毁灭的状态——毁灭发生这种失衡的事物(或存在),甚至连同他与他的周围。

比如,一种不间断的\textbf{行动}状态,不与充分的\textbf{不行动}期、休息期交替,便导向衰竭、疾病与(极限处的)死亡——这恰恰是近来于我格外切身的事\fnote{45}{(**)}{对此可参见第XI行列的前两条注(n\textsuperscript{o}~98、99),“亡故者(仍未死去……)”。}!但反过来,过度的无行动状态,又导向身体或心灵(视情形)的官能与机能的衰弱与硬化,极限处导向毁灭。在我那桩“患病事故”里,我恰好有一个\textbf{两}种失衡并存的例子:精神的行动过度,身体的无行动(外加对两者都没有充分的休息……)。

就这个个案而言,这份对阴阳之平衡-失衡“哲学”的“解释”,还停留在表面:因为它不曾触及一份积习已深的文化成见——它抬高阳的一端、行动,却\textbf{对立地}看待阴的一端、不行动。不行动被感受为一种“消极”的东西:不出产,从任何角度看都乏味;勉强被容许,作为一种无奈的权宜——它(唉)甚至强加于世上最善良的意愿,因为人总得时不时歇一歇,才好继续投入行动(否则,正如我方才解释的,就有操劳过度与天知道什么的风险……)。总而言之:不行动被看作行动的寒微婢女——(唉)不可或缺,除此之外不配注意、不配敬重。

\dmark{}\origpage{494}当然,这样一种对行动的“官方”抬举、以牺牲不行动为代价,随即有一份后果:在人身上发动抵抗机制(它们常常保持隐秘,或至少极度混乱),其表达是一种\textbf{相反}的价值化:行动,顿时显为生存的严酷必要所强加的东西,总之是那份“苦工”——在办公室、在工厂乃至在田间,烦得不能再烦,即便不那么烦也累得半死。行动的真正存在理由,是挣面包与住处(这是不可或缺的),而在此之上尤其,是要有美滋滋的闲暇(在有生之年),以及日后一份体面的退休、一份惬意而永久的休闲——那时人们将免除“苦工”这份可惜的义务。这一回,被(或多或少有意识地)抬举的是不行动(别名“闲暇”),而行动成了它的寒微婢女。于是有了一场角色的\textbf{颠倒},却始终是同一份失衡:即当事人(在文化制约的推动下)在自己生活的两个本质侧面或两极之间设立的那份\textbf{敌对};这敌对,经由其一面的专制支配与另一面的奴役状态,表达并延续着自己。

我觉得,最常见的是:这两种态度与价值化在同一个人身上叠加,一种在意识的层面占着街面,另一种在无意识的层面。而这两种相反失衡的叠加,分明并不产生平衡!平衡倒是自然出自对行动与不行动之真正本性的领会(即便这样一种领会停留于纯然“本能”,直接经由一种平衡的行为显形,而非经由一份形诸言语的“所知”)。\textbf{在完整意义上的行动之中,也有不行动}——我是说,它就在\textbf{行动的当下}在场,而不只是“事后”(因为事后总得歇一歇嘛!)。行动之中的这份“不行动”、这个“阳中之阴”,宛如一种深深的宁静,充当着发生于表层的一个运动的立足。比如,它显形于一只行进中的猫科动物所散发的那份完美松弛的印象——无论是随便哪只街猫,还是一头体格雄壮的母狮……而同样,\textbf{在真正的无行动之中},哪怕它全然彻底,\textbf{也有行动}。比如,\textbf{睡眠}富于它的梦:那些梦向我们诉说着我们自己;经由它们,我们过着另一种\textbf{生活}——一种更强烈、更精微的生活;对它,我们常常太沉睡、太怯懦,以至于无从在醒着的生活里去活。而且只需端详一个熟睡的婴儿,或仅仅从酣睡中被唤醒,就会感到:即便没有梦,真正的酣睡本身也是一种\textbf{工作}——以它自己的方式:某种全然吸收我们的事物,说到底,是把一份已然四散的能量“重新充满”——我们刚从其源头把它重新汲回……这又是那个“阴中之阳”\origpage{495};没有它,阴本身也会成为毁灭性的。

\dmark{}同样的方向的反思,想必也可以为醒着的不行动——睡眠时间之外的——发展出来。只消对人们知觉为“无行动”的这样或那样的状态,就实际情形作一番留心观察即可。人们会发觉:无行动之中有行动——哪怕是一个思想继续原地打转的贫瘠饶舌,尽管它已经停止了工作。但说实话,把那个纯然机械的运动称作“行动”是不恰当的:它仅凭惯性而延续——靠的是那台机器停不下来!而这份内在的躁动,当然不会给“无行动”带来一份使之有益的阴阳和谐。相反,倒可以是这样的:那些用来充实闲暇的种种活动(即便人们把闲暇体验为一种不行动)。甚至在完全休息的状态中——比如疗养期——也可以有行动;不然,这份休息或“不行动”就会沦为瘫软懒散,显然不利于疗养(也就是恢复被扰乱的平衡!)。比如,休息可以使人留心自己的身体与近旁的环境(它仿佛构成第二层皮肤……),由此生出认识,乃至一种相通;这本身便具有真实的“行动”性质。毫无疑问,学习确实是一种行动(因为它有一个无可争辩的效果:认识的出现……)。

逐一察看我纳入“行动——不行动”组的十四对对子(想必还可以找到更多自然归入其中的),可以发现:或许只有一个例外,其余每对都是第一端、即“男性”的一端,按我们文化传递并自童年灌输的反射性态度,被赋予声望与“价值”。这仍是我们文化中那种积习已深的失衡的标志:唯独抬高阳——我在前面已有机会提到它\fnote{46}{(*)}{见“阳葬阴(1)——或肌肉与肚肠”一注,n\textsuperscript{o}~106.}\dmark{}\origpage{496}。我遇见的几乎所有阴阳对子,也都可以作同样的观察。这一点实在触目;此前我从未花工夫如此细致地核对过。

在我方才写下的这些对子之中,唯一一个依我看堪称例外的,是\textbf{热情——宁静}这对子:因为在惯常用法里,“热情”一词常与一场爆发、一种暴力的形象相连;或不然,与一种\textbf{放纵}相连——那放纵,遗憾地贴近环绕“\textbf{堕落}”这样的词的那团联想之云。说来凑巧:放纵与堕落,指的正是以\textbf{阴}的、女性的过分支配为特征的心理失衡状态!而对称地,循着同一套按钮机制(它们揭示的是我们时下的制约,而非“宁静”这样一件事物的本性),“宁静”一词(与“热情”相对)便与\textbf{自制}的形象相连——即一种恰如其分地属于\textbf{男性}本质的品质。(实际上,“控制”的阴的对应者,绝不是“热情”,而是“交付”。)

于是,这里发生的是:由于人心对某些事物之本性的一份普遍混乱——它表达为人们用以指称它们的某些词的用法的同样混乱——“热情——宁静”这个阳-阴的对子,与另一对观念

\begin{center}松弛——控制\end{center}

混为一谈了;后者两端是阴-阳(却并不因此构成一个“对子”:在场的两端根本无意成婚!)。因此,依我看,所谓价值化规则的那个“例外”……

恰恰是一份格外有趣的确认!至于我辑得的其他几个例子——那些在一个阳-阴的对子里似乎被抬高的是“阴”端者——想必也是同样情形。

况且,我全然没有把握:我在这所谓“西方的”文明里所确认的这份世界观上的扭曲——它出自那份系统偏袒男性、对立于女性的成见——在中国传统里、乃至在今日的中国世界(或更普遍的“东方”世界)里,就真是轻微许多\dmark{}\origpage{497}。在日常生活的层面上,没有任何征兆能使我作此设想:无论经由我的东方朋友们的日常,还是经由中国或其他远东国家的传统与今日生活传到我耳中的回声——恰恰相反。依我看,倒不如说,一份对阴阳动力学的精微知觉,几乎排他地局限在某些\textbf{技艺的实践}之中——比如书法、诗、烹饪艺术,当然还有医学之艺\fnote{47}{(*)}{(10月21日)我在清单里忘了\textbf{占卜之艺}——《易经》或“变化之书”;它如今在欧洲与美洲的某些圈子里享有盛名。《易经》占卜语言的基本“词”,是那64个“卦”:它们不是别的,正是六位“爻”(阴爻与阳爻)的$2^6$种可能组合——从纯阴(阴的六次重复)到纯阳(阳的六次重复)。那里似乎有一种关于阴阳组合的极其精微的炼金术;据说它曾令Jung心醉神迷。这门炼金术的趣味(尤其作为“一部原型集”)依我看,先天而言,独立于它在占卜中的使用、以及人们对这类使用肯给的信用。}。

尤其后者,顶着“中医”之名、并借着针灸的几桩惊人成功,在过去二十年里终于在我们这里取得了公民权,并被赋予了威望。然而仍有许许多多的人不知道:在中医里,对身体、对气的循行及其扰动(它们构成我们称为“疾病”的病态)之把握的阿尔法与欧米伽,恰恰正在一份极其精微的阴阳辩证之中。这份辩证“行得通”这一事实——既然以它为根基的“中医”是有效的(包括在许多西方武库无能为力的个案里)——可以被视为阴阳这些“本原”或“侧面”或“模式”(把握的或存在的模式)之实在的一种“证明”:它们不是某些哲学家与其他诗人(不说是江湖术士)从帽子里掏出来的纯粹思辨。

人们当然会问:这样的证明有什么意味,乃至关于某一种世界观之效验的任何“证明”又有什么意味。且假定证明果真说服了人(也就是说,当事人肯被说服),甚至再外加一份厚礼:所论的视野深刻、因而是有益的——世上最好的证明,依然无力\textbf{传达\dmark{}\origpage{498}一份视野},更遑论一种世界观。死死“被说服”于一份依然陌生、依然未被理解的视野,对你毫无裨益。干脆说:那甚至没有意义——或者更确切地说,那份“信念”的真正意味,当事人并不比他所佯装修进其沉重视野的这份视野懂得更多。

当视野被理解、被消化之时,“证明”这个问题本身,便显得出奇地离奇——有点像当人们明明看见天是蓝的,却还要证明天是蓝的;或者像证明一朵所爱的花的香气是香的……

\subsubsection*{(c) 一半与全——或裂缝}
\tocsubsub{(c) 一半与全——或裂缝}
\label{ch18:local:h14}

\notehead{112}{}
\phantomsection\label{ch18:local:n13}

(10月17日)我关于“女性”与“男性”双面的最初反思,出自一场对我自己的反思。那是在1979年初;那时我还不知道“阴”“阳”这两个中国词,也不知道中国文化传统里有一套关于阴阳之不断游戏的精微“哲学”。这东西我是同年末才知道的,我想,是经由我的女儿、尤其我的女婿艾哈迈德——他那时开始对中医感兴趣,并在随后几年里深深咬住了它。他对我说的东西,大部分与我抵达的视野相互印证;这毫无使我惊讶之处。倘若说有惊讶,倒是在几个“对子”的个案里:其“自然的”阴阳角色,在中国传统里似乎是被颠倒的。我的反射(在这个案里,是强“阳”的反射!)曾是一种浮于表面的确信:这份“颠倒”必定出于一种文化扭曲——而且我不曾凑近细看\fnote{48}{(*)}{这份对着一个千年传统——它本可以教我更谨慎些——的武断自信的反应,与儿时的那个反射如出一辙:当年它使我驳回书本里教的那个(说实话够复杂的!)公式 $\pi=3.14\ldots$,而偏信我以自己的手段确证的 $\pi=3$。(见“化圆为方”一注,n\textsuperscript{o}~69.)诚然,关于这段阴阳的故事,我早已有充分的机会觉察:“女性”与“男性”之本性及其相互关系的知觉,被积习已深、力量可观的文化扭曲败坏着。但当时我还不曾觉察:对这份关系的精确而精微的知觉,在某些中国传统技艺的实践中是本质的东西,而且被推到高度精微的地步。}\dmark{}\origpage{499}——那是一个我的男女阴阳练习显得十分遥远的时刻,因为我那时正投入一场个人得多的禅观:关于我父母的生平、关于我的童年。想必是数月或数年之后,经由几路旁证,我才发觉:在某些个案里,我对这样那样“对子”中阴阳角色的把握,仍有点肤浅;我曾相当匆忙地,把一些本性不同的局面装进了同一个口袋——而中国的阴阳辩证,是小心区分它们的(注112$'$)。如今我明白:在我这里,对阴阳的把握,仍然相对粗糙、相对静态——尤其若与某些中国传统技艺(如医学:它与食疗及烹饪之艺血肉相连)的实践所要求的那种精微相比;在那里,这份把握最终成了第二天性。

我已有不止一次这样的印象:在这些技艺的修习者与从业者身上——无论东方人还是欧洲人——这份把握的精微停留在断片的形态:它在很大程度上,小心地只局限在该技艺的练习之内。在日常生活中,它更像一个寻常的“所知”那样起作用:干脆叠加在文化(及其他)制约的“所知”之上,而对后者多少形同虚设。换个说法:我有这样的印象——在这些“博识之至”的人们身上,世界观与自我观、知觉实在的压抑机制,与常人身上的毫无二致。

这份印象,又与我读过的两三份文本相互印证——它们的作者号称“在局中”,担着向人概述中国传统阴阳哲学的差事。(其中一位是法国的著名东方学家;名字我一时想不起来。)打动我的,是在这些文本里,阴与阳被呈献为“\textbf{对立}的”(或“相反”的)乃至\textbf{敌对}的本原(后一个词在其中一份文本里反复出现),而不是\textbf{互补}的。这份“对立”或“敌对”,据说有其典型表达:即在人类社会之内、在社会所设立的伴侣之内,女人与男人之间的那份敌对。

夫妇之间的敌对,确实是实在——东方与西方皆然。它深深扎根于文化,以至于有时俨然人类境况的一个(有时令人错愕的!)方面,乃至人身上或人类社会里冲突的根源。这份敌对的实在无可辩驳;它当然超出那些尽力驱邪的通行老生常谈。\origpage{500}这份“社会性的”实在,是一份不可追忆的制约的产物;它早早便在成形中的“自我”里扎根,并结构着它。然而,在这份实在之外,还有一份更深的实在——它来自更远得多的地方——它在性爱的冲动本身之中起决定作用。那是两性之\textbf{互补}的一份深刻而本质的实在:其中没有任何“敌对”的位置。这也是在一切物种之中分明显现的实在——唯独我们的物种例外;在我们这里,它在很大程度上被文化敌对所遮蔽,即被人与人类社会特有的那种\textbf{分裂}状态所遮蔽。

时下浪漫的老生常谈——“我俩”体——支配着大部分文学与媒体,把一份廉价赝品的“互补性”捧上云端,同时对男女之间那份令人不安的敌对侧面,罩上一块遮羞布;或(至多)把它当作一桩有点刺激的意外事件来打发——正好给一顿本来太过寡淡或甜腻的饭食加点辣。一旦越过这类安人心的套话,人们立刻直面这男女敌对的实在——一个显然普遍、且兼有一份无懈可击的顽强的实在:狗牙根般的顽强!但以这份无处不在、无可辩驳的实在为出发点,去设立一种阴阳、“女性”与“男性”的宇宙敌对,便是把人类社会与个人那份撕裂的、深刻分裂的状态——我们这个物种特有的一种疾病——投射到整个宇宙。这同时也是延续自己对自身之内\textbf{另一种}实在的无知(它接上那宇宙性的互补和谐之实在)——一份同样顽强(或不如说,不可摧毁)却更为隐蔽的实在。它逆着那些心照不宣地设立敌对的制约:无论在女人与男人之间、妻与夫之间,还是在我们身上作为“女人”的东西与作为“男人”的东西之间。

老实说,这种\textbf{二元}的或\textbf{好战}的宇宙视野——其中事物的一个侧面,与一个同样本质的“对称”侧面恒常交战——这份视野绝不是\textbf{反思}的果实:仿佛它会“出发”(如我方才所写)自人类伴侣与人类社会里冲突的实在,继而把它“推导”(或按我更准确的说法,“设立”)到整个宇宙。它不多不少,正是文化制约的忠实表达——简直可以说自动的表达——并且顺应该制约的一项本质功能:\textbf{维持人身上的冲突、分裂}。\origpage{501}分明,倘若我肯以那双我出生时领得的眼睛去凝望宇宙,并确认:处处——除了(据说……)在我自己身上、在我的同类之中——“女性”与“男性”都是彼此不可分解的互补者;正是它们的联姻与结合,诞生了造化中一切活物与“死物”里的和谐、创造之力与活的美——那么,维持我身上“女人”与“男人”之间那份既立的敌对,想必是无从谈起的事;或者不如说,这份敌对,\textbf{自那一刻起}便已化解。相反,倘若我自命在宇宙中处处“看见”并非在场的“对立”与“敌对”(即便如此做时,我循着一份可敬的、几度千年的传统),那绝非因为我用过了我的眼睛,而是因为我不过像人人一样,\textbf{重复}了自洪荒以来代代相传的东西;无论如何,是服从了文化共识那份无声而强制的指令——正是它,在我的身上牢牢设立了一份分裂、一场冲突;而我还要把它合理化(并由此延续它),称为一种“宇宙的必然”。

关于伴侣之中的敌对,更普遍地,关于女人与男人的敌对,当然大有可说——我信得过我的同类:关于它,包括中肯的东西,已经写下了许多。这里不是铺陈这个极为有趣的主题的地方,尤其不是铺陈这份敌对在我们父权社会里所采取的特殊形态的地方。我觉得,在那些分明见过其存在的人们中间,有许多人认为:社会的结构——它映照并落实男人对女人的支配——应为这份敌对负责。他们肯定有理;而且我疑心,在一个母权倾向鲜明的社会里,想必能找到一份类似的敌对,以或多或少对称的方式显形。我只想补一句:这份因果在我看来却是\textbf{间接的};它似乎经由一份更隐蔽的因果而行使——今天的反思已拂过其表面。伴侣中分裂的这份更隐蔽、更本质的缘由,是\textbf{个人内部}的分裂状态——无论女还是男——对着自己的冲动(尤其性的冲动)与自己的官能。我在其中看到男人与女人之间敌对的真正根源,一如看到他们在精神层面上那份\textbf{相互依附}的真正根源:我是说他们双方那种\textbf{内在自主的匮乏}\dmark{}\origpage{502}。

这一在自身之内的分裂,在于一种亲身的、隐秘的信念——无论男还是女——相信自己只是一个\textbf{一半}。这份信念的一个征兆,是那种弥散而阴险、从不被检视的\textbf{裂缝}感,或许还有\textbf{残缺}感:仿佛唯有另一性的伴侣,才能解救我们——至少临时地。在那些“大男子气”或“喀耳刻”的应景门面背后(还有许多别的门面),每一个人——男人一如女人——面对潜在的或现实的伴侣,都摆出一副\textbf{乞求者}的姿势:等着对方的(或多或少)好心,给他一场昙花一现的解救——他盼它彻底,它却总是跛脚——解救他那份可怜的破缸状态,甚至不如说碎缸状态:总之是\textbf{半只缸},正在寻找另一只,好与它多少重新粘合(而且粘得坏多于好,人们猜得到……)。

这份裂缝感,或者说,这份对我们真正本性的\textbf{无知}——对我们那超越与性别相连的生理特殊性的根本\textbf{统一}的无知——我身上的这份深层分裂,依我看,是纯然社会制约的产物。至少在婴儿的最初几天与几个月里,人们察觉不到它的踪迹。况且,这份制约绝不限于对“男性”的抬举、以牺牲“女性”为代价,或反过来。说到底,倘若我感到、接纳并且被接纳为\textbf{同时}是两者、“男人”\textbf{与}“女人”——带着一份可以随我人格的这一侧面或那一侧面而变动的主音;那份主音,绝不限于在生殖器官层面占支配地位的那个(诚然极其重要的)主音——那么,我周围抬举的是“男性”还是“女性”,就不再那么要紧了。在我的性冲动的层面上,我的“个人价值化”反正会往往投向与我相反(抱歉,我是想说:互补)的性别,而不会因此感到自己在这具\textbf{不同}的身体面前低人一等(也不高人一等)——一股强而深的冲动把我引向它。况且,无论与性别相连的价值化,还是任何别的价值化:社会共识所借贷的“价值”或威望(对自己或对他人)所占的分量,在一个不大(或全然不)被我所说的“裂缝”感触及的人身上,是相对次要的,甚至微不足道——即在一个身上活着那份自发的\textbf{安然确信}的人:那确信既非僭妄也非门面,而是一份对自身本性的完好无缺的认识的显形。\dmark{}\origpage{503}

“裂缝”或人身上的分裂\fnote{49}{(*)}{我在这里避免使用时下相当流行的“阉割”一词——一个暴力之至的词(这一回真是超级阳的!);它还有另一重不便:暗示一份无可补救、不可逆转的残缺的图像,从而刺激一份慌乱、反叛或认命的反应——这类反应,只会加强一种受阻状态,而无助于其朝渐进化解的方向演化。}并非仅仅是价值化的产物,还有一个征兆:这份分裂在男人身上与在女人身上同样肆虐——也就是说,肆虐于那份自命“抬举”他的共识的所谓“受益人”身上,而这份共识(在某种意义上)却把他连同其伴侣一起压断了腰。人们看到:这份分裂愈发尖锐、愈发暴烈,恰恰是当对一种性别(以“惠及”另一性别为名)的压制愈发强硬、愈发无情之时。\origpage{504}人们简直可以说:“社会”(压抑的源头与工具)在安置这些“压制”之时所遵循的原则,是让男人否认其阴的一面,让女人否认其阳的一面:“分而治之”!而这份由共识造出来、用以压断并奴役男人与女人的“分裂”,还同时在两块板上运作。最显见的一块板,是伴侣之内的分裂:它经由设立一种或多或少的专制的支配——一方性别支配另一方,男人支配女人,或反之——而获得\fnote{50}{(**)}{至少在表面上。不过,如前面所提示,往事情深处走,人们会发觉:伴侣中这份由男人支配女人所维系的分裂,有一份更深的“根”——几行之后我还要回到它。}。一方被认定统治另一方——而双方到头来都成了奴隶\fnote{51}{(***)}{况且,还是这样的奴隶:无论如何不肯与他们的锁链分手;那锁链,对他们比生命还珍贵……}。因为,当妻子或丈夫遭轻蔑时,被交付轻蔑的是双方——有时是旁人的轻蔑;但更深、尤其,是自我的轻蔑。

而在这里,我们接上了分裂游戏的“第二块板”——那更隐蔽的一块:即人本身的分裂。它是伴侣之分裂的隐秘弹簧;伴侣的分裂使它加剧,却不把它归结掉;而它绝不由对一种性别的抬举、以牺牲另一种为代价所单独造出。它宁可是一份无声而不间断的强制的产物:自我们最初的年月起,周围环境便对我们行使这份强制。它逼我们否认我们人格的整整一个“坡面”(“阴”的坡面,或“阳”的坡面\fnote{52}{(*)}{原则上,除意外情形,强制的方向,推着男人否认其阴的一面,女人否认其阳的一面。女人的处境更为微妙:她被认定要否认的,恰恰是那些被社会共识赋以威望、因而她想必感到受怂恿要去栽培的特质。她于是受着两股相反方向的压力;无意识要结构一份“可操作的”身份认同,其任务也随之倍加复杂。}),否则就有被拒斥之虞——被当作可笑,或不合体统;无论如何,当作不可容忍。

\subsubsection*{(d) 原型认识与制约}
\tocsubsub{(d) 原型认识与制约}
\label{ch18:local:h15}

\notehead{112$'$}{}
\phantomsection\label{ch18:local:n14}

\fnote{53}{(**)}{本注出自前一条注的一条脚注(见其首段的指涉)。}比如,在母体-胚胎与阴道-阴茎这两对里:阴阳角色的分配毫无疑义,阴的一端环绕并包含阳的一端。这曾使我匆忙断定:在“容者——所容”这对子里,“所容”是阳;而形-底、外-内、周-心这几对(在其中,正如我的确感到了的,前一端分明是阳,却同时是“容者”)却不曾使我警觉。实际上,在母体-胚胎与阴道-阴茎这两对里,我把重点错放在了在场的两端之关系的“几何的”或构形的侧面——一个次要的侧面;而在本案里决定角色分配的,是主要的侧面:\textbf{滋养者}相对于\textbf{被滋养者}(后者为阳)是阴;\textbf{进入者}相对于\textbf{被进入者}(后者为阴)是阳(同样,\textbf{给予者}相对于\textbf{领受者})。

我关于阴阳的反思,尽管如此有限,却在我身上立稳了一份亲身的信念:在人们对阴阳角色分配(或对一个给定的人身上“阴”或“阳”的“主音”)的种种个体领会差异之外——那些领会极易染上“文化扭曲”——这样一种分配(或“主音”)的“自然的”样子,确实存在。它有一份同样无可辩驳、“宇宙性”而不可改移的实在(就那些普遍本性的对子——如迄今所论者——中的角色分配而言),一如一条物理定律或一个数学关系;即便它既不能经由经验(按自然科学实践中该词的意义)来“确立”,也不能经由一份“证明”甚至一份“论证”来确立。\dmark{}阴阳的这份实在\origpage{505},是经由一种直接的知觉把握的;这份知觉可以(除别的途径外)经由一场足够深入的反思而发展、纯化。

我觉得,这样一场反思的一个主要效果,恰恰是使我们超越那份由周围文化在我们身上程序化的老生常谈式的反射,重新接上与实在本身的接触。实在,依我看,早已在场于心灵的深层——作为一种“原型认识”:它置身于文化制约的射程之外。反思的角色,是使人得以重新接上这份已在的认识,并细心把它从“所知”的表层——即文化制约——中澄清出来。

我在这方面着手的工作,对理解世界与我自己,并由此对我的日常“做”与我生活的航向,曾是重要的。这项工作(一如在许多别的场合)在我看来,像一次最初的突进,像一扇我刚刚推开的门:门外是一座宏阔的全景,尚待勘察。做成它的手段我样样在手——但我不知道有朝一日会不会去做\fnote{54}{(*)}{一如我不知道我在这里看到正在敞开的这类工作是否已经有人做过(说到底,是对宇宙中事物的诸品质及其把握方式,在阴阳互补的和谐之光照下,作一种“局部的与全局的”“地图”式的研究)。况且,这是一个纯然次要的问题:这里要做的,不是提交一篇这个或那个的博士论文,而是深化对世界与自己的领会——那领会,只能是个人工作的果实。}。姑且把数学放在一边,同样“多汁”、且更个人、更灼热的反思主题并不缺,它们无疑会先于对阴阳的一场更一般的反思之深化而占据优先……

\subsection{(4) 我们的母亲,死亡}
\label{ch18:local:h16}

\subsubsection*{(a) 行为}
\tocsubsub{(a) 行为}
\label{ch18:local:h17}

\notehead{113}{}
\phantomsection\label{ch18:local:n15}

(10月21日)三天没有写注了。我的日子被别的任务与事件占去。其中之一,是皮埃尔的来访——带着他的小女儿娜塔莉,昨晚到的。他打算住到明晚,并在此期间读一读葬礼已写成的部分。对我花了将近三个月写成的一份文本,这恐怕稍嫌短促……

我得以分给反思的时间,用来继续\dmark{}\origpage{506}摆弄“阴阳对子”与它们构成的组。这个题目自有其迷人的地方:它兼具探查一个数学“结构”的那份独特兴味——其本性在工作的进程中渐渐明确——与一场关于世界、关于存在的反思的兴味。每一个主要的阴阳对子,都代表一把“锁孔”(在无数别的锁孔之中),透出世界的一个侧面,或世界的一个角落。我迄今辑得的那些对子的“组”,似乎对应于把握宇宙事物的不同可能方式,宛如一扇扇向宇宙敞开的门,从各个不同的角度把它显示给我们。每一扇这样的“门”都有大量的锁孔——甚至数目无限的锁孔——可供窥视;且待哪天,或许干脆把门推开?眼下我只限于侦察出其中相当多的一批锁孔(我找到的足有两百来个),并逐一凑上眼睛看上片刻;同时每一次都意识到:这里有的是值得久久观看而不虚掷光阴的东西……

不过,我更急的,是先去瞧瞧这里或那里另一把还可窥视的锁孔,并且把这些门统统巡视一遍,多少弄清它们相互之间如何布列,或许还弄清:使人们得以侦察出这一扇或那一扇门之存在的那些锁孔,是循着怎样的“模式”布列的……

终于,一周多前侦察出的那十八扇“门”,又添了三扇,共二十一扇;它们布列成一幅图(我曾称它“大致呈圣诞树形”):如今带一根“主干”,由九个“顶点”(或“门”,或“组”,或“角”)构成,以垂直的“棱”或“纽带”相连;主干两侧各带六个别的顶点,它们与主干相连、彼此相连,构成“树枝”\dmark{}\origpage{507}。

\fnote{55}{(*)}{(10月24日)要预言会不会终于出现一些自然地插入我迄今辑得的任何一个组之外的阴阳对子——也就是说,是不是还有别的组或别的向世界敞开的“阴阳之门”,乃至数目无限?——我实在窘于断言。况且,我找不到别的组,绝不意味着不可能存在无限多个别的组——甚至无限多个逃离人类经验、逃离我们把握宇宙之手段的组。这使我记起,近几年里我不止一次被这样一份直觉打动:从蚂蚁或微小的蚜虫,直到已经与我们相当接近的哺乳动物,每一个动物物种都有一些逃离一切其他物种(当然包括我们)的把握宇宙的手段与知觉方式;以致就(不妨这样说的)感知把握我们周围事物的模式之丰富而言,我们的物种不“覆盖”也不“包含”任何别的物种,一如没有别的物种包含我们。而我顺势贸然说出的那句“一如没有别的物种包含我们”,在我看来又嫌仓促,乃至僭妄:鉴于在纯然感知的丰富与精微的层面上,我们物种的演化反倒往往走回头路、退化。只是在智力的层面、在心智图像的精微——尤其与语言相连的那些——上,我们胜过别的物种,依我看。自发来到我注意中的大部分阴阳对子属于这个特具“人性”的音区,而只有寥寥几把带有(此外)显见的感知意味——如影——光、冷——热、低——高,以及别的几把——这绝非偶然。}相当滑稽的是:最近几天出现的三个“新”组里,有一个恰恰是所有组中最显眼的、最原初或最原始的:即对应于对阴阳的那份头一个直觉——“女性”或“雌性”,与“男性”或“雄性”——的那个组。我觉得,它以“\textbf{父——母}”这个原型对子表达得最为醒目(胜过“男——女”——后者属于同一个组)。这个组带着强烈的性意味,显形于诸如“生育——受孕”或“阴茎——阴道”这样的对子;而这些对子本身,又属于环绕那个卓越地的行为——原型之\textbf{行为}——的联想之云:那创造性的拥抱,经由孩子的诞生——经由出自这行为的\textbf{作品}——(至少在潜能上)把女人变成母亲、把男人变成父亲。

与性爱冲动相连的这些意味,五年前的那场反思里始终占据前台。它们而且还在那部大名鼎鼎的“诗作”——当时那场反思凝聚而成的、约130页的东西——的通篇,享有一份近乎不间断的抒情强调;其效果,连最宽容的读者也会读得发腻。在近日反思所依凭的那份独一无二的参考文本里,我干脆“忘掉”了那组当然开道于游行队伍之前(而且不无道理)的对子,肯定正是对这份双重“蓄意”——诗的与Erotique的——的恼怒反应\fnote{56}{(*)}{形式上的这份蓄意,映出一种内在的态度:一个传讯使徒的角色之选定。对此可参见“非古鲁的古鲁——或三足马”一节(n\textsuperscript{o}~45)的结尾,及与之相关的注n\textsuperscript{o}~43.}。

那部著作的题目“乱伦颂”也有些挑衅的味道,\dmark{}\origpage{508}足以使人对其意图与其“讯息”产生误解。况且,这意图与讯息在写作进程中演变颇剧:诗的轭具并未妨碍一项深化的工作继续进行,一份澄清也确实发生了。一个最初的、主要的用意,是探测性爱冲动的某个侧面(我感到它深刻而本质)——就它经由我自己的经验所知的样子。所以说的,首先是男人身上的性爱冲动,或更确切地说:那份“阳”的冲动——它在爱的游戏与爱的行为中对应于“男性角色”,但它以可变的强度\fnote{57}{(*)}{(10月24日)这份冲动的在场,常被力量强大的压抑机制或多或少地彻底遮蔽。我有这样的印象:在男人身上,这份阳的冲动往往比与之互补的阴的冲动占优势;在女人身上则相反。但文化制约,以及对这些制约的种种“正面”或“负面”的内化方式,如此严厉地(而且常以复杂的方式)干扰着原初冲动的交互作用,以致有时很难从零星、短暂且往往变了质的表现背后辨认出它们。}在女人身上一如在男人身上在场。很久以来——或许一直如此——我知道这份冲动按其本性是“乱伦的”:它也是“回归母亲”的冲动、回到原初怀抱的冲动。

这场宏大的回归,在爱的游戏中被“排演”并重新经历,终于在一个寂灭、一个存在的熄灭、一场死亡之中达到顶点并完成。以全部的丰满活出爱的行为,也是活出自己的死——如一次“倒转的诞生”,使我们回到母亲的怀抱\fnote{58}{(**)}{此外,我深信:这份阳性的爱欲冲动的内容存在于一切生物物种,甚至存在于生物界之外;它对应着宇宙万物同一种深层的动力:每一个创造的过程(或“行为”),都是阴与阳、“母亲”与作为孩子的厄洛斯的一次拥抱;孩子回归母亲,消融于她之中。从这孩子回归母亲的“死亡”(或“倒转的诞生”)里,创造行为的果实——“作品”——仿佛从滋养的母胎中涌现。这是“孩子”、新事物的出现,由赋予它生命的“旧物”的死亡与更新之行为而来。在这个宇宙的维度上,原初的性冲动自古就在,远在人类出现之前,甚至早在我们的星球上出现(生物学意义的)生命之前。}。但这同时也是一并僭越两道威力可观的禁忌:乱伦的禁忌——它把“母亲”排除出爱欲对象;以及那份(至少在我们的文化里)把生与死分开、把它们对立起来如不共戴天之敌的禁忌——生与死\dmark{},诞生与死亡。然而我早已分明懂得:爱的行为既是一场死亡——它在高潮的痉挛中完成——又是一次诞生、一个存在的更新:出自这场死亡……宛如一株新苗从滋养的大地里婀娜升起——而那大地本身,由那些倾身于她的存在们的创造性分解构成……

正是在五年前那场关于爱之行为之意味的反思中,我终于懂得:“死亡”与“生命”是同一场紧紧相拥的伴侣中的妻与夫;生命永恒地诞生于死亡,又永恒地倾身于死亡。\fnote{59}{(*)}{(10月24日)奇怪的是:在其后几星期我辑得的一些阴阳对子中,竟没有“死——生命”这对子。也许是因为它与那对相近的“死——生”(或更好:“亡——生”)相混了——后者在册;于是前者便俨然与后者重复了。}或者不如说:生命永恒地倾身于\textbf{死亡},以便永恒地从\textbf{她}——那多产而滋养的母亲——重生;而她自己,又被她的无数儿女对她的永恒回归所不断滋养、不断更新。

而人类的伴侣——妻与夫、恋女与情郎——当其充分地活出把彼此相吸的那股冲动时,便宛如生命与死亡那无尽联姻的一则寓言:每夜爱之终了,情郎倾身于恋人之中、死于恋人之中,以便与她一同从这场死亡中重生——在他们共同的拥抱里……

在同一场反思之初,我曾把人格分裂的一个本质侧面,视觉化为一种“切分”、一道“水平”的切分:即由乱伦禁忌所设立的那道切分——它把孩童从母亲那里“切开”,一如它把生命从其母亲死亡那里切开,一如它也把一代人与生养它的那一代切开。

倘若我起先看到的是这道切分,那无疑恰恰因为我是免于此切分的人。然而,我的生活——一如每个人的生活——曾被另一道大切分深深打上印记;那是我在后来的反思中看到的、并称之为“竖直切分”的:它把每个存在之中“女性”与“男性”的两个“一半”分开来、互相对立,只容忍其一而排斥其二。这正是这一两周来我所投入的这场关于阴阳的长篇枝蔓所谈论的切分。\dmark{}\origpage{510}

我觉得,如今看来,这一(“竖直”的)分裂比另一道(“水平”的)更为关键;在某种意义上,后者为它所蕴含、为它所“包含”。说到底,把孩童与母亲分开、把生与死分开;把一份污秽感、一份排斥或羞耻感——加到死亡上、加到那把孩童与母亲连起来的冲动上——这分明也是在把两者彼此切开、使它们互相对立:把夫与妻——在那两场不可分解而原初的宇宙伴侣之中:母亲——孩子,死亡——生命——彼此切开\fnote{60}{(*)}{这里我按“自然的”阴阳次序书写这两对,以阴的一端——“原初的”一端——起头。关于“母亲——孩子”这对子,应当注意:“母亲”一词也出现在前面已谈过的另一个重要的原型对子里,即在一切之中最原始的“母——父”对子里;它把名字借给了它所描述的那个组。(此外,“母亲——孩子”这组的组名属于别的对子,即我以“因——果”对子之名所称的那一个。)而且,“母亲——孩子”这对子的阳端“孩子”,本身也属于另一个原型对子“老人——孩子”——它与十分有趣的“成熟——天真”这对子相邻。这两对插入我称为“高——低”的组里;那个组,是我迄今侦察出的所有组中(单就数目而言)最丰富的一个,含着许多别的可观的对子,如衰落——腾升,亡——生,毁灭——创造,忘却——学习,终——始……在列举这几对时,我几乎是强制了自己:违背积习,按阴阳的次序为它们命名。粗看之下,这个新次序显得有点滑稽,甚至离奇——总而言之,世界颠倒了!可凑近一看,人们发觉:这个不惯常的次序,向我们显出两端之关系的另一个侧面,一个对惯常侧面的补充:比如(照惯常)“生”先于“亡”——而我们方才分明看到,在更深的意义上,“亡”先于“生”。

我那场反思的总名字“收获与播种”,也是一样——它无疑是一个阴阳对子(这正是我此刻发现的!)。它的命名仍旧循着与惯常的阳——阴次序相反的次序:收获理应跟在播种之后,而非相反。然而这个名字是毫无歧义地自己强加于我的;任何时刻,甚至不曾闪过“这名字本可以是反的——‘播种与收获’”的念头。那是因为:一次次面对不合时宜的收获,总归把我的注意力引回到它们所自出的那些播种——仿佛收获的深意与功能,恰恰是执拗地把我带回到我那双早已被遗忘的亲手播种……}

有意思的是,这最后两对并不曾在《颂》里所辑的那些对子之列。而“死——生”这对\dmark{}\fnote{61}{(**)}{要注意:在“死——生”这对子里,“死”一词的意味,与“死——生命”对子里的不同:在前者它指一个行为(“辞世”的同义词),在后者指一个状态。德语里有两个不同的词:“Sterben”(不带“辞世”那份略显轻佻的意味)与“Todt”。法语里,依我看,最好以“亡——生”来称这对子;这就消除了“死”一词意味上的歧义。},则与我的恋爱经验有更直接的联系,是在册的。“母亲——孩子”与“死——生命”这两对,只是在最近这些天的反思里才出现——在同许多迄今逃脱我注意的别的对子之中;其中最有兴味的一个,是“恶——善”。那是可以称作“艰难的”对子中的一个(一如“死——生命”):那些威力极大的制约使我们把两端知觉为敌对的“相反者”,而不是不可分解的互补者。分明,五年前写《颂》时,这些制约在我身上比今天更强。不过,《颂》里已经有相当一批“艰难的对子”,其中有“混沌——秩序”与“毁灭——创造”……\origpage{511}

回顾起来,对各色阴阳对子之本性——作为不可分解的互补者构成的和谐整体——的一份多少深入的领会\fnote{62}{(*)}{我是指一种不停留于纯然智性的领会;它具体显形为一种改变了的、对他人、对世界或对自己的关系——因为为人之道变了。},如今在我眼前,俨然我们那通往发现世界与发现自我之旅中一级级有待越过的“门槛”。这样一个“门槛”,越是在所论对子“艰难”之时,便越是可观;也就是说,把它作为“对子”来把握所撞上的内在抵抗——文化制约的表达——越强。

\subsubsection*{(b) 挚爱者}
\tocsubsub{(b) 挚爱者}
\label{ch18:local:h18}

\notehead{114}{}
\phantomsection\label{ch18:local:n16}

(10月26日)昨天的反思\fnote{63}{(**)}{这里说的是前一天的反思,记在我排在今天这条注之后的注n\textsuperscript{o}~116里。}起步时有些费劲。这无疑是因为最近几天频频的打断。不过从前一天起,有一样东西在我身上还热乎着,我急于把它托付给纸面——哪怕只写几行。事后我大为赧然:它竟半路丢失了,被一些寻常杂事挤掉了!今天,我不忍心就这样过早地、仿佛出于一场误会地与它分手——在真正与它相识之前,干脆说。

我翻阅了《Zupfgeigenhansl》新近再版的版本\fnote{64}{(***)}{Wilhelm~Goldmann出版社(Wilhelm Goldmann Verlag),1981年。}——这部德国古老民歌的经典集子,约在本世纪初编定出版。据说它已无处可寻,但几位途经我家的德国朋友\dmark{}\origpage{512}带了一本来。前天动工之前,我匆匆翻了一页,有点像路过时与一位老朋友握握手。我撞上了那首歌《今天与明天》;我扫了一眼,没有真正停留——当时我也急着回到等我的工作。但这并不妨碍某样东西“咔嗒”了一声。我分明感到:这些如此单纯、外表如此天真的歌词,轻轻地触到了我身上某种深的东西——而且,这东西离我三天前勉力想唤起的那个,又如此之近。我正准备把这方面的笔记誊清。或许我模糊地感到:我方才扫过的那几节歌,比我那些以断然的简短、匆匆写就(又赶着去做别的事、仿佛顺带)、当日情绪尚且缺席的笔记,更是我想传达的东西的更忠实、更有说服力的信使。

今天早晨起身时,我试着把这几节歌译成法语——它们的曲调我并不知道,但两天来它们一直在我体内歌唱。这无疑是一种更好地找回它们的方式,更好地让它们的滋味与旋律渗入我体内。令我惊讶的是,我在另一门起初显得桀骜不驯的语言里,找回这部德语文本的节奏与音乐的一些部分,并没有太费劲;而字面的意思,始终贴得很近。以下便是这七节歌词,我尽力将它们译了出来\fnote{65}{(*)}{(10月29日)下面的版本,是随后三天里修订过的版本。那天晚上大家唱了歌,我得以学会这首歌的曲调。初版的改动,大部分是为了照顾歌唱文本的节奏与重音的要求。必要时不妨把音节在曲调的音符之间妥善分配;这支曲子可以用法文唱出,而不至于在任何时刻违背重音(一如近来某些新酿的法语歌不幸常有的情形)。}:

\begin{quote}\raggedright
“今天还在,明天\\
我将在你身旁;\\
可一到第三天,\\
我便要离开。”\par\medskip
“可你几时回来\\
我的爱,我温柔的意中人?”\\
“等红玫瑰落雪,\\
等新酒落雨!”\par\medskip
“玫瑰不落雪,\\
酒也不落雨;\\
就这样,我的爱,我温柔的意中人,\\
你也不再回来!”\par\medskip
在父亲的园子里\\
我躺下,睡着,\\
一个美丽的梦落到我身上,\\
白雪纷纷落在我身上。\par\medskip
\dmark{}而当我方才醒来,瞧——\\
纯然的空,纯然的无——\\
\origpage{513}却是那些娇美的红玫瑰\\
在我身上开花……\par\medskip
小伙子回来了,轻轻走过\\
那座美丽的园子,\\
头戴玫瑰花冠,\\
手捧一杯酒。\par
\newpage
他用脚轻轻一踢\\
那座小小的坟茔,\\
倒下——玫瑰落雪,\\
新酒也落雨……
\end{quote}

当我摸索着试图还原我所读到的东西时,我心中有一种喜悦、一种幸福:随着时刻流逝,这首歌仿佛变成了我的一部分。那里有这份剥落而温柔的美的东西,既宁静又切肤,一种由亲密相拥的喜悦与悲伤构成的重而深的美。我相信,很少有人不对这样一首歌多少有所触动——即便他们会抗拒这种触动;人们也常常抗拒一份不期而至的情绪,当我们身上某种我们不自知的深的东西,忽然进入共鸣,并向我们无声地诉说我们宁可不知道的事情。

首先正是梦,有能力在我们身上使那份必须保持隐藏、保持不被知晓、必须保持沉默的东西发出共鸣。也许唯有梦的语言,有能力触到我们身上这些隐秘的弦,并不由分说地使它们歌唱。而当你曾容它们歌唱片刻——哪怕是一首痛苦或沉重的悲歌——你忽然感到轻盈、仿佛焕然一新——像被大水洗过,仿佛一股丰沛的水流穿过你的存在,溶解并带走你身上一切纠结、坚硬而陈旧的东西……

当诗人准备拨动其中一根弦——它的歌会引发内心的洪水——他凭本能会借用梦的语言:既清澈又充满神秘——一种以图像与寓言言说的语言,以其表面的荒谬令理性错愕,又以其隐秘的分明径直抵达它要触动的地方!

这里无须“死”这个词出口,无须任何一个对醒着的理性暗示它的词。然而她在场;她那张迷雾的脸,正是挚爱者的脸。\origpage{514}那沉睡而遥远的挚爱者——你早已离她而去,却又近在咫尺——既是雪,又是落雪而降、又从雪中诞生的玫瑰……那把你吸向\textbf{她}的力量,宛如一股极深极强的波浪,一股来自那位召唤者的波浪、又把你带回她身边的波浪。而这召唤是切肤的悲伤,这回归是低声吟唱的喜悦;喜悦与悲伤本是一体,正是那股把你载向挚爱者的波浪——带着一场分娩的不容置辩之力的波浪。

而且无须提一个字,就已是那缠绵与渴望的冲动——你,孩子——被挚爱者唤入她体内的那个“小伙子”。只消一个梦讲到那位睡在父亲园子里、梦见落雪、醒来见玫瑰的她,你身上那股久被遗忘的波浪便也苏醒,应答着那位做梦又醒来、召唤又等待的她那缠绵的呼唤……

\subsubsection*{(c) 信使}
\tocsubsub{(c) 信使}
\label{ch18:local:h19}

\notehead{114$'$}{}
\phantomsection\label{ch18:local:n17}

这首古老的西里西亚歌谣,只是许多或老或新、歌咏挚爱者与死亡那神秘而切肤的交融的情歌中的一首。我刚刚抄录的这一首,或许多亏其含义饱满的意象之繁复、及其所激起的联想之丰富,而显得不同寻常。这里的主旨,不是把它们一一追索下去——在唤起其中最触动我的一两样之后。昨天与前天,当我的思绪回到那几节匆匆读过的歌时,并不是要深化那份起初还停留在表皮的情绪。倒是这份情绪,提请我注意到:“爱”与“死”的主题——或挚爱者与死亡的主题——显出何等紧密的联系,仿佛被某种神秘的咒语拴在一起!而且,越过那张挚爱者面孔的死亡主题,它们又接上诞生的主题——从睡—雪中觉醒—玫瑰;两者神秘地合一于那幅切肤的图景:玫瑰落雪,落在那位在父亲的园子里睡着、同时做梦又醒来的她的身上。

禁忌尽管把对死的排斥、死与生命及爱的不相容灌输给我们!但不得不相信:它逆着一份扎根极深的认识、或一股既强大又隐秘的冲动——以致那必须不惜一切代价分开的东西,竟以这样的顽强想要重聚,借符号与梦的迂回之路,穿过一代又一代、一个世纪又一个世纪传唱的歌与神话。

关于这些令人不安的交融,写就的渊博卷帙想必不少——为的是多少驱一驱邪。尽管有这些努力,肯定,在我们每个人身上的“某个地方”,这些顽强的联想的深意分明被知觉着——至少在那些我们没有刻意对身上的情绪关闭心门的时刻:它们接待那些信使,用梦那闪烁而有力的语言,向我们诉说着我们自己。

这份“深意”又一次向我们揭示——直接地、以一种基本元素般的力量——经由恋爱的经验:只要我们敢于充分地活出它、倾听它那分明的讯息。它便向我们诉说死亡与诞生的奥秘:二者在传递生命、更新恋人的那\textbf{行为}之中,密不可分。

无疑,我不是第一个:这份“扎根极深的认识”从它长期流放的幽暗深处上浮,变得全然有意识,并由此更加强烈地浸透我与死亡、与生命、与世界、与我自己的关系。然而我有这样的印象:见证这份有意识层面的认识的、已刊的书面文字,想必罕见。迄今我所知的仅有例子,是老子《道德经》里的三四节诗\fnote{66}{(*)}{(10月30日)我撞上《道德经》的这几段,是在1978年末。那是一份触目的、全然意外的确认:确认一些我感受强烈的事物(有些由来已久,有些新近才……),而且我似乎曾是唯一这样感受的人。这场“相遇”被经验为一种大的欢喜、一种无声的欢呼。这份欢喜、这份欢呼,承载了其后六七个月《乱伦颂》的孕育与写作。孕育,是在相遇后的几天或几星期里完成的。以一个更谦逊的音高,近日我又感到一份相似的欢喜:当我“认出”曾鼓舞一位匿名诗人(已死去几个世纪)的那份情绪——他歌唱那些落雪的玫瑰:荒诞地、奇迹般地从“纯然的空无”中诞生;或者不如说,经由我自己亲身的内在经验,重新找回那份\textbf{同一}的情绪——同一份认识的标志。人们也是在《道德经》里,跨越四千多年,重新找到这同一份认识——只是有一层分别:在中文文本里,这份认识以一个高度觉醒的意识的意象化、却绝非符号性的语言表达,而不是以梦的语言(那也是心灵深层的密码语言)表达。\hspace{2em}况且,我在这几节《道德经》诗里认出的内容,分明也逃过了我手头五六种不同译本(法、德、英)的译者们。我毫不惊讶。这类讯息——它们表达一种逆着千年制约的领会——只把其真正的意味(在用来表达它的词与图像之外)传递给这样的人:已经凭自己消化了的东西认识它的人;或那份消化的工作正在其中继续、且已然临近的人。}。

从另一方面(且有点悖论地),我同样有这样的印象:“爱——死”的交融,想必在某个时刻终于沦为一种浪漫的老生常谈、一块稳赚眼泪的“奶油馅饼”——专从最不情愿的眼里抠一滴讨好的泪。事实是:这套手法终于落得声名狼藉——(唉)到这种地步,以至于甚至在那些情感细腻的人们中间,有时也有一种倾向:把纯金与马口铁的粗劣赝品混为一谈。有人在这些地方看见过时的乃至可笑的调调,而那里分明有一份对一种隐藏实在的鲜活精微的知觉、一份远离一切“时尚”的精微表达。一个“好品味”的共识在此援助那些四面八方的内在抵抗:它们自动屏蔽一切鲜活而真实的情绪的闯入——无论是喜悦还是悲伤,是享受还是折磨——凡来搅扰熟悉的老一套者。也是同一机制,屡屡给爱的游戏的原初力量、及其高潮的完成设置壁垒。所幸,这些鼓舞性爱冲动的原型,单单保持被遮蔽、被放逐于意识领域之外,丝毫无碍于它们仍然在场——使该消失的消失、熄灭,以便爱之游戏的意味得以表达并完成,使最终的行为成为一个创造的行为、一次更新。\dmark{}但常常,也有一份隐秘的恐惧,给人们自以为在追求的那份“快乐”本身设置壁垒:人们被一股近在咫尺、既陌生又可畏的力量吓住了——它(倘不留神……)会像扫麦秸一样,扫走我们身上那位不惜一切代价要保住“控制权”的那一位。这样的恐惧,断不容许快乐靠近那个切肤的强度阈值——在那里,快乐同时又是折磨,二者在一场漫长而难忍的拥抱中结合,寻求解脱,以便最终消解、倾身于高潮的空无……\fnote{67}{(*)}{(10月28日)正是这份恐惧——显形为一种对快乐的拒斥——同时又推着人把快乐从恋爱经验的全体中孤立出来,把它化约为恋爱经验的全部,并奉为目的(有时心照不宣,有时明白宣示)。“爱”于是被化约为“寻欢”——说到底,是两个伴侣之间的一种礼尚往来,仿佛互邀去四星级的餐馆用餐,甚或去福利贝热歌舞厅(Folies Bergère)。这份惶恐地拴着皮带的“快乐”,与那股原初的冲动毫无共同之处,一如从大师亲手所绘的画上刮下的干漆屑之于那幅画;一如吹风机之于浩瀚的海风——那带着海与大地的芬芳的远风……}

(10月27日)我想我已经懂得了像《今天还在,明天……》这样的歌与梦的隐秘讯息——就它们所共有的本质而言。于是留下这个问题:究竟是怎样一股力量\dmark{}\origpage{517},如此执拗地推着人给这份“扎根极深的认识”——它想必比我们这个物种还古老——赋以声音;顶着那位乖戾而狭隘的检查官的警惕,逆着所有人把它表达出来,逃出家门、在梦那资源无穷的符号语言里自由驰骋?

倘若神话、歌与梦以无数的口吻不倦地向我们低语同一份讯息,那么同样真确的是:它们所诉说的那个囚徒,听它们也听不厌!诚然,那是一个自愿的囚徒;他小心不肯听。他被剥夺了空气、空间与光,却由四面墙安心地护着——那墙围起一种没有大意外、没有大奥秘的生存;若说有奥秘,或许只有尽头那场死亡,无限遥远……他的监狱保护他,避开墙外的那个未知——他佯装不知的未知。它既使他害怕,又使他着迷。正因为墙外的那个“彼岸”使他害怕,他的这座避难所式的监狱才对他比生命还珍贵。然而它又使他着迷、吸引他——违背他的本意;一如那些偶尔前来向他诉说彼岸的信使们吸引并迷住他。有时他向这份不寻常的吸引让步——只要能背着检查官——那位总监——行事:他一面不动声色地侧耳,一面却依旧处于“暂停”状态——他什么也没听见,尤其,什么也没听进去!

我方才向自己提出的问题,仿佛消失不见了——被一幅有说服力的图像藏掖了。一旦我记起那份讯息的效果——那种来到讯息之前的情绪、以及这份情绪的益处——问题便重新出现。

但老实说,任何触到一根深弦的情绪,都是四墙之外的彼岸的信使、是辽阔的信使。即便下一刻我们就努力抹去它的一切痕迹,它仍是有益的;它已经留下自己的痕迹,如一丝精微的香气——仿佛那些阴沉的墙挪开了哪怕一点点;或仿佛经由一道不曾料到的开口,有一缕哪怕极微弱的林野芬芳,透进了这消过毒的空气里。\origpage{518}

\dmark{}(10月28日)近十五天来,反思走进一个全然没有预料的方向——它与葬礼的主题没有明显的联系,甚至(似乎)与我自己的这个人也没有——这多少是违背我的本意的。我心底清楚:事情并非如此,我继续被卷入这些注中,不比以往少,甚至更多。这并不妨碍我在两种欲望之间摇摆:一边想“了结”,一边想开掘日复一日隐约瞥见的东西、追索那些最紧迫的联想——后者又与另一份关切合流:不放过任何能照亮我这场葬礼“侦查”的东西。显得最远的,有时恰恰也是最切身的……

总之,这十五天来——倘若不是自患病事故后重新提笔以来——我有一种(有时颇不好受的)印象:我在“赶工”、匆匆忙忙;仿佛每一条新注都是我又一道打开的括号(在一个会喊饶命的假想读者面前),而我理应尽快把它关上!肯定正是这些心境,或许更甚于最近几星期里相当多朋友在我这里的走动,造成了一份同样匆忙、有时略显杂乱的书写。我不得不随手把近来写的大部分注陆续誊清。这又拖慢了进度,吊着我那份眼看工作进展的急切胃口!

诚然,这些我有时佯装要趁热打铁、当作“众所周知”只为求个心安、并为一位刚刚“下船”的读者之益而费心显豁的主题——它们过于精微、分量过大,经不起如此漫不经心的态度。我忍不住在页页进程中觉察到这一点,并“拨正船头”——我是说,调整我的内在态度,顶着我本打算偷偷摸摸应付的东西的那份(不妨这样说的)重量!

这使我想起:这场我已投入将近四个星期、且远未终结的阴阳长篇反思,说到底,不过是在显豁一份瞬间的直觉——一份在我看来再单纯不过、甚至可以说明明白白的直觉;一份在5月12日的次日“闪”来的直觉——那时我刚写完关于某篇“葬礼悼词”的第一条注。一个月前\fnote{68}{(*)}{见“肌肉与肚肠(阳葬阴(1))”一注,n\textsuperscript{o}~106.}我接续那条注时,打算循着这一条联想(而非另几条在我看来趣味稍逊的联想)走下去;\dmark{}\origpage{519}我预计这顶多再花五六页。而那里,我已经越过了六十页的大关……

昨天,我在这个问题上停了笔:在爱与死之间、死与生之间、生命与死亡之间的联系的象征性唤起,其意味何在——以及这样一场唤起在我们身上激起的那份情绪,其意味又何在。在神话、歌与梦之中起作用的力量是什么,推着它们“以无数的口吻不倦地向我们低语同一份讯息”?而在我们身上——那些令人安心的监狱里的自愿囚徒——又是什么力量,屡屡以那份情绪回应它们,迎着那场唤起,表明它“命中了目标”、触到了它想触的地方?还有:梦的语言的那份奇特的力量从何而来——那种暗示而不点名的语言,传递着任何别的语言都不会传递的东西?

追索这些问题,也是更深地探测性爱冲动的作用——一如梦的作用——以及连接两者的深刻纽带:两者彼此滋养又被对方滋养,各自表达着自己,并经由一种它们共有的、逃避检查官的语言与对方交流。这同样是更深地探测原型与符号在性爱冲动中的作用、以及该冲动那些“符号性”满足的作用。

分明,这一切把我引向远远超出我所能合理指望“塞进”这场阴阳“枝蔓”的东西的限度——那场枝蔓(该提醒我自己了!),正延续在一场葬礼仪式的正中间!我觉得,是时候放下这根新“线”,回到三天前\fnote{69}{(*)}{见“失乐园”一注(n\textsuperscript{o}~116),它排在“本注”(注n\textsuperscript{o}~114)之后。}悬置的另一根“线”:当时它把我带回我自己的这个人。

\subsubsection*{(d) 安格拉——或告别与再会}
\tocsubsub{(d) 安格拉——或告别与再会}
\label{ch18:local:h20}

\notehead{115}{}
\phantomsection\label{ch18:local:n18}

\dmark{}\origpage{520}(10月30日)一两天以来,有几行诗在我脑际萦回,出自三年前写的一首诗。我先是用德文写成,次日又用法文重写。浮上心头的是头两节——第三节即末节仿佛从记忆里被抹去了,只除第一行“Ein~Kreis~schliesst~sich”——“一个圆环合拢”。(还有末行也一样,它重复了第一节的末行。)昨夜醒来时思绪又回到它;我终于起身,翻检纸堆。没费什么事就找回了那首诗——整齐归置毕竟有好处!它在此。

\begin{quote}\raggedright
果实沉沉\\
熟而重\\
我的生命俯身\\
为了归返\\
她之中\par\medskip
甜而稠的汁液\\
浸透了我\\
曾开花\\
脆弱的奶白的花\\
已成果实与酒\par\medskip
一个圆环合拢——\\
自我的怀抱\\
升起甘甜\\
画出它的环弧\\
又以弱音俯身\\
为了归返\\
她之中……
\end{quote}

我想,这是我写过的诗中唯一一首死亡之思\fnote{70}{(*)}{我更应当写:对我自己之死的思。1957年——我母亲去世那一年——写下的两首诗(各有几行),浸透着对这场死亡的预感。}分明在场的诗。它在这里以“她”之名现身。在前一晚的初稿里,它由德文词“Erde”——大地——唤起。顺便说,这三节的德文“译本”远非直译;第一节当初是这样来的:

\begin{quote}\raggedright
Voll und schwer\\
reife Frucht\\
neigt sich mein Leben\\
gen Ende\\
der Erde zu\par\medskip
Die süssen Säfte\\
die mich durchtränken\\
haben geblüht\\
weiche Blüten und wurden Frucht und Wein\par\medskip
Ein Kreis schliesst sich\\
aus meinem Schoss\\
steigt Süsse\\
kreist\\
und neigt sich\\
gen Ende\\
der Erde zu……
\end{quote}

\dmark{}\origpage{521}最终,刚刚把德文初稿重写一遍时,我竟不由自主地把它一气写到了底——那后面两节仿佛自发地从第一节里流出来!这三节,对我来说,是一首情诗(此外我写的诗,也几乎全是情诗)。倘若它致予的不是我自己,那便是致予她——致予那位在沉默中等待、准备迎接我的她……

同一天,我还写了两首别的诗,一在之前,一在之后。它们致予一位有血有肉的“挚爱者”:安格拉——“安琪儿”——一位高大、金发、苗条的姑娘,浑身洋溢着生命;一星期前,在那条夏天热浪蒸腾的路上遇见她,她正搭车。一两个钟头里,我们来得及彼此说了许多;然后就此别过。我本想把这些她所引发的诗送给她——包括在她被遇见的当天晚上写的另一首,以及还有一首(仍旧用德文,我们的共同语言):它诞生于那“三首(几乎)一挥而就”的次日。我也本愿我们相爱……但我失去了她的踪迹,一如她想必也失去了我的。

这场相遇所引发的诗的一个共同点:每一首都或者极强地“阳”,或者极强地“阴”。它们属于我写过的最炽烈的诗之列,每一首都一挥而就,几乎不加修改——仿佛它们早已齐备在那里,只等这场相遇的信号,好取得有形的言语之身。

乍看之下,在这些满载炽烈情欲张力的诗中间,竟有这另一首秋声阵阵、准备进入冬日长眠的诗,似乎奇怪。但只有感觉不到情欲的冲动与死亡之感之间那条深深纽带的人,才会诧异。在那些孤独的日子里,有一种对生命的炽烈知觉——由情欲、由支撑着它的原型意象的繁富所放大;同时,又有一种充分活过、正走向终点的生命所特有的宁静的超然:准备“归返她之中”。

与死亡——我们沉默的母亲、被感受为朋友而近在咫尺的她——契合的这样的心境,肯定被身体的一种大疲惫所促成:它把我们带回那些简单而本质的事物:我们的身体,爱,死……当时,我刚走出一场“数学的长期亢奋”——关于它,我已在《收获与播种》的导言里谈过\dmark{}\origpage{522}\fnote{71}{(*)}{见“梦与完成”,尤其第(iii)页。所论的“亢奋时期”,自1981年2月延续到6月;那也是“穿越Galois理论的长征”的时期(见“Galois的遗产”一节,n\textsuperscript{o}~7)。它通向一场关于我与数学之关系的漫长禅观(见“扫兴的老板——或高压锅”与“非古鲁的古鲁——或三足马”两节,n\textsuperscript{o}~43、45),自1981年7月19日延续到12月。献给安格拉的诗(以及献给“她”的那首),写于7月8日和9日(只有最早的一首作于7月1日)。}。我刚刚才开始从一场身体的衰竭中往上爬——那场颇有些疯狂的时期把我抛进的状态。它刚刚结束(来得与去时同样突然),借着一场碑文般有力的寓言梦的冲击;我当时肯倾听它的讯息\fnote{72}{(**)}{见注n\textsuperscript{o}~45的开头——前一条脚注已引用。}。那是可听可看的日子、倾听的日子——两波之间的一段“敏感时期”:身后,是一道长而宏阔的“数学”波;眼前,一道同样宏阔的“禅观”波已在酝酿……十来天后,它将开始它的冲刺——随着另一个梦:它的叙述开启《收获与播种》的导言,那幅“如我所是”的我自己的视野。

那是内在工作强烈的几个星期、无声的孕育、转变的几个星期。而这些情诗——音色不同于我此前写过的所有诗——正是这份强烈、这份丰满的果实与见证。

它们也是我写下的最后几首情诗。也许我身上有一种预感:这将是我最后一次恋爱,最后一次为挚爱者燃起歌唱的盛大焰火!一份预感:这些致予一位陌生姑娘的诗——我未曾相识却强烈地感到她的美——同时也是一场告别:向情歌、向我爱过的女人们告别;向我那正在这束光焰四射的花束中燃尽的爱的热情告别——它即将离我而去。而更隐秘、更深地:这也是一场告别(或许是一场再会):向着所有的女人——她们交相融合、在一个新的面孔下化为一个“她”。一张也许更遥远的面孔,淹没在雾中,在路的另一端——却同时非常近、非常温柔……

\subsection{(5) 拒受与接纳}
\label{ch18:local:h21}

\subsubsection*{(a) 失乐园}
\tocsubsub{(a) 失乐园}
\label{ch18:local:h22}

\notehead{116}{}
\phantomsection\label{ch18:local:n19}

\dmark{}\origpage{523}(10月25日)\fnote{73}{(*)}{(11月1日)本注早于前两条注;它们写于10月26至30日之间,并且是“行为”一注(n\textsuperscript{o}~113,10月21日)的直接延续与深化。本注倒是系于更早的10月17日那条注(n\textsuperscript{o}~112)的结尾,即“一半与全——或裂缝”。从这里起,反思遂分成两条平行的道路:一条关于死亡之感及其与爱欲冲动的联系,在依次列出的注113、114、115中继续;另一条则由本注n\textsuperscript{o}~116开启。}又过去了三天,我没有找到时间乘势写下去。第一天,星期一,主要被皮埃尔携女儿(两岁的)娜塔莉的来访占去;晚间很迟我才送他们去奥朗日赶夜车。再过几天,盘点这次来访带给我什么,还来得及——那是一次我已不再指望的来访……眼下,我更愿意断断续续地继续阴阳反思的线头。

这场反思,俨然一段哲学枝蔓,突然闯进某场它本无用武之地的侦查——只除了一点:它是从围绕某篇葬礼悼词的几缕联想中,不声不响地冒出来的……然而我分明感到:恰恰是借着这场“枝蔓”,我方才开始超越对构成葬礼的那些“原始事实”之总体的揭露阶段\fnote{74}{(**)}{(11月14日)这句顺口说出的话,并非深思熟虑,只有部分根据。至于一份更周详、更细化的概观,可参见“一场禅观的回顾——或一幅三联画”一注(n\textsuperscript{o}~127)。},终于稍稍逼近那些在种种似乎离奇出格的行为与举止背后起作用的力量……肯定并非偶然:正是经由这场“枝蔓”,我在自己毫无预料的情况下,被引至一种比《收获与播种》的任何其他时刻都更深的意义上牵连我自己的这个人的地步。这是最近的患病情节的一个意外果实——它恰好发生在我准备轻装把七周来的侦查带向其近在眼前的结论之时……

因此,这场“枝蔓”——有人会在其中看到一场亲身的忏悔,另一些人会看到一场玄思的形而上学——对我而言(比《收获与播种》的任何其他部分都更),处在葬礼的正中心、冲突的正中心。改变的只是眼光、看这件事的“视点”——而这一变如此苛烈,以致方才还在检视的那件事仿佛骤然消失!\origpage{524}我想,我们很快就会重新接上那看似半路丢失的联系——与“市井小事”即葬礼的联系。

不过,人们也可以把那桩市井小事忘了——它的主要功劳,毕竟是引发了这场“枝蔓”……昨天的一天,用来誊清前一条注的草稿——它写于四天前,我最终给它命名为“我们的母亲,死亡——或行为”。那份草稿有相当一部分涂改很重:一个征兆,表明表述还欠清晰;一些重要而精微的主题,当时是有点“从旁边”、在赶往别处的半路上混进反思的。老实说,起笔这条注时,我主要打算接上前一条注的线头——它名为“一半与全——或裂缝”,写于整整一周之前。可最终那根线头仍旧悬着;该是终于把它拾起来的时候了。

也是为那条注,我同样重誊了文本的相当一部分,缘由本质相同:沿途修正笨拙与晦涩。那是一场关于伴侣之内的分裂的反思的开端;它与个人内部的分裂血肉相连——更确切地说,与我(在四天前的“行为”一注里)所称的“竖直切分”血肉相连:它把我们身上原初之“全”的两“半”之一——阴或阳——“切开”,或削去。

在一个如今仍属于直觉领会、未形诸言语的层面上,我“懂得”、它对我“分明”:人本身内部的分裂(它俨然是全然由制约造出来的)是人类社会中无处不在的冲突的深层缘由。无论伴侣之内或家庭之内的冲突,还是更大群体内部的冲突、或把一些群体相互对立起来的冲突,直至各民族各国家之间的武装对抗。伴侣内的冲突——它把两位典型的、分明可辨的对手相互对立起来——满可以名正言顺地显为人类社会之内冲突的基本寓言、其不可再分的元素个案。“裂缝”那场反思的“要点”,主要是把伴侣内冲突的个案,归回到那个更根本、更“元素”的个案:每\dmark{}\origpage{525}个人自身内部的冲突——它使人自身的“一部分”与另一部分对立。

在七天前那场反思的眼光里,首先想到的,自然是我们身上“阴”与“阳”两“部分”之间的冲突——其中一方被接纳并妥善地抬举、吹胀,另一方被拒斥、被或多或少彻底地压抑。然而我心里明白,个人身上还有别的对抗,系于别的禁忌,而非性别只能有单一面向的禁忌。诚然,后一种禁忌与乱伦禁忌同样强,却更阴险,因为它带着一副明摆着的面貌——仿佛正因其明摆着,便免去了把它表述或命名的心思!尚未一步一步亲自检证,我便有这份印象(自《颂》的反思以来):这个禁忌是所有禁忌中最关键的一个;它在人身上设立的分裂或“切分”,是人身上积习已深的分裂之诸多侧面的终极根源。细心弄清事情在多大程度上确实如此,无疑会是一个最诱人的起点,足资开启一场“探索冲突之旅”。然而这里不是启程的地方——更不必说,在我面前摆着的、为我而设的那些旅程之中,我看到的起点有比这更灼热的……

誊清“一半与全——或裂缝”那条注的文本时,我还发觉:写它时,我不曾想到哪怕稍稍显豁一下,为什么我把个人内部的冲突视为伴侣冲突、社会冲突的深层缘由。这件事,我方才说过,属于那些我“懂得”过(却至今无须向自己“解释”过)的事情:它们由千百件日常小事那无声而雄辩的语言教给我、并向我证实,经年累月\fnote{75}{(*)}{这份“领悟”或信念,在我看来,并未被下面这桩我屡次得以作出的观察真正驳倒:由母亲与父亲结成的伴侣之内的分裂,以及表达这分裂的那些对抗姿态,会在孩子身上留下深深的印记,并常常支配着成年后的姿态与行为。至少在很大限度内,我们身上的分裂,是童年岁月里使我们母亲与父亲相互对立的那场分裂的印记与遗产——这样说无疑是站得住的。于是,要裁决个人内部的分裂究竟比伴侣之内的分裂更根本、更“元素”,抑或相反,这个问题有点像先有鸡还是先有蛋!
然而我确信:在一桩伴侣里,倘若夫妇中有一方是“一”、不与自己为敌,那么即便另一方对他怀着对抗的姿态,冲突也\textbf{不会}传递给这对夫妇的孩子。我相信这份确信的缘由是:在那种情形里,孩子会被他的父母之一\textbf{全然}接纳。幼童身上分裂的出现,在我看来不多不少,正是他的存在的一部分遭到周围环境\textbf{拒斥}的效应——而拒斥,首先是他的\textbf{双亲}。}。我并不是说:在这里花几页、或许卷帙浩繁的几册,来显豁或“解释”其“所以然”与“如何然”,毫无趣味。无疑,关于这个话题的几页,置于这些注中,并不比已在此安家的任何其他关于阴阳与冲突的篇幅更“不合适”。我肯定会在其中学到许多东西,一如继续那另一条反思的主题——我们身上被设立起来的阴阳冲突、作为我们内部分裂之终极缘由——也会教我许多。\dmark{}况且,这两个主题之一分明延续着\origpage{526}另一个,这使得它们各自愈发诱人!然而,眼下我毫无意愿朝那个方向走,哪怕一步。那不是我一周以来最想拾起、却始终悬着的“线”。

一周前,当我为那条注的反思收尾时\fnote{76}{(*)}{即“一半与全——或裂缝”一注(n\textsuperscript{o}~112)。},我忽然满心欢喜、精神一振:反思出人意料地重新接上了某样重要的东西——前几日我已有点失去视线的东西:接纳。这一接,是经由否定的路径恢复的,靠着那个为反思收尾、宛如一场出人意料的持续强音的词——“不可接受”。正是由于我们人格的整整一个“坡面”被周围环境当作“不可接受”而拒斥——首先由给我们定调的父母(或当父母缺位时,由代行其职者)——正是经由这份\textbf{不接纳},冲突才在我们身上安营扎寨。我们身上的冲突、分裂,不是别的,正是我们对自身一部分的退位——那一部分遭否认;是我们那未分裂的本性的退位。这份退位,是我们为被周围环境多少“接纳”而付出的、必须付出的代价。况且,那样的“接纳”并不是完整意义上的接纳——不是对我们真实所是的接纳。它宁可说是一份\textbf{报偿}\dmark{}\origpage{527}——报偿我们对某些\textbf{规范}的服从、报偿我们按规范自我塑形与浇铸——说到底,是一份\textbf{变形}、一份对我们存在的\textbf{戕害}的报偿;那变形与戕害,与我们的周围人自幼年所承受的如出一辙。

在前几条注的反思里,已有两次谈到接纳;两次,接纳都显为一件关键的事。第一次,在“天真(阴与阳的联姻)”一注(n\textsuperscript{o}~107)里:我在其中重提一份上溯到四年前一场禅观的确认:我身上一股未分裂之力得以绽放并充分舒展,竟发生在一个被冲突与潜伏的憎恨撕裂的家庭的语境里,\textbf{唯独因为我被我的父母、被我的环境全然接纳};冲突进驻我的存在,是在其后——五岁之后——在一个比我原生家庭“平静”得多的环境里。当然,亲人之间的冲突(至少在我那时候)远未达到我原生家庭那种加剧的(尽管蒙着面纱的)强度。然而,在那原生家庭里,我自己的这个人始终置身冲突之外。即便我偶尔表态站在哪一边,那也不是撕裂,而是一个未分裂的存在之自发的表达——一个从不曾尝过被自己人拒斥、以及对拒斥的恐惧之咬噬的存在。

如今,带着半个世纪的回顾,我明白:在我当时的新环境里,我身上这股天真的力量同样散发一种辐射,一种我愿称之为着魔的引力;宛如一座\textbf{失落的乐园}——无限遥远,令人一生怀恋——忽然借一个孩子的声音与目光向你发问。它当时为我赢得了强烈而持久的感情,追随我直到成年,直到那些如此爱过我的人一一谢世\fnote{77}{(*)}{我记得有七个人曾这样把他们的感情给了我,其中今日只有一位尚在人世。}。但\textbf{与此同时},理所当然:这股力量\textbf{不可能被容忍}——一如在一座修剪得笔直的观赏园子里,人们无法容忍某棵繁茂而恣肆的树或灌木:人们自以为爱它,却执拗地把它剪成方块、圆锥或球形……

按我对事件的复原\fnote{78}{(**)}{对我童年这些重要事件的复原,作于1980年3月。},这股力量坚持了或许两年、两年半,然后深潜下去,\dmark{}被贬入地下——那是在我终于决意“跟人人一样地做与是”之后:\origpage{528}十足十的肌肉与大脑(肚肠活该!),好落个清净!我终于步了后尘:凭着我周围所有成年人毫无破绽的共识,我\textbf{拒斥}并否认(以忽视的方式)了一切该被拒斥、该被忽视的东西。而且也凭着我父母自己的共识——他们到头来几乎不再给一点活着的音讯,在与他们的孩子尽可能远的地方过着他们的伟大的爱……

\subsubsection*{(b) 循环}
\tocsubsub{(b) 循环}
\label{ch18:local:h23}

\notehead{116$'$}{}
\phantomsection\label{ch18:local:n20}

(11月1日)我拾起恰好一周前中断的线头——10月26日,我曾出人意料地投身一场关于死亡之感在爱与情歌中的“诗的枝蔓”。

我刚重读了10月25日的前几页,并把其中最后一页誊清。我仿佛看见一个圆环正在合拢:它的轨迹两周前已经开始,随着“力量的绽放——或联姻”一注(n\textsuperscript{o}~107)。这轨迹在前面几页里收束:它们重拾并放大了10月17日那条注“一半与全——或裂缝”(n\textsuperscript{o}~112)的最后的“持续强音”。这个持续强音——或为当日反思收尾的“临终词”——归结为最后一个词的绝对律令:那个词,“不可接受”。

这个收尾的词,在形塑了我们生命的、形形色色令人错愕的大量制约之中,在我看来恰如其分地圈定了\textbf{那个}使我们身上发生分裂的决定性缘由:它就是我们生命最初年月里对我们的这个人的\textbf{不接纳}、\textbf{拒弃}\fnote{79}{(*)}{在这一点上,我本人的情况是个例外:直到六岁起,我才开始遭遇身边亲近之人这类姿态。}。它落实为对我们身上某些力量与冲动的失于接纳、予以拒斥——那些力量与冲动,是我们的存在、我们认识与创造之力的一个本质部分。对这些力量的压制,被一位忧心忡忡、毫不容情的内心检查官接了过来、据为己用,乃是对我们身上这种力量的一种戕害。其效应,往往无异于我们的创造官能一场真正的瘫痪\fnote{80}{(**)}{(11月2日)它往往还更显豁地表现为种种“卡壳”的效应:既无力在我们身陷的某种处境里“运转”,又无力从这毫无出路的处境中脱身……}。

\dmark{}这种不被接纳的力量,或这些“官能”,不是别的,正是那谦卑的能力:做我们自己。\origpage{529}这同时也是说:凭我们自身官能那份谦卑而完整的运用,活我们自己的生活,而不是过一种千篇一律、被编程的生活——那样一种生活,主要(且往往单单)靠\textbf{重复}、\textbf{模仿}的反射来推动。这些反射把我们封闭起来、隔离开来,宛如一副沉重的甲壳,僵硬而不透水,我们无论何时都无法从中脱出\fnote{81}{(*)}{睡眠与梦的时辰不计——那时甲壳会变轻:有时甚至会消失……}。

这副甲壳自我们幼年起便逐渐成形,年复一年,愈结愈厚。它最初的功能,无疑主要在于保护我们免遭亲近之人(往往出于好心)的侵凌,为我们从他们那里挣得一份多少还算善意的容忍。然而,这副甲壳不仅保护我们抵御外部世界——它还有另一层功能,或许更深、也更根本:使我们隔离,保护我们抵御\textbf{我们自己}——抵御我们身上那份知识与力量,它们被我们周围立法定规的那些无声共识宣布为“不可接受”、没有存在的理由。这在我们幼年便已如此,并年复一年愈发如此:一副\textbf{两面}的甲壳,一面朝“外”,一面朝“内”。它们保护着“我”、那位“东家”:一面抵御他惧怕的来自外部世界的侵凌(而且他一年比一年更胆怯!),另一面——也是\textbf{最要紧}的一面——抵御“工匠”那些扰人而不见容的奇想与乖张;说得直白些,抵御那个\textbf{野孩子}——难以逆料而又无所不能,纵然隔着三层厚厚的角质,耐火又防水、保证牢靠,他依然叫人心忧……

(11月2日)继“天真”一注(n\textsuperscript{o}~107)凸显出我最初几年被身边亲近之人接纳所起的作用之后,还有第二个时刻,“接纳”与“不接纳”处在反思的中心。那是在“接纳——阳中之阴”一注(n\textsuperscript{o}~110)里,我对自与孩童王“重逢”之日起在我身上发生的种种变化作了一次部分盘点。它们朝着一种逐步“回返”“童年状态”的方向。

\dmark{}这回返绝非“倒退”回某种先前状态——仿佛它有本事抹去我这旅人身上所留的我走过那条路的痕迹。\origpage{530}唯有凭\textbf{成熟}——一场内在工作的果实——我们才能重新接上那仿佛已经消失的天真,重新接上我们身上那个仿佛早已死去、早已下葬的孩子。而凡有成熟,无不同时也是些许的回返——回返孩子,回返质朴,回返孩子的天真。正因如此,一种活得完整的一生,宛如又一个正在“臻于圆满”的圆环;那是老境重新找回童年,是成熟重新找回天真——并或许终结于一场死亡,而这死亡准备着一次新生,一如冬天准备着一个新的春天……

在这份对一条尚未走完的回返之路的“盘点”里,显出:那“临终词”曾是\textbf{接纳},一如我的决裂之路、出发之路的临终词曾是\textbf{不接纳}、拒斥、拒绝。我的成熟不是别的,正是这样一个过程、这样一场内在的工作:藉着它,我逐步接纳了、迎纳了那些我长年以来拒绝、尽力清除、置之不理的东西。

这绝非一次“折返”,不是把一条已经走过的路再反着重走一遍;因而不是“倒退”——借用方才的说法。它毋宁像一个圆的上弧:延伸并接续那条已经画出的下弧,\textbf{诞生}于下弧,成了它滋养的根基、一次新跃的跳板……

(11月3日)昨天的注,收束在一幅意想不到的意象上——它不召自来,从反思中涌出。起初我带着几分迟疑接待它,生怕这意象旋即暗示出的那幅现实图景失于造作;生怕这意象“逼我就范”,逼我说出“牵强附会”的话。可一旦最后几行落笔,我在其上驻留片刻,便知道:我刚刚把手指按在了某一现实的意想不到而重要的侧面之上;这侧面也许我本已认识,却未充分消化——一个我素来倾向于忽略或遗忘的侧面。

\dmark{}多年来(118),我倾向于抬举朝着“接纳”方向去的东西,\origpage{531}而相反,主要以负面的眼光看待朝着“拒绝”方向去的东西。尽管这事情未必总是被清晰地表述出来,我却把这两类姿态——接纳与拒绝——感受为一对“对立”、一对“相反”:其中一个于我自己、于所有人都是“好的”,另一个则是“坏的”。

在这份未曾成形于言语的领会方式里,我始终是那套老生常谈的“二元论”世界观的囚徒(当然,我自己并无察觉)——那也就是我先前称之为“战争式”的世界观:它把一些事物当作对手对立起来,而一种更深的视域却向我们显明,它们是同一现实的\textbf{互补而不可分的}两面。当(10月25日,即十天之前)我开始眼下这场关于接纳与拒绝的反思时,我刚意识到:它们正是这一个月来——自这场关于阴阳的“枝蔓”开端以来——不断谈起的那些著名的阴阳“对子”、或“宇宙”对子之中一对里的妻与夫。因此我预计,反思将转向事物的这个侧面。两天来,它看上去似乎正在远离这个侧面。可昨日反思结尾的那几行——同一个循环的两道弧相互延续的那幅意象——出人意料地把我带回了那最初的直觉:它一直未曾形诸言语。

我倾向于以(即便不是全然)主要\textbf{负面}的眼光,看待从我八岁到四十八岁支配了我一生的那些\textbf{拒绝}:视之为一副有时令人不堪重负的\textbf{重担},我在我生命的四十年里一直拖着它,而在逝去的八年里,我终于甩掉了它(毋宁说,\textbf{开始}甩掉它)。这“重担”的那一面,是在禅观的发现之后、在与我身上的“孩子”“重逢”之后,才开始向我显露的。恰是在那一刻,我开始发现我生活中拒绝的进程——它表现为一种“超级阳”的循规蹈矩。事物的这个侧面绝非虚构。在先前只有一片“空白”、一片全然虚空之处知觉到它,是这八年里持续进行的成熟所结的果实之一。但与此同时,同一现实还有另一个侧面,同样真实、同样重要:它作为“\textbf{强大的行动原则}”的“正面”。这个侧面头一次出现(而且十分含蓄),是在10月5日那场禅观“阳葬阴——或超级父亲”(n\textsuperscript{o}~108)里,当时我写道:

\begin{quote}
\dmark{}“‘我将像他们一样’(而不是‘像我自己一样’)也意味着:我要把宝押在‘头脑’上——它在我身上并不比在任何人身上更差——\origpage{532}用他们自己的武器打败他们!”
\end{quote}

正是这股动力,宛如我1945至1969年间在数学中那场无度投入的活力——那股在四分之一世纪里滋养了一场发现之冲劲的力量\fnote{82}{(*)}{更确切地说,是这股冲劲中利己的成分,是这一“活力”的\textbf{利己}“因子”。}。无论人们选择以“正面”还是“负面”的眼光看待这样一场投入,清楚的是:那里确确实实有过\textbf{冲劲}、有过强烈的\textbf{行动}。在向生活学习的一侧,有着那副“有时令人不堪重负的重担”,从未被检视过——且不必说全然的停滞;而这同一副“重担”,恰恰在同一个时刻,却又滋养着一股认识的冲劲,赋予它活力。

自1970年我“出发”以来,我倾向于贬低、有时甚至否认这样一股冲劲——朝着对所谓外部世界作“科学”发现与理解的方向——所应得的“价值”。在《收获与播种》里,我已有好几次尝试去圈定这样一种发现与自我发现之间的共同侧面,以及反过来说,它们又在何处不同\fnote{83}{(**)}{尤其参见“欲望与禅观”“禁果”“孤独的历险”诸节(n\textsuperscript{o}~36、46、47)。}。这样说无疑是有根据的:朝着科学方向(哪怕是生物学,或“心理学”……)的发现冲劲,使我们远离我们自己、远离对自我的领会。当这样一份领会的角色被充分理解时,人们或许因此会倾向于在科学发现的冲劲中(以及在任何使我们“远离我们自己”的冲劲中)看出一种“恶”,或至少,看出成熟的一个“障碍”——并由此,看出我们自身充分绽放的一个障碍。(至少在我长时期所处的那种情形里:这股冲劲动员了绝大部分、乃至全部的心理能量。)然而,同样真实的是:我们所经历的\textbf{一切},都是我们学习生活、学习自我的原料。\textbf{材料}就在这里,任由我们让它转化为认识——只要容许一场成熟的工作在我们身上发端并继续。这也是为什么,我对自己经历过的一切并不懊悔,最终看到“其中一切都是好的,没有什么可抛弃”;也包括那些漫长的灵性停滞时期的荒漠——那是我为一场吞噬性的热情作无度投入所付的代价,我付出时不曾吝惜(而且闭着眼睛……)\dmark{}\origpage{533}。如今我看到,连这些荒漠也曾有东西要教给我,或许唯有它们才能教给我。这份代价我本无法省去——顶多,或许早在几年之后,我就本可以启动这循环的“第二道弧”——它的期限,我曾拖延了几十年。

也是在这一天,同样显出:对我自己、对他人的接纳——它诞生并成长于我成年的岁月——曾“滋养”于那些拒绝,它们标示了我一生最长的那段时光——即昨天唤起的那道循环的“下弧”,及其“滋养的根基”。诚然,在我生命的最初六年,我身上确有一种对自我的全然接纳,它无需任何先行的“拒绝”便能存在、舒展、确立。恰恰相反,它的绽放之所以可能,\textbf{正}因为它不曾被对抗、不曾被某一把拒绝的剪刀裁剪。但我童年身上的这份“接纳”,与我成年时的那份,不是“\textbf{同一份}”。它欠缺一个维度——那份仅凭我童年周围之人对我这个人的接纳所无法给予的维度。那是一份对拒绝的认识——对被他人、或被我自己拒斥(或我的某一部分被拒斥)的认识。这份认识,是经由拒绝的经验来到我这里的,也是经由轻蔑的经验——它是拒绝的众多面孔之一。

或许有些人生来就带着对拒绝的一种认识、一种领会,使他们能够在童年所暴露的种种拒绝面前依然保持\textbf{一}、依然天真而有知。我清楚地知道,我不是这样的人。他人的、和我自己的拒绝与轻蔑的经验,作为一片土壤,孕育出对拒绝与轻蔑的领会(无论多么不完善)——这份经验我无法省去。

\subsubsection*{(c) 伴侣——或“恶”之谜}
\tocsubsub{(c) 伴侣——或“恶”之谜}
\label{ch18:local:h24}

\notehead{117}{}
\phantomsection\label{ch18:local:n21}

我刚刚探测了我自己生活里拒绝与接纳之间关系的一个意想不到的侧面——它在昨天的反思中不期而至。不过,这里所说的“拒绝”并不是完整意义上的拒绝;我是说,一份被完全承担下来的拒绝——还差得远。这份拒绝也曾是一场对被拒之物的\textbf{长久逃遁}。它在于\textbf{不去看它}、\textbf{把它忽视},并由此在某种限度内,使它从我有意识的领会之场消失,也\dmark{}从他人可见的场中消失。\origpage{534}它曾是一种失谐、失衡状态的缘由与发条——在这里,是一种“超级阳”的失衡:它标示了我的成年,而它的若干关键机制至今仍在起作用。所以,这“拒绝”在此绝非以对称、甚或阴阳互补的角色,显于方才所谈的那种“接纳”(对我自己与对他人的)对面。相反,后者落在一场认识我自己的工作之中,朝着恢复一份被扰乱的和谐行进。所以,那里的接纳,是“充分知其所以然”的接纳,是完整意义上的接纳——而绝不是与方才名为“拒绝”的那场逃遁方向相反的又一场逃遁。

然而,“拒绝”与“接纳”之间,有一层比方才所探测的更明显的关系。它显现在两者都被“按完整意义”来取的时候。那时,它们是同一份和谐、同一种被完全承担下来的态度的\textbf{同时}而互补的两面。(而方才谈的,则是一段路程、一种进程的\textbf{先后}两面:途经一种失衡、失谐的状态,走向一份更新的平衡。)在这个视域里,没有那种把拒绝排除在外、对拒绝闭门的“真”接纳。也没有那种不诞生自接纳、不是接纳之一份可触的显现的“真”拒绝;没有那种不是同一件不可分之物的两“面”之一——阳的那一“面”——的拒绝:那物本有两面,其“阴”的或“母”的一面,便是接纳\fnote{84}{(*)}{有趣的是可注意到:接纳-拒绝这对偶子里,阴阳角色的这种“自然”分配(在法文里,由这对偶子两项一阴一阳的性属来表达)——在前一天反思结尾自发浮现在我心中的那幅意象里,却\textbf{颠倒}了。会有这样的互倒,没什么可奇怪的——正如在一对情人与情郎之中,其爱的关系并非凝定,爱的游戏里必不缺少角色翻转的时刻:好让活在情人身上的“阳”性爱欲冲动、与活在情郎身上的“阴”性爱欲冲动,畅行无碍。何况,我在“接纳(阳中之阴)”一注(n\textsuperscript{o}~110,该注第一部分末段)里,正谈到这种偶发的角色互换之重要。}。

一份把拒绝排除在外的“接纳”,不是接纳,而是一份讨好(对他人,或对自己,或兼而有之);或者——当涉及对他人的“接纳”时——一份串通、一份默契。全然接纳一个存在,无论这存在是自己还是他人,绝不意味着对其一言一行、种种习惯与癖好,作\dmark{}无条件的赞许。\origpage{535}这样一种无条件赞许,本身便是一场\textbf{逃遁},是对一份(往往雄辩的)现实拒绝去认识,而绝不是接纳。它非但造不出一个有利于更新、有利于与一份被遗忘的太一重新接上联系的“力场”,反而强化一种惰性,助长把人滞留在车辙里。

一个不同时是一次敞开的拒绝——它不同时也像一只向他人伸出的手(或“一根撑竿”),或像一个在自己与自己的关系中标出断裂与更新之点的跃起——这样的“拒绝”,真正是一场切断:它既“切断”又隔离拒绝者与被拒者。这仍是逃遁——逃开一份令人不快、甚或扰人的现实,一份沉沉压在我们安稳生活、我们种种便利之上的威胁——对这份现实,我们自信一刀便能躲开:“这没有存在的理由……”然而,\textbf{它有}!而且我们那道命令式的“拒绝”,丝毫不能拦阻事物如其所是,哪怕它们惹我们不快。恰恰相反,一如一场自动赞许的讨好,这样的拒绝强化着抵制创造之更新的种种惰性;它宛如一道\textbf{判决}:你是不可接受的,而且永远如此……

我并不自许在我这个人身上实现那份被完全承担下来的接纳与拒绝之和谐。恰恰相反,我知道远非如此——而我不敢肯定曾遇见过一位实现了这份和谐的存在。实现它,也就等于在你自己这个人身上,解开了“恶”的大谜:不义、谎言、恶毒、卑怯、轻蔑——以及那些被击中而喑哑无声的人们的苦难。这肯定也等于:充分领悟那份“善”——它就在那内在的跃起如此经常地给我们指为“恶”的东西之中。

拒绝战争,同时看到并接纳战争无处不在、在所有人身上;看到连我所爱的那些人也在他们身上携着它、传布着它,一如我自己也曾重拾它、携带它、传布它、传递它。拒绝战争,同时接纳它的存在,同时爱它那数不清的盲目的士兵。确定无疑地,唯有这个,别无其他,才也意味着:走出战争,走出冲突——停止传布战争。

\subsubsection*{(d) 阳扮阴——或主人的角色}
\tocsubsub{(d) 阳扮阴——或主人的角色}
\label{ch18:local:h25}

\notehead{118}{}
\phantomsection\label{ch18:local:n22}

\dmark{}(11月4日)\fnote{85}{(*)}{本注由“循环”一注(n\textsuperscript{o}~116)的一条脚注生出。参见11月3日诸注开头处的参见号。}这份“倾向”的出现\fnote{86}{(**)}{即抬举“接纳”、以与“拒绝”相对的那份倾向。},是在七十年代初——也就是说,在我从数学舞台“出发”之后的那些年头。\origpage{536}在一个与从前大不相同的氛围与朋友圈的影响下,我当时所信奉的那一整套“价值”,发生了一个严苛的转弯。带着回顾,我可以把这个转弯描述为:从一个“超级阳”的或“父权”的价值系统,过渡到一个几乎相反的、“阴”占强势主导的系统——一个“母权”系统。在这场颠倒中起过作用的影响里,还有对Krishnamurti的几次零星阅读——关于这一点,参见“Krishnamurti——或变成羁绊的解脱”一注(n\textsuperscript{o}~41)。

我当时之所以放任这些影响起作用——它们把我引向这样一场“意识形态”的转弯——无疑是因为(尽管我当时并未察觉)在我身上有一种深刻而紧迫的更新之需;而首先,是一种解脱的需要:从那份积习已深的“超级阳”姿态的重负下解脱出来。这同一份需要,无疑在1969年已经起过作用:那一年,正当数学活动强烈而多产之时,我忽然从数学上“脱钩”,转而关心生物学\fnote{87}{(***)}{我起初是从“分子生物学”这一头对它发生兴趣的,是受我的生物学家朋友Mircea~Dumitrescu的影响——是他把我领进这个迷人的世界。};继而在次年,离开(不存回头之念的)数学舞台,乃至科学研究。当时发生了一场突然而严苛的氛围与活动的改变,对此我在《愚妄与更新》(《收获与播种》的第一部分)里已屡有机会提及。

然而,把这些环境、活动乃至“价值”的引人瞩目的改变,看作一场“更新”、一场“解脱”,会是不确切的,或只部分为真。关于这一点,在“与Claude~Chevalley相遇——或自由与好心情”一节(n\textsuperscript{o}~11)里,我已表白得相当清楚。在眼下这场关于阴阳的反思的更透彻的光照下,我可以说:那个无疑显得在一切之中最有意义的改变——阳的价值被清除出去(甚至还没在我自己身上被辨认、更未被检视),让位给阴的价值——却丝毫没有改变“我”的(超级阳的)结构,至多稍稍缓和了由之而来的种种姿态与行为。诚然,我对外部世界的领会\dmark{}已发生相当大的转变,朝着一种突然拓宽的方向;但这一转变依然是零散的,几乎排他地限于智性的层面、“抉择”的层面。\origpage{537}事情不可能别样,只要这场转变恰恰只限于我的“外部世界”之景——在那景中,我自己的这个人并不出现,或只是偶然地、浮面地出现,主要经由我的“社会角色”及其种种含混与矛盾。与过去一样,我当时毫无些许的疑心:\textbf{我自己}身上也会有含混与矛盾!恰恰相反,我被一种不可动摇的确信驱动着:\textbf{我}这个人,是无矛盾的(尽管我当时已开始在四周处处辨认出他人身上的矛盾);尤其确信:在我有意识的欲望与有意识的事物之知、与我的无意识之间(倘若我的情形里竟有一个无意识的话——倘若它不是我意识的一件如实复制品的话……),有着完美的契合。

这份确信的第一道裂缝,直到1974年春才出现:我终于懂得,某种东西在我\textbf{自己}身上也必定出了毛病,而不只在别人身上——它是我与所有亲近之人的关系那场无情恶化的缘由(整个成年期间,我的生活似乎就缩小到了这关系之上)。当时,由于缺少一份对我自己的真正的\textbf{好奇}——它本会以钻进去为乐,看看背后有什么,并由此看一座沉重的建筑坍塌下来:一座由从未检视过的、天花乱坠的幻觉筑成的建筑——这道益人的裂缝的效应,便仍然有限。

这份对一种自然好奇的顽固阻塞,无疑主要缘于:我从未在他人身上遇见过这样的好奇——它本可让我起疑:在生活中,一如在数学里,凡有一个问题呈现,就有可看的东西,而看着,便能学到许多意想不到、极其有用的东西——换句话说:世上有\textbf{自我发现}这样一件事。我当时读过一些Krishnamurti,并得以确认:他说的某些东西是真的、深的、要紧的。于是我倾向于对他全盘照单收下。除枝节而外,我已在事实上默默采纳了Krishnamurti的世界观\fnote{88}{(*)}{(11月5日)这样一份“采纳”——它变成一种文化行囊——在我生活中的效应,始终极其有限。我的注意力被引向现实的某些此前完全逃过我知觉的侧面,但并未因此启动一场具有更新之力、深入的清理与消化的工作。倘若说,在1970与1976年之间(在我的数学舞台“出发”与禅观的发现之间),Krishnamurti在我的行程里曾是重要的,那与其说是因为我向他借来的那份“行囊”,远不如说是因为他已变成(当然,是在我不自知的情况下)一个默默的\textbf{模型}——我遵从着它,又不肯露出遵从的样子:总而言之,“非古鲁的古鲁”的模型,那位自称不是大师的大师。}。\dmark{}在我所说的那个时刻,这份行囊确确实实起了一场“羁绊”的作用——羁绊着一场真正的解脱、一场完整意义上的更新。\origpage{538}关于这一点,我恰好在那条已援引的注里作了说明(我刚刚重读过),我在其中力图圈定(Krishnamurti的)“教诲”在我自己的行程里所起的角色。

完整意义上的第一次“觉醒”,只在两年半以后才到来,随着禅观的发现。那也是对自我发现的发现:\textbf{有一个未知之物},它就是“我”,而我有力能穿透这物、认识它。这一关键的发现,是在一切教诲(无论有无大写)都已被遗忘的时刻完成的。那也正是“建筑”头一次坍塌的时刻——它由种种现成观念与形形色色的“教诲”筑成,靠一份巨大的惰性维持着;也是一份主动的好奇头一次出现的时刻——它往往顽皮,而始终满怀善意。

正是在这个转折之后,随着一份好奇在我身上绽放——首先朝向我自己的这个人,继而是“生活”,作为自然的果实——我才得以用一双新的眼睛,同时看Krishnamurti与他的讯息。带着回顾,我得以赏识那份讯息的丰富,同时辨认出它的限度与匮乏,以及那位\textbf{主人}(“the Teacher”,他的信徒与追随者如此称他)身上某些根本的矛盾。这些匮乏与矛盾中最重的一个,在我看来,正是我方才又一次擦过的那个:主人自身之中全然没有好奇。他的著述里没有任何东西让人起疑:在遥远的岁月里,这份景是否曾\textbf{诞生}于一个\textbf{人}——一个如你我一样陷在现成观念之网与从未被辨认的矛盾之中的人;这份景是否曾在一场强烈的、有时痛苦的、逆着巨大惰性之流的\textbf{工作}进程中,\textbf{从错误里澄清}出来;这项工作的诸阶段,或这些劳作中跨过的“门槛”,是否是一场场意想不到的\textbf{发现},每一场都撼动一整套积习已深的、\dmark{}由模仿与重复的\origpage{539}普遍机制所延续的观念\fnote{89}{(*)}{(11月5日)这些机制,分明是心灵的基本机制的一部分,在人与动物身上皆然。它们先于一切制约、一切学习(如幼童的语言学习,以及日常生活里几乎全部行为的学习)而在——没有它们,那些学习根本无从确立、无从展开。而在那位年轻的未来主人身上,它们并不比在任何人身上少在、少奏效。}。

所有这些事,那个孩子有一天都曾懂得,甚至曾认识——因为曾强烈地经历过。但主人把它们忘了,而且无意记起。他不愿做那个激情地\textbf{发现}、\textbf{学习}、并在发现中被转化的孩子;他要做一个一成不变的\textbf{主人}——凭天赋之知而不移地\textbf{知}着,并把一生献给传播他的\textbf{教诲},以造福凡俗大众。

他变成了他的追随者与信徒、那些信他的人要他成为的那一位:一份静态的、重复的、因而是抚慰人心的讯息之化身,一种新意识形态的使徒。总而言之,一位\textbf{非古鲁的古鲁}——一如我自己当年也曾是(或许正是仿效他的榜样\fnote{90}{(**)}{(11月5日)说真的,这个“或许”所带的存疑口气,根本不合宜!关于这一点,参见今天写下的倒数第二条脚注。})……

(11月15日)我把前一条注(11月4日)命名为“阳扮阴——或主人”。正如一场对我自己的禅观所应有的做法,这条注的主名关涉我自己的这个人:它指的是某一场“戏”——自1970年我离开科学世界之后\fnote{91}{(***)}{此外,1976年10月禅观发现的时刻,也标志着这场戏的陡然衰落;它好歹继续着,只是调门更低调,直到1981年——那一年,它终于被察觉、被拆除引信。关于这一点,参见已援引的“非古鲁的古鲁——或三足马”一节(n\textsuperscript{o}~45)。},我曾扮演了几年。至于第二个名字“主人”,它可以无差别地被解读为:或指涉我自己的这个人——指我在“阳扮阴”这场戏里所扮的角色、所摆的姿态;或指Krishnamurti——他曾默默充当我的模型。

实际上,从Krishnamurti的书里透出的价值,几乎是排他地阴的价值。在我头一次读Krishnamurti时(1970或1971年),我是头一次看到这样的价值被摆到前头来,并看到我自己那阳性的世界观(以及“人人”的世界观,出入无几)的限度与裂隙被透彻地圈定。这无疑就是那几章的阅读在我身上留下极深印象的缘由。六七年之后,我又有机会读到吕滕斯夫人写的那部Krishnamurti的漂亮传记。它印证了从他的书里已然透出的对他这个人的某种印象(尽管他在书里从不亲自出现)。今天我会这样表达:他气质里的基调,是\dmark{}强阴的\origpage{540}。

还须补上:透过他的全部著述,人们看到一条恒常的主导动机:把带阴色的品质、姿态与价值摆到前头,而(明言地或以沉默)贬低阳调的品质、姿态与价值。

所以,Krishnamurti的一生与教诲,实现了“\textbf{阴葬阳}”这一相当例外的心态——它与迄今为止最常见的那种心态、“阳葬阴”的心态正好逆行;而我自己的生活(至少直到我四十八岁)恰是后者的一份同样极端的例证。Krishnamurti的“超级阴”抉择\fnote{92}{(*)}{这些“抉择”无疑上溯到他的童年,更确切地说,上溯到他与他那些神智学监护人的最初接触。}有一个大功绩:逆着周边文化的基本价值而行。但这并不妨碍它们在我看来同样具有压制性(以他这个人的一部分压另一部分),一如我曾经的那些抉择。

然而,Krishnamurti的一生里,有一个非常触目而惹眼的“阳”的侧面——它无疑首先是被强加给他的:是他那些显赫的神智学监护人在他尚是孩童时,就为他定下的那号(未来的)“精神主人”、旗舰人物的角色。其后,在他生活的大转折——以那些翻天覆地、彻底撼动其世界之景的发现为标志(这些发现在其后成了“教诲”)——之后,这个“主人”或“引路人”的角色,似乎被他整个地内化了,被他自己接了过来,连同一份他个人(而非承自其神智学老师)的学说之传布。这份传布代表一种强烈乃至耗竭人的活动。它看上去很难朝着阴阳的一份\textbf{平衡}去,毋宁显为一份\textbf{强制}:由一个“我”强加给一种彻头彻尾禅观性的气质——这“我”在主人身上,与其在任何人身上一样强、一样具有侵夺性。在这个光照下看,眼下这条主要谈Krishnamurti的注“阳扮阴”,同样可以叫作“\textbf{阴扮阳}”。

这样,在我的一生里,有两次、以两种不同的方式,我玩过一些“戏”——它们恰似一种\textbf{颠倒}:颠倒那支配了某个人一生的种种姿态——那个人,在我行程的某段时期,本该成为我(同样默默的)形象招牌的默默模型;也颠倒我身上与之相应的某些姿态与架势。但透过这两种彼此相反的表达风格,我今天认出一种显然的亲缘。其一,在于一份\textbf{压制}的存在(无意识的,不言而喻),它\dmark{}生成阴阳自然平衡的\origpage{541}破裂\fnote{93}{(*)}{在这一份亲缘里,我们想必不乏人数众多的同道!}。其二,在于对一个\textbf{角色}的选择,在于\textbf{这角色之重}——它的刹车效应,乃至阻塞效应,落在一场绽放、一场成熟、一份领会或认识的进程之上。这个角色(或这份架势),在我身上与在我那位模型身上,是同一个——我也许不过是把它原封借来。那就是\textbf{主人}的角色。

\subsection{数学的阴与阳}
\label{ch18:local:h26}

\subsubsection*{(a) 最“大男子气”的艺术}
\tocsubsub{(a) 最“大男子气”的艺术}
\label{ch18:local:h27}

\notehead{119}{}
\phantomsection\label{ch18:local:n23}

(11月5日)想谈谈数学里的阴与阳,已有一阵子了。一项数学工作之中、或一种通向数学的进路之中的阴阳两面,只是在这几周关于阴阳的反思进程中,才向我显现。我预计:在这些注里对这个双面稍作探测,将是“言归正传”的最自然的方式——这些注,本应对“一个数学家的往昔”构成一场回顾。

从一开始(五年前)我关于阴阳的最初反思起,对我就已十分清楚的是:“做数学”,或许是已知的一切人类活动之中\textbf{最阳}、最“男性”的一种。老实说,一切全然智性的活动——尤其科学研究,以及更一般地,一切通常被称作“研究”的活动——都是一种阳占绝对优势的活动。我本要写:“带有强烈的阳失衡”,而事实上,当这种活动竟至吸收一个人的几乎全部能量时,情形正是如此。这种阳的优势(或失衡),经由列举许多阴阳对子而显现:对这些对子,分明是阳的那一项——若不说排他地——才“在”一项智性工作之中“在场”。我只限于举出其中几个;它们全都属于同一个“组”(或同一扇“通向世界的门”),我称之为“模糊——精确”组。(按:在这个对子以及接下来的那些对子里,是阴的一项放在前头。)

\begin{itemize}
\setlength{\itemsep}{0pt}
\setlength{\parsep}{0pt}
\item 感性——理性(或智性)
\item 本能——反思
\item 直觉——逻辑
\item 灵感——方法
\item \dmark{}视野——融贯\origpage{542}
\item 具体——抽象
\item 复杂——简单
\item 模糊——精确
\item 梦——现实
\item 未定——已定
\item 未言——已言
\item 无形——有形
\item 无限——有限
\item 无界——有界
\item 全(整体)——部分
\item 全局——局部(或细部)。
\end{itemize}

我刚把我的阴阳清单巡了一遍,又捡出一大批对子,它们都让人尝出纯粹智性活动那份超级阳的品格。只再举方才已想到的那些之中居首的一个: \textbf{身——心}对子。

有鉴于此,我觉得,在各种智性活动之中,代表着终极极阳的,是数学工作。这无疑首先缘于它极端抽象的品格:它在极大的限度内,不依赖任何借助感官经验、借助对外部世界——即我们生活于其中、我们身体活动于其中的那个世界——的理性观察的“支撑”。这种抽象中的极端品格,把数学同一切其他科学区分开来,把数学工作同一切其他智性工作区分开来,使它成为一门“纯理性”的科学、一份“纯理性”的工作。与实验科学和观测科学相对,它也是唯一一门其结果经由最严格意义上的\textbf{证明}来确立的科学:证明循着一条严格编码、原则上不出错的\textbf{方法}——所谓“\textbf{逻辑}”方法——行进,达到不容任何怀疑或保留的\textbf{确定性},不容许有逃过迄今所察案例的例外。这些都是在数学工作中、且唯独在数学工作中汇齐的极阳特征。

当然,这些特征自童年起就自然而然地吸引着我——我这个从头到尾选定了“头脑”、选定了极阳的人!\fnote{94}{(*)}{军事与战争的那个变体且不算:检阅、制服、挺着胸膛的立正,以及组织得无懈可击的屠杀与尸坑……}尤其在经历了战争与集中营的考验、直面种种似乎向最起码的理性也发起挑战的歧视与偏见之后,数学活动中(经由我中学年头所能认识的那么一点\dmark{})最让我着迷的,是它所给予的那份\textbf{权能}:凭一个简单的证明,赢取哪怕最不情愿的赞同;总而言之,\textbf{强使}他人首肯,无论他乐意与否——只须他同我一道接受数学的“游戏规则”。这些规则,早自我1940年在芒德中学最初接触学校数学之时(我得以去那里上学,尽管被拘禁在离那儿五六公里的里克罗营地),我就仿佛本来就认识它们、本能地感到它们,好像我从来就认识它们似的\fnote{95}{(*)}{这些最初的接触,就发生在我童年的“化圆为方”思考之后不久——见注~n\textsuperscript{o}~69。}。\origpage{543}当然,我比老师本人把它们感到更好:他毫无信念地背书给我们当时通行的那些老生常谈——“公设”(亦即欧几里得公设,唯一一条他同我们都有幸听说过……)与“公理”的分别,或三角形全等“三情形”的“\textbf{证明}”——一路照着教科书念,活像一个初领圣体的学生跟着念他的日课经。

五年之后,我被原子物理骤然显赫的名望所吸引,却先在蒙彼利埃大学注册学习物理,心里想着去入门物质结构与能量性质的奥秘。但我很快就明白:倘若我想入门奥秘,靠听课是不行的,只能靠自己的手段、孤独地去工作,有书没书都一样。既然我既没有那份天分、也没有那套装备,以这种方式学物理;我便把这事推到更合适的时日。于是我着手做数学,同时“远远地”听着几门课——没有一门能让我满足,或能给我超出通行教科书的东西。但考试反正总得考……

\subsubsection*{(b) 陌生的美人}
\tocsubsub{(b) 陌生的美人}
\label{ch18:local:h28}

\notehead{120}{}
\phantomsection\label{ch18:local:n24}

(11月6日)刚才重新翻阅昨天的注,我得以确认:我留了神,不曾重新落入某种混淆——把数学\textbf{工作}(一种阳占绝对优势的活动)同“\textbf{那}数学”混为一谈。法语如同德语,指称它的那个词是阴性的——这肯定不是偶然:一如涵盖着它的“\textbf{那}科学\dmark{}”,或更宽的词“\textbf{认识}”\origpage{544}\fnote{96}{(*)}{相反,“\textbf{知}”是阳性的,而在“认识——知”这对阴阳偶子里,“知”恰是“夫”。德语在此不那么分明,因为~Kennen~与~Wissen~两词(作为名词化的动词)都是\textbf{中性}的。},或“\textbf{实体}”。对本来意义上的数学家——我是说那个“做数学”的人(如同他“做爱”)——他与数学的关系之中角色的分配,毫无歧义可言:数学,即那未知的实体,他与之结交的实体,他穿透它而认识它。那时的数学,同他认识过或仅仅渴望过的任何女人一样是“女人”——他曾感到过那神秘的权能,那把他拉向她的吸引,带着那既极温柔、又无可回驳的力量。

我头一次察觉:把我拉向“女人”的冲动与把我拉向“数学”的冲动之间,有着深刻的同一——那是在《道德经》的几节诗章触发我写《乱伦颂》的几个月之前(顺带,也触发了我对“阴柔”与“阳刚”的第一场系统反思——它们的中国名字“阴”与“阳”,我当时还不知道)。那是六年之前,我在写一篇两页的短文,题目叫“以代纲领”,言外之意:为(C4的)“研究入门”课程而作——这篇短文构成那门“课程”的引言,或更确切地说,一场关于这门“课程”之精神的宣言。写完这篇笔底最自然不过流出的文字之后,我惊于其中意象之繁富——它们彼此相生,满载着爱欲的意味。我分明意识到:这既非偶然,也非单纯的文学故意的产物——那是支配了我成年生活的两种激情之间深刻亲缘的一个不含糊的信号。当时,我虽未想着以一场系统反思去深究此事(它只在几个月后,借着写《乱伦颂》才出现),我相信甚至也不曾把那骤然知觉的东西对自己清晰地表述;但我想我可以说:在那一刻,我悄无声息地学到了某样要紧的东西——我“发现”了某样东西\fnote{97}{(**)}{那是一场“阴”式的、“女性”式的“发现”——它经由我们对一份新知的迎纳而完成,怀着对来到我们心中之物默默敞开的姿态。这样的时刻,我相信,在我一生中是稀少的。无论如何,凡我留有记忆的发现时刻,几乎全带阳的、“男性”的调门。},某样此前整个逃过我知觉的东西。

当然,同大家一样,我听说过弗洛伊德、力比多的升华等等,但这毫不相干。纵然成吨的精神分析书籍、随你什么书,也替代不了这样的时刻:那时,一切理论、一切“行囊”都被遗忘,而忽然有什么\dmark{}“咔哒!”一声对上了\origpage{545}。正是在这样的时刻,我们对事物的认识得以更新。这与读书、听报告毫不相干——也就是说,与增益一份\textbf{知}\fnote{98}{(*)}{这份确认,并不被下面这件事驳倒:很可能,甚至大概,这次“意识到”(即一件在无意识中被知觉的东西上升到意识层面)曾得益于弗洛伊德共识的存在——我听说过它,但它不曾真的令我或热或冷。一份知可以助一份认识之绽放;但在我看来,更常见的情形,倒是它把一切绽放的企图扼杀在萌芽里——一如那些现成的“答案”,把一个(好)问题的绽放扼杀在萌芽里……
令人惊讶的是:尽管“人人都听说过”性爱冲动在(且说是艺术的或科学的)创造性中所起的作用,但在那些我一度置身其中的圈子的共识里,却看不到一丝痕迹。然而,惹眼的事实并不缺,它们本可早就让我有所警觉。比如,直到三年前,我一生中创造强烈的时期,尤其内在更新的时期,也都以一股强大的性能量之涌入为标志。不过,我的数学活动从来未曾伴随有意识的、性的意象或联想。但我记得,五十年代,在一次Bourbaki工作会议上,我曾有点错愕:一位同事兼朋友当着我的面,像谈世上最平常的事一样,谈起他数学工作中的一个特点:当他抵达一件艰难工作的尽头,他便感到一种强迫性的做爱的欲望(有无伴侣均可)——而且他对自己刚做的东西越满意,这欲望便越强。}。

当我想到“那数学”,我想到的肯定不是那可以被称作“数学”的全部\textbf{知}——从上古到今日,写进种种出版物、预印本、手稿与通信里的那份知。即便剔除重复,它想必也有几百万页排得密密麻麻的文字;也许十来吨的书,或几千册厚厚的卷帙,够装满一座宽敞的图书馆:这肯定不会让人血脉贲张——恰恰相反!说“那数学”,唯有在一幅\textbf{视野}、一份\textbf{领会}的语境里才多少有意思——而这些,在本质上是一些属于个人的事情,绝非集体的。有多少数学家,就有多少“数学”:其中每一位都对他那份数学有某种个人的经验,或广或狭;那经验的一个果实,就是他自己对“那数学”(即他所认识的那份)的领会、他自己的视野——它总归或多或少是零散的。这有点像“\textbf{那}女人”:在某些人眼里,她也许显得只是一个抽象、一个空洞的公式;然而她却有一种深刻的、强大的、无可辩驳的“实在”(至少对我而言),而所遇见或所认识的每一个女人\dmark{},都是她的一个化身、代表她的一个侧面\origpage{546};而同\textbf{一个}女人在另一个人的经验里,无疑又代表另一个化身、另一个侧面。

我在此的用意,绝不是要直面那份困难:把这一大堆互异的关于“那数学”的经验、领会、景“整合”为一个总体、一个统一体——何况,还正是在这样一个时代:依我看,数学成果的产出正经历着一场猛烈的“发散”,而恐怕没有一个数学家敢自夸,哪怕只在大轮廓上,认识我们这门科学里迄今完成的那份实质工作的全部或要害。我的用意倒是稍许检视阴阳在数学\textbf{工作}中的戏——也就是说,也在数学家(或某位数学家,从我自己开始)与“那数学”的关系中的戏。所以,被检视的东西是“数学家”或“某位数学家”(在其与数学的关系之中),而不是“那数学”本身。

\subsubsection*{(c) 欲与严格}
\tocsubsub{(c) 欲与严格}
\label{ch18:local:h29}

\notehead{121}{}
\phantomsection\label{ch18:local:n25}

(11月7日)在我们的智性官能、理性的层面上,“认识”一物,在一切之先,便是“\textbf{理解}”它。而在一场落在我们官能这个层面上的发现工作里,鼓舞着我们身上那个孩子的认识冲劲(独立于“我”、“东家”所特有的那些动机),便是\textbf{理解的欲望}。这或许正是智性的认识冲动区别于它姐姐——爱的冲动——的主要差别。这份理解的欲望,先于一切“方法”而在,科学的或其他。方法是一件工具,由欲望所塑,用来达成欲望的目的:穿透理性可及的未知,以达理解。认识诞生于认识的欲望——当想要认识的是理性时,便诞生于理解的欲望。\textbf{方法},欲望的器械,本身无力分娩任何认识——正如医生的产钳、甚至助产士那双娴熟的手,都不会分娩。但有时,当时机成熟、而它们懂得及时赶到,它们便能有益地助产那个新生儿……

中学生与大学生——即便不是全体,也必定是大半——应当都感受过数学里的那份\textbf{严格}:它由一些阴沉着脸的老师反复灌输给他们,显得像一种全然外在于他们谦卑个人的先天之物,不可理解、专断横行,由一位专断而\dmark{}不容情的上帝,颁给一位欧几里得——他被擢升为总大检查官\origpage{547},领的使命是叫无数世代的学生在课业面前面无人色,好歹把那大写C的文化整个吞下。我想必是少数几个例外之一:在与学校数学的关系中,我从未经历那一阶段——从第一次相遇起,便在一本六年级数学教科书那逼仄的框架里,本能地感到了严格的原始功能与意义:它是一件柔韧而效力惊人的工具,服务于对这些所谓“数学”之物的领会——那些理性凭一己便能全然认识的东西。这份“严格”,也像我前天反思中所谓的“数学游戏规则”、方才我所谓的“方法”的灵魂与神经。仅仅对它们惊鸿一瞥,便仿佛我从来就认识它们——仿佛是我\textbf{自己的}欲望,温柔地、爱恋地把它们塑成这样一把钥匙:它有力能使一个未知而神秘的世界向我开启;那世界被预感到的富饶,将显得永不枯竭……而恰恰是我自己的欲望,在我的中学与大学岁月里继续打磨这件工具——在遇见任何人能使我起疑:世上某处竟存在我的\textbf{同类}——那些如我一样,以探测未知为乐的人——之前;这把显然无人(包括我的老师们)知晓的钥匙,唯有它有力能把那扇门推开一道缝\fnote{99}{(*)}{然而,我在中学与大学学到的那么一点数学,本该也够让我明白:至少在过去,必定存在过像我一样的人——实际上,就是人们称作“数学家”的那些。而且,苏拉先生(我在大学的一位老师)还同我谈起勒贝格,说他解决了数学里最后一批未解决的问题,包括测度论里的(自1945年离开中学起,我就在做测度论)。但在那些年头(1945--48),我凭自己的手段澄清\textbf{我自己}向自己提出的那些问题的欲望,是那样的排他,以致排除了一切对过去或当下的数学家之存在、之著作、之其人的好奇。}。

\subsubsection*{(d) 上涨的海……}
\tocsubsub{(d) 上涨的海……}
\label{ch18:local:h30}

\notehead{122}{}
\phantomsection\label{ch18:local:n26}

(11月8日)三天以来,我的反思原则上一直“论数学中的阴与阳”,而我有种印象:它老也起不了步,因为我部分地被别的营生与事务吸住了。铺垫做了这么多,我一直没来到那个一开始就想来到的点上:在我自己的数学工作\dmark{}里,\origpage{548}占主导的,是\textbf{阴}的音符、“\textbf{女性}”的音符!

我是几周前察觉的,在这场阴阳反思的边边上,并联系着那场“由三部分葬礼悼词引发的联想”——它正是这条长枝蔓的出发点。(见“阳葬阴(1)——或肌肉与肚肠”一注的开头。)老实说,这场联想(我日后还有机会回到它)多少基于这样一个直觉:我的数学进路阳占强势。这个直觉相当自然,因为正是我的超级阳抉择,促成了我在数学中那份旷日持久的投入。但这并不妨碍这个直觉——或更确切地说,这个观念——是错的;只须容我自己稍事检视它,便会发觉真实情形恰恰相反。要说出乎意料,这就是:一场意外!我当时没有“趁热”在注里谈到它,为的是不打断当时反思的线索——那时我正试着圈定我感知阴阳的方式、以及由之透出的那份哲学。可如今,我们总算来到这里了!

这个关于我的数学进路之品格的错误观念,想必是在五六年前——我开始留意事物阴阳两面的那时——未经检视、仿佛天经地义地溜进我身上的。它想必是我那份阳性的、男性的形象招牌的一个残迹——残迹因纯粹的惰性而一直滞留在那儿,怪我只懒得起身,往那个角落里扫上一帚……

读者或许会觉得我在哄他:就在三天前,我还长篇大论地说,数学工作是一切超级阳活动中最超级阳的——在与数学的关系中,数学扮“女人”,数学家扮殷勤进取的情郎——可转眼之间,我竟提出这样的问题:就我这微末个人而言,我的工作或我的“进路”究竟是阴还是阳,并且(像世上最自然不过的事一样)断定:是阴!谁想得到呢!倘若这里有一副表面上的混淆,那缘于对下面这条普遍事实的误解:在一切事物之中——哪怕它是世上至阴或至阳的——都上演着阴与阳的动力学,借着两股原初力量的联姻。比如火,万物中最阳者、阳本身的象征,在其某些侧面却是阴的(此即“阳中之阴”)\dmark{};反过来,水,阴本身的象征,在其某些侧面与功能中却是阳的(此即“阴中之阳”)\origpage{549}。无须在此展开这两个格外有教益的例子——凡被这些论断(它们或许显得武断,或许像谜语)勾起好奇的读者,只须亲自追随系于火与水的那些联想,便能在这两个个案里,自己发现阳中之阴、阴中之阳的实在。而如果他是个数学家,或哪怕只是个对智性工作不陌生的人(纵使他既非数学家、也非科学家),他也将毫不费力地辨认出:面对任何一种智性工作——无论与其他较零碎的活动类型相比,它显得多么“阳”——都存在着互补的阴与阳的进路模式。

一个可能的出发点,是重拾三天前那场反思开头标出的那十五来个阴阳对子\fnote{100}{(*)}{见“最‘大男子气’的艺术”一注(n\textsuperscript{o}~119)。},当时我曾确认:拿这样的智性工作与其他活动类型相比——如做爱、唱歌、画画(画一幅画,或者,画一面墙也行)、侍弄园子等等——对其中每一个对子,都是阳项占优势(而这在数学工作那里尤其如此)。但这并不妨碍:倘若我们留在一种确定的活动之内——比如说做数学(其全然属阳,自不待言)——我们仍能分辨出一个(或有时是一种失)平衡:或阴或阳的特征,因数学家而异;有时在同一位数学家那里,也因工作而异。比如,在某些工作里,被摆到前头的是所展开理论的\textbf{逻辑}结构;在另一些里,则是\textbf{直觉}的侧面。有一种失衡,它显现在读者或听者身上,是那份大家熟得不能再熟的\textbf{不适}之感(有时作者身上也有):当两个必不可少的侧面之一被粗暴地忽视,以便让另一个独占“好处”。(而当两者都被粗暴忽视时,人们就把书扔进垃圾堆,或摔门离开教室!)当两个侧面都强有力地在场——或明言,或藏在字里行间——那便显现为另一种同样熟得不能再熟之感:和谐、美、平衡、满足之感。事情便是这样,与那支配着所循进路的“基调\dmark{}”无关——无论那基调朝“逻辑”的方向,还是朝“直觉”(或同样朝“结构”)或“实质”的方向\origpage{550}。无疑无须展开这个有教益的例子,去描述——比如说——当两个侧面之一被忽视时,症结何在(即圈定方才唤起的那种“不适”);读者凭自己的经验,对此早已了然!同样方向的论断,必将从三天前考虑过的大多数阴阳对子里透出。甚至或许对所有的对子都如此——尽管其中某些更为微妙,无疑需要比直觉——逻辑那一对更深入的检视,方能被充分把握。

现在,我该试着稍稍显豁这个事实,或者毋宁说“把它传递过去”:在我做数学的方式里,领舞的,是我那些阴的、“女性”的特征,而不是我那些“男性”的特征。倘若要做的是把这份印象追到底,在尽可能多的侧面下检验它,那个自然的想法(昨天它确实在我心头掠过)便是:在我所知的阴阳对子之中,检阅那些能代表一种智性工作的某一侧面或把握模式的对子(我猜想总有五十来个),并逐一查看,这两个“配偶”之中哪一个在我身上占优势。我预计:在所有情形里,经检视,两者之中确有一个将显出优势。

这样,在直觉——逻辑那一对里,我一眼便看出:两个侧面在我的数学工作里都强有力地在场。这便是平衡、和谐的标志之一,与其他同方向的标志并列。一如阴阳对子所应有的:对我而言(我是说在我的工作里),两个配偶真正不可分——一个理论的逻辑结构,与对它所处理之事的一份\textbf{领会}之深化,并步共进;也就是说,与对它一份越来越精细、越来越完整的\textbf{直觉}之展开,并步共进。或许在我已刊的著作里,按数学家行当的规矩,最显眼、对读者最醒目的,是阳的侧面——“结构”或“逻辑”或“方法”的侧面。然而,我分明知道:在我的工作里引领并支配着的、它的灵魂与存在的理由,是工作进程中形成的、用以把握数学事物之实在的\textbf{心灵意象}。

当然,为了借助数学语言,尽可能一丝不苟地圈定这些意象、以及它们所给予的把握,我从不吝惜气力\dmark{}\origpage{551}。或许正是在这持续不断的努力之中——把尚未言说的表述出来,把尚且模糊的弄精确——有着数学工作(或许也有一切创造性智性工作)特有的那种动力学:在多少还无形的\textbf{意象},与那个赋予它形式、并一路催生新的多少朦胧的意象的\textbf{语言}之间,一场持续的辩证——新的意象深化着先前那些,而它们同样呼唤着表述,好在轮到它们时得到形式……况且,正是这份永不停歇的工作——用语言去圈定那最初显现为不可名状而无形的“预感”、显现为未成形的“感觉”、显为一幅没入雾霭的意象的东西,尽可能精确、尽可能完美——正是这份工作,自我的童年起、直到今天,是数学发现工作里最让我着迷的东西。不过,倘若这里的“努力”似乎总是朝“语言”一侧——即朝表述、结构、逻辑那一侧,它们构成数学\textbf{方法}的关键成分;倘若(情势使然)一篇数学\textbf{文本}——它理应复现一项数学工作(或至少其果实)——的可见方面,主要也落在那里;这一切并不妨碍(至少在我身上):数学事物之领会的灵魂,并不在那个侧面,数学工作里起作用的活力或动机,也不在那里。我相信,在我的著作之中,这样的关系被颠倒了的工作——我曾仅凭其内在逻辑、凭融贯性的要求或形式主义自身的其他侧面来引导,而不是由一个内容、一种实质——一种以“几何”性的意象与直觉来显现的实质——所引导——必定极为罕见。

无论如何,我一生都无力阅读一篇数学文本——无论它多么平淡、多么简单——倘若我不能以我数学事物的经验为它给出一个“意义”;也就是说,倘若这文本不能在我心中唤起心灵意象、唤起将赋予它生命的直觉——一如鲜活的肌肉与器官的肉身赋予一具身体以生命,没有它们,身体便只剩一具骨架。这份无能,也把我和我的大多数数学同行区分开来;并且(如我有机会提过的)正是它,常常使我在Bourbaki小组的集体工作里难以容身——尤其在共同读书的场合:一连几个钟头,所有其他人都从容跟进,我却往往被晾在一边。

\asterisms

我刚刚追随了几场关于我数学工作的联想,系于“直觉——逻辑”这对子,以及紧随着它自行进场的那几个相邻的对子:无形——有形,未定——已定,未言——已言,模糊——精确,灵感——方法,视野——融贯……逐一检阅(如我曾盘算的)一切可以想见的、与智性工作相关的“对子”,并逐一探测两个配偶各自以何种方式、在多大程度上在我的数学工作里在场,以及两者之一是否显得“定着调子”、是哪一个——这无疑会很有教益。\origpage{552}它非但能让我更精微地把握\textbf{我的}数学工作的特质,也肯定能让我深化对一般数学工作之性质的理解,同样也深化对如此检阅过的每一个对子的把握。但这样一项系统的工作,分明会把我带得太远,超出眼下这场反思的合理边界。我觉得更自然的,是试着在这里重拾、并(若办得到)“传递”那些联想与意象——它们使我确信(无须再走更远):在我的数学工作里,确是我存在中那些“女性”的特征,倾向于悄悄定着调子;并由此为它们在我生活其他领域所受的压制,找到一场出乎意料的“报复”(在人们最想不到的地方!)。

拿证明一个尚属假设的定理这项任务来说吧(在某些人看来,数学工作似乎就缩小到这件事)。我看出两种极端的进法。其一是\textbf{锤与凿}的进法:摆在面前的问题被看作一颗大核桃,又硬又滑,要去够着它的内部,那被壳护住的滋养人的果肉。原理很简单:把凿刃抵住壳,用力敲。必要时,换几个不同的地方重来,直到壳裂——于是心满意足。当壳上有棱有凸起、可据以“下手”时,这条进法尤其诱人。有些情形,这种下手的“茬口”一眼便见;另一些情形,则须把核桃在各种方向细细翻转、细细勘探,才能找到攻击点。最难的个案,是那壳浑圆而均匀地坚硬:任你用力敲,凿刃打滑,只在表面上留下浅浅的划痕——人们终于厌倦而罢手。但有时,凭肌肉与耐力,终究也能成。

我仍可拿开核桃作比,来说明第二种进法。方才浮上心头的第一个寓言是:把核桃浸在一种软化液里——就用清水,何尝不可;不时搓一搓,好让它渗透得更透;其余的,交给时间。\origpage{553}几个星期、几个月过去,壳渐渐变软——时机成熟时,手上一压就够:壳便开了,像一只熟透的鳄梨!或者,把核桃留在太阳底下、雨水里,或许还有冬天的霜冻底下,让它熟。时机成熟时,一根纤嫩的幼芽,从那丰腴饱满的果肉里钻出,像玩儿似的穿透了壳——或者更确切地说,壳已自行开启,给它让路。

而几周前浮上我心头的意象,又自不同:那有待认识的未知之物,显现为一片泥土或紧实泥灰岩的地界,不肯让人穿透。你可以抡镐、使钎,甚或动用风镐:这是第一条进路,“凿”的进法(带锤或不带锤)。另一条,是\textbf{海}的进法。海悄然无声地向前推进,仿佛什么也不会破裂,什么也不移动;水还离得远,几乎听不见……然而它终归围住了那块倔强的实体;后者渐渐变成一座半岛,继而一座岛,继而一块礁石,终于也被淹没——仿佛它最终消融在那一片一望无际的海洋之中……

对我的某些工作多少不陌生的读者,认出这两种进法中哪一种是“我的”,将毫无困难——在《收获与播种》第一部分里,我已有机会就此表白,语境稍有不同\fnote{101}{(*)}{见“梦与证明”一节(n\textsuperscript{o}~8)。}。那是“海的进法”:凭淹没、吸收、消解——那种进法里,倘若你不是十二分留心,任何时刻都仿佛什么也没发生:每一事物在每一时刻都如此显然,尤其,如此自然,以致人们常常几乎不敢把它白纸黑字记下来,生怕显得自己在搞名堂,而不像大家那样抡着凿子……然而,这恰恰是我自幼凭本能实行的进法,从来不曾真正学过。

说到底,这也曾是Bourbaki的进法;而我与Bourbaki小组的相遇,在这一点上简直是天赐:它印证了我、鼓励了我\dmark{}\origpage{554}在这个自发属于我的“风格”里走下去——不然,我多半会发觉自己这一类的人,多少形单影只\fnote{102}{(*)}{在这份极阴的进路里,我的走法,甚至比Bourbaki的大多数朋友所愿意走的更远。这无疑是我终于在五十年代末离开那个小组的缘由之一。}。诚然,“形单影只”的处境,久已为我所熟悉,也并不怎么叫我难受。至于我的直觉性数学工作进法会不会“有效”——即首先(按当时的通行标准,尤其按评判一个初出茅庐的数学家的标准)我有没有能力解决那些尚无人能答的“未解决问题”——这我无法预知,也未曾过分操心。我自然的取向,更倾向于给自己提出我自己的问题,而不是想去解决别人提出的问题。而我的数学著作之所以证明为多产,确确实实主要凭新\textbf{问题}的发现、同样也凭新\textbf{观念}的发现,或凭新的\textbf{视点}、乃至新的“\textbf{世界}”,胜过凭我对已被提出的问题所能给出的“解决”。这股把我推向发现好问题、而非发现答案,推向发现好观念与好命题、远甚于推向证明的强劲冲动,也是我的数学进路里若干打着强烈“阴”印记的特征\fnote{103}{(**)}{况且我有这样一种印象:在我这里,别的任何研究工作,情形并无不同——尤其我所谓的“禅观”。}。这也是为什么,无疑地,当我看到我献给数学的最好的东西,被我从前的某些学生——即恰恰是它们的第一批受益者——漫不经心地或轻蔑地打发了的时候,我格外敏感。

无论如何,我只是事后才得以确认:我的自然数学进法,当我感到被别人提出的一个问题吸引、被它启发时——总之,当它“咔哒”一声对上了、而那问题旋即变成“我的”之时——也同样“行得通”。倘若我试着就这类情形开一份多少穷尽的清单,我猜想那清单会相当长。粗粗一望,便有四桩这样的情形,以其分量“出类拔萃”\fnote{104}{(***)}{我此时想到的问题,按其解决的时间先后,是下列这些:\par\noindent
1.~Riemann--Roch--Hirzebruch~公式在任意特征下的有效性。\par\noindent
2.~任意特征的代数闭域上一条代数曲线的“与特征互素”基本群的结构。\par\noindent
3.~有限域上有限型概形的~L~函数的有理性(它构成“Weil猜想”的一部分,也是通往这些猜想之证明——由Deligne完成——的重要一步)。\par\noindent
4.~定义在离散赋值环的分式域上的阿贝尔簇的半稳定约化。}。在这\dmark{}四桩个案里,那个假设性的定理终于大体得到了证明,凭的是“上涨之海”的进法——被某一套或大或小的理论淹没并消解,那理论远远超出当初要建立的那些结果\origpage{555}。况且我得以确认:我在这些情形(或同样在其他情形)里发展出的观念、想法、公式、方法,早已进入数学“众所周知”的领地——“人人”都知道、随意取用,而不问它们的来历\fnote{105}{(*)}{我自己也常常实行这份对所用的“众所周知”不问来历的轻心——例外是:那些我第一手知道来历的情形,或因为我多少亲历了它们的诞生,或因为我自己就是父亲。一如我在逝去的这些年里、尤其在我对葬礼的反思中所屡次确认的:这份起码的体贴,在我从前的某些学生、或数学世界里的近友们身上,却常常付之阙如——即便事关他们除我之外无从学得的东西、即便他们对那来历不容有任何怀疑。关于这一点,参见“掘墓人——或全体教众”一注(n\textsuperscript{o}~97)中的反思。}。

\subsubsection*{(e) 九个月与上涨的海}
\tocsubsub{(e) 九个月与上涨的海}
\label{ch18:local:h31}

\notehead{123}{}
\phantomsection\label{ch18:local:n27}

(11月9日)昨天唤起的那四桩个案——被“上涨之海的进法”所解决(或毋宁说“消解”)的未解决问题——还有一个共同点。那便是\textbf{J.-P.~Serre}在这四桩个案的每一桩里所扮演的角色。那首先是一个“雷管”的角色:促我“启动”这些问题——借用《导言》里一条提到这角色的脚注的说法(见“一个秘密的终结”,《导言》第8节)。实际上(如我当时所确认),情形显出:在我1955至1970年间——即从我离开泛函分析转向几何,直到我离开数学世界——所展开的那些主要的力量观念与大工程的发生学里,Serre都扮演过这样的角色。

我可以说,几乎没有夸张:从五十年代初直到1966年前后——即十五年间——我在“几何”(取极宽的意义,涵盖代数或解析几何、拓扑与算术)里学到的一切,要么是从Serre那里学的\dmark{},\origpage{556}要么是我在数学工作里自己学的。我想是1952年,Serre来南锡的时候(我在南锡待到1953年),他开始成为我一位得天独厚的对话者——而且在好些年里,对泛函分析之外的种种主题,他甚至是我\textbf{唯一}的对话者。我想,他跟我谈的头一件事,是Tor与Ext——我原把它们想得何等了不得,而他说:瞧,简单得跟请安问好一样……还有内射分解与投射分解、导函子与卫星函子的魔法——那时Cartan--Eilenberg那头“diplodocus”尚未出版。从那时起把我引向上同调的,是他刚同Cartan一道展开的、关于Stein解析空间的“定理A与定理B”。我想此前也听说过它们,但正是经过一两次与Serre的面谈,我才感到这些上同调式的命题——它们如此简单——所蕴藏的全部力量、全部几何的富饶。在他同我谈起它们之前,它们起初整个儿从我头顶掠过——那时,我还没有在一个空间的层上同调里“感到”几何的实质。我欣喜若狂,以致好些年里我一心打算:一旦把手头泛函分析里的工作(反正我分明不打算在那里久留!)收尾,便去研究解析空间。我之所以没有真正兑现这些打算,是因为Serre其间转向了代数几何,写下了他那篇著名的奠基之作“FAC”——它使先前在我看来乏味到极点的东西,变得可懂而极其迷人:它甚至迷人到这种地步,我没能抵住这份魅力,于是转向了代数几何,而不是解析空间。

倘若不加克制,我便会顺藤摸瓜,从这里讲起我同Serre关系的来历——那也几乎就是1952至1970年间我数学志趣的来历。此处不是地方。我只想补一句:理所当然地,把我“领进”上文唤起的那四个问题的“水”里的,正是Serre。这当然不是指把问题的确切表述给我指出来,如此而已。要紧的是:每一次,Serre都强烈地感到一个表述背后的丰富实质——冷不丁看到那表述,它多半不会叫我或热或冷——而且他能把这份对一份丰富、可触、神秘的实质的知觉“传递过去”;这份知觉,同时又是对这实质的\textbf{认识的欲望}、穿透它的欲望。这或许正是一切发现工作里最关键的时刻:那“咔哒”一声对上的时刻\dmark{}——尽管人们此刻对从哪儿下手接近未知、从哪儿进去,还没有任何想法,再模糊也没有。\origpage{557}这真正是“受孕”的时刻——从这一刻起,一场孕育的工作能够开始;倘若境况相宜,它便开始……

倘若说Serre在我的工作与我的数学著作里扮演了一个重要角色,那么依我看,他扮演得更重要的,是在那些关键的时刻的出现里——火花闪过、黑暗而不可见的劳作开始启动的时刻——而不是凭他那为我所不识的技术手段(他不时在恰当的时刻提供给我),或凭我在工作后继阶段向他借来的那些观念。

无疑,Serre之所以扮演这个特殊的角色,缘由之一是:我生性不爱靠读书去打听数学的时新,甚至也不爱靠读书——读那些讨论某套“众所周知”的理论的书或论文——去学那种理论的ABC。但凡可能,我喜欢经由“局内人”活生生的言语来打听。我有这份运气:从最初接触数学界(1948年)直到1970年我离开,从不缺一位内行而好意的对话者,给我讲那些可能使我感兴趣的事情。这也许造成了对这些对话者的依赖,但我从不曾这样感觉\fnote{106}{(*)}{头一个也是唯一的例外,出现在1981年——那已是我离开数学世界很久之后。那是在“穿越Galois理论的长征”之后,我求助于Deligne,把他当作我的阿那贝尔式(anabéliennes)反思的当然对话者。那时我分明感到,他有意利用这唯一对话者的处境,把我“折腾得团团转”——于是我便断绝了一切数学层面的关系,直到今天。关于这段情节,参见“两个转折”一注(n\textsuperscript{o}~66)。}。老实说,“依赖”的问题几乎无从谈起,只要我的对话者与我,就他所教给我的东西,怀着同一音高的兴趣。教一个渴望认识的人,对双方都有益,而且是“教者”自己学习的机会,与被他教的人一样。

方才给出的那个“缘由”,很好地解释了对话者在我数学家生涯中的重要性,却解释不了Serre所扮演的异常角色——依我看,它远远超过我所有其他“对话者”加在一起!确定无疑的是:Serre与我,彼此互补得妙极了。我们有强而多的共同志趣;我在他身上感到那同样的要求、那同样的严格——我在自己的工作里也是这样要求自己的。除此之外,我们的工作循着极不相同的“风格”进行。我有这样一种印象:我们对数学的进路\dmark{}与我们的工作相互\origpage{558}\textbf{补足},而几乎从不相互侵越。我所做的那类工作(以及我做它的方式),与Serre的那类工作大不相同。他会在一篇五十来页的文本里,为一门理论打下最初的地基,甚至肯花一年工夫,写一部中等篇幅的书,优雅而简洁地阐述某个令他心动的题目——但他肯定不会把五年生涯的大好时光、甚至十年或更多,用来经年累月、卷帙浩繁地展开一整套新的语言(此前人们没有它也过得极好),以便为比如说代数几何的某种新而多产的进法奠定基础。他引入过许许多多新的、多产的想法与观念,却不让自己被拖去把它们“怀”到足月、怀到底。而恰恰不止一次,这些想法与观念成了我的出发点,引向一项规模宏大的工作——它对我再合适不过,而要让Serre本人投身其中,则无从谈起。

在这里,一场联想不可抗拒地涌来。借着近日反思的光照,我看待我与数学工作、与我的“著作”的关系,更多是“\textbf{母性}的”,而非“父性的”。受孕的时刻,无论多么关键,在我这里只代表那场“工作”里微末的一小部分:在那场工作里,孕育中的东西、那个将要来的“孩子”,生长并发育。这份工作,正像一位孕妇的妊娠:它自受孕起启动,延续整整漫长的九个月……那把一个胎儿怀到足月、把它\textbf{分娩}出来所需的时间——也就是说,把一个\textbf{孩子}生到世上来,一个活的、\textbf{完整}的孩子,而不是一个婴儿的头、或一具躯干、一具婴儿的骨架、或天知道什么。母亲的角色,分明与父亲的角色(哪怕是世上最好的父亲……)大不相同:后者几乎只消播下一粒种子,便扬长而去,忙别的营生。显然,Serre的数学工作、他对数学的进路,是阳占强势主导的、“男性”的。他对付困难的进法,宁可说是凿与锤的进法,而极少是那上涨而淹没的海的进法、那浸润而消解的水的进法。而且他似乎乐于播下种子,并不过分操心它落到哪里,或它会不会触发受孕与分娩,甚至不操心由此可能生出的孩子会不会像他、会不会承他的名字。

\dmark{}一幅意象能帮我们把握某一现实的重要侧面,却穷尽不了那现实。\origpage{559}现实永远比任何想要表达它的意象更复杂、更丰富。我所得到的这些意象——它们不召自来,为表达对数学的两种不同进法:Serre的进法与我的进法——也是如此。Serre也曾完成过需要长劲的工作;我也曾播撒过观念的种子,其中有些发了芽,并由别人而非我怀到了足月。一如我的数学进法并不缺乏“男性气”(尽管底色是“女性”的),Serre的进法也不缺乏“女性味”,平衡着他那“男性”的底色。

在对一份未知实质的创造性进法里——无论数学的或其他——不可能不如此:没有阴与阳这两股原初的能量与冲动在同一存在者之中的联合而不可分的行动,便没有发现,没有认识,没有更新。恰在这两者的内在融合之中,存在着一个存在者、或一部著作的\textbf{美}——那份精微而难以把捉的品质,它借着那份和谐与满足的特殊感觉,向我们显现自己。这份品质,在我所认识的Serre的全部著作里都在场——无论当面聆听,还是经由他写下的文本。我认识的数学家之中,这份品质如此恒常、带着这份力量而在场者,极少。

\subsubsection*{(f) 阴的葬礼(阳葬阴(4))}
\tocsubsub{(f) 阴的葬礼(阳葬阴(4))}
\label{ch18:local:h32}

\notehead{124}{}
\phantomsection\label{ch18:local:n28}

(11月10日)昨天与前天的反思,远远没有穷尽我数学工作里那些打有强烈印记的、属阴品格的特征。顺着眼下这场关于数学中阴阳的反思,再往深处探测它们,对我也是一次绝好的机会,去深化对一般数学工作之性质的理解。数学中的阴阳这个题目,我原以为一天反思便可巡完,却已经连着花了五天,还觉得刚刚碰了碰边儿。它方才显出真面目:又是一个外表平淡的题目——愈是走近它、进入它,它便愈发广阔、愈发深邃。断断没有道理:我趁乱草草穷尽这个多汁的题目(甚至也不“巡完”它,哪怕是跑步通过),而葬礼仪式正进行到半中间——我可不愿意让它拖沓得漫无节制!

我只想再指出(不加评论,我保证!)我数学工作里这些朝着“阴”的、“女性”方向的“强烈印记”中的两个\dmark{}\origpage{560}。其一是对\textbf{一般}的偏爱,胜过对特殊者的偏爱(后者与前者“配对”或“成偶”)。另一个特征,我觉得还要更强,或者说得更准些,更要害、更要紧,也更宽广(就它\textbf{包含}第一个而言)。倘若有一场“追寻”贯穿了我整个数学家生涯——从十七岁(刚出中学校门)直到今天;一场不停歇的追寻,自我著作的发端起便打下了我全部著作(已刊或未刊)的印记——那便是对\textbf{统一}的追寻:穿过数学事物那无限的多样性、以及通向这些事物的种种可能的进法。在多样性之外辨认出、发现出这份统一——那份常常令人错愕的富饶(而丝毫不削减这份富饶);在差异与不似之外,认出共同的品质;一路走到类比与相似的根子跟前,去发现那份深刻的亲缘——这便是我的激情,我的一生。连那些差异本身——那无限而难以把捉的多样性的表达——也终于显为同一棵枝繁叶茂的大树的枝与杈:枝杈无限分岔,而每一根、每一枝、每一杈,都给我指明通往它们共同的树干的路。凭本能、凭天性,我的行程一直是\textbf{水}的行程——水总倾向于\textbf{往下}走:走向那树干、那些根须的行程。倘若我也爱在途中流连,那很少是在树梢,去探查叶片与纤弱的细枝;多半是在粗枝上、在树干与几条主根上,去认识它们的肌理,透过树皮,感到那滋养人的树液上涨的涌流\fnote{107}{(*)}{我自信能在这场穿过多样性而对统一的追寻里,辨认出打下了我一生印记的那三种激情的共同特征——爱欲的激情与禅观也在其内。甚至或许,在一切激情之外,这在我身上是把握实在的一个\textbf{模式}:我倾向于主要去看那些共同的品质与亲缘,把注意力与分量放在它们上面,而不是放在差异上(却也不曾起过把差异变没了的念头)。我注意到,远为常见的倾向是相反的倾向,即阳的倾向。它常常走到这一步:忽视或否认深刻的亲缘。(超级阳的倾向,我们文化的特征。它常常伴着这样一种反射:想把差异碾平,把一切对齐到同一个号称“完美”或“高等”的模子上,以成全一份虚假的“统一”——那是极度的贫瘠化,同时也是一场暴力。)对话者与我之间这些重音的差异,常常是聋子对话的缘由:两场平行的独白各自展开,却永不相交……}。

\asterisms

\dmark{}老实说,对这个不久前才发现的新事实,我还不知道该拿它怎么办,不知道怎么给它安一个位置:在我的数学进法里、在我“做数学”的方式里,那基本的调门,在我身上是强阴的、“女性”的。这与我已经暗示过的某一直觉同向:我存在深处的那个基本调门——我是说我身上“孩子”或“工匠”的调门,即那创造者的、在制约之外的调门(也就是说,在“我”、“东家”之外的)——它同样也是“女性”的,而非男性的。或许此刻我手中已样样在手,足以查清实情究竟如何:细细检视所有这些或朝这一方向、或朝那一方向的征象\fnote{108}{(*)}{我身上若干打有强烈印记的阳的特征,在我看来是\textbf{习得}的特征,来自制约,更确切地说,来自上溯到我童年的那份超级阳的形象招牌。这些特征之中,有对行动的一场无度的投入;朝向未来的极强的投射——即朝向完成我的任务;对一种首先是智性的发现工作的偏爱,以及思想的侵夺性角色;对一切不显得与我当下任务直接凝聚者关闭心门的姿态,尤其我对风景、季节等等的漫不经心。然而有一个阳的特征,我觉得是天生的而非习得的:那把我与\textbf{火}系在一起的极强的亲缘关系——不同于我与水的关系,水分明不是“我的元素”。况且据说,我的星盘也打有一个极强的阳失衡的印记:进入它的所有星座,都是“火象星座”,而没有任何一个水象星座。},去认出每一个的分量,以及从它们的总体里透出什么。倘若经由这样一份工作,我达不到一个“是”或“否”的可触的结果,它也肯定不是无用的:它将帮我更好地圈定我的无知——这份无知,此刻依然模糊、尚未定位,因为我不曾对它禅观。或许我会做这份工作,在《收获与播种》的工作完成之后,并顺着它的势头。但在这里,这同样不是地方。

不过,我之所以被引向这场关于阴阳的反思,是在一场反思的进程中——在那场反思里,我尤力图理解我自己与其他一些人(尤其那些曾是我学生的人)之间的某些关系。所以,在此最使我感兴趣的,是这个方才显出的“新事实”可能有的反响:对我与他人的关系、以及他人与我的关系。而这也正是我在“安放”、运用这个事实上的为难所在。这为难或许系于:这样一件事,除了我,恐怕无人曾有所察觉——至少不在意识的层面、不在表述出来的层面。无论如何,就记忆所及,我从不曾收到过任何可作此解的回声——除了一例不算——正如我也不记得收到过什么回声,把一幅“阴”的影像从我自身反射给我;尽管自我童年起(婴幼儿期不算)我所扮的那个角色,一直是强阳的\origpage{562};以致时至今日,这“男性”的品格似乎仍像一种第二(?)天性,仍以种种方式支配着我的生活。

诚然,某个特征在某个人身上(此处即我)不被意识层面所知觉这件事,并不必然拦阻它作用于与他人的关系。而这个特征确确实实被数学世界知觉到了——在我那些对我的著作多少熟悉的数学同行之中;并且这份知觉,在一个分明广阔得多的数学公众之中“油渍般扩散”开来——这在我毫无怀疑。当我在“葬礼悼词(1)——或恭维”里写道:“那支在这里细心料理我葬礼悼词的无名之笔,把今天已被交给轻蔑的东西,加倍再加倍地赏给了我”,我当时还无从一言道破:在那些我珍视的东西之中,被数学时尚“今天交给轻蔑”的,究竟是什么。但第二天,借着那场我日后必须回头的“联想”\fnote{109}{(*)}{见“肌肉与肚肠”一注(n\textsuperscript{o}~106)的开头,那场联想在那里头一次被唤起。},我已经感到(或许还不曾对自己表述,而它也还不像如今这般清晰):“这某样东西”不是别的,正是那一切被(往往在未成形的层面)承认为做数学的一种“阴”式的、“女性”方式的东西——那方式被默默等同于“瞎拼凑”(bombinage)、“胡话”(借用我的学生兼朋友Pierre~Deligne的恭维——针对那构成他全部著作之根基的文本)、“花架子”、“省力”等等。

当然,在那篇葬礼悼词(由同一位朋友皮埃尔所致)里——包括把我与他相提并论的那一段\fnote{110}{(**)}{见“葬礼悼词(2)——或力量与光环”一注(n\textsuperscript{o}~105)。}——恭维是题中之义!那里谈的绝不是什么胡话或瞎拼凑,而是“\textbf{泰坦般的}气象”、“二十卷著作”、“剥出了\textbf{本质问题}”、“最大的\textbf{自然普遍性}”(原文如此)、一个被他传达观念的那份\textbf{慷慨}所“\textbf{滋养}”的学派、“\textbf{传说般深邃}的理论”、“更新了\textbf{基础}”、“开辟了\textbf{新的应用}”、如此\textbf{自然}的观念,“以致我们难以想象它们曾花费的气力”(更不说它们“容易”了——不过这一点,我留神亲自说清了\fnote{111}{(***)}{见“陷阱——或省力与衰竭”一注(n\textsuperscript{o}~99)。})\dmark{}、“对术语的\textbf{高度关注}”(更不说“瞎拼凑”了)、代数K理论的“\textbf{鼻祖}”、在\textbf{一般}基域上引入的topos、“由Grothendieck\textbf{提示}的类比”、“那些\textbf{猜想}……依然同样无从下手”、“一如Grothendieck所\textbf{梦}想的那样”……\origpage{563}

在这些引文里,我给关键词划了着重号——它们全是一些指称对事物的阴式进法的词。这场由“配量得当的恭维”所行的葬礼里,那“完美的手法”就在于:对着这些品质——它们一面“被交给轻蔑”,一面却是真实的、在我这里价值连城——系统地使用夸张;而与此同时,对那些互补的侧面——今天独享荣耀的那些“男性”侧面,它们在我的著作里,也如在任何人的著作里一样强地在场,几乎无一例外——则拿橡皮整个儿、彻底地擦了个干净。

况且,在那篇谈Pierre~Deligne的文本里,被摆到聚光灯下的,恰恰是这些“男性”的侧面与价值——带“女性”味的一丝半点也被排除在外:既凭那几顶选定的形容语(“众所周知的\textbf{困难}”、“令人惊讶的\textbf{结果}”、“使$\ell$-adic上同调成为一种\textbf{强有力}的工具”、“\textbf{第一步}”、“\textbf{出奇地有用}”、“\textbf{敏捷}”、“\textbf{透彻}”、“对每个问题都作出\textbf{启人心智而富有建设性的回应}”、“\textbf{耀眼的发现}”),也凭对可触的成果的一份周详的罗列(而在给我画的那幅速写里,我的成果一项也没有被提及,甚至也不曾暗示:我的成果或许曾对Deligne的成果起过作用)。

我不后悔费神做了这份形容语的快速汇编——它的效应真正触目惊心!倘若在一份有结构的知识层面,对阴阳略有所知的人依然寥寥,那么只好相信:在我朋友皮埃尔的无意识里、如同在那位为他执笔者的无意识里,有着一份毫无破绽的知觉。它在这里被用于某一项事业:把该交给轻蔑的交给轻蔑,并给人群指定一位可供仰慕的英雄。

况且我很怀疑:我方才重读的这三篇短文,有过许许多多的读者。但读者或多或少,依我看是个次要问题。对我而言,这些文本所诉诸的,并非什么假想的潜在恩主(毕竟,替他的机构寻找恩主以筹措经费,不是我朋友皮埃尔操心的事),而是“\textbf{大会全体}”——它出现在与那篇同名的注(别名“掘墓人”,n\textsuperscript{o}~97)一同进行的那场反思里\dmark{}\origpage{564}。它们所传递的讯息,如同一份触目惊心而手法老到的速记,浓缩了无数同方向的讯息——它们出自我的朋友皮埃尔、出自那些曾是我朋友或我的学生的人中的另一些、或许还有其他一些人——并被这同一个大会全体接收、认可。倘若存在一个集体无意识(而此刻我相当倾向于相信它存在),那么毫无疑问:在这个大会全体(别名“数学共同体”)的无意识里,一如在那位主持我隆重葬礼的大司仪的无意识里,有着这份同样毫无破绽的知觉:什么是阴(受够了!),什么是阳(脱帽!)。

而这葬礼,忽然以一种新的、意想不到的光照显现在我眼前:在这光照里,我自己的这个人变得无关紧要,变成了那个“应当被交给轻蔑”者的\textbf{象征}。这不再是一个人的葬礼,也不是一部著作的葬礼,甚至不是一场不可容忍的异见的葬礼——而是“数学中的阴柔”的葬礼——而或许更深地,在每一位为葬礼悼词鼓掌的众多与会者身上:是\textbf{那被否认的、活在他自身之中的女人的葬礼}。

\subsubsection*{(g) 超级妈妈,还是超级爸爸}
\tocsubsub{(g) 超级妈妈,还是超级爸爸}
\label{ch18:local:h33}

\notehead{125}{}
\phantomsection\label{ch18:local:n29}

(11月11日)破例地(下不为例……)今天早上我醒得早,睡了不过四五个钟头。昨天反思那出人意料的收束,旋即启动了一场强烈的工作:去“安放”、消化这个方才显现的新事实——那是在我热一份丰盛的汤、睡前吃点东西的当儿,彼时已过凌晨三点。而一大早,这同一份工作又把我从睡梦里、从床上拽了起来……

倘若我说“出乎意料”的收束、“新”的事实,却还须补上:自这条关于阴阳的没完没了的“枝蔓”的最初几天起,我身上就有一种仿佛克制的等待:等一个“了结”,或至少等一场“会合”——它迟早要与某支队伍会合:那支在一场葬礼仪式中聚拢起来的队伍。看上去,我离葬礼的场地越来越远,甚至仿佛葬礼已被彻底遗忘——然而并非如此:它们一直在那儿,像弱音器里、像水印里。我从不曾真正离开它们。它们无声的在场,显于这份含蓄而恒常的等待、这份紧张与悬而未决之感——它把我推向那个尚且云雾迷蒙的点:那“会合”终须在那里完成\origpage{565}。我本能够预感到这个会合点的大致所在——它绕着某一场“联想”(我已不止一次唤起它,却始终未曾表述):它正是这场穿行阴阳、穿行我一生的意外旅程的出发点、最初的动因。这场旅程,总而言之,又将像一个大的循环,回到(或多或少……)它的出发点;或毋宁说,像沿着一条下行螺线转了一圈,把我带到那被探测之物更深一层的地方,带到这些葬礼的“正中心”(倘若我的预感不曾骗我)。

然而,正当我开始准备“降落”,并在一条还十足是“枝蔓”甚或“回锅”的“注”的最后一个段落的转弯处,我忽然发现自己已经登临葬礼仪式的正中,而且实实在在在它的正中心——有点像一个外星人,被弹射到恰好落在披着祭披的神甫面前、落在信众大会跟前;更糟的,像一个被认为已经死去、(几乎)已经下葬的亡人,忽然掀开棺盖(于是花圈与感人的墓志铭纷纷滚落!),以真身现身,裹着白色裹尸布、目光闪亮,活脱脱一只小魔鬼,在人们最不提防的时刻,从盒子里跳出来!

这样,昨天反思的收束,同时便是我方才所说的那份悬疑的了结——一份十分特别的悬疑,在“如铺展的海一般”的工作方式里,它是我熟得不能再熟的,无论数学工作或其他。但就在这场长悬疑之松弛的尾流里,旋即显出一份\textbf{困惑}。我相信,自那时起主要吸住我的,正是它;而且正是在不合时宜的钟点,是它把我从床上拽到打字机前。有困惑,本不足怪——差不多每次,当一个情势忽然显现在一种新的光照之下、而这光照初看似乎与旧的景相矛盾时,事情便是这样。于是头一件非做不可的工作,便是细细探测这些矛盾:检视它们在多大程度上是实在的,或仅仅是表面的——也就是说,是一种心智惰性的表达,那惰性不肯在不同的两道光照之下认出“同一个”东西。这份必不可少的工作便算完成:当一切不谐之音都消解于一份新的和谐(哪怕这和谐本身仍是暂时的),消解于一幅把先前诸幅局部景涵盖并统一起来的景——必要时修正、调整它们,并剔除那些终于显出根本为假的。在这样一幅焕新的景里,那个诞生了它的“旧”\dmark{}——即融入其中的那些更为零碎的景——本身也获得一份新的意义\fnote{112}{(*)}{可与“孩子与老天爷”“错误与发现”两节(n\textsuperscript{o}~1、2)中的反思对读。}\origpage{566}。

回到我的“困惑”。它在此:那个“了结”或“新的日子”,在于一幅忽然显出的意象——“数学中的阴柔”的“象征”的葬礼,大操大办,以我这个人而为化身;同时,它又是每一位葬礼与会者身上“被否认的女人”的投影;换句话说,这是一场象征性葬礼的意象,葬的是一位\textbf{超级母亲}——总而言之,一位赎罪的牺牲,代替那个女人——却极少是母亲——而她,正在每一位前来为葬礼鼓掌的与会者的黑暗地窖里,萎顿着。这幅意象似乎与\textbf{另一幅}相反的、尚且朦胧的意象相矛盾——后者在6月之前的反思进程中逐渐成形(在“掘墓人——或全体教众”一注达到顶点):那是一位\textbf{超级父亲}的意象,又受敬慕又遭惧怕,又有吸引力又遭憎恨,被他的孩子们“屠杀”,他被残害的躯体,就在这“同一场”葬礼中被交给嘲弄。两幅浓墨重彩的意象并排一放(倘若还需要的话),似乎要在滑稽与谵妄的边上打转;我很容易想象,这些按精神分析路数来的幻象,会激起怎样一场剥头皮的舞蹈——倘若竟有读者有那份气力,一路跟到这里!

我乐意把他们留给他们的舞蹈——它将给这场不同寻常的葬礼添一支最有效果的异国情调的插曲;而在此期间,我宁肯追随一场昨夜便已呈现在心中的联想——我相信,它能把这两幅所谓对立、乃至不可调和的意象或侧面调和起来,甚至使它们相爱、成婚。

\subsection{阴阳的颠倒}
\label{ch18:local:h34}

\subsubsection*{(a) 颠倒(1)——或激烈泼辣的妻子}
\tocsubsub{(a) 颠倒(1)——或激烈泼辣的妻子}
\label{ch18:local:h35}

\notehead{126}{}
\phantomsection\label{ch18:local:n30}

(11月12日)我原打算在注里继续这场联想——昨天注的结尾谈到过它,它能把从我葬礼中透出的那两幅表面敌对的意象“调和”起来、让它们“相爱”。正当我准备朝这个方向起笔,我感到一阵迟疑,而我不愿越过它。

这联想事关我母亲与我父亲的关系,事关1933年那个家庭的解体之意涵——那是凭我母亲的意志、带着我父亲的首肯(起初勉强而尴尬,继而急切而彻底)而发生的。这段关键的情节,在我父母结成的那对\dmark{}里打下一个颠倒的印记\origpage{567}:在这对偶子里,我父亲一直扮演男性价值的英雄化身——那份化身被张扬地崇拜着;而我母亲(倘若曾有一个专断而好支配的性格的话,便是她)则把自己打扮成被征服而乐在其中的女人的颜色,底下却是被持续不断的交锋打满了印记的日常生活。对孩子之牺牲的首肯,标志着\textbf{上帝与英雄崩塌}的时刻;继之而来的,是“胜利的轻蔑”的一场真正的放纵——在那位直到前夜还在扮演神魂颠倒的崇拜者的人身上;而她从此取代了那位跌落的、被阉割的而乐在其中的英雄,把他贬进那个受轻蔑的“女人”的角色——恰好就在同一时刻,她自己从这个角色里被解脱出来……

我说的这寥寥数语,如此梗概,我担心如此浓缩,以致它引起的误会恐怕数不胜数,而非帮助人理解某场葬礼的隐秘发条。然而我感到,此处不是把我方才以几句话勾勒的东西稍加展开的地方。要以最低限度的精微,复现一份复杂的实在——那实在被两位当事者着意搅浑——须要一场新的、长长的枝蔓,其规模并非语境所能辩解。我此刻不觉得自己有跳进去的兴致;尤其因为,这涉及一份牵涉他人的情势,而我自己的责任(作为共同行动者)在我看来并不真正牵涉其中。我自己、我的姐姐,在其中不是作为行动者出现,而是作为我母亲手里的\textbf{器械}:用来打倒那位被热烈敬慕与嫉妒的英雄,好取而代之,把他变成嘲弄的对象。

倘若这场五年前被耐心揭开的脚本\fnote{113}{(*)}{关于这一点,参见“表面与深度”“写作之颂”两注(n\textsuperscript{o}~101、102)。},是我所认识的这一类里最极端、最暴烈的,我此后却仍有的是机会,在其他伴侣之中察觉到全然相仿的脚本。对我父母生活所做的这份工作,大大帮助我睁开眼睛,看见那些此前整个逃过我知觉的事情。但在当时,我目瞪口呆——而也有得可呆!今天,我倾向于相信:除了那份色彩之暴烈与众不同,我在我父母结成的那对偶子里揭开的这种对抗关系,或多或少是伴侣关系的典型——或至少极其常见。因此,那位像我一样终于动用官能、去探测伴侣对抗(或男与女对抗)之隐秘发条的读者,对我在这里说的寥寥数语,将不至于特别惊讶(乃至震惊)。

倘若我试着撇开一件件个案各自特殊的东西,而从我有机会近距离看到、并有所领会的那些男与女的对抗里,提炼出共同点\dmark{}\origpage{568},那么得到的是下面这些。

\textbf{1.}~在女人身上,有一种朝向男人的钦慕与嫉妒的姿态,缘于罩在他身上的威望(常常言过其实)——那威望来自他的处境(尤其作为雄性),以及为之张目的种种品质(真实的或假定的)。

\textbf{2.}~常常,其中混有一份怨怼、甚或仇恨的成分,缘于男人(比如情郎或丈夫)与父亲之间的一份(理所当然是无意识的)混同。母亲对父亲的对抗关系,被女儿接了过来据为己有——她(或多或少充分地)与母亲认同。还常常加上一些(对父亲的)更直接的怨怼缘由(他的暴君姿态、缺爱、不关心、不体贴等等)。此后,这些“现成可用”的对抗之情(以及别的),便原封不动地投射到那位伴侣(现实的或潜在的)身上,无论他有没有“对得上号的长相”。

所以,方才(在第1条里)我写道,女人朝向男人的姿态(尤其钦慕与嫉妒)是“缘于威望等等”,这只是部分为真。我觉得,最常见的情形是:这些姿态里的\textbf{活力}\textbf{源自对父亲的关系}(哪怕父亲早已死去、早已下葬),而它进入行动,只在有限的程度上取决于伴侣个人的品性。

\textbf{3.}~作为对她那些劣等之感(不用说,完全出于主观)与蒙着面纱的对抗、乃至敌意或仇恨的补偿,有一种执念缠着她:要对伴侣行使权力(而按那份或多或少默契的普遍共识,被认定握有权威的,恰恰是他)。女人对权力的行使,凭她手头一切手段进行(其中最有力的,是她的身体,而尤其,是孩子\fnote{114}{(*)}{然而,那最主要的共同“手段”在此被略而不谈,因为它天性更为精微,不易用几句话说清。它在于某一套通行无阻的“战术”——那在关于阴阳的反思的后文“天鹅绒中的利爪”(注~n\textsuperscript{o}~137--140)部分里得到检视。}),而且它几乎总是隐秘的。伴随着它的那份犒赏,因此多半是无意识的,但它并不因此少了实在、少了分量。常常,这权力的游戏变得吞噬一切:它变成女人生活的主要内容,吸收她几乎全部的能量;其余的一切(包括爱的冲动与孩子)都从属于它,甚至被毫不犹豫地牺牲给它。

\textbf{4.}~\dmark{}最极端、被撕裂得最厉害的,是这样的个案:对雄性的钦慕与嫉妒——她要支配他,又显得臣服于他\origpage{569}——伴随着对一切属于阴柔性之事的轻蔑、乃至厌恶与仇恨:对她自己身为女人的处境。然而,恰恰只有玩弄她自己的“女性味”,她才能指望降服男人,或至少把他摆布到她的心意!这样,为满足她最强烈的利己冲动——“让男人听使唤”的冲动(甚至制服他、乃至毁掉他……)——她竟不得不全身投进一个她所憎恶的角色,一个她自觉可鄙、配不上她的角色。正是在这拒绝自身处境与天性的极端个案里——一份超级阳的、反阴的抉择——她将动用自己全部这些力量,去求一条摆脱她身上冲突的幻想的出路:一场\textbf{角色的颠倒}:她自己取代男人——那位英雄与主人,昔日的钦慕与嫉妒之对象,而今的跌落者——把他贬进她长久以来像一件卑贱号衣一样穿着的角色、那个她终得解脱的受轻蔑的角色……

我方才勾画的这幅速写,同样只是梗概,至多能在那种人身上“唤起”某一现实——他已在自己的那边、此地彼处知觉过它,却或许还不曾试图像这份粗略的描述那样,多少把它圈定;倘若我想给它一些立体感,我至少该试着把这整套相互敌对的感情与意愿所上演的各个\textbf{层面}(几乎全部无意识)弄清楚。此外,在这场无情利己机制的缠结里——爱的冲动在其中似乎被严格排除在外——也该试着给这冲动安一个位置:看它在多大程度上、以何种方式,助长这永世打转的转圈(也许就像风力,被一座精巧风车的翼板捕获,叫一副沉重的磨盘永远转下去……),以及在多大程度上,齿轮也会时而停下、归于沉寂,好让\textbf{别的东西}畅行无阻。

最后,我完全忘了谈在他身上、“伴侣”或主角身上上演着什么——仿佛他只是相对于她而存在,仿佛他只是迎面那位女人的吸引与厌拒、钦慕与嫉妒的\textbf{对象}。这份遗漏的缘由之一无疑是:在这对偶子的旋转木马里,扮演主动角色、全身投入、并往往在其中找到(退而求其次的)真正存在理由的,正是\textbf{她};而\textbf{他}却什么也没看见,他忙着别的事,而且\dmark{}天真得没人比得上\fnote{115}{(*)}{(11月23日)当然,倘若旋转木马转得起来,那是因为(他再“天真”也罢)\textbf{他}也在其中捞到他的好处,一如她——而把这些看牢,正是她的本职!依我看,她借此“拿住”他(而她自己也借此被拿住……)的两个主要“挂钩”,是虚荣,以及对一份情感与爱恋之安稳的需要——由一位稳定的伴侣作保。还有孩子……},\origpage{570}一次接一次地起反应,而不试着理解;而且(更要紧的)实在不曾理解——甚至(依我看)在无意识的层面也不曾。至少,这是自我开始留意这对偶子的旋转木马以来,我一贯得到的印象!不过同样真确的是:我对男人的角色知之甚少,因为我只能在我自己这微末个案里真正贴身观察它;而女人的那一侧的角色,我却不止一次有机会从第一排包厢里看到。

反正,即便我肯花上十页、乃至整整一卷的工夫,来充实我这份稍嫌太梗概的描述,对那位在这门事上还不曾“动用他的官能”、从未见过或感到过此类事情的读者,也终究是白费。至于那位多少“在局中”的读者,我说的这寥寥数语,尽管笨拙含混,肯定够把他重新浸回他凭自己早已知觉的那些事情里,并在他心中唤起不亚于我写下这份速描之时背后所存的那些意象与联想。

依我看,无须更多,便可看到那“缺失的环节”显形:它连接着对“超级父亲”的对抗——它在所说的那场象征性葬礼中找到表达——与对“阴柔”的轻蔑、拒斥,而更深地,与对自己身上“女人”的否认(后者或许将在一位“超级母亲”的“葬礼”的象征中找到表达,在一大堆滥情的、一语双关的颂词之下……)\fnote{116}{(**)}{(11月23日)这句“无须更多”,后来显出有点匆忙:一周之后,这条结论与这个“缺失的环节”已被忘得一干二净!至于要到达一个更有说服力的“缺失环节”所缺的那一步,参见昨天那条注“阴阳的颠倒(2)——或反叛”(n\textsuperscript{o}~132)。}。

\subsubsection*{回顾(1)——或一幅画的三联}
\tocsubsub{回顾(1)——或一幅画的三联}
\label{ch18:local:h36}

\notehead{127}{}
\phantomsection\label{ch18:local:n31}

\dmark{}(11月13日)现在时机似乎成熟了:试着以几笔大轮廓\origpage{571},勾出对葬礼的一幅更清晰、也更有层次的景——它(如我前天所写)“把先前诸幅局部景涵盖并统一起来,必要时修正、调整它们……”。粗粗一望,我看见三幅这样的先前的景:该把它们认作同一个\textbf{全}的三个局部侧面。

显出的第一个侧面,最明显也最简单化的,是“\textbf{对一场异见的报复}”的侧面——它是“掘墓人——或全体教众”一注(参见注~97)里主要摆在前头的侧面,即患病情节之前的最后一条注。在I至X各行列(事故之前的那些行列)的注里,依我看,它也是把那些\textbf{集体动机}圈定得最深的一条——“掘墓人”别名“(几乎)大会全体”的集体动机。

我刚刚重读了这条注。第二个侧面——我可以称之为“\textbf{超级父亲}的(远不止象征性的)\textbf{屠杀}与(象征性的)\textbf{葬礼}”——在那里头并不出现。这或许是因为:葬礼动机中的这一成分,并不真正涉及“大会全体”——它当时是我注意的焦点——而主要(若不说排他地)涉及“那些曾是我学生的人”。诚然,这些人——即便把他们那位无可争议的头领、我的朋友皮埃尔放在一边——在葬礼的操办中扮演了头等重要的角色:没有其中一些人的积极贡献、没有所有人的首肯,葬礼便无从举行。(关于这一点,参见“沉默”一注,参见注~n\textsuperscript{o}~84.)所以,主要是经由他们,“超级父亲”的侧面才对我显得对理解葬礼至关重要。

第一个侧面、“报复”的侧面,是自1976年Yves~Ladegaillerie的挫折以来,就强加给我的注意的\fnote{117}{(*)}{见“谁也拦不住进步!”“棺材2:锯成段的切割”两注(n\textsuperscript{o}~50、94)。};此后我倾向于遗忘这个侧面,但在随后几年里,它不时重新唤起我的记忆。它终于越过了那种只被“感到”的无形阶段,变成一份清晰而有层次的理解的实质——在那条谈“掘墓人”的引注里。第二个侧面、\dmark{}或“超级父亲”侧面,只是在《收获与播种》\origpage{572}的反思进程里\fnote{118}{(*)}{(11月29日)老实说,这个侧面早在许多年里,就已以一种皮肉直觉的形式,存在于我同Deligne的关系之中;只是在《收获与播种》的反思之前,我从不曾在此驻足。}才开始显出,而且起初\fnote{119}{(**)}{见“敌对的父亲(1)(2)”两节(n\textsuperscript{o}~29、30)。}与葬礼本身无关——葬礼本身,是我要在随后几个月里才发现的。在整个葬礼反思的进程中,这个侧面渐渐走出迷雾,终于在“屠杀”“躯体……”“……与身体”(87、88、89)诸注里,以惊人的形态现身。这些注写于5月12、16、17日;“掘墓人”那条注写于5月24日;患病情节6月10日登场,把诸注的写作中断了三个多月,9月22日才重新接上。至少很可能:倘若这场(实在不是时候的!)情节不曾在一个我正准备对整体作盘点、并画一条最后界线之时出现,我对葬礼的景,本会止步于5月12日至24日那两周里透出的景——一幅“两联”的景,各自停在各自的角落,而我想不到试着把它们拼起来。

然而,一直有一种弥散的感觉,像一层几乎察觉不到的薄雾,觉得那最终的答案始终没有被真正把握住;那是“在黑暗中摸索”的人的感觉(这个说法,肯定在我的葬礼诸注里出现过一两次)。“掘墓人”那条最后的注,多少起了点薄雾里一阵微风的作用:它能给人雾已散尽的幻觉,而其实它们只是稍稍挪了挪地方。或者说:这条注所重拾的那个侧面,在其中显得如此清晰、如此有立体感,以致对这一侧面的一份可触、透彻的理解之印象(绝非幻觉),连同伴随它的满足之感(在注的结尾,这感觉肯定显而易见)——这份印象与这份感觉,造出了一种亢奋,一种自觉马上就要够着目标的亢奋,使我或多或少忘了另一个同样有分量的联——“超级父亲”的侧面,它还“记在账上”呢!

第三个联,只在三天前才显出(患病那个不合时宜的情节出现后整整五个月)。那便是\dmark{}“\textbf{阴柔}”的(象征性)葬礼与(实实在在的)下葬这一侧面\origpage{573}——这“阴柔”被可视化成一位“\textbf{超级母亲}”,而她本人,正是由我这微末个人所化身的!这个侧面,是在一场全然未曾预料的、关于阴阳的长“枝蔓”的尽头显出的——那场枝蔓,终于落实了一份努力:把从某篇“葬礼悼词”里生出的一场“联想”以可懂的方式表达出来,而那篇悼词,本该为葬礼仪式收尾。这场著名的“联想”或“直觉”(我在“肌肉与肚肠(阳葬阴(1))”一注的开头处首次暗示过它)至今仍未被显豁——但万事俱备,而且我许诺要谈它,已有些时候了!

无论如何,一路走来,显出了大量的事实与直觉——其中一些对我而言是新的、意想不到的——而它们全都让我有效地重新接上了我生活的重要侧面,一如接上一般存在的重要侧面。这些事实之一——我数学工作的“基调”是“女性”的——看来与那场还在等它时辰的联想所依据的直觉之一相矛盾:那个直觉说,作为数学家(在其他方面亦然),我是一个十足\textbf{阳}的人物;一个系于葬礼“超级父亲”侧面的直觉。而这同一个事实——它看来与那场联想(整场阴阳反思都由它而生!)相矛盾——也一转眼便催出了那个此前一直逃过我的第三联:“超级母亲”的侧面。与此同时,(到头来)与那场似乎已被忘了近百页的“葬礼”的“会合”,也完成了。

要说“上涨的海”,这就是上涨的海——只须希望:最终的结果,我是说那幅我正准备从冥府里拽出来的“景”,配得上手段,配得上整整一海关于阴阳的哲学-弗洛伊德式的枝蔓……潮水已经启动(随着发令的那条注“肌肉与肚肠”),是10月2日;关键的“新事实”在随后几天便已显出\fnote{120}{(*)}{我记得,就在两天后,在“天真(阴与阳的联姻)”一注(n\textsuperscript{o}~107)里,所说的事实便已显出,并构成那条注里所谈的“种种迹象”之一(对它们未再多加说明),那些迹象“不止一次使我起疑:……在我身上占支配地位的,是那些‘女性’的品质……”。},\dmark{}正当我一天又一天地准备把那场著名的\origpage{574}“联想”终于写成白纸黑字——它出现于五个月前,5月12或13日,就在“葬礼悼词(1)——或恭维”一注的反思之后,与那条关键的注“屠杀”同日。但这个事实,直到五天前、11月8日,才在注里“被揭开”——在那之前,是连着三天关于数学中阴阳的三条预备的注。那就是“上涨的海……”那条注(122)。两天后,11月10日,随着“阴的葬礼(阳葬阴(4))”那条注(124),“超级母亲”登场了(但这个词,是在次日那条注“超级妈妈,还是超级爸爸?”(125)里才说出口的)。于是,这便是葬礼的“第三联”!

我并不是有意要投入这场回顾——借着当下的势头,从葬礼三个主要侧面相继显现的角度,对它作一番回顾(依我此刻看事情的方式)。在一部长程的禅观之中,这样偶一为之的回顾,每一次都被证明极其有益:它们给出反思进程的一幅总览,同时又给出其中某些主要“成果”的一幅新视景\fnote{121}{(*)}{这种回顾,在我看来,在一份数学工作里十分罕见;我自己也是从去年春天动笔写“Pursuing~Stacks”以来才实行它。倒是有一项惯常的工作实践,对一份进行中的数学工作的观念与成果,能起类似的“新视景”效应:把所展开理论的全部观念与命题,按当下此刻的理解看来最自然的次序,从头重述一遍。这样一份乍看纯乎例行的工作,常常导致理解的实质性深化——比如,借着新的编排对内在融贯的要求,显出一些同样“自然”的、此前未被见到的观念、性质、关系等等。有时,借着显出某些假设的偶然或人为品格,或整个出发点语境的逼仄品格,“重述”的工作会通向初旨的一场意想不到的拓宽,赋予最初展开的理论以新的维度与分量。}。或许最让这场回顾的假想读者吃惊的,是我绕了这么远的一场枝蔓,而不是马上来到那场著名的“联想”(它还欠着呢)跟前,把它谈完,好终于抵达葬礼之下那条著名的“最后界线\dmark{}”;\origpage{575}那条界线,我在9月29日“葬礼悼词(2)”那条注里,是那样急于画出——彼时我刚把6月悬着的反思的挽具重新套上!况且,三天后我开始下一条注“肌肉与肚肠”时,正是抱着这样的心态:那条注以对这场联想的暗示开头,却不给它任何具体说明。

我之所以当时没有给出说明,并且一天天、一星期又一星期地把它推后——已推了一个月零十天——绝非出于什么或早或晚会显形的故意。倘若我试着探测它的缘由,我会说:我当时想必凭本能感到——甚至无须对自己说破——在我当时所处的那个点上,劈头写出那场联想,将毫无意义;那会像一份单纯的“表述”,纯粹形式的、口头的,而被那些词句所覆盖的丰富实质——它们只是凭纯粹记忆的效应来到我笔下——将依旧被忽略、不被知觉。读者倘若是个数学家(或科学家,倘若不是数学家),肯定屡次认识过这样的情势,以及它所激起的不适:人们就这样面对一个表述——可以毫不费力地确认它完全精确,其中每个词的意味我们也多少知道——却分明感到它的“意义”与实质整个儿逃过我们。这样的情势,或许在那些\textbf{非}技术性、却表达着一份被作者强烈知觉的可触实质的文本里,还要常见得多;只是有一点不同:读者相当清楚地觉察到自己所读之物的意义正逃过自己,这种情形反倒少见得多。就眼下这个个案而言,还有更多——因为连\textbf{我自己}也一样:几个月来,我早已不在那篇葬礼悼词、以及系于它的那些联想的“水里”了;几年来,我也不曾真正“潜入”阴阳的实在(虽然一路走过时步步擦过它……)——即便是\textbf{我},我当时所能写下的、用来“说出”这场联想的东西,也会是口头的东西,并不真正被感到、被知觉。迫使自己屈从于此,或者更确切地说,强制自己去做,会是一种纯粹形式的做法,一种凭良心交差的做法:总算清偿了一份义务,总而言之“封了账”,同时留心“给足分量”、不在路上丢失那场“联想”——它(我记得清清楚楚!)曾经多汁而滚烫,却早已在记忆的一角里凉透、发霉了!

\dmark{}倘若我记得的东西,确确实实要用来深化一份依然零碎\origpage{576}的领会,那么对我而言清清楚楚:这上百页的“枝蔓”我当时省不得。它们构成那场贯穿整部《收获与播种》的反思里最深的部分。

我还无从预言:我正准备顺着它们的尾流析出的那幅葬礼之景,会留给我一份圆满满足之感,还是仍会剩下一些晦暗的角落或一些不谐之音——我或许会放弃去照亮、去解决它们,至少眼下如此,或在《收获与播种》之内如此。但无论如何,一如在我的数学著作里,我知道:这上百页中的每一页,一如至此写就的《收获与播种》六百页文本(出入无几)中的每一页,都有它独一无二的位置、它的讯息与它的功能,而我一页也省不得(有没有读者肯一路跟到那里,且由他去!)。纵然所追逐的目标已远(甚至全然被忘掉……),这些页中的每一页,都带给我它自己的收成——唯有它能带给我的收成。

\subsubsection*{回顾(2)——或死结}
\tocsubsub{回顾(2)——或死结}
\label{ch18:local:h37}

\notehead{127$'$}{}
\phantomsection\label{ch18:local:n32}

(11月17日)我刚度过了相当难受的四天,身边纷扰不断。谈不上顺势前进;我在诸注上的工作,仅限于一点后勤事务:重读那部分要送去打字誊清的文本,改正已誊清的部分。每条注的“初稿”——写下一注之前重读一遍——与送去复写的定稿誊清本之间,我至少要读三遍,三遍都用心,至少在前两遍里随时调整措辞。这样下去,《收获与播种》的文本我快要熟透了!但尤其要紧的是,我做了必要的一切,以确保送去复写的文本,确实是我真正拿得出的最好的东西——包括它的形式。除葬礼的一条注之外,《收获与播种》中凡我写就并重读过的各节各注,末次阅读都给我一份圆满满足之感。我感到,每一次我都已抵达:以我力所能及的清晰与有层次,说出了我要说的话——既不隐瞒写时于我已清晰、已懂、已知的任何东西,也不隐瞒那依旧晦暗、模糊、未懂、甚至全然神秘、未知的东西……

\dmark{}唯一的例外,是10月17日那条注“一半与全——或裂缝”——自那条注起,禅观的“线”分成了两股,沿着我所命名的两个主题(在“阴阳之钥”诸注的续篇里作副题):“我们的母亲,死亡”与“拒受与接纳”\fnote{122}{(*)}{把关于阴阳的“枝蔓”诸注用副题归组,这个需要几天前才感到。这也使我重新调整了当初给这些注起的名字;因此,这些注在有的地方被引用时,所用的名字与其定名略有出入(但号数照旧无误)。与此同时,也冒出了这组注的当然之名:“阴阳之钥”。}。所指的是那条注的最后部分,即那两三页:其中我谈个人内部的分裂,说它是伴侣之内、家庭之内、人类社会之内分裂与冲突的终极根源。这份直觉,最先显现在我“出发”离开科学世界之后的最初几年,并在岁月中发展、印证、深化,直到今天。它对我已变得如此“分明”(尽管我从不曾费心把它细细检视、面面检视),以致它有点像天经地义地溜进了反思,毫不费心以某一个“抓手”把它提出,好多少显出这份“分明”。可是,倘若这几页的阅读留给我一份模糊、不满足的印象,那肯定不单是“表述”笨拙的问题。毋宁说,我感到我当时是想一步跨过一个关于这个复杂主题的实质反思——对这场反思,我分明觉得要素样样在手,可以去做,但它终究没有做!在10月25日那条注(“失乐园”(116))里——它直接系于17日那条注(由它出发展开“拒受与接纳”的主题)——我先多少试着“补救”前一条注里我注意到的缺漏;但最终并没有说出多少更多的东西,只说了这么一句:至于一场可能的“探索冲突之旅”,“眼下我毫无意愿朝那个方向走”——罢了,那只好留待来日!

\asterisms

\dmark{}在四天前那条注里,我巡了一圈葬礼图景至此透出的三个侧面、或“三联”\origpage{577}。事后我想起:在葬礼反思的进程中,已有两个时刻,我感到、并写到:我触到了冲突的“死结”。那是在“死结”与“葬礼悼词(2)——或力量与光环”两条注里。这两条注,与《收获与播种》最早几节里(表面“十分一般”的)一些反思接上了头——“他人之无谬与对己之轻蔑”(第~n\textsuperscript{o}~4节)。正是\textbf{对己之轻蔑}、对我们身上那份给我们认识与创造之力的力量的不识,它同时也是\textbf{轻蔑他人}的源头:那永世不改的补偿反射——把自己抬到他人之上,好向自己“证明”自己的价值;比如,动用那贬低或压服的可怜权力,或干脆叫人受苦、叫人受害。

写那条注时,例子当然不缺。当时在我心头最在场的,是Pierre~Deligne:我曾多次见他动用他那使人灰心、甚至羞辱人的权力,其方式常常让我感到无法解释。只是在这条注写成两个月之后,我才开始发现“葬礼的全部辉煌”——4月19日的诸注(“一个梦的回忆——或Motif的诞生”与“下葬——或新父亲”(51)(52))可为佐证。同样逐步地,我发现了我的朋友皮埃尔作为我葬礼与出殡之大司仪的角色。六月之前关于葬礼的诸注(I至X行列),大部分以他这个人为中心。他也是我掌握着无可比拟地更丰富、更个人的材料的那一位——比其他众多与会者中的任何一位都丰富。所以,我感到“触到冲突心脏”的那两个时刻,处于我注意中心的还是他——他也是唯一一个与我保持经常联系、直到今天的人。

\subsubsection*{(d) 父母——或冲突的心脏}
\tocsubsub{(d) 父母——或冲突的心脏}
\label{ch18:local:h38}

\notehead{128}{}
\phantomsection\label{ch18:local:n33}

(11月18日)昨夜睡了十二个钟头——前几夜都睡得太短,正需要!我感到把一份已开始有点散架的能量重新打足了气——比昨天更有劲,可以从头拾起那条著名的“线”了。

\dmark{}在昨天谈到的那两个时刻里,我身上有过一种“闪光”,如此清晰、如此强烈,我压根想不到去怀疑它——我是说,怀疑它向我启示了某种真实的东西、在我这个个案里外在于我个人的东西;怀疑它不是某种纯粹主观之物,(比如说)一份单纯故意的产物:想看到某套我心心念念的“心理理论”得到应用——总之,不是那只被捕蝶人天助地网进网里的“蝴蝶”\fnote{123}{(*)}{这幅意象,见“孩子与海——或信与疑”一注(n\textsuperscript{o}~103)。}!\origpage{578}怀疑这样的征象——无论在禅观中、在数学中还是在别处——便干脆是放弃我的认识与发现之力。我有幸认识这份力量;倘若有一件事是我全然信赖的,那便是它。

我可以考虑在这个“闪光”里、在它所教给我的东西里,看出葬礼图景的第四个“联”——叠加到另外三个之上(它们已在11月13日那条注里巡过)。但我一上来便看到:它与“超级父亲”和“超级母亲”两个侧面亲密地系在一起——而这层显然的联系,远远超出我朋友这个人的范围。昨天我重新唤起的这份对我们身上“认识与创造之力”的不识,不是别的,正是对我们根本统一的不识——那统一,是我们存在之中“阴”与“阳”、“女性”与“男性”的品质、能量与力量的联姻之果。因为我们身上的“男性”,单凭它自己,并不能使我们能够认识、能够创造;我们身上的“女性”,单凭它自己,也同样不能。有认识与创造之力的,不是我们存在中那虚假而可怜的\textbf{一半},而是\textbf{全}、是我们存在的\textbf{总体}。它有这份力,不是作为一场追寻、一段长路、一个生成的终点——仿佛我们在一种暂时的无能状态里跋涉,一路慢慢积攒“权力”;而是:这份力,凭我们的天性便是我们的,自我们出生之日起,便作为无偿的赠礼授予了我们\fnote{124}{(**)}{而且无疑,早在我们出生之前很久……}。

而这份“对己之轻蔑”或“对己之不识”,不是别的,也正是那对着这份赠礼的\textbf{拒绝}:对这份根本统一的拒绝,对与之不可分的力量的拒绝。或者毋宁说,它像这份拒绝那不可分的影子:它是一份\textbf{无能之知}\dmark{}\fnote{125}{(*)}{正如我在一行之后要说明的:这份认识是“搅浑的”,就其本质内容而言,它依旧无意识。然而人们常常看到它露出小小的一角(像冰山的峰顶,其底座被小心翼翼地浸在水下……),借着某种\textbf{无能的信条}——它们不止一次让我目瞪口呆。它们以一份不容回驳、不许还嘴的\textbf{断言}的口气说出;在其背后,人们感到一种激烈的、桀骜的封闭——仿佛这份被当作不可触犯、神圣“事实”来申认的无能,是最贵重的财产,任何代价也不肯让渡……},由这份拒绝所设立的无能之知\origpage{580};一份固然怯懦、搅浑、未被承担的认识——它留神停在已知的东西上(那已知,又知得何其之坏……),生怕潜得更深,生怕见识那份被藏起、被这份蓄意培植的无能所阻塞的未知力量。

在我们这个超级阳的社会里,这份对我们统一的拒绝所取的最常见形式,便是日复一日、时时刻刻地埋葬我们身上的“阴”、“女性”。这恰恰便是那个“超级母亲的联”,别名“‘女性’的葬礼与下葬”——而尤其且\textbf{首要}地,是\textbf{自己身上}的女性的葬礼。

但我分明感到:对己之轻蔑与“超级父亲的联”——别名“父亲的屠杀与下葬”——之间,同样有一层直接而深刻的联系。这份被强烈预感到的联系,正是我现在想试着圈定的。换一种说法来表述这份“预感”、这份直觉:在我们身上的分裂,与对父亲的对抗之间,必定有一层直接而深刻的联系。

当然,这份“对抗”既有机会对着生身的父亲表达,也有机会对着童年里代行其职者表达,或对着任何其他这样的人表达:他在某一时刻、因某种缘由,充当了或多或少象征性的“替补父亲”,而那些原初的对抗冲动,便投射在他身上。所以,我的用意是圈定这些对抗冲动与姿态的深层缘由——它们如此常见,以致人们有时会倾向于把它们看作普遍的;这份缘由,须比一堆具体的怨怼更深——那些怨怼,固然常常全是可触的,是人们可以对生身者怀有的。我不止一次得以确认:这些怨怼,常常更像是一份说得通而又受欢迎的\textbf{合理化}\dmark{}\origpage{581}——为一桩对抗所作的合理化,而那对抗的真正的根、其暴烈与顽韧的缘由,却在别处。

我还可以换一种方式,以它自发呈现给我的形式,表述这份我正试着圈定的直觉:那就是,我怀有一种贴身的信念:在一个是“\textbf{一}”、未曾分裂的人身上——在一个以其存在的总体接纳他自己的人身上——对父亲的冲突、或对母亲的冲突,已经解决。他是\textbf{自主}的,对他的双亲都“\textbf{自由}”。那根继续把我们系在父母身上的脐带——它系得远比童年与青春期长久(而且最常见,贯穿整个成年,直到死亡)——在他身上,已经剪断。缆绳已经解开:不久之前,它们还拽着我们,不许我们真正出发,去\textbf{我们自己的旅程}——去发现我们的母亲,那世界\fnote{126}{(*)}{奇怪的是,在法文里,“世界”“宇宙”与“太空”三个词都是阳性的。德语里相应的词——die~Welt、das~All、der~Kosmos——则分属阴、中(中性在德语里往往是一种“超阴性”)、阳三性。依我看,这更契合这些词所指之物的本性。人们谈“太空”(Kosmos)时,其意味(撇开晚近发明的空间舱与外星人不论)是一份受法则支配的\textbf{秩序}——这观念恰与阳性相应(在这一点上,两种语言一致)。相反,“世界”与“宇宙”暗示的是一个\textbf{全}的意象:我们自己与一切别物,都是它的一个\textbf{部分};而且,这还是某种有待我们去\textbf{发现}、去\textbf{穿透}、去\textbf{认识}的东西。就这些在我看来是要害的方面而言,这两个词所指之物,是“阴”的、“女性”的——尤其相对于我们而言。我实在难以看出,法文为何仍然给它们派一个阳性。
顺带指出另一桩奇怪的“反常”(?)——这一回似乎是在德语里:“太阳”与“月亮”在德语里作~die~Sonne、der~Mond,其性属恰与法文相反,而法文的那一种看上去倒更“自然”。太阳直接让人联想到热与火,那是典型的阳性之物。或许这桩“反常”在北欧诸语里是共通的,因为在那些寒冷的国度,太阳的热从不被感到炙人灼人,而是被期待为一种恩惠、一种生命之源;在那里,太阳(连同大地)被感到像一位滋养的母亲,把热量布施给造物——他们凭这热量“滋养”自己,一如凭来自大地的食物……}。

这份贴身的信念,并非“一厢情愿”之想;它不是一个愿望的投影(为应景改名为“信念”)。它的源头确凿在我的亲身经历之中,而且首先\dmark{},在我与我自己父母的关系里所能确认的东西之中\origpage{582}。我这里指的是:八年前那个转折之后的岁月里,我与我父母的关系中所发生的深刻转变——那转折,以我身上那场“阴的苏醒”为标记,继之以其后几个月里禅观的发现,再继之以两天之后与我童年的“重逢”\fnote{127}{(*)}{我生活中的这些关键情节,我在“重逢(阴的苏醒(1))”与“接纳(阴的苏醒(2))”两注(n\textsuperscript{o}~109、110)里、以及在“欲望与禅观”一节(n\textsuperscript{o}~36)里谈到。}。我意识到,这个转折打上了一份\textbf{自主}的印记——它与先前对种种现成观念、尤其被采纳的观念的依赖,恰成对照。所有这些依赖中最深的,是对我父母的依赖:他们的价值与抉择,塑造了我的价值与抉择、塑造了我自己的世界观;而他们对自身、对彼此结成的伴侣关系、对他们与孩子的关系所抱持的那幅理想化图景,我也整个“成块地”、原封不动地接了过来——几乎分毫未改。自我童年起,我便在这套价值、抉择、意象上“运转”:它们全然不是我自己生活经验的果实、不是我对此经验消化工作的果实,而只是一份“行囊”。这份行囊,好大一部分由陈词滥调与自我陶醉的幻象构成,是我凭“信任”从我父母那里接来的;它们在我一生中,屡屡取代了对周围事物的一份直接的、活的知觉、一份创造性的知觉。

诚然,我所说的这份“自主”,是随着禅观之力的发现立刻显出的。凡我留神检视的东西,它(我相信)都是全面的。但这并不妨碍:许多现成观念——尤其、首要地是那些来自我父母的——起初仍凭纯粹的惰性留在原地,因为还不曾被检视。要看的东西太多,不可能一下子全看!更不必说,经过几个月强烈的工作之后,我让自己被“继续着的生活”分了心——尤其是几段恋情,自不待言\fnote{128}{(**)}{1976年禅观发现之后的那些年头里,我的爱情生活,比我一生其他任何时期都强烈,也更动荡。它肯定代表着一种分散、一场岔开——相对于禅观最初的那股冲劲而言;那冲劲,直到1979年8月,才(以它应有的规模)被重新拾起:那是一场关于我父母生活的长程禅观。(关于它,见“表面与深度”与“写作之颂”两注(n\textsuperscript{o}~101、102)。)然而,带着回顾,我意识到:这份分散,我当时还省不得——必须有某一份激情、某一份饥饿在我身上烧尽;而一路走来,我继续借着那些我做她们情郎的女人,学习那些我在已往的生活里不曾学全的东西。到了我当时的那个地步,我怀疑单凭对这段往昔的禅观,它本会教得会我。}。在其后将近两年里,我的禅观仅限于几场应景的、分量极有限的反思——当我面临某个尖锐冲突的情势、急切地需要看个分明之时。直到1979年8月之后(禅观发现后将近三年),才开始对种种现成观念作那场“大扫除”——尤其关于我父母、关于我自己的;它们继续堵塞着我,挡住我看这个我生活于其中的迷人世界。对我父母生活的工作,吸住了我七个月,直到次年3月。那时,我五十二岁在望。正是凭这份工作,那份我谈过的自主——它在某种意义上,三年间只停留在“潜在”——变成了全然现实的、完整的、不可逆转的。也正是凭这份工作、且唯凭它,我才能够\textbf{爱}我的父母,取这个词的完整意义;也就是说:也\textbf{接纳}他们所是的东西、或他们曾是的东西,连同这一切所牵涉的(而我那时开始隐约望见的)、尤其牵涉到我——他们的儿子——的东西。

倘若我感到需要做这份工作(注~128$_{1}$),并且能够做成它,那是因为三年之前,我已懂得接纳生命的这份赠礼——它在我出生之时授予我,却被我拒绝了四十年——我自身统一的赠礼。或者说:因为我已懂得\textbf{接纳我自己的天性}。正是凭对我自己的接纳、对我自己的爱,我才能够接纳、爱我的父母\fnote{129}{(*)}{这与“接纳(阴的苏醒(2))”一注(n\textsuperscript{o}~110)结尾处的反思相接。}。

我还可以说:正是唯凭这份工作,“对父母的冲突”才得到“解决”——这份冲突,几年前我甚至不曾疑心它的存在,而那时\dmark{},我的父母双双谢世已二十多年\origpage{584}。诚然,自我婴幼年起,我对父母姿态的底色,曾是一种敬佩的尊重、一种抬举、一种毫无保留的认同;而在他们死后,是对他们这个人、对他们音容的一种默默的尊崇。这不是人们惯常用“冲突”一词所指的关系——“冲突”暗示一种对抗、敌意的底色。在这份由我这个人带给他们的抬举里,我的父母当然自得其乐:他们觉得这样极好、合乎秩序——而天下的父母,不愿处在他们位置上的、或不为此自庆的,想必寥寥!只是在这份对我父母的工作之后、尤其在随之而来的对我童年的工作之后,我才能充分地、充分知其所以然地看清:我与父母之间那场田园诗般的关系,在多大的程度上曾是\textbf{假的}、造作的,不“\textbf{真}”。它之所以能维持,只凭着在一幅感人的图画上执拗地拿橡皮擦去一大堆“不入画”的东西:包括那些难受的时期(恰恰是尖锐的对抗,常常被感到是一场\textbf{撕裂}),或那些慢性的“越轨”,它们在我母亲与我之间的关系里,以同样毫不容情的规律性重演(即便频率低些),一如当年她与我父亲之间。更不必说那些在意识层面整个逃过我知觉的事情:比如那道我在八岁上、给父母画下的“大十字”——那时我在异乡已住了两年,母亲一年三四封匆匆的来信,便是他们二人全部的生的音讯……

而那深层缘由、那份\textbf{真正的}缘由——它使我把1933年夏(五岁)到1979/80年冬(五十一岁)之间我与父母的关系称作“冲突性的”——并不在于:这四十六年间有过一些冲突,把我和他们之一、或他们二人对立起来——无论这些冲突频繁或稀少、暴烈或潜伏、有意识或无意识。而毋宁在于:这份关系不曾被\textbf{承担},也不可能被承担(照它当时那样,我是说,不经深刻的转变)。它之所以能照我所活的那样被活、照我所看的那样被看,只凭对我的认识官能与领会官能的一份恒常而顽强的\textbf{压制};凭一份\textbf{拒绝}的执拗\origpage{585}:拒绝去认识这份关系的真正本性,或至少认识它的若干要害的侧面——那些侧面,本质上牵涉着我的双亲与我本人,牵涉着我所喂养的关于我们的那幅影像。换一种说法:这份关系当时所取的形态,是凭一份执拗的、不停歇的\textbf{逃遁}而自我延续的——逃开一份十足可触的实在;而那实在,同样执拗地一次又一次要让我认识它,可在我父母生前,我从不曾真正从中汲取种子。那清晰而无可否认的、把我和他们之一对立起来的冲突的那些情节——有时是撕裂性的——只是“我与父母关系之‘冲突性’”(即这份发生\textbf{在我自己身上}的压制与逃遁)的诸多或隐或显的征兆中的几桩。

换一种说法:与他人的“冲突性”关系,取这个词的深层意义,便是那份“分裂了的”关系——它凭一个压制、逃遁实在的进程而原样自我延续,又反过来助长这些进程在自己身上延续。关系中“冲突”、“分裂”的征兆,既可以是对抗品格的,也可以是效忠品格的;既可以是批评、乃至低估或轻蔑的故意,也可以是赞许或钦慕的故意。

就这样,未经寻索、也未可预料,我又回到了人们或许会称作我的哲学“老毛病”上:人与人之间的冲突,不过是各当事人身上冲突的“征兆”;或者说:社会中冲突的“源头”,是个人身上的冲突、分裂。(而父母,在这整个论说里,终究销声匿迹、不留痕迹了!)

这幅世界之景,看似全然忽略了那幅更简单化、也远为流行的景:两人之间的冲突,是双方各自“利益”或欲望冲突的结果——那些欲望“客观地”相互对立,即一方之满足,唯有损及另一方方能达成。这是举世公认的看待方式,无论涉及两个不同个人之间的冲突,还是同一个人内部的冲突。比如(在第一种情形里),这些不相容的“欲望”,在双方身上,可以是支配的欲望、定调子的欲望、掌舵的欲望——这当然是最常见的个案之一,包括在父母与孩子之间(也全然一样,在妻子与丈夫之间、或情郎与情人之间\dmark{}\origpage{586})。当然,我并不否认这种看待方式有任何实在性、任何效用——至少在某些个案里。但我看到:它只涉及一份浮面的实在,而一份更深的实在整个儿逃过它。为在这层意义上举一个例子,我指出:支配的欲望(或炫耀的欲望,或一般地,把自己抬到他人之上的欲望),它的根恰恰在那份“对己之轻蔑”、那份“对己之不识”里——方才已谈过;人们试图以种种足以\textbf{搅浑}、\textbf{补偿}这份对自己秘密的低估的姿态与行为,从它那里逃开。这样,在对立的欲望之“客观”冲突之外,人们在这个个案里看到:个人内部的冲突显出轮廓,作为如此这般欲望的缔造者——那些欲望,除了唤起并滋养对他人的对抗之外,别无他能。

当然,凭这几句评注,我并不打算穷尽冲突的两个侧面之间那份微妙而重要的关系的问题——我想把它们称作“浮面”侧面与“深层”侧面——而这无疑也不是地方。毋宁说,我感到需要回到对父亲的冲突、对父母的冲突这个主题上——我方才正要离开它。我一度可能给人(甚至让自己也一度信了几分钟!)这样一种印象:对某位父母的冲突,同对皮埃尔或对波莱的冲突,是一码事。然而我分明知道远非如此!我分明知道:\textbf{对父亲的冲突、对母亲的冲突,正是我们自己身上冲突的心脏}。

方才我已在这个意义上谈到我的“贴身信念”(我也乐意称之为我身上的一份\textbf{认识}、一件透彻懂得的事):在一个自身未曾分裂的人身上,对父母的冲突已经解决。这份认识,我说过,首先(我相信)来自我自己与父母关系中冲突之解决的经验\fnote{130}{(*)}{关于这一点,见紧随其后的那条脚注。}。换一种说法,便是:\textbf{对我们父母的接纳}(即对父母冲突的止息)\textbf{属于对我们自己的接纳的一部分}。他们(相对于我们而言),既是我们的\textbf{本源}\dmark{},也是我们的\textbf{制约}(或至少是其中好大一部分)\origpage{587}。前者(我们的本源)与我们的这个人不可分,无论我们的行程与命运如何;后者(我们的制约)深深扎根于我们,因而一如我们的本源,属于我们的这个人。否认我们母亲或父亲的真正实在——无论这拒绝以对抗表达,还是以效忠表达——也就是否认我们自己的一份要害的实质、否认我们的一生(就我们记忆所及的而言)……

还不止于此。正是经由我们的母亲与我们的父亲,先于一切他人,那曾在他们各自身上的冲突传递到了我们身上。(这正是片刻之前那个凝练的词“我们的制约”所表达的!)他们由此与我们身上的冲突相连,比世上任何别人都更贴近。而这冲突在我们身上的第一份外部的投影、一切投影中最古老也最关键的,便是对我们母亲与父亲的冲突。于是,在我看来,我们身上的冲突,与对我们双亲各自的冲突,不可分解地系在一起——它们宛如同一个冲突。方才我表达过那份“贴身信念”:当我们身上的冲突得到解决时(或至少,当它在根子上、在“阴对阳”的分裂里得到解决时),我们对父母的冲突也同样得到解决;或者换一种说法:我们身上冲突的解决,要经过对父母冲突的解决。但我确信,反过来说同样为真:一旦对父母的冲突得到解决,我们身上的冲突也随之解决\fnote{131}{(*)}{在这里,我或许给人造成一种印象:仿佛我在摆“那个已在自己身上解决了冲突的人”的架子。千真万确,我是毫无保留地说:对我父母的冲突已彻底解决。同样真确的是:我这个人身上的冲突,仍以种种方式让人感到,它没有消失。这件事,在《收获与播种》的每一页里,想必都显而易见;而我也不止一次有机会,在这个或那个具体个案里指出它。这看来便与这条脚注所评注的那句断言——“一旦对父母的冲突得到解决,我们身上的冲突也随之解决”——相矛盾。然而,在某种意义上(即我写这几行时所着意的那种),“冲突已在我身上解决”千真万确。至少,这份冲突里的某个本质的东西、在它根子上的东西,确确实实解决了——凭着这份对我自身统一的认识、这份对我自己的接纳。倘若把冲突比作一棵根深干壮的树,那么可以说:当根被切断、或已枯干时,树已经死了,只是凭习得的惰性,树干与几条主枝还立在原地,等待慢慢干枯、慢慢崩解。我能分明感到这些年来冲突这份渐进的“干枯”——它像一份曾经强而顽的掌控,在渐渐松脱。《收获与播种》的写作,在我看来,是这个进程中的一站,是过去八年里许多站中的一站。另一幅试图描述同一实在的意象,是一份徐徐铺开的深沉的宁静——像深海的宁静,不为搅动水面的涡流所动。关于这一点,我在“重逢(阴的苏醒(1))”与“接纳(阴的苏醒(2))”两注(n\textsuperscript{o}~109、110)里有更周详的表述。}。正是由此,我在与我们父母的关系里,看到一份在我们灵性历险里的\textbf{枢纽角色}——一份独一无二的角色:它不属于任何其他亲近之人,无论是配偶、孩子,还是朋友、老师或学生。

\asterisms

\notehead{128$_{1}$}{}
\phantomsection\label{ch18:local:n34}

\dmark{}(12月1日)\fnote{132}{(*)}{本注由前一条注“父母——或冲突的心脏”(n\textsuperscript{o}~128)的一条脚注生出。}“去认识我的父母”这件事之于我的重要性,是一个梦向我启示的——那梦来到1978年10月28日。那是关于我父亲弥留的一个梦。这场弥留,在痛苦挣扎的日日夜夜中迁延,周围是他身边人忙碌的冷漠——凭所有人的默契,他已被当作“已死之人”——“这宛如一道判决,它斩断一切怀疑,把他的死变成了事实。”我醒后把梦记述了下来;但在随后的三个月里,我回避了对它的一切反思,以致让它沉入半遗忘的昏影。总之,我当时把这个梦里对我谈到的、我父亲之死给“埋葬”了;而在这个梦里(它唤起我醒觉生活的一个关键侧面),我正“埋葬”着还活着的父亲。对这个梦的讯息——它那样清晰、那样透彻,美得撼人心魄——曾有过相当强大的抗拒。抗拒是在第一夜禅观结束时消解的,那是次年1月31日,以一场对梦之意义的执拗禅观为始;其后的三周里,又有四场禅观相继。

这个梦使我懂得:我与我父亲、我母亲的关系,是一份凝定的、“死去的”关系——它与一份活的实在隔断,对那份实在的知觉被压抑着——一如(在梦里)对那场被宣告为“子虚乌有”的弥留的知觉被压抑着,连同由之而来的那个自发的行动:去救助那个痛苦地、被所有人抛弃地、为活命而挣扎的人。

\dmark{}终结我内心这份隔绝的第一件事,便是去认识我的父母\origpage{589}。当时我丝毫没有料到这项任务的规模:我幻想“几个小时”便能“直捣题心”!经由我的童年、尤其经由它来认识我自己的想法,当时连我的边儿都没擦过。这份需要是后来才感到的:它将自发地从那场我正准备启程的旅程中流出。旅程只在六个月之后才开始,1979年8月——缘于“乱伦颂”情节那场长长的枝蔓(它在许多方面绝非无用)。(关于它,见“行为”一注(113)。)

与1976年10月18日的那个梦(它触发了“重逢”)一道,这个关于我父亲弥留的梦,是对我生活进程起过最强烈作用的两个梦之一。对它的讯息的抗拒,依我看,要强大得多。前一个梦的讯息,在醒来后的几个钟头里便被接收;而后一个的,被推拒了好几个月。它只是在九个月之后才开始兑现:凭我动身去作一场发现之旅——它延续到今天……

直到最近这几天,这个梦的意义与我正在眼下这场反思里力图穿透的葬礼之实在之间的关联,才来到我心头。这场葬礼——我在其中扮演“头号亡人”——不久前曾向我显为一场“事物的回转”(见同名注(73))。这一回,我又看到一场“事物的回转”,却是从一个全然意想不到的角度。因为在葬礼里,我相继扮作“父亲”与“母亲”。我从不曾想到:我自己也曾处于一个儿子的类似境地——“埋葬”活着的(哪怕是象征地、或凭默契)他的父亲或他的母亲——恰恰相反!而且我当时有十足的理由相信相反的情形;那些缘由,我在“屠杀”那条注的结尾头一次唤起(诚然,是在\textbf{屠杀}父亲的语境里,而非下葬的语境)。(在“天真(阴与阳的联姻)”一注(107)里,我更周详地回到它。)在“屠杀”那条注里写下关于我最初童年的这两段时,我肯定给了人(甚至我自己当时也处在)这样一种印象:我与父亲的关系,一生都没有冲突。对这份关系浮面一看,也会这样暗示。但已经在这条被评注的注里——“父母——或冲突的心脏”——在那里我并不止于这类皮肉的印象,事情分明显出\dmark{}:全然不是如此\origpage{590};那幅世界之景(它确确实实是我的,直到1979年1月31日)是我在大半成年生活里乐于喂养的诸幻觉之一。这份幻觉,从我终于费心检视那个关于父亲弥留之梦的意义的那一刻起,便清晰显形——那是迄今为止生活赐予我的所有梦之中最\textbf{美}的一个。这个梦以惊人的写实,呈现了冲突加于我与父亲关系之上的掌控——它也让我活出了这份冲突的\textbf{解决}。冲突之解决,凭的是我身上与那道宣告我父亲之死的共识之间的一场\textbf{决裂}——决裂旋即为\textbf{别的东西}敞开了门;也凭我父亲的一个爱的动作,向我示意:他已听见我那被扼住的喉咙无法向他迸出的呼喊……

这个梦的亲身经历——它是那份凝定的父母关系(它旋即重新活了过来……)的惊人寓言——与我已探测了将近九个月的葬礼之实在之间的深刻亲缘,此刻以一份明证之力呈现在我眼前。惹人注意的是:在这整场漫长的反思里,直到最近几天,这份亲缘的念头竟不曾擦过我的心。我是凭最大的偶然才“撞上”它的:当时我打算在一条脚注里,顺便指出某一个\textbf{梦}——在八年来宛如天助灯塔照着我路的众多梦之中——又一次(在对父母的反思的启动中)所扮演的角色。这个打算的效应,是把我与这个梦的亲身经历及其实质稍稍重新接上——我离穷尽它们还远得很。一旦这份联系重新接上,在此语境里,它与葬礼的亲缘便不可能不显形。

诚然,这份亲缘,眼下只涉及某个“死结”而已;而在这个梦里、在它所转录的实在里,既有死结,也有它的解决。何况那份解决——梦曾让我活出它,当夜我便知道了它的滋味与力量——要使它也成为一个在我醒觉生活里、在我与父亲及母亲的关系里被活出的实在,那是我、唯有我的本分。我可以去做,也可以不做——而有几个月,我选的恰是第二条路!今天——那场解决之后五年——在这份某种意义上对称的情势里,情形肯定依然如此:我在其中牵涉,而我扮的是被一道共识-判决所下葬的父亲——想当年,我可是那个恭恭敬敬活埋自己血肉之父的儿子\dmark{}\origpage{591}!而或许这一回,仍将是一场对我的亲身经历之意义的禅观——在此即对这场葬礼之意义的禅观——将解开我身陷的这个新的死结,并或许再消融我过去之重的一份余量。

至于这场禅观会不会对我之外的什么人有用处——对这场葬礼的某位当事人也许(在这场葬礼里,被下葬的不止我一个,而赶赴仪式的下葬者更是不计其数)——这轮不到我操心;我看到的别人身上的这个或那个死结解不解脱,也轮不到我操心。那是他的本分,我的本分已经够多了!但倘若趁我在世,它果真解脱了,我肯定是头一批得到消息的人之一,而且我会为此高兴……

\subsubsection*{(e) 敌对的父亲(3)——或阳葬阳}
\tocsubsub{(e) 敌对的父亲(3)——或阳葬阳}
\label{ch18:local:h39}

\notehead{129}{}
\phantomsection\label{ch18:local:n35}

显然,在前面的几页里\fnote{133}{(*)}{即注~n\textsuperscript{o}~128的几页,本注是它的直接延续。},我对\textbf{对父母的冲突}这个主题,不过刚刚擦了个边;甚至连对父亲的冲突——我的出发点——也谈不上。我从那里出发所追随的种种联想,似乎把我带离了它,而非深入它。在我方才关于对父母之冲突所说的话里,母亲与父亲的角色可以互换;同样无差别的是,这几页里所说的“我们”,指一个男人还是一个女人。然而,在我们与父母的关系里,母亲与父亲扮演的角色远非对称;而他们各自扮演的角色,关键地取决于“我们”是男孩还是女孩(后来长成男人,或女人)。

在这个个案里,对父亲的冲突(它以对父亲的象征性下葬、甚或以对他的屠杀来表达),首先令我感兴趣的,是那些我知道曾积极参与我下葬的人的个案——他们\textbf{全是男人}。于是,在我之结构的历程中,父亲是那个人们\textbf{与之认同}的人、那个人们\textbf{以他为模型}的人——在他的与他人的关系里(尤其与女人的关系里),也在他与自己的关系里。这份认同\dmark{}很少有不带一份相当大的“越轨”便能完成,而对父亲的对抗\origpage{592},便是它最顽的痕迹之一。此处不打算遍数这些偏差,以及哪怕最愿意以爸爸为榜样的小男孩也常会遇到的种种不顺;也不打算探究它们在与父亲的关系中往往如何表现。况且,我在这方面的亲身经验,实在太不合典型,以致做这样一份清单,我或许是最不称职的人:凭我自己的亲身经历,我并不贴身感到那些主要个案中任何一个的来龙去脉与特殊的“滋味”\fnote{134}{(*)}{可与“屠杀”一注(n\textsuperscript{o}~87)结尾处的反思对读。}。我这里的经验,主要是间接的:凭我对周围的观察,而首先,凭我的孩子们与我的关系。

在“越轨”的特殊品格之外、在由此汲取的对父亲的怨怼与积愤之外,却有一个共同的侧面,我在许多场合强烈地知觉过——那时,任何“解释性”的故意全然缺席。这便是:男孩或男人对父亲的对抗——父亲多少做过他的模型,而他正复现着这个模型,以“正”的或“负”的方式(凭模仿,或凭反着来),无论他愿不愿承认——这份对抗不是别的,正是一份对着\textbf{他自己}的对抗的格外雄辩、格外关键的侧面。更确切地说:经由对父亲的(或多或少明白说出的)拒斥,它是\textbf{他对自己一部分的拒斥}的外部征兆——是那样一些东西的征兆:凭着它们(在他不知不觉中,或违背他某些有意识或无意识的抉择),他像极了他所否认的那个模型——他的父亲。

这样一来,我稳稳落回了双脚——我看到那份被预感到的联系清晰起来:在“对己之轻蔑”(或“对己的拒斥(或不识)”)与“对父亲的对抗”之间——但我是从一个意想不到的一侧落回来的。我正准备在对父亲的对抗与自拒之间,找一层或多或少直接的联系——自拒取\textbf{他自己身上}女性之拒斥(或“下葬”)的形式。结果却似乎是:我落在了对\textbf{男性}的拒斥上——按“好逻辑”,这本该在我意料之中!然而我分明知道:这份拒斥,在男人身上不如他身上女性之拒斥那么显眼、藏得更深(我主要有机会谈的是后者),却几乎同样不罕见,压在他身上的分量,同样沉甸甸。它常常叠加在另一份之上,以致无论一个人的我以什么方式结构——以阴的颜色还是以阳的颜色——都注定不容于他自己!或者换一种说法:这份对父亲的拒斥、或对自身之中使我们像父亲的“男性”“阳刚”之物的拒斥,常常\textbf{与}毫无保留地采纳一套“阳”的、“大男子气”到底的价值系统\textbf{结伴而行}(既然“阴”的配重已被否认)\fnote{135}{(*)}{(11月29日)至少,就我所知的那些个案而言,这是远为常见的情形。}。

一个念头来到我心头:这份矛盾(一旦白纸黑字说出、写出,着实骇人!)无疑也是那场毫不容情的\textbf{竞争}里真正的\textbf{神经}——那竞争,是我们这个超级大男子气社会的特征之一(无论在科学的高层,还是在任何别的地方……)。因为,倘若“往上爬”与“超越”是十足十的超级阳价值,这些价值却不至于以这样的烈度被内化、其付诸实践也不至于以这样的暴烈进行(哪怕在“高层”里,那暴烈是压低了嗓门的……)——倘若在那位占了上风的对手身上——我们正要超越乃至排挤的对手——我们没有同时看见“父亲”那副可畏的影子显出轮廓:他既受钦慕又遭嫉妒,还被秘密地仇恨;他先我们而在,而他的存在本身,就我们记忆所及,一直是我们生活里的\textbf{大挑战}。

\subsubsection*{(f) 箭与浪}
\tocsubsub{(f) 箭与浪}
\label{ch18:local:h40}

\notehead{130}{}
\phantomsection\label{ch18:local:n36}

(11月19日)我又感到急不可耐,要接着把反思继续下去。老实说,已经一个星期了(自11月12日那条注“激烈泼辣的妻子(阴阳的颠倒)”(126)起),我天天感到自己就要进入“题心”了——就要来到我应承过自己的那幅葬礼总图跟前:它将把反思进程中析出的那些局部“联”统一起来\dmark{}——也一个星期了,那个“点”却一天天被往后推\origpage{594}。每天写完注(时候不早了,总得停下上床),我分明感到自己做成了一份非做不可的工作,“前进”了一步——却同时觉得,我想抵达的那个“点”也退了同样远!这里明显的诱惑,是一口气写下去,直到抵达那著名的“题心”。但经过这三年里的几次“健康事故”,我也分明知道:这正是该避开的失着。

况且,在心底我分明知道:我已经全身在那个“题心”里了。只是,我急不可耐,恨不得立刻把它整个摸透。这份抵达任务尽头的焦躁,这股朝着某个“点”、某个“题心”的冲劲——它在前方被我强烈地知觉着,或近或远,反正一样——这份“目标”对我的吸引,把我向前抛射,如一支直扑靶心的箭——这个在我看来最强烈地“\textbf{阳}”的侧面,刻画着我的\textbf{工作之外}的存在方式。它是“\textbf{东家}”的一个惹眼的侧面,是我身上受制约的、习得的东西的侧面。在我幼年所知的一切里,没有任何东西能预示这个品格——它在我的童年后期才显出,却如此强烈地打下了我整个成年生活直到今天的印记。

而在工作本身之中,这个侧面似乎几乎消失。我有这样一种印象:它残存的些许痕迹,此地彼处,不多不少,正是“东家”在工作中偶一为之的、不显眼的插手之标志(说实在的,他在这里无事可做!)。工作本身,顺着那匠人——他经我的双手,按他自己的节拍劳作——的心意,循着一股全然不同的气息进行。那急躁的烈性,退到一份宁静、平和而执拗的东西面前。不再有箭,急匆匆扑向靶心;只有一道浪,铺向极远,不知往哪里推进——那推动它的、活动的力量把它带向哪里,它便去哪里——一道浪跟着一道浪,再跟着又一道……这个运动里没有迟疑:在每一处、每一刻,都有一个它自己的方向托着它、把它往前拉。每一刻都有一场推进——说不上朝着什么;有一种“工作”在做着,在一种不知努力为何物的运行之中——却没有目标。“目标”这个念头本身,在这里显得出奇地乖谬——人们想把它放在哪里呢?!目标消失了,一如\dmark{}箭消失了。\origpage{595}倘若还有箭,那也不是一支颤动的箭,呼啸着扎进靶心、插在靶上折损自己——而是在这片相互追逐的浪的活动之体的\textbf{每一}处,都有一个毫不含糊的运动与力量,都有一场推进之中的方向,其精确与分明,不亚于一支箭——一支看不见却威严的箭,标示着这方向、这力量、这运动。

这样,依我看,在我的工作里,我同样“阴”,同样“海与运行”,达到了可能的极致。我相信,我一生中一切发现的工作、一切我满怀激情投身其中的工作,都是如此——首先,是我的数学工作与禅观的工作。而今,这幅意象不期而至,威严而骤然,描绘出我如何感受这份工作;我觉得,这幅意象同时也描绘了\textbf{我生活的运动}——自与我自己的重逢之日起,或许更早,或许早自我那场“有益的拔离”——拔离一个安乐的巢\fnote{136}{(*)}{见同名注(n\textsuperscript{o}~42)。}。至少,它描绘了我生活深层处的“如何”——即我昨天那条注的一条脚注里(就在几个小时之前)谈到的那个“宁静”:一份不为水面上的纷扰所动的宁静。在这份深沉的宁静里,有运动、有推进,却没有目标——目标消失了。

而今我也记起:3月份,正是这同一幅意象来到我心头,当时我谈到我那两种激情——禅观与数学——的显现,说它们像“一道道相互追逐的浪的高低起伏,像一份宽阔而平和的呼吸的气息……”\fnote{137}{(**)}{见“我的激情”一节(n\textsuperscript{o}~35)的结尾,这几行便摘自那里。}。八个月后的今天,我相信在这些意象里认出了我存在的自发运动——在它最自发的东西里、在我身上真正原初的东西里——在那个渴望认识的孩子身上,在“显得体面”的操心、与“求变成”的饥渴还未触及他之前……

\subsubsection*{(g) 冲突之谜}
\tocsubsub{(g) 冲突之谜}
\label{ch18:local:h41}

\notehead{131}{}
\phantomsection\label{ch18:local:n37}

\dmark{}(11月20日)昨天几乎整个晚上,都花在重读前天的注上\origpage{596}:一路修改,重打一页实在写得太挤的纸,写脚注(前天就已打算好的)——而已经午夜了!可我还急着当晚再往前走一点,哪怕一点点;于是又坐回打字机前,接上昨晚中断的“线”。接着,来的却是全然另一件事——箭与浪的意象。长久以来,我在箭的意象里认出自己,而浪的意象,在我看来说的是一种与我大相径庭的气质。这正是这场关于阴阳的反思进程中显出的意外之一:恰恰是浪的意象,以最触目、最贴切的方式,表达了我的存在里占上风的那个“基调”——当“东家”不在场、或至少退到别的东西后面时。意象升了起来,仿佛它早已在那儿备好,只等着那些终于使它成形的词。词不慌不忙、毫不犹豫地来了,而我所做的,不过是尽力\textbf{描述}——尽可能忠实,不省略、不歪曲——那尚处于弥散之感状态的东西。描述完成,已近凌晨两点。当夜我重读了这两页,几乎无须修改。最棘手的一段,是我试着描述那份对“箭”之连续无限的直觉——它们合拢起来,宛如一个力的“场”。这是一个强有力的观念,却不肯轻易让语言把它唤出。可我感到,这是整幅意象的一个重要侧面:“阴中之阳”的侧面。浪里有“箭”,有一股托着它向前的\textbf{冲劲}——循着一份它特有的运行,而不是\textbf{一支}箭的运行,毋宁说是整整一簇、一\textbf{连续}簇的运行,它柔韧地复现着浪的那份运行。我也分明知道:在我的工作里,我\textbf{同样}是“箭”;只是我之为箭,循着一种与我迄今想象不同的方式——因为我从不曾费心以某种注意去看这份工作、让自己浸入它,仿佛它是另一个人,以便知觉它特有的调门。八年来我一直在禅观,却至今不曾这样做;这无疑是因为,我在不知不觉中,始终是一个积习已深的故意的囚徒:认同于我身上的“\textbf{东家}”,而不是认同于那个匠人-孩子;也就是说,当我说“我”时,\origpage{597}首先(甚至常常排他地)想到的,是那个“东家”在前台时的那个我。而那些时刻,差不多也就是我工作之外的时刻。

教学的必要与波折(以及其他缘故),终于在禅观发现之后,把我的注意力引向我工作的\textbf{某些}特征——即那些我感到具有普遍品格的特征:它们应当存在于一切创造性的工作、一切发现的工作之中\fnote{138}{(*)}{我想,头一次写出、并唤起这些特征中某些特征的文本,是1978年10月的那份“以代纲领”(11月6日那条注“陌生的美人”(n\textsuperscript{o}~130[原书如此,当为120])里暗示过它)。那份文本之后,我没有费心把这方面的观察白纸黑字显豁、深化,直到今年《收获与播种》的反思。它的最初八节,基本上就献给这个主题,还不算反思进程中此地彼处的许多别的评注。}。但在这场阴阳反思之前,我还不曾想到去辨认我自己的工作里那些独特的特征——它们使我的工作区别于任何别人的。这些特征之一,依我看是所有特征中最要紧的,终于在11月8日那条注“上涨的海……”(122)里得到圈定。那条注里最初唤起的意象——在一桩有待证明的猜想这个典型语境里——在昨天那两条注里被重拾,却换了一道光照,脱开一切特殊语境。

我终于接上反思的线,从前天中断的地方续起。我起笔时\fnote{139}{(**)}{见“父母——或冲突的心脏”一注(n\textsuperscript{o}~128)。},本打算试着圈定对父亲之对抗的深层缘由——越过人们对父亲可能怀有的那些具体怨怼。顺着强劲涌现的种种联想,我离开了这个打算,主要被引去谈对父母的冲突——父亲或母亲,不分彼此。这份“冲突”,既可取效忠的形式(我的情形如此),也可取对抗的形式。自对我父母生活的工作以来,这份“\textbf{对父母的冲突}”向我显为真正\dmark{}“处在冲突的心脏”——我们身上的冲突的心脏。解决后者,我确信,不多不少,正是解决对父母的冲突,也就是说:\origpage{598}从他们那里自由,在灵性上充分自主,走\textbf{自己的}旅程……

回到男人身上对父亲的对抗,我重新接上了一份直觉——近年来它屡屡强加于我:这份对父亲之对抗的深层意义,是对我们身上使我们像父亲的那些东西的拒斥——是对我们这个人“\textbf{阳刚}”的方面与特征的拒斥。我把昨天反思(其实\fnote{140}{(*)}{实际上,这不是昨天那条注,而是前天那条注——我正准备接上的正是它。}是前天的)的最后这部分,单独抽出作一条注,命名为“敌对的父亲(3)——或阳葬阳”——这名字同时也暗示它与“敌对的父亲(1)(2)”两节(n\textsuperscript{o}~29、30)的联系:那“敌对的父亲”的主题,正是在那里头一次出现。

这样,前天反思开头所谈的葬礼的那个侧面——“对己之轻蔑”、或“对己之不识”、或“自拒”——显为一种联接的痕迹,或更好地说,一道“\textbf{铰链}”:铰接在前面的两联之间——“超级母亲——或‘女性’的下葬”之联,与“超级父亲——或父亲的屠杀与下葬”之联。这铰链的品格,一经清楚地知觉到下述两点,便显出来:在前一联里,“女性”在一切之先,是“\textbf{我们身上}的女性”(11月10日那条注“阴的葬礼(阳葬阴(4))”——“超级母亲”之联登场的那条——里,确实已被知觉);而且,“父亲”首先是“我们身上的男性”的象征性替补。于是,这两个相关侧面,显为完全对称的两联,对应着“自拒”的两个显然的“图式”:拒斥我们身上的“女人”(别名母亲),拒斥我们身上的“男人”(别名父亲)\fnote{141}{(**)}{我提醒一句:这两种“对称的”拒斥,叠在同一人身上的情形,绝非罕见。鉴于我们社会对阴的贬抑,无论如何,阴的拒斥以或多或少显著的形式在场的情形,想必相当常见。所以,我倾向于在对父亲的对抗里,看到一份对阴\textbf{与}阳的双重拒斥的(至少是盖然性的)征象。}。而对父母的冲突这个主题\dmark{}——它是对母亲的冲突与对父亲的冲突这两个不同主题的某种合取或叠加——也显为一道铰链。\origpage{599}或者更好地说,一如昨天反思里所看到的\fnote{142}{(*)}{见倒数第二条脚注。},这个主题显为与自拒的主题不可分:两者是同一份未分裂的实在——\textbf{我们身上的冲突}——的两个不同侧面。

在这一切之中,最初的那个打算——“圈定对父亲之对抗的深层\textbf{缘由}”——似乎始终悬着。我可以说:对父亲的对抗,是对抗自己、自拒所取的\textbf{形式}之一。于是,最初的问题仿佛一分为二。一方面,在哪些“缘由”之下,自拒在某些个案里取这份特殊的形式?探测它,也就或多或少要进入若干不同类型的典型情势——那些易于引发这份对抗的情势。

另一方面,我们回到那个更深、也更关键的问题:自拒的“\textbf{缘由}”,也就是我们身上冲突、分裂的缘由。我相信自己至少把握住了那份共同的\textbf{机制},冲突凭它代代相传:我们对自身的拒斥,不是别的,正是我们最初几年里周围环境对我们的拒斥之内化——至少,是我们身上某些侧面与某些冲动之被拒斥的内化;那些侧面与冲动,构成我们原初存在、我们创造官能的一个本质部分。我在“阴阳之钥”的“拒受与接纳”部分里触及事情的这方面(以及其他),尤其在最初两条注“失乐园”与“循环”(116)(116$'$)里。

把握住冲突传递的这份共同“机制”,却丝毫并不等于:已经懂得我们身上、并(经由我们)在人类社会里冲突的\textbf{缘由}。\textbf{为什么},古往今来、无远弗届(凭着世代传到我们手里的一致证词),“社会”竟不能容忍它的成员是一些\textbf{完整}的存在?也就是说,是一些充分\dmark{}握有\origpage{600}自己创造官能的存在——他们不必花费高昂的代价,压制自己之所是的一部分:那份被看得如此可耻(或如此可畏……)的部分,以致最好假装不知道它的存在,并默契地判定它\textbf{不存在}……

这在我,是存在的伟大奥秘之一,或许是最伟大的奥秘\fnote{143}{(*)}{这份提议纯属主观,它只反映这样一件事实:在“存在的伟大奥秘”之中,正是这一个,我以格外强烈的方式感到——那方式,超出单纯的智性好奇。它是唯一一个在我身上唤起一份\textbf{欲望}的:探测它、认识它、认识它的“最终答案”(在我这有限的官能所能认识的限度之内)。区别正如数学里的情形:一边是我“很有感觉”的未解决问题(我可以当场就投身其中),另一边是我按术语的意义“理解”的问题——我(浮面地)知觉其分量,它们却“不叫我或热或冷”。Riemann假设属于后者(无疑归因于我对解析数论的极端无知);“费马大定理”直到几年前也属于后者。是我的“阿那贝尔式”反思改变了我对后者的心态,尽管我对它所催生的那些工作,无知一如从前。}。

曾有一段时光——几年前还如此——我对压制与冲突这一普遍实在的姿态,是一种战斗的\textbf{反叛}:反叛那把“\textbf{利剑}”——它声称要把凭本性当为一、\textbf{本}为一的东西剖成两半。五年前我写《颂》时,我的心态犹是如此\fnote{144}{(**)}{这一情节,《收获与播种》里谈到过几次,最后一次是在“行为”一注(n\textsuperscript{o}~113)里。}。是其后那场关于我父母生活的长程禅观工作,改变了这份心态。凭这份工作——它日复一日让我贴身接触我父母身上冲突的种种显现,并让我耐心地由显现上溯其意义与缘由——我终于感到了冲突的\textbf{奥秘}。反叛的姿态消失了,仿佛从不曾有过。它曾是一种皮肉的反应,一份单纯的能量耗散。一场反叛——反谁?不是反某个人、或某群人,不是反那著名的“他们……”!我们全在同一条船上,而人们在这儿已经待了一百万年或两百万年了……反“上帝”反叛?就差这个了。

\dmark{}在心底,我分明早就知道(我甚至说不出从何时起,尽管有很长一段时间我装作不知道……):这世上的万事万物,都有它存在的充分理由;甚至,倘若人们懂得事物的底里,万事万物就其本身而言,肯定都是\textbf{好}的\origpage{601}。死亡与死亡的“彼岸”(倘若真有这样一个彼岸),也属于这些事物。这是一个奥秘;倘若我在这件事上有什么“\textbf{信}”,它绝不在关于彼岸之存在(或不存在)及其种种细节的“信条”里,而只在于这份单纯的确信:事物如其所是,已然圆满——包括一切涉及死亡的,也包括一切涉及诞生的,而诞生同样神秘。

许久以来,我却把“冲突”排除在这些事物之外——我把它当作一种“越轨”、一块不可容忍的污渍、创造的大合唱里一个顽固而乖谬(乃至令人义愤)的“走音”。只消我终于贴身去认识冲突,而不是浪费自己去装模作样地与它搏斗,我与它的关系便深刻地转变了。

死亡与“死后”的奥秘、诞生与“生前”的奥秘,并非我们这个物种所独有。它们所引出的问题,对一切活物都有意义——或许对一切事物都有意义,从电子到星云。冲突的奥秘,却依我看为人所独有,为人类这个物种所独有\fnote{145}{(*)}{(12月3日)人们或许会(有理地)反驳我:冲突——以个体之间或个体群之间的攻击性与交锋的形式——也存在于我们之外的其他物种之内。当我在这里谈“冲突”时,我想的是它在人类社会里所取的那种特殊形式,尤其是它与个人内部的\textbf{分裂}与\textbf{压制}之间的深刻联系——对其存在之大半的压制,尤其对其知觉实在之手段的压制、以及知觉本身的压制。种种形式的压制,在我看来,都扎根于那在我看来最要紧的压制之中:所谓“性”的压制——它灌输对自己身体的羞耻、对身体的功能与冲动的羞耻(或至少,对其中某些功能与冲动)。就我所知,这些机制是人类物种之外所不知道的。我或许不该把“冲突”“分裂”“压制”三个词当作近乎同义词来用,或至少当作指称同一实在之不同侧面的词。关于“冲突”一词对我所取的意义,我在“父母——或冲突的心脏”一注(n\textsuperscript{o}~128)里略有说明。}。它向我显为\textbf{那个}大奥秘:关乎\textbf{我们这个物种}的特殊意义、特殊命运的奥秘。民族学家与心理学家所给出的那些“解释”\dmark{}——至少我所听过的那些——分明不过是些\textbf{合理化}\origpage{602}:为那被承受、被内化的压制辩护,说它是社会良好运转乃至社会存在本身所不可少;一如在一群残臂者或独腿者的社会里,必不缺几位显赫的理论家,摆出一套看似无可辩驳的论证(无人想到反驳),证明:一个人们可以使用双臂(或双腿)的社会,无论如何也运转不了\fnote{146}{(*)}{一如奴隶制社会的时代,对“最优秀的头脑”(他们同样役使奴隶)一如对其他人,没有奴隶便没有社会,乃天经地义。据说,柏拉图非得“有幸”意外沦为奴隶,才开始换一副眼光看事情。}。这里说的是一些一望即穿的辩护,它们竭力以“科学”面孔的解释,把一个奥秘变没了。实际上,人类社会之内冲突(或压制)之起源与意义的问题,仍是纯粹的修辞——只要装模作样提出这问题的人,不曾经历一场强烈而深入的、对\textbf{他自己身上}冲突、对\textbf{他身上}冲突之起源的认识工作。没有这样一份自我认识,这问题(一如关于自由之本性、爱之本性、创造性之本性那些问题)便是中世纪那著名的“天使的性别”问题的现代翻版——不过是一场文体练习,好把反正得塞进去的东西塞进去。这个问题,严格说来,并不是一个“科学”问题——不是一个其检视不以一份\textbf{成熟}为前提、而只以某份预备的知识、某个智性力量或灵敏的程度为前提的问题\fnote{147}{(**)}{(12月3日)冲突之意义的问题不属于科学的管辖,这或许使人期待:人们能在神话与宗教里找到答案的要素。可依我看,全然不是如此。就我所知,神话与宗教的一个要害功能——若不说其主要功能——便是设立一份“法”:那法,就其实质而言,是一包禁令,压制便在一个个特殊的社会里凭它而落实。这法被宣称为神圣本质,便无须自我辩护,无须解释自己的“意义”,更无须解释它与其他社会的“法”所共有的意义。}。

在这个个案里,我要做的不是猜谜式地设想:压制在人类社会里是凭哪些机制设立的\dmark{}——也就是说,不是要为压制这个事实找一份\textbf{解释}\origpage{603}。即便人们达成了一个说得通的、乃至有说服力的脚本,我也不会因此觉得自己前进多少。它或许会照亮奥秘的某个有趣的侧面——“机械的”侧面——却穿透不了它。正如古生物学与分子生物学的周详成果、甚至达尔文那些深刻的思想,也真正穿透不了生命之出现、及其在三四十亿年里在地球上创造性绽放的那个奥秘。在冲突的奥秘里,使我感兴趣的,不是机械的、科学的侧面——一个\textbf{外在于我这个人}的侧面,一如那著名的“费马大定理”——而是冲突之\textbf{意义}的问题。这份意义\textbf{与我直接相关}、本质地相关,一如它与无数代以来相互撕扯、相互屠戮的无数男人与女人每个人直接相关,他们把从父母那里接来的冲突,又传给了自己的孩子。

冲突必定有一份\textbf{意义}、而我能够多少认识这份意义——这肯定属于我方才谈到的那份“信”。对我是显而易见的事——而那份熟得不能再熟的“奥秘之感”:此处有某种深的东西可探测——同时也告诉我:这“某种东西”恰恰\textbf{便是这意义}。所说的那份“信”,与我对自身官能的信彼此重合:当我的官能在此毫无一丝怀疑地向我启示,我面前有一个待发现的“意义”时。

或许有一天,这份意义将变得显然,仿佛我一直认识它!这个奥秘在我看来丝毫不遥远、不可及。它以切近之物的样子呈现给我:更贴身地认识它,只在我一念之间。而且肯定地,我此刻已望见一条接近它的路,或毋宁说一个侧面——它已经在向我友善地招手。毕竟,冲突有许多要教给我的东西,而且它已经教了我许多……

\subsubsection*{(h) 颠倒(2)——或暧昧的反叛}
\tocsubsub{(h) 颠倒(2)——或暧昧的反叛}
\label{ch18:local:h42}

\notehead{132}{}
\phantomsection\label{ch18:local:n38}

(11月22日)这已是连续两条注,我眼看自己登程于十足计划外的远足——这一回\dmark{},我留神头一件事就从那\textbf{预定}的东西开始,破例一回\origpage{604}。我想检视前一条注里唤起(未加说明)的“类型情势”之一:易于引发对父亲之对抗的情势,而更深地,引发对自身阳刚特征的(或多或少彻底的)拒斥的情势——这份拒斥在对父亲的拒斥中找到它的象征表达。那份情势,早在11月18日的反思里——以“敌对的父亲(3)——或阳葬阳”那条注收尾的反思——我就记起了。我当时的心愿,是在这个“类型情势”里(至少),把手指按在\textbf{对男性之拒斥与对女性之拒斥之间的一层直接联系}上。

离我最近的、我又曾长期下过工夫的个案,是我母亲的个案。她一生都怡然于一份几乎不加掩饰的、对一切女性之物的轻蔑;她以过度的男性价值为自己塑形;而与此同时,自她的青春期起,与男人的关系,对她一直是一份“彻头彻尾”对抗的关系\fnote{148}{(*)}{与对女性之轻蔑相反,这份透过一份激烈而动荡的感情生活透出来的、彻头彻尾的对抗,终生无意识。我只是在我1979年8月至1980年3月的工作中才察觉。}。我有这份大运气:我母亲曾极自由地同我谈起她自童年起的生活;我又掌握着极详尽的、直到她与我父亲共同生活最初年头的自传笔记,还有卷帙浩繁的通信。这是——加上与她接触时我自己的亲身经历所归还给我的——一份异常丰富的材料;况且我远未穷尽它。但我对它下的工夫已足够,使我无可置疑地感到:她身上那份我刚唤起的双重拒斥——对女性的拒斥、对男人的对抗——根在一份与父亲撕裂的关系里。她父亲,在许多方面是个讨人喜欢的人,慷慨、正直、重情,却在战后德国(我说的是1914--18那一战之后)一场漫长的社会坠落里变得刻毒——那样的坠落,当时何其多。老实说,这坠落开始得更早:从一位高车驷马的富人的地位,一路跌到走街串巷的擦鞋匠。在忧愁与失意的刺棒之下,他那易怒的脾性有时变成家庭暴政,而他那位体弱多病的妻子,首当其冲\dmark{}\origpage{605}。我母亲,对父亲一如对母亲依恋至深,被这些父权的暴政场面所激怒——那些场面,她母亲默默承受,有时实在撑不住了,却从不诉苦。这孩子激情地与母亲认同——那父权专横的牺牲品;而与此同时,她母亲所扮演的角色(牺牲品的角色、被动的角色——“女人的角色”……)在她看来难以容忍。一边是与母亲的认同,表达为一场反叛、一份对父亲彻头彻尾的对抗;\textbf{与此同时},又有那“我永远不做她那样的人(那个不反抗而忍受的人)”的跃起——这跃起必然同时意味着:“我永远不做女人们那样的人。”

但更深地,还有对那份权力的嫉妒——父亲的、男人的权力,它使他得以随心所欲地支配。我母亲的一生,被这份吞噬一切的支配欲所支配、所摧残;而首先,是支配并摧毁\textbf{男人}的欲望——恰恰是那在她身上唤起如此狂怒反叛之跃起的人、那凭其天性本该支配她的人——一如她父亲曾支配过她母亲:她默默承受着,脸色苍白,全然无力。

我本要在这里写:这份反思,如今“接上”了11月12日那条注“激烈泼辣的妻子(阴阳的颠倒)”(126)里进行的反思。由于对那条注的记忆已不很清晰,我刚刚重读了它。奇怪的是,我忘了那条注(一如今天这条)也是由我母亲的“个案”引发的。十天前,我对稍许展开这个个案还心存迟疑。今天我卷土重来、越过了这份迟疑(它后来我也忘了!);这无疑是因为:在被检视的情势里,有一个侧面当时还搅浑着。我也忘了:今天这条注的出发点——“把手指按在……对男性之拒斥与对女性之拒斥之间的一层直接联系上的心愿”——本就是十天前那场反思最初的动因:它自然接续前夜那条注“超级妈妈,还是超级爸爸?”(125)结尾处的那份追问。实际上,12日那场反思的最后一句:

“无须更多,便可看到那‘缺失的环节’显形……”,似乎是说,我当时自信当日之责(建立\dmark{}这样一层联系)已经完成。倘若我全然忘了自己早已揭示过\origpage{606}这层联系、甚至早在四天前那条注(今天的反思正接续它)之前就已提出过这个问题——那无疑是因为:我对刚才引用的那份辉煌结论,尚未全然信服——它写成于“敌对的父亲(3)——或阳葬阳”那条注之前不过六天。引全那句话,情形便更清楚:

“无须更多,便可看到那‘缺失的环节’显形:它连接着对‘超级父亲’的对抗(它在其下葬中找到象征表达)与对‘女性’的轻蔑、拒斥,而更深地,与对自己身上‘女人’的否认(它或许将在一位‘超级母亲’的‘葬礼’的象征中找到表达,在一大堆滥情的、一语双关的颂词之下……)。”

这份结论里缺了一步,使它显得匆忙:那便是“对超级父亲的对抗”与对“男性”的拒斥之间的联系——这层联系,直到11月18日那条引注“敌对的父亲(3)——或阳葬阳”里才在反思中显现。对父亲的对抗,那时向我显为一份远更要害的实在的象征表达:对自己身上阳一侧、“男性”一侧的拒斥。在对女性之拒斥的“对称”个案里,象征表达与其深层意义之间的联系,早在“超级母亲之联”显形时——11月10日那条注“阴的葬礼(阳葬阴(4))”(124)——便已被知觉。于是,11日那条注“超级妈妈,还是超级爸爸?”里显出的两“联”相对——父亲的下葬与母亲的下葬——前天才被看作自拒(或对己之轻蔑)的对称的显现,取\textbf{对自身男性之拒斥与对自身女性之拒斥}的双重面目。

况且,在18日那条注“敌对的父亲(3)——或阳葬阳”里,我把自己限于“主语”为\textbf{男人}的个案——然而,我所知道的最极端的个案,却是我母亲的!她在那份反思里被整个忘掉了,甚至已被忘了十天(若不算在11月17日那条注里,她藏在“我父母”这个词底下)。

正是我对我的孩子们、及其与我的关系的认识,四天前使我感到对父亲的对抗与自己身上对男性之拒斥之间的一层联系\dmark{}\origpage{607}。老实说,对我的五个孩子中我有机会相当贴身认识的那四个,近年来我不止一次感到:在他们对我——他们的父亲——积习已深的对抗姿态背后,有一份对他们自身阳刚一侧的拒斥,而尤其,是对他们身上那股冲劲的拒斥——它把他们掷向世界,却使他们像一个被否认的父亲!我从不曾问过自己这是不是一条普遍的事实;或者毋宁说,我身上有一种未说出口的预设:事情想必如此——但在四天前的反思之前,我从不曾感到需要把这件事对自己清晰表述,更不曾费心稍加检视。老实说,这类“普遍的”问题,根本不是我在禅观里向自己提出的那一类:禅观的用意要切近得多:理解我自己——首先经由我与他人的关系;并由此也多少理解“他人”——即那些与我发生关系的人。

当然,四天前的反思里,当我暗示应当有这层联系——对父亲的对抗是一份更深的冲突的表达,即对自身“男人”的拒斥——这还只是一个预设,由我极其有限的经验所暗示。这层联系在我看来至少说得通,尤其在男人身上;但我不自许“看见”这层联系在一般情形下成立。关于它,我没有那份我如此经常选作可靠向导的“贴身信念”。比如在我母亲的个案里,我看得分明:对父亲的对抗,是一份对\textbf{男人身上}阳刚特征的隐秘而恶毒的对抗之源,却全然不是对女人身上这类特征的——恰恰相反。诚然,单单充分推崇阳刚特征、在自己身上过度培植它们,或许并不必然意味着:人们全然接纳自己存在的阳一侧;那毕竟还意味着:也接纳那份“阳中之阴”——它自发地居于每一个以阳为“主导”的特征之中;而这当然不是我母亲的情形。

但反思正在这里取一种略嫌辩证的转向,它不叫我放心!我宁肯诉诸我对母亲这个人的直接知觉——它已由我对她、以及对我父亲生活的反思所磨细。我不记得曾有过这样的感觉:她身上有什么根本属“阳刚”的东西被拒斥\origpage{608}。相反,我在她身上强烈地知觉到这份矛盾、或毋宁说这份\textbf{撕裂}:她在自己身上培植(如一件件\textbf{武器})、并爱逾生命的,恰恰是那些在男人身上在她身上唤起如此烈性、如此狂暴的求战与求毁之饥渴的特征;而她的一生,正是被这份热病所磨蚀(并过早燃尽):不停地遭遇与迎战、并不遗余力地在他人身上压服那\textbf{同一份}力量——她押上全部赌注的力量,它摧残着她自己的生活,一如它摧残着所有她所疼爱之人的生活。

\subsection{主人们与仆人}
\label{ch18:local:h43}

\subsubsection*{(a) 天鹅绒的脚掌——或微笑}
\tocsubsub{(a) 天鹅绒的脚掌——或微笑}
\label{ch18:local:h44}

\notehead{133}{}
\phantomsection\label{ch18:local:n39}

(11月24日)前天那条注的反思里唤起的个案,并不是我所知的、印证这份预感的仅有的个案:父亲的超级阳失衡(无论这失衡取不取专制的形态),会在孩子们身上回响为一份对阳的拒斥,而这份拒斥又可以千姿百态的面目表达。男孩身上,就我所知、此刻在我想象中的个案而言,这份拒斥取对自身阳刚一侧(或多或少彻底的)压制之形——而这份拒斥,肯定将跟随他一生(除非一场深刻的更新,那当然极其罕见)。我母亲的个案使我确认:在女儿身上,情形不总是如此——除非我母亲身上也有过一份对自身阳刚一侧的拒斥,以更精微的方式表达,而至今逃过了我\fnote{149}{(*)}{一份邻近的情势,是一位禀性专断、侵夺性的母亲——超级阳失衡的征象。在我贴身知道的两个个案里,这在女儿身上表现为对她身上“阳刚”特征极深入的压制。}。而在她的个案里,相反,触目的是那极端相反的效应——她身上阳刚特征的过度发育(外加对一切女性之物的厌恶)。况且我还在别的方向知道一些个案,在\textbf{男人}身上(比如我母亲的父亲)——那是一场对父亲的\textbf{反叛},表达为一份强阳人格的发育,足以“势均力敌”地迎战父亲。既然我不曾有机会贴身认识这样的个案,我倾向于相信它应当更罕见。但这到底无关紧要。

\dmark{}倘若在或近或远我所知道的全部个案里有一个共同点,那便是:父亲的超级阳失衡,在孩子们身上回响为一份\textbf{失衡}——它可以朝阴的方向(或许最常见的情形),也可以朝阳的方向\fnote{150}{(*)}{当我在这里谈“朝阴方向的失衡”时,这并不意味着他的阴特征的(或许过度的、单侧的)发育,而毋宁是对阳特征的一份\textbf{压制}——这根本不是同一回事。在相反的、称“朝阳方向的失衡”的个案里,说的确是阳特征的“过度发育”,它常常伴着对某些阴特征或多或少深入的压制。}\origpage{609}。在我脑海在场的所有个案里(我并不打算在此对我所知道的个案作一份系统的清点),这份失衡伴着一份\textbf{对父亲的对抗关系}。我有这样一种印象:它同样伴着一份对另一些人的彻头彻尾对抗的姿态——那些阳特征打有强烈印记的\textbf{男人}——至少当那些特征不曾被互补的阴特征所平衡之时;也就是说,对着这样一些男人:在他们身上占上风的,是一种使人想起父亲的超级阳失衡。

这样一份超级阳的失衡(一如相反的失衡),固然易于在任何人身上唤起一份\textbf{不适},如我已有机会确认的\fnote{151}{(**)}{见“超级父亲(阳葬阴(2))”一注(n\textsuperscript{o}~108)。}。但这份不适并不必然表现为一份自动的对抗姿态——比如,它并不罕见地经由一份臣服、一份或多或少无条件的钦慕、或效忠的姿态而消解(或至少从意识的场域里消失)。

这里,一场联想来到心头:在数学世界之内、对我这个人(有着光环的)的关系里,最常见的,肯定正是这些调门——至少在那些同事(或学生)之中:他们(如我在别处所写)“并不觉得自己被一份相当的声名所保护”,或者(我要在此补上一句)那些人:某种内在的平衡、对自己力量的一份自发的认识,也未能使他们免于这类失衡。但无疑,这样一种“效忠”关系的本性,是它暗藏一份对抗;这份对抗,一俟合适的机会出现,便显形(公开地,或仍以隐蔽的方式)\dmark{}\origpage{610}……

我刚刚追随了几场联想,它们重拾并补足前天(在上一条注“阴阳的颠倒(2)——或反叛”里)的反思,并由此,也补足11月18日那条注“敌对的父亲(3)——或阳葬阳”的反思。它们使我明白:一位父母身上的某种阴阳失衡状态(在此即父亲的阳失衡),与它在孩子身上所取的那些回响之间的关系,绝无一义性——我先前说得匆忙。毫无疑问,父母的失衡(在此即父亲的)传递给孩子时所取的形式,必取决于许多别的因素:既有家庭环境的(尤其母亲这个人与她的姿态),也有孩子天生的禀性\fnote{152}{(*)}{这样,我确认:我母亲的三个弟弟(都比她小)各自走着一段与她(她多少像鸭群里的天鹅)全然不同的演进,也不同于其他兄弟。}。

但老实说,方才起笔反思时,我想走的方向并不是这一边。毋宁说,我想继续一场全然不同的联想——它自11月12日的反思起便在场,或头一次把阴阳角色的\textbf{颠倒}之动力学引入反思(在同名那条注里,“——或激烈泼辣的妻子”(126))。或许读者已在他那边做了这份联想——反正,当我在11月12日、随后前天22日唤起这个问题时,在我脑海的某处,确确实实像弱音器一样,有着另外两次“颠倒”已被谈起的记忆:在这场葬礼反思的进程中。第一次,是在V行列那条同名注里,“我的朋友皮埃尔”(4月28日的注(68$'$))。第二次,在9月30日那场反思里的一条脚注中,它属于“葬礼悼词(2)——或力量与光环”那条注。甚至还有第三次,但藏在字里行间:在两天后那场反思的开头,它开启了“阴阳之钥”的反思。(那是10月30日的注“肌肉与肚肠(阳葬阴(1))”(106)。)\dmark{}这里涉及的,是那场著名的“由三部分葬礼悼词引发的联想”的内容\origpage{611}——人们暗示过它;恰恰是它,当天便触发我踏上这条关于阴阳的枝蔓,我已沿着它走了将近两个月。或许,现在是时候揭开底牌了——自我谈起它以来;更不必说,自5月12日之后——“葬礼悼词(1)——或恭维”那条注之后——第二天起,我便想着它,已有六个多月。

这三桩情势的共同点,是它们涉及我的朋友兼昔日学生皮埃尔与我之间角色的一场“颠倒”。在那两桩被明白表述的情形里——方才重提过——我显为我昔日学生的“合作者”(甚至干脆是他的学生!)。第一次,是作为那个曾为“强有力的工具”之发展出过力的人——那$\ell$-adic上同调的“强有力的工具”,归在我那位才华横溢的前辈兼朋友名下(出力的方式当然乱糟糟,但人们也承认,有时有点意思)。第二次,当我们被相提并论时(因“以‘跨学科’手段联结了拓扑、代数几何与数论……”),同一份实在的颠倒,凭一个狡猾的手段被暗示出来:一场“排字”的遗忘——仿佛出于天大的巧合\fnote{153}{}{一如我在“屠杀”一注(n\textsuperscript{o}~87)里先前得以确认的:只要排字工人和搬家工人掺和进来,巧合常常把事情办得极好!}。这份颠倒的意味,当人们留意那些颂扬形容语的选择时,便变得比一个单纯的先后问题更别有用心(在此,是在一个我曾是唯一与Dieudonné一道、在科学层面上“启动”了它、却早已离开的机构之内):一边是“传说般深邃的理论”,另一边是“耀眼的发现”——后者还配得上着重号,人人都配,唯独我不配。这份意味,在“阴的葬礼(阳葬阴(4)”((124),11月10日)那条注的反思里,“触目惊心地”亮了出来——阴阳的反思凭它忽然“降落”在葬礼仪式的正中央:一边是(时而滥情到极点的)阴的、超级阴的形容语的堆积,另一边是阳的、超级阳的……

这正是早在5月12日“恭维”那条注的第二天,就已经打动过我的东西\dmark{}——早在两周前得以那样周详地显豁它之前\origpage{612}。按照我当时感受事情的方式(我日后还须在此重新检视),那里有过一场实在的真正\textbf{颠倒},或更确切地说,一场“颠倒”——被推向漫画式的极端——而其底下的实在,我感到是有层次的、平衡的。我把自己看作一个“阳”占强势主导的人,甚至超级阳——至少在我那些最显眼、最明白的特征上,尤其那些对他人彰显的特征上\fnote{154}{(*)}{而且,在我“出发”之前的那些年月,更甚于如今。}。相反,我在朋友皮埃尔身上感到一种基调属阴的气质——分明比我平衡得多,在我们常见面、他还是学生身份的那些年月,比当时的我要平衡。

况且我相信,这份对实在的把握在本质上是正确的。这些年里,有时——乃至最近\fnote{155}{(**)}{见“箭与浪”那条注(n\textsuperscript{o}~130,11月19日)。}——我偶然预感到我身上有一份原初的“阴”底色;而依我看,我是头一个也是唯一一个感到它的人——他人始终是经由我那些“阳刚”的特征(它们往往已经相当具有侵夺性)来把握我的\fnote{156}{(***)}{对我自己,亦然。},无论在意识的层面还是无意识的层面——至少就个人关系而言。而个人关系(恋情除外)所牵动的,主要乃至排他地,是我们身上的“东家”,那受制约的东西。阴阳反思进程中新显出的事实——\textbf{在我的工作里},我把握事物的方式阴占强势主导、“女性”——并不真正与我另外知道的东西相矛盾。它使后者更有层次,修正了其中一处我曾默默“一锅烩”的地方。细细权衡,我觉得:我心中那份骤然而强烈的印象——一场实在的漫画式“颠倒”,或更确切地说,一场蓄意\textbf{颠倒}的\textbf{意图}——这份“直觉”在本质上也是正确的,尽管粗糙。我现在想细细开掘的,正是这份直觉所未曾完全把握的实在。

\subsubsection*{(b) 兄弟与配偶——或双重签名}
\tocsubsub{(b) 兄弟与配偶——或双重签名}
\label{ch18:local:h45}

\notehead{134}{}
\phantomsection\label{ch18:local:n40}

\dmark{}(11月25日)我首先该试着更贴切地圈定这份印象——它于我显而易见:朋友皮埃尔这个人身上的“底色”,是一份\textbf{阴}的底色\origpage{613}。依我所知觉的,无论在“我”的层面上——我尤其在他与我的关系、与其他人的关系里看到它的表达——还是在他的工作里,也就是说,在他身上认识的冲动、创造官能的层面上,都是如此。

就第一个方面而言,分明地,他与我是气质\textbf{互补}的,还加一层细节:我身上那份过度的、“超级阳”的东西,有时似乎让他有些错愕。我相信,尤其是我那份朝着完成任务之实现的恒常的前投、那份对一切与此无关者的\textbf{隔绝},在他身上唤起一种难以置信的惊异——我在其中感到一缕温情的惋惜:那份惋惜,我曾在我母亲身上多次感到,当她看到我与周围事物的美如此隔绝之时\fnote{157}{(*)}{我母亲与我父亲一样,直到生命尽头都保持着一份与自然交融的能力,同时对周围的一切有敏锐的观察感——这两样,至今都是我所欠缺的。这或许是她存在中唯一未被她压制、得以自由绽放的“阴”的方面。另一方面,就“朝向目标的投射”而言——它是我“我”的主导特征之一——这或许也是我这个人唯一成功的方面:比我的母亲更阳!}。

在他身上,这不是本来意义的“不适”,不是拒斥某份实在的征兆。至少,我不记得有一次在他身上感到过对我的不适,也不记得有过拒斥、疏远、或哪怕我们之间的顶撞的印象。而我毫无怀疑:这绝不是他身上什么“外交式”的故意——决意什么都不表露。相反,他有时会表达我暗示的那份“惊异”,不带一丝局促,也不带一丝恼怒。分明地,我们关系里的基调——它至今\dmark{}从不曾变过\fnote{158}{(**)}{(11月26日)基调固然还是同情、吸引,但这并不妨碍:自我出发之后,年复一年,这份关系越来越凝定、僵化,失去了赋予它生命质量的东西。我感到自己面对一副“甲壳”,密封得如此完好,再没有什么能从任何一个方向通过。关于这一点,见“两个转折”与“墓”两注(n\textsuperscript{o}~66、71)。}——是一份温情的同情,没有一丝阴影从中穿过。

\origpage{614}这在我始终是一件奇怪的事,而且我相信,在我离开IHES的情节之前,任何人都不可能从中猜到什么(甚至那时,也只是在面对面的接触里“透”出来的层面)——即:自我与他相识的最初几年起,他与我的关系里就有一份深刻的、根本的含混:一份暗藏的对抗、至少一份想与我这人划清界限的愿望,以及排挤的愿望。这后一份,在我离开IHES的情节里(见“驱逐”一节(63)),以格外粗暴的方式显形(当场叫我目瞪口呆),尽管方式上无比地压低了嗓门。我的朋友,不久前刚被遴选为IHES第五位“终身教授”——主要凭我热心的奔走。在我们之间进行的那次“解释”(也许有过不止一次,我已说不清)里,他时时刻刻都保持着那份完美的、微笑的自然——带着善意亲切的一切表象,那是他如此讨人喜欢的地方。他向我解释,不带一丝犹豫或为难——更不带对抗或敌意、或暗藏的得意——说他早在青年时代就决定把一生与全部精力献给数学工作;说他这份对数学的奉献,无论好坏,对他而言必须置于一切之上;说我之所以等待同事们、尤其他本人的声援(以要求取消来自军部的经费),其缘由在他看来与数学全然无关;他当然惋惜:这对我是一票否决的情势,由于“公设”与他不同的生活观,我将离开IHES,为一个在他看来无关痛痒的缘由;但令他深感遗憾的是,他不能声援这份对他陌生的请求——其结局他毫不关心(注~134$_{1}$)。

我上面给出的,是我朋友那番话的“表层”内容、明白说出的内容,一如我的记忆所归还给我的,而我并不曾费心同时去复原一种表述的风格、或一场谈话的氛围——它的细节,除我此处所说的,我也不曾记住分毫。这段情节发生之时,我还没有丝毫疑心:在一番话那十分平淡(有时怪诞地荒谬)的表层内容背后,往往以弱音\dmark{}、而清清楚楚地,表达着一份全然不同的讯息\origpage{615}。这份讯息,肯定在无意识的层面被知觉着,却被惊惶地拒斥、逐出意识的场域。一如我在引注的“驱逐”那条注里所暗示的:要把一份其实足够刺眼的讯息清场出去,想必需要相当大的能量!而恰恰是在那条注里——写在十四年之后——我才头一次费心把这段情节交给有意识的注意,并把那份被否认了如此之久的意义清晰地表述出来。

我刚刚追随的,是呈现在我面前的联想线索中最强的一条。我追随它时,怀有某种迟疑——仿佛经由这场“枝蔓”,我正离开我的主要打算。可事后我明白:全然不是如此。无疑,一个人的影像与气质,当它自发地从他被牵涉其中的具体情势的描述里透出时,比一份号称能圈定它们的“特征”罗列,更活、更有说服力。但与其投入其中,我宁肯记下又一场联想,并踏上又一场枝蔓:把这里检视的关系,同Serre与我的关系相比。在我们两个人的关系层面,占上风的印象,绝不是与皮埃尔在一起的那份“\textbf{互补}”,而毋宁是两个气质之间的\textbf{亲缘}——双方都强“阳”。在十八年紧密的数学交往里,这份亲缘不止一次经由偶发的摩擦显形,表现为一时的冷淡,没有一次持续长久。依我的记忆,这些情节,缘于Serre那些不拘小节的焦躁举动,它们与我的敏感不“合拍”。Serre有时会恼火于我执着一个想法的顽固——当它在我看来重要时,我顶着逆风硬干;我一有机会就把它重新端出来,不管它“通不通过”,我深信(这份深信很少骗过我)自己握有“\textbf{对的}”观点。不知为什么,Serre对我的那些上同调“大家伙”发展出一种反感——也许他干脆过敏,一如André~Weil对一切“大家伙”过敏。另一方面,当我开始展开“我的”上同调yoga——在五十年代后半期——Serre实际上是我唯一不时的对话者;所以一开始就很难顺风顺水!我相信,他是在$\ell$-adic上同调自1963年发展起来之后,才谨慎地开始对这份工作产生兴趣的;同年,又有我对L函数有理性的证明草样(“三下五除二”)\dmark{}\origpage{616}\fnote{159}{(*)}{我记得的另一处摩擦点——无疑更带偶发性——是我坚持把代数群与形式概形中的取商理论(五十年代还被人误解)系于平坦等价关系之“有效性”的问题,乃至(后来)系于fpqc层的取商语境。这些观点,先由加布里埃尔与曼宁重拾,如今已是代数几何里几乎处处通行的通货,甚至在别的领域亦然。我觉得,Serre的迟疑,自我终于费心(因为似乎没有别人肯干这苦活)把第一条有效性定理——对平坦而有限的等价关系——白纸黑字证明出来那一刻起,便消散了。},才开始明白它们通向某个地方。

依我看,Serre与我的关系,是一份典型的阳-阳亲缘,与Deligne的关系正相反,那是一份阴-阳的互补。而在数学工作、数学进路的风格层面,情形恰恰颠倒过来。一如我在前面一条注里(“九个月与五分钟”(123))已有机会说过的:我感到Serre的进法与我的进法是\textbf{互补}的,取阳-阴互补之义。而恰恰是这份互补,成了那些偶发摩擦的缘由——摩擦缘于他与我双方都强阳的气质。

Deligne与我,在数学进路上的关系,则全然不同,这毫无疑问。我可以毫无保留地说:正是在Deligne身上,胜过任何别人,我体验过那份完美的\textbf{亲缘}——在我们彼此都感兴趣的数学问题之看法与进法上。每当我们的数学对话发生,这份体验便重新来一遍。我清清楚楚:这绝非一种偶然的情势——比如说,不是由于在那些关键的学习岁月里,我确确实实给过他影响。这份亲缘并非在长年的熟悉中发育——恰恰相反:它从我们最初接触起便在场,正是它,才是那份力量,几乎一夜之间造出那样一份系于我们共同激情的纽带\origpage{617}。这是两种数学进法之间的一份深刻亲缘\dmark{}——先于我们的相遇而在,而它(我深信)表达着双方原初气质的一个重要侧面:把握事物、发现事物时那份“阴”的“基调”\fnote{160}{(*)}{(11月26日)本注的反思,与“上涨的海”“九个月与五分钟”两注(n\textsuperscript{o}~122、123)的反思相续,似乎意味着:对任何人来说,都有一份“双重签名”、一份\textbf{双重}“基调”在:其一(无疑最显眼)关乎“东家”,即“我”的结构与支配它的机制;另一份关乎“匠人”,别名“孩子”,也就是说,也关乎认识的冲动、发现世界的冲动、创造的冲动(当然,也包括爱的冲动)。(诚然,把东家当作匠人、把匠人当作东家——也就是说,把鱼目当珍珠——是世上最平常的事,不过那是另一段故事……)
如此,在我身上,这双重基调是阳(东家)--阴(孩子);在Serre身上,是阳-阳;在Deligne身上,是阴-阴(对此,我心中没有一丝怀疑或犹豫)。在双方彼此好感的底子上,正是这份“记号”(或“调门”)的“分配”,使得在人与人的关系层面,我与Serre的关系是亲缘、我与Deligne的关系是互补;而就我们对数学的进路而言,恰相反。
四种可能的“分配”里,只剩双调阴-阳还“记在账上”。鉴于阴在我们大男子气社会里的失宠——这失宠倾向于主要作用在第一个调门(“东家调门”)上——我推测,双调阴-阳应当比阳-阳少见。可我至少认识一位知名的数学家,依我看正对应这份签名。当然,第二个调门、“原初调门”更难圈定,因为它常常被外部影响、被“跟人人一样”地做与是的心思搅浑。}。

这样一份贴身的信念,谈不上“证明”——正如我不曾想过要“证明”我自己数学工作里(姑且这么说)的基调是阴的、“女性”的。对这类事情,至多有时可能把一份感受从一个人“传递”给另一个人,在别人身上触发对某样东西的认识——那样东西,人们此前不曾对它留意;它逃过了他有意识的注意,却其实早已以弥散的形式被“登记”在某个地方。情形肯定被搅浑了——一如常常如此——被当事人为按时兴价值、按“阳”的“男性”价值自我塑形所作的那些努力。而我分明看到:他的数学著作、以及他所施加的那份(可观的)影响,都深深地打有他与我的关系那份暧昧的印记;我却怀疑:那些为抹去一份与我的底色相近、却被否认的底色而作的努力——是否终于成功\origpage{618}。诚然,严谨的那些要求——在我“出发”之前,它们尚未在他身上起作用——早就使他不再屈身(至少在打算刊布的文字里)理会那些远低于他的东西、或那些今日被视为该绝的东西。可是我觉得,在他刊布的东西里,他终究没能不追随那份自发属于他的进法风格。至少,这是我翻阅他那几份为数不多的抽印本时得到的印象——在我“亡故”十五年之后,他还肯把这些越过坟墓寄给我。

当然,我对Deligne数学进法的把握,首先汲取自我“亡故”之前的那些年,1965至1969之间。整整五年,我们强有力地都接在同样的事情上,数学交往不间断(除了他在比利时的那一年),而且比我与任何其他数学家——包括(依我看)Serre——的交往都强烈。我已有机会不止一次唤起这些年头\fnote{161}{(*)}{尤其见“孩子”“下葬”“驱逐”“授职”“死结”诸注(V行列,“我的朋友皮埃尔”内),以及“继承人”一注(IX行列,“我的学生”内)。},那是双方都创造强烈的岁月。在我朋友身上,它们以一个令人印象深刻的启程为标记——它却并不叫我惊讶,因为一切在我看来如此理所当然!那时,他对实质的那份极可靠的感觉——对最抽象的表象背后、或在最“一般的胡话”的表述之中的可触之物——还不曾被一种自满、也不曾被那后来才出现的葬礼综合征所遮蔽。他那时为这些主题(我可以说,极阴的主题)作过许多贡献——后来的共识(经他毫无保留的祝福)早已把它们开除出“正经”数学的行列\fnote{162}{(**)}{(11月26日)我顺带提醒:这些数学的一部分,曾在1981年的“反常层会议”上被大张旗鼓地开棺出土——而我连名字都不曾被提起;次年,又借着那“值得纪念的LN~900卷”。关于这一点,见“不义——或归来的意义”“赊账的学位论文与全险”“一个梦的回忆——或Motif的诞生”三注(n\textsuperscript{o}~75、81、51)。译按:原书此脚注角标作“18.5.4.4**”,当系排印之误。}:topos的形式主义、上同调的“大家伙”……我在SGA~4的引言里,怀着明显的快意,把它们逐一清点、郑重表彰\dmark{}\origpage{619}\fnote{163}{(*)}{(11月26日)这些评注,是在SGA~4彻底重编的第二版里补入的(重编主要关乎site与topos的部分)。它们会给人这样一种印象:仿佛Deligne曾参与平展上同调与$\ell$-adic上同调那份“强有力的工具”之主要观念与主要成果的绽放。于是,我在那里为Deligne、为我其他那些上同调家学生的水磨添了水——十年之后,他们瓜分一位亡故师父的躯体!}。还有其他这样的贡献(其中有些还更“带劲”,足以让他一上来就跻身“大明星”之列),在我1968/69年的那双份报告里——见“授职”那条注\fnote{164}{(**)}{我提醒:这双份报告已重印于《数学反思》第一卷。}。

\notehead{134$_{1}$}{}
\phantomsection\label{ch18:local:n41}

(11月26日)\fnote{165}{(***)}{本子注(系于前一条注“兄弟与配偶——或双重签名”~n\textsuperscript{o}~134)由该注的一条脚注生出。(参见该注第三段末尾处的参见号。)}典型的细节:那笔军费——当它关乎要我离开时,谁也不肯为它抬一抬小指——就在我离开的那一年,被悄无声息地取消了!谁也说不准,没准儿它会叫某位对这个章节有点计较的贵宾心里不痛快……况且,那笔经费在IHES的资源里只占一小部分(如果我没记错,是5\%)。不曾有人串通,可我那四位IHES的同事(不算所长)之间,却有一种漂亮的全体一致:抓住一个机会,把我甩掉(几乎同时,顺带把所长本人也甩掉了)。而我,还以为自己不可少、还被人爱着呢!

(12月6日)IHES的两位物理学家——Michel与吕埃勒——不满意:IHES的“物理”分部,与由Thom、Deligne和我(其中有两位“菲尔兹奖”!)代表的数学分部相比,显得有点像穷亲戚。Deligne的遴选,又加剧了这个失衡(那次遴选,本就是在Michel与吕埃勒毫无保留的同意下进行的——IHES科学委员会事实上全体一致,唯Thom例外)。物理学家与数学家之间有过协调,要向所长Léon~Motchane施压,以求在可能的限度内恢复两个分部之间的正当平衡\dmark{}\origpage{620}。不过我猜,我的物理学家同事们,看到这失衡被有效地补偿——而且比他们所敢望的早得多,借着我突然要走的前景——心里未必不乐意。

至于Thom,他怒不可遏:Deligne的遴选,竟是顶着他的正式反对进行的。他曾把我那份辉煌的“授职”报告里列举的Deligne的各项贡献——全未发表,分明全从他头顶掠过——称作不过是一些“练习”!Deligne登上IHES“终身教授”之位、与他平起平坐,叫他震惊的是:年轻的Deligne——当时25岁——竟尚未浑身挂满荣誉。按Thom的看法,登上这样的职位,只应作为“一项事业的加冕”而来。不到十年之后,便离那些英雄的年月那样远——当年,我把尚不知名的Heisuke~Hironaka迎进临时凑合的屋子……总之,Thom的怨气如此之深,以致他当时(据他亲口对我所说)琢磨着离开IHES,回到他在斯特拉斯堡的教授职位——那职位他留了一手(比我当年离开CNRS进IHES时更谨慎)。我热心地为Deligne作保,便成了他受挫的第一位的、主要的缘由;我猜Thom在心底里必定觉得:我不过咎由自取——我竟敢如此无礼,在把我那位辉煌的“被保护人”引进IHES不过几个月之后,自己便不得不离开!

至于所长,在他被终身教授们一致的意愿逼到墙角、催他下台的时刻,他便(凭一套他耍得炉火纯青的惯用战术)玩起了“分而治之”的把戏:拿军费问题当作一个方便的手段,声东击西,同时把最碍他的那位终身教授甩掉。(局势的颠倒,手法老到!而他对这笔经费之存在所守的那份秘密,在我看来,反倒成了催他下台的一条追加的、强制性的理由!)这并不妨碍:我走之后,事情到底没再拖多久,他的离任紧跟着我的离任——我,与他一样,自IHES最初那些寒酸而英雄的年头起便在其中,而正是我,与他一道、各凭自己的手段,保证了它的信誉与长存\dmark{}\origpage{621}。

\subsubsection*{(c) 阴,作为仆人——与新主人}
\tocsubsub{(c) 阴,作为仆人——与新主人}
\label{ch18:local:h46}

\notehead{135}{}
\phantomsection\label{ch18:local:n42}

(11月26日)在我出发前的那些年月里,Deligne与我之间的许多亲缘之中,有一份是这样的快意:他与我一样,喜欢在需要的时候,展开我所谓的“大家伙”。我作为数学家的能量,若不说全部,至少绝大部分,都献给了这样的任务。倘若说是盖一座房子,“做大家伙”便意味着:不止于画一张诱人的房子草图、甚至从不同角度画上两三张,也不止于画出带尺寸的详图;而是把用来盖房子的石头,一块一块地搬来、凿好;砌成墙,架上梁、椽,铺上瓦或石板瓦;安上门窗、盥洗盆、水槽、管道与檐槽;并且(倘若真打算自己住)一直安到窗帘与墙上的画。房子可以轩敞,也可以是只一间的小屋——做工里的精神却是一样的。而只要住进去,任凭你做得多么彻底、多么到头,你很快便发现:工作永远做不完,总会有新东西冒出来——至少,当“大家伙”——对不起,是房子——够大的时候。

1955至1970年间,我作为数学家的最大部分的精力,献给了启动并彻底展开四项\textbf{大}“大家伙”——当然,没有一项做到了头,见前文。按时间先后,它们是:上同调工具、概形、topos、Motif\fnote{166}{(*)}{“上同调工具”不曾等我到来便已存在。这里说的是某一种个人的进法,它尤其引向“平展上同调的掌握”(在我看来,它是Deligne所完成的Weil猜想之证明里主要的技术与概念的要素)。二十年之后,我在《追寻栈》(A~la~Poursuite~des~Champs)里重新沿这条进路走下去,朝着“非交换上同调”(或“同伦上同调”)的方向。至于“交换上同调”的方向,我在“我的孤儿”一注(n\textsuperscript{o}~46)的开头处对这份进法给出若干说明。这里所说的四项“大家伙”,大体对应引注里的五把“关键观念”——出入只在于:“上同调工具”对应其中\textbf{两}把观念(即导出范畴,与“六运算”的形式主义)。
耐人寻味的是:四项“大家伙”(或主要研究主题)之中,唯一在我的葬礼悼词(见注~n\textsuperscript{o}~104、105)里被点名的,是topos。像特意安排的一样,它也是被我的那些上同调家学生下葬的三项之中、在葬礼悼词的时刻还不曾被换别人爹妈“开棺出土”的那一项。(葬礼悼词在1983年;导出范畴1981年于反常层会议开棺,Motif则1982年,在那“值得纪念的LN~900卷”里。)}。这四项主导主题,彼此亲密地系在一起\dmark{}\origpage{622}——如同同一座农庄、同一个村落里的几座建筑,共同成就一份蓝图。而每一项“大家伙”,都强制地——我毫不曾寻索——把我引向展开另一些明显小一号的“大家伙”:有点像盖一座大房子、乃至整个村落,人们便被引去砌一座石灰窑、建一间木构与细木工场,等等。比如,每年都会重新感到需要给范畴的观念与构造的武库,添上两三件(小号的)“大家伙”。那些晚了十年、二十年才来的人——他们发现一切都做熟了,便舒舒服服住进现成的房子(还有另一些人,心底里知道该有个分寸)——对着那么多读不通的“胡话”(Deligne语)、对着头发里分出的四股(“Spitzfindigkeiten”,一位对我满怀善意的著名德国通信者这样称呼),耸起居高临下的肩膀\fnote{167}{(*)}{我的通信者和颜悦色地向我担保——为博我高兴——他自然明白,我的著作“在很大的限度内,没有这类毛病”(“weitgehend~frei~von~diesen~Übeln”)。他所说的“毛病”,指的是一旦有人胆敢在尚属猜想的地基上发展一门理论(如我就Motif所建议的)就免不了要掉进去的那些毛病(如形形色色范畴论者的“Spitzfindigkeiten”)。这里我们又见到那份对“数学之梦”的彻头彻尾的拒斥——见“被禁的梦”一节及随后三节(第5至8节)。这又是数学里对一切“阴”的、“女性”的进法或行径之自动压制的侧面之一。}。这些人对在光地上盖一座房子是怎么回事,毫无概念;而他们,恐怕永远不会去盖,只满足于在别人凭双手、凭整颗心为他们盖好的房子里,扮演房东。

\dmark{}我方才有点急躁,看上去把朋友皮埃尔也塞进了那些“对盖房子毫无概念”的人的口袋\origpage{623}。他非但见过我做工,而且他自己也满心欢喜地在盖他的房子,仿佛他来到世上二十年,不曾做过别的事。况且,这“大家伙”与盖房子的故事(倘若读者还没看出来……),又是另一个侧面、另一幅意象,用以圈定我此前曾试着以“上涨的海”的意象、继而以一道道相互追逐的浪的意象,多少把握住的那个东西\fnote{168}{(*)}{见“上涨的海”“箭与浪”两注(n\textsuperscript{o}~122、130)。}。说的是把握实在的“阴的样式”、“女性”的样式,以及与之相应的行径:浸润其中,从中析出一幅意象,柔韧而忠实地复现这份实在。于是,经由我自己的这个人绕了个弯,我又回到了最初的打算:把我心中那份强烈的知觉“传递过去”——Deligne的数学进法与我的之间,有一份亲缘、一份本质的亲缘。可在Deligne身上这个我方才试着以意象圈定的方面,自1970年我的“出发-亡故”之后,依我看,有过一场全面的“搅浑”:我相信,在他“之后”的出版物里,“大家伙”全然缺席。当然,他也无法合理地动用他那被申斥的师父身上的这个特征,来贬损师父,却又容许同一个特征照他自己的天性在他身上绽放。

诚然,倘若要做的不是追随一份内在的需要——一份原初冲动的表达——而只是靠积累“立得住”的\textbf{成果}来抬举声名,那么我的朋友,实在没有任何兴趣继续让(或大或小的)“大家伙”来缠身。就在我的时代,在Bourbaki小组之外(小组本身正投身于一项相当规模的“大家伙”!),这已经是一件颇不受待见的事。这毫不奇怪:毕竟,我们这个社会、科学世界的共识里,那份“超级阳”的马眼罩,不是昨天才戴上的。或许这正是主要的缘由:我所乐意盖的那些房子,在漫长的岁月里无人居住,除了那位匠人-泥瓦工自己(他同时又是建筑师、木匠等等)\origpage{624}。而直到今天,即便我工作中早已成为公共遗产的那部分(甚至在没有别的参考文献、唯我的文字可依的地方),对某些人而言(至少对那些不属于那个上流圈子、也不以居高临下为本分的人而言)仍环绕着一圈近乎令人畏惧的光晕,仿佛踏进去便需要近乎超人的官能。这是真的:那常常很费时间,也不可能不如此——既然一切确确实实是手工的、细部的,从头至尾,甚至在每一章的转弯处都有说明,告诉你走到哪儿了\fnote{169}{(*)}{我想,只是随着岁月的流逝,我才意识到纳入这类说明——它们往往纯乎启发式——的必要:尽可能给读者传达一份“方向”感与用意感;写的时候,它在我身上强而在场。今天,依我看,这比一丝不苟地写出关键的证明要紧得多:一旦读者感到去向何方、而那“何方”吸引着他,那些关键证明,他自会津津有味地复原、甚至凭空造出。}。我的学生们当年与我一道工作的时节,进入这水里,并不显得格外吃力。但那是在“可触成果”已经挣得数学建制背书的时节,我的学生们怀着稳操一张“保险牌”的自信工作。我有这样一种印象:从那以后,倒不止一个人乐得给“读不通”的版本续命\fnote{170}{(**)}{这件事,只有Deligne一人说得有鼻子有眼——他最近的来访中,还当面重提过。说的是SGA~4(其一大半以极度的细致展开topos的语言),被我的朋友判为“读不通”,以此为他那天才的“SGA~4$\frac{1}{2}$工程”作辩护。},顺应一股今日比我那时暴虐得多的时风。

可是,即便撇开时风的种种要求不谈,当人们计算收益与“回报”时,肯定会像躲瘟疫一样躲开“大家伙”。展开一项“大家伙”、并把它交到所有人手边,这是向科学共同体提供的一项\textbf{服务}——而这个共同体,往往不情愿地接受它。对这份完全可以理解的推拒,我向来不怎么为难;我分明知道,我手里有“好家伙”,人们或早或晚,总免不了要来的\dmark{}\origpage{625}。而即便他们来了,以“信誉”计的“收益”,也只能微薄。倘若我作一份数字的盘点——不是盘点我在1955至1970这十五年里引入并发展的那些观念、问题、想法(它们要么进入公共的、匿名的遗产,要么无声无息地下了葬,等候被大张旗鼓地开棺),而是盘点人们可称作“大定理”的东西——我怀疑连十个也凑不出。直接花在证明它们上面的时间总量,大约几个星期,至多几个月。1957年之前,一个也没有(Riemann--Roch--Grothendieck定理)——然而我知道,此前三年我并没有虚度。说不定,倘若这十五年里,我没有执拗地追随我身上那份理解的激情、信任它向我口授的那份进法——不管它“有利”与否(按这样那样的要求来算)、在上流世界里体面与否——眼下这些“大定理”竟一个也不会得到证明(而这本非我主要的操心)。这份进法,每一次都在于:从一份强烈的出发直觉、或一小把这样的直觉出发,把它们当作一根结实的、经得住一切的线,由它拽着我走向未知;而一边走——换个意象说——我总按捺不住,随着那未知渐渐为人所知,如一些粗粝的石头,人们一边凿它一边“认识”它,盖起一座座房子,很大的与不那么很大的,全都宜于居住——一座座房子:其中每一个角落与旮旯,都注定要成为不止一个人的亲切而熟悉的居所。门窗端正,开合顺滑,不夹不吱嘎;屋顶不漏,烟囱抽风。未必是巴黎圣母院,每一只面包箱里也不都藏着一个“大定理”——不过是些该盖的房子,我为居住而盖的。我盖它们、盖得漂亮而宽敞,从中找到我的快乐;我分明知道:我做的工作——独自或与人同做——是必须做的,而且每一刻都做到我当时所能做到的最好。

这份精神,我在五十年代的Bourbaki小组里也找到了,它使我在其中如鱼得水、“在家”一样——尽管有氛围与文化的差异、尽管有我在该提的地方提过的那些偶发的困难。至少在那时,我在其中找到的,仍是一份\textbf{服务}的精神\dmark{}\origpage{626}:为一项\textbf{任务}服务;而在任务之外,也是为其他人服务——那些像我们一样渴望懂得大大小小的事物、并要懂得个透彻到底的人们。这种“服务”并不呈现为严苛的义务或苦行。它自发而欢快地出自一份内在的需要;它表达着一种共同的东西,把这些如此不同的人系在一起。

而我在Cartan研讨班里也认出这同一份精神——那么多法国数学家在那里初试锋芒;后来(六十年代)又在我自己的研讨班里(它的代号为SGA,“布瓦玛丽代数几何研讨班”)。两个研讨班的一个差别是:我的那些,强有力地以方才唤起的“大家伙”的展开为中心(因此是“\textbf{我的}大家伙”),它们从来嫌人手不够;而Cartan逐年跟随的主题,更为驳杂。可依我看更要紧的,是两个研讨班共同的东西,尤其,是它们似乎曾有的那个要害功能、它们的\textbf{存在的理由}。老实说,我看到两个。这些研讨班的功能之一,与Bourbaki的用意相近,是准备并向所有人提供易上手的文本(我是说,实质上完整的),周详地展开那些重要而难入门的主题\fnote{171}{(*)}{“难入门”,或因为这些主题仍被理解得不完全,或因为它们只有寥寥几个入室者知道,而讨论它们的零散出版物只给出一幅不称职的图景。}。另一项功能,是构成一个\textbf{场所}:在那里,胸怀热望的年轻研究者——即便不是天才——也确保能学到数学家的手艺,就着正当时的种种问题,在高明而善意的长者身边。学手艺——也就是说,把手插进面里,并由此找到让人认识自己的机会。

看上去,我1970年的出发,标志着“大研讨班”的终结——至少在法国:那些\textbf{持久的}场所——年复一年,当代数学的某些大主题在其中开工;也是\textbf{善意的}、激励人的场所,为一切前来把手插进面里的人\dmark{}\origpage{627}。我不知道世上别处是否还有(或许在莫斯科,凭I.M.~Gelfand的推动?)。确定的是:这样的场所,分明不合时宜,一如那些白纸黑字、细致入微地写下来、为的是交到\textbf{所有人}手边的“大家伙”。

并非偶然:如今几乎再没有人写那种精心结构的、(暂时)穷尽的阐述——论那些熟透了十年、二十年的主题,分明至关重要,而在此期间却只有一撮“局内人”能够问津。属于数学“上流世界”的人,倘若不同时属于那一“撮”,一旦需要,尽可以毫不费力地让那撮人里的一位给他讲个明白——人家求之不得。至于其他人:见鬼去吧!六十年代,我看到一大摞书,声嘶力竭地要求有人来写。我自己倒很愿意写,可我不能事事同时做。据我所知,这些书至今一本也没有写出来\fnote{172}{(*)}{(11月28日)这里我该为我推动之下写出的那些论文作例外。鼓舞我的那份精神——我相信,至少在他们与我一道工作的时期,它也传给了我的学生——正是鼓舞我自己的工作的那份;也就是说,打个比方,“盖那些分明需要盖的房子”——尽管某一幢特定的“房子”需要不需要,常常只有我一个人感到。我有这样一种印象:一般而论(一个例外除外),这份感觉终于传到了学生身上,使他“挂上”某个题目,继而强有力地认同所选的题目。把Verdier除外——他不屑于把我们约定好的基础工作交到所有人手边,那工作至今等着有人来写——凡与我一道完成国家博士学位的学生的学位论文,都成了所谓的“标准参考文献”。都是些宜于居住的房子,而且没有任何两幢重复……}。然而,不止一个人(且不说别人,单在我昔日的学生们中间)曾足够“在局中”、有那份感觉与那份手艺,本可以毫不费力写出某一本来就该写(而且至今仍该写)的书。而从某些人后来的工作所传回我的那一点点来看,我并没有这样的印象:是他们那些更个人的工作之繁多与困难,使他们……不肯(“真遗憾,实在没有工夫!”)向那著名的“数学共同体”提供这项服务。而且对不止一个人,大可以打赌:写一本被阅读、被引用的书(即便书中所阐述的并非都出自他自己——但“怎么写”绝不是可以忽略的量……),比他那一叠或厚或薄的抽印本,更会使他知名\dmark{}\origpage{628}。分明地,叫这个那个以一种令人印象深刻的全体一致,拒绝把那依旧只属于少数人的东西交给所有人的,绝不是单纯的“没有工夫”——或者,不肯采取一种“服务”的姿态(哪怕此地彼处,花点时间写一本书也好)。这里,一场联想不可抗拒地涌来:1965/66年的SGA~5研讨班——被埋没了十一年,而埋没它的,恰恰是它最初也是唯一的受益人:我的朋友皮埃尔领头,还有我其他那些上同调家的学生!诚然,那个案子里有一具躯体要瓜分,动机因此有点特殊。但我也想到别的个案:已提供的服务填补了明显的空白,却被在位者们一挥手扫到一边\fnote{173}{(*)}{我这里当然想到Yves~Ladegaillerie的工作,以及Olivier~Leroy的工作——前面四条注与节里谈到过它们(“谁也拦不住进步!”“棺材2——或锯成段的切割”“科学院报告——或新伦理”“棺材4——或没有鲜花没有花圈的topos”,注~n\textsuperscript{o}~50、94、第33节、注~n\textsuperscript{o}~96)。}。人们会说,这些仍是一路特殊的个案:被瞄准的是我这个人,而明摆着,那些工作正是由我所启发的。然而,我在这一切之中分明感到一种“时代精神”,它超出一切个案。

我此刻正多少圈定着的“时代精神”的这个侧面,便是\textbf{笼罩服务姿态的那份污名}——我经由一大堆汇聚的征兆知觉到它;对我而言,它是明摆着的事实。谁都可以否认它,正如谁都可以亲自检视它、确认它。我这里的用意,不是向一位不情愿的读者“证明”它,而是试着把握它的意义。在眼下这场反思的光照里,有一个意义一望即知。服务的姿态,典型地是一份“阴”的、“女性”的姿态;它落在被贬抑者的份里,毫不足怪。我屡次仿佛知觉到的那层细节是:这样一种姿态,只配给那些没有本事摆“主人”姿态的人;照这份精神做成的工作,是\textbf{下等}的苦役,是给那些驾着伟大观念与“耀眼发现”之高车驷马的人打下手的小卒干的\dmark{}\origpage{629}。

然而我也知道,不止于此——因为不然的话,人们为什么要想方设法,拦住一伙(倘若居然找得到)有心的小卒,不许他们在自己的角落里安安静静地做那份合乎他们身份的低贱苦役:终于供给一些结实的参考文献,好叫人不必再满足于(当人们还肯说点什么的时候……)“人们知道……”或“人们可以证明……”,或更罕见也更老实的“我们将承认……”?!

八年前,Yves~Ladegaillerie为“安放”他的论文而屡遭挫折时,我头一次直面这份扰人的问题\fnote{174}{(*)}{关于这一点,见前一条脚注所引的两注(n\textsuperscript{o}~50、94)。}。老实说,那正是我对数学、对数学家世界的兴趣最边缘的时刻。我有点目瞪口呆,却不曾试着去澄清这个谜的意义。此后若干年,直到今年2月、随着《收获与播种》里继续的那场反思,我的态度几乎没变——出入无几。然而,由于不断地接收征兆——甚至并非有意——我终究渐渐不由自主地也接收到它或多或少的意义,或者毋宁说,它的诸般意义。我确实看到两个。一个关乎我这个人——那是针对我的葬礼综合征,我对它的巡礼还没有全然做完。另一个与任何具体的人无关。那是一份\textbf{对“信息”之占有与控制的独占姿态}——在“科学”的建制之内占上风的姿态,它使这建制变成所谓“科学”“共同体”之内一个君权神授的统治种姓\dmark{}\fnote{175}{(**)}{(12月6日)人们会注意到:支配欲是一份\textbf{超级阳}的失衡,而且是这类失衡远为最常见的形式。它对应于“主人——仆人”这对阴阳偶子(或“支配(或掌握)——服务”这对偶子)里,阴的、“女性”那一项的泯灭;它邻近“掌握——服务”这对偶子。}\origpage{630}。

这是我在“职业道德共识——与信息的控制”那条注里(只擦了一擦,只擦了一擦)轻轻触及的主题,在“年轻人的势利——或纯洁的卫道士”((25)(27))里也谈了一点。我怀疑:这是科学世界里的一桩\textbf{新事实},在逝去的二三十年里,蹑着狼步悄悄安顿下来的。我不相信我曾属于传播、迎纳这份“新伦理”——“两套砝码,两套量法”的伦理——的人\fnote{176}{(*)}{我不知道,在我上一辈或同辈的同事之中、乃至在更年轻的同事朋友之中,有多少人看见过它。我怀疑,在“那些曾以兄弟之情迎接我、走进那个变成我的世界的世界”的人们之中——《收获与播种》正是献给他们的——除Chevalley或可算一员,连一个人也没有。这当然属于那些我本愿与他谈谈的事情之列——可他已经不在,不能告诉我了……}。倘若我在这份新伦理的问世之中也有一份共同的责任,那毋宁是因为:我不曾看见它的到来\fnote{177}{(**)}{我不知道,在我上一辈或同辈的同事之中、乃至在更年轻的同事朋友之中,有多少人看见过它。我怀疑,在“那些曾以兄弟之情迎接我、走进那个变成我的世界的世界”的人们之中——《收获与播种》正是献给他们的——除Chevalley或可算一员,连一个人也没有。这当然属于那些我本愿与他谈谈的事情之列——可他已经不在,不能告诉我了……}。在最近的这几年之前,我从不曾疑心:那份全方位的信息——自1948年我最初接触科学世界起,我便自由地享有它——已在不知不觉的岁月里(我说不清从何时、如何),变成一份惊人的\textbf{特权}:我与一小撮伙伴共享的特权——一份\textbf{阶层的特权},借一个被用得有点滥的词说;而这词在此,依我看,恰恰表达着一份十足可触的实在。

但我的用意,不是给数学世界作一份“阶级分析”、罗列它的“力量关系”与“权力的手段”,也不是作一幅“风俗画”。该回到一个更有限的事情上了:在几位主要的当事人身上,就其要害的发条,理解我那场提前举行的葬礼的“市井小事”!

\subsubsection*{(d) 阴,作为仆人(2)——与慷慨}
\tocsubsub{(d) 阴,作为仆人(2)——与慷慨}
\label{ch18:local:h47}

\notehead{136}{}
\phantomsection\label{ch18:local:n43}

(11月28日)前面两条注,基本上是绕着Deligne与我之间那份阴-阴亲缘的主题、在数学工作与数学进路的层面上——所作的枝蔓。我不知道它们是否有助于把我对这份亲缘及其本性的知觉“传递过去”——对我而言,它无可怀疑。

我在别处写过:“在我的工作里,我同样‘阴’,同样‘海与运行’,达到了可能的极致。”细细想过之后,我要说:这话按字面并不真——“可能的极致”还可以更甚,因为(依我所知觉的)Deligne比我更甚。或者至少,那份“阴中之阳”,在我看来,在我身上比在他身上更显豁。我身上的那份烈性,在他身上取更沉着的姿态。在我大胆向前扑的地方,他不止一次停在一份审慎的观望上——而那观望,往往是有根据的。只要我有一星想法的苗头、一个可以进去的“抓手”,我便毫不犹豫地扎进一潭我感到有实质的数学泥沼,不操心先把出发的想法凑近端详端详(德语所谓“ihr~auf~den~Zahn~fühlen”……),也不预料这场混战的结局。有时,那想法站不住,缘于某个一眼可见的缘由——而我因一头扎进“汁水”里而把它放过了。我终究会明白过来——有时我觉得自己蠢透了;可我很少后悔自己扑了出去。正是这样、也只有这样,我才与一份未知的实质建立联系:蹭进去,不管“得当”与否。

我的朋友呢,先探测、先检视——然后才扑出去:当他感到稳当时才扑,未必稳在对岸——那要求过分——但至少稳在有着陆之处、稳在不至于空手而回。我从不曾在他的工作里有过任何\textbf{能量耗散}的印象——那份耗散,我自己常有——而毋宁是:在他那里,\textbf{招招中的}\origpage{632}。从这个观点看,他工作的风格带着一份\textbf{成熟}的印记,而我的,宁可带着一份\textbf{青春}的印记——烈性所至,有时乱糟糟。可我们初次相遇时,已近不惑的是我,他二十岁。而且不止一次,我在他身上感到对我的一份微笑的宽容——有点像一位善意的成年人对一个他喜爱的孩子:当他看到我又一次登上一项(小小的)“大家伙”、什么都不怀疑地……
此处唤起的这些侧面,想必难以在“誊清的”、已刊的工作里察觉——那些工作呈现的是一份反思的最终阶段、或至少相当靠后的阶段。我在工作中的要求,不比他低;这些注在手稿打字员或印刷者手中之前,我几乎从不交出,除非它们已达到满足我心中完全清晰之需的阶段。相反,在《数学反思》里(尤其《追寻栈》里)我所追随的那种书写风格中,工作里原初的行径,每页可见。读者可以在其中确认大量“失手”。它们全都幅度很小——多半在次日或隔日、甚至当天就被发现,并在其后几页里得到纠正。(事情如此,连我自己也惊讶——这就是我在别处谈过的、我数学工作里那份出奇的“轻捷”的一个征象\fnote{178}{(*)}{见“陷阱——或省力与衰竭”一注(n\textsuperscript{o}~99)。我觉得,这份“轻捷”,如今甚至比当年——我“出发”之前——更大。这在我看来,系于一场在逝去的十五年里在我身上完成的成熟;它在我的数学工作里、一如在别处,让人感到。}。)“小失手”在场的缘由之一,当然是我对一个七八年不曾触及的题目不再熟悉;而随着工作的推进、随着失去的联系渐渐接上,这些失手也愈发罕见。但这并不妨碍:这种一次次把我那份相当混沌的记忆所归还的东西、我当年多少知道的东西,不加迟疑当作“现钱”收下的做派,正极好地例解那份“闯将”的、有时乱糟糟的侧面——它(与其他侧面一道)构成我数学工作(或非数学工作)里“阴中之阳”的侧面。我深信:一篇同样自发的文字,倘若出自Deligne的笔,会更接近人们通常所认为“可刊”的东西——甚至,接近按他那套严格的标准的“可刊”。

\dmark{}我之所以在这里着力强调我朋友的工作风格与数学进法里那份“成熟”的品格、“阴之又阴”的品格,绝不是要借此暗示:他的工作里有任何失衡,也就是说,这份工作会打有“阳”的、“阳刚”的品质之缺乏或缺席的印记。倘若如此,他的著作便不会像Serre的著作、像我的著作一样,每页都带着那份精微而不会骗人的\textbf{美}的印记。但此处不是地方——一如在Serre的个案里、在我自己的个案里,我都不曾这样做——去逐项追索那阴阳、“女性”与“男性”的精微和谐:在他那份我所认识的已刊著作里,以及在我与他将近二十年的个人接触所让我认识的那部分工作里。

也不该以为:我对阴阳之均衡所作的这份确认,是一种老生常谈,会立刻适用于任何一个以这种那种身份扮“大数学家”的人。我方才唤起的那份对美的知觉,并不是在每一个在自己时代的数学上留下持久印记的数学家的著作面前,都同样在场、都同样程度地在场。这些人之中,我认识两位,如Deligne一样,在我看来无论工作还是人格都以阴为主导;而他们的工作,却从未在某一时刻给过我那份内在平衡的印象、那份从不叫人对美失望的印象。这两位同事之一,那份阴的失衡取如此极端的形态,以致他似乎全然无力——哪怕清晰地、正确地表述任何一个定义、任何一个命题(更不用说一个想法了……)——尽管他对许多事情有深刻的直觉,并且引入过许多重要而多产的想法。它们每一次都靠别人的工作才获得形体。分明地,他身上有一份罕见有效的压制,压着“阳”品格的力量——在他的工作里、在他的存在方式里都如此。这份压制取一种真正的无能之比例,包括在他的工作里:凭自己的手段,他绝无能力把最小的一件事做到底。他用一份自大狂的姿态补偿这份存在的无能,同时内化那些他乐于在自己身上培植的毛病,仿佛正是\textbf{凭着它们}\dmark{},他才得以设想\origpage{634}那些(在他眼里)使他成为\textbf{本}千年大智者的想法……\fnote{179}{(*)}{我在这里谈的,是我出发之前那些年月里得以确认的姿态与存在方式——那时我有机会与这位显赫的同事非正式地见面。自那以后,某种东西变了,并非不可能(尽管那会是罕见之至的事……)。}。

我在朋友皮埃尔身上感到一份反方向的压制:清除某些“阴”的特征,把他(或多或少成功地)引向一份超级阳的影像自我塑形。这份压制,当然离我方才唤起的那极端相反的个案还远。它不曾达到那样的地步:把那份美的感觉、那份不带一丝不适余味的满足之感——真正领会的征兆,时时刻刻给明晰与幽影、给奥秘各留其分——从读者或对话者身上抹去。这就是说:我朋友所选的那份“超级阳”的招牌,想必不曾侵越他的工作本身——我是说工作的时刻;在工作的时刻,“东家”的在场,想必也如(我相信)在Serre身上、在我身上一样地淡\fnote{180}{(**)}{关于这份匆忙的印象,我将在本注的子注~n\textsuperscript{o}~136$_{1}$(12月4日)的结尾处回头再谈。}。

可是,恰恰在选\textbf{题}的层面上,依我看,“东家”的角色才变得重要、乃至具有侵夺性。有那份固定不变的念头:与我这个人划清界限;由此,便拒斥追随他自己天性的某些癖好——它们在他身上与那被否认的主人影像系得太紧。于是,倘若他偶尔——如任何一个手段高强的人——证明了困难的定理(乃至“困难尽人皆知的定理”),甚至引入了好想法并把它们展开,他却决计不会“重新思想”——以他自己的方式、哪怕只在大轮廓上——一整门科学(比如拓扑,它分明需要这个……);更不会凭空创造一门新的科学,“把新的世界拖到光天化日之下”(如我在别处所写)(136$_{1}$)。然而,倘若说有谁,我对他的这份能力毫无怀疑,那正是他。倘若有一样东西,至今于他有所缺,那便是\textbf{慷慨}——真正的慷慨:它同时是一份沉静的确信,使我们追随我们自身天性的冲劲,任它把我们带向哪里,不操心有无勉励、有无“回报”\dmark{}\origpage{635}。

但也有那份单纯的快乐:盖大大小小的房子,让别人去住——不必非得是“一整门科学”或“一个新世界”的规模;搬石头、架梁,像随便哪个泥瓦工、木工,不怕因此被当作这个、像极了那个;或把(在某些人眼里)本该是极少数人禁脔的东西,交到所有人手边。这是一种服务的姿态、一份谦逊——它仍是方才唤起的那份慷慨的表达、那份对自己天性的忠实的表达。我的朋友把它换成了另一份自满的姿态(“我——做这样的工作!”)与一份种姓的姿态\fnote{181}{(*)}{这份“阶层”的姿态——在我的朋友身上、在那个“数学上流世界”里——首先出现在我3月份的反思里:在“职业道德共识——与信息的控制”与“年轻人的势利——或纯洁的卫道士”两条注(n\textsuperscript{o}~25、27)里;上周,它又在“阴,作为仆人——与新主人”那条注(n\textsuperscript{o}~135)里重现。},在那些号称“可接受”的工作主题的拣选上。

最后,还有第三份姿态或力量,“东家”凭它压在我朋友对工作主题的拣选上、压在他给自己圈定的探测对象上——一份给他划定强制性边界的力量。那便是“师父的下葬”的综合征,或\textbf{掘墓人的综合征}。这里,不只是止步于不提名那个必须被无视的人。还要把他的著作本身下葬,或更确切地说,把它干净利落地“锯断”——用\textbf{电锯},在他自己的工作里一如在别人的工作里,在从一具强健的树干上迸出的每一根主枝的水平上\fnote{182}{(**)}{我头一次直面“电锯”的实在,是5月19日,在那双注“继承人们……”“……与电锯”(n\textsuperscript{o}~91、92)的反思进程里;随后是四条棺材注(它们与“掘墓人”一道,构成葬礼的“灵车”,即X行列),5月21、22日(注~n\textsuperscript{o}~93--96)。}。一如我前天(在上一条注“阴,作为仆人——与新主人”里)还提醒的:在我“几何学家”时期——1955至1970年——所析出并展开的四个大主题之中,只有一个被我那辉煌的学生兼继任者公开“取用”;其余三个都被“锯断”了——当然,是以弱音\dmark{}\origpage{636}。其中一个主题1981年得到一份极局部的开棺,另一个次年——像几株孱弱的幼芽,从被截断的主枝那结痂的残桩上重新发出,人们还为这场合给它围上斑驳的花环与刺眼的霓虹,好掩人耳目……

\subsubsection*{注~136$_{1}$}
\notehead{136$_{1}$}{}
\phantomsection\label{ch18:local:n44}

(12月4日)\fnote{183}{(*)}{本子注(系于上一条注“阴,作为仆人(2)——与慷慨”~n\textsuperscript{o}~136)由该注的一条脚注生出。(参见该注结尾前第三段中的参见号。)}我自己的行径,不断地把我引向“重新思想”——彻底地——我在数学家路上遇到的东西:无论它是看似最微末的事物,还是有着“一整门科学”规模的事物。诚然,我和大家一样只有两条胳膊,我不可能每一次都把一项“彻底重做一整门科学”的工作纲领,实现到我在代数几何的情形里所达到的地步——那是从概形观念周围几个极简单的力量观念出发的。即便在这个个案里——我一口气投入了十二年数学家能量的很大一部分——我也远没有“收尾”预定的纲领:为此,我还得再要整整十二年!(而我出发之后,没有人操心去继续这项任务——它想必(错误地)显得吃力不讨好……)

作为我重思过一门科学、却肯定没有走得这么远的其他个案,我举出\textbf{同调代数}(交换的与非交换的——后者在我1955年最初的反思时还不存在),以及\textbf{拓扑}:随着\textbf{topos}观念的引入——它至今仍在等候它的时辰,好成为几何拓扑学家的日常面包,与人们今日惯于使用的种种“空间”“簇”观念平起平坐\fnote{184}{(**)}{可与3月底那条注“我的孤儿”(n\textsuperscript{o}~46)第二部分里的某些评注、及其子注~n\textsuperscript{o}~46$^5$至46$^7$对读。}。无疑,今日拓扑的某些重要部分,将不大被拓扑学中topos视点的系统展开所触及\dmark{}\origpage{637}。于是,这个视点依我看,宁可说是“凭空创造一门新科学”里的关键要素——这门科学实现了代数几何、拓扑与算术的综合(当我在五十年代乍到之时,它还全然出人意料)\fnote{185}{(*)}{见前一条脚注。(1985年3月11日)“全然出人意料”一语想必过分,因为这样一份综合的先见,已见于Weil猜想——它们曾如一股强大的灵感之源起作用。}。在新代数几何的奠基之外、并经由“平展上同调的掌握”(及随之而来的$\ell$-adic上同调的掌握),规划并施工这项尚在生成的新科学、为它发展坚实的技术基础——依我看,这才是我对我的时代的数学的主要贡献。那“Motif的yoga”——它仍然……尚属猜想,依我看像这份新科学的灵魂——或至少像它最要害的部分之一:它如此宏大,以致至今我还不曾想过给它一个名字。或许可以称它为\textbf{算术几何}:这名字暗示一幅“几何”的图景——人们将在“绝对基”$\operatorname{Spec}\mathbb{Z}$“之上”展开它;它容许向各种特征的“代数几何”之特化,也容许向“超越”的几何观念(在$\mathbb{R}$、$\mathbb{C}$或$\mathbb{Q}_\ell$等基域之上)的特化——经由“簇”(或更好地说,\textbf{多重簇})的解析或刚性解析的观念、及其种种变体。

我还看到又一门“新科学”——六十年代我就已隐约望见它,它汲取自1955年开始的同调代数反思。那是一场宏大的综合:综合来自同调代数的观念(它在与代数几何的需要——或更好地说,“算术几何”的需要——的接触中发展)、来自同伦代数、来自“一般拓扑”的topos版本,以及(自六十年代起尚在冥府的)(非严格的)$\infty$-范畴理论——或者,按我如今更愿意说的:$\infty$-叠($\infty$-champs)。我曾以为理所当然:这场综合将交由我那些上同调家的学生中的一些人接手——领头的是Verdier,他那著名的论文\fnote{186}{(*)}{关于这一点,见“赊账的学位论文与全险”一注(n\textsuperscript{o}~81)。}恰恰被认为朝这个方向走。依我看\origpage{638},发展一份令人满意的共同语言——具有一切可希求的普遍性与柔韧性——理应是几个年头的事:由一小撮心怀热望的研究者,做一场肯定引人入胜的工作。在我那些上同调家的学生作出一些极零碎的开端之后,我1970年的出发,敲响了放弃这项工作纲领的信号——连同许多别的我引以为心的纲领。这就是为什么,在1975年与Larry~Breen的书信往来里,我回到自己的某些想法,盼着能看到一幅世界之景重新活过来——我分明感到它“在路上”,而“人人”每每一与它照面,便小心翼翼地绕着走。在我致Larry~Breen的信里(重印于《追寻栈》第一章),我提议把这门尚在孕育的科学称作\textbf{拓扑代数}——一二十年以来,唯有我一个人隐约望见它\fnote{187}{(**)}{至多Deligne例外:我曾自信把一幅景传递给了他,而在我出发后的翌日,他便赶着把它与其余的一切一道埋葬。在《收获与播种》里,我屡次暗示我这整套“全方位几何”之奠基纲领中最古老的部分——尤其在“做梦人”(第~n\textsuperscript{o}~6节)里,以及在“我的孤儿”“本能与时尚——或强者的法则”“同伙”诸注(n\textsuperscript{o}~46、48、63$'$)里。}。终于,兴味索然、眼看无人肯接手这项二十年来燃着要做的工作,我便于1973年2月动起手来,写下《追寻栈》,至少以大轮廓画出我心目中该做之事的总设计。

分明地,方才所谈的“算术几何”,与拓扑代数,没有公度可言——而拓扑代数,依我看,它的一个主要角色,正是这门新几何发展中的“后勤支撑”。要使这门几何达到完全成熟的阶段——比如说,达到对Motif观念的掌握,堪比我们对平展上同调的掌握——想必须预期:好几代几何学家都得投身其中,比我目睹在场的那些更有活力、更大胆;更不必说,在\textbf{阿那贝尔式代数几何}的层面上,达到一份堪比的掌握——它(连同Motif)在我看来是算术几何当下便可辨认的两处“要害”之一\dmark{}\origpage{639}\fnote{188}{(*)}{(阿那贝尔式代数几何的若干主导观念,见《纲领概要》第2、3节。)凡言“要害”,我指这门“算术”几何的一个部分:它为这门几何带来直觉、引线与问题——相对于六十年代的所得,全然是新的。(这份“所得”,其实质在于一个框架、一门语言、一份同调与同伦的形式主义——为算术几何所涵盖的三门学科所共有。)或许,还该在上述两处“要害”之外,补上第三处:它与Motif亲密相连,即Langlands式的\textbf{自守形式}理论。我之所以对它略而不谈,是因为我对自守函数理论那份说来遗憾的、从来的无知。(我不知道机会会不会出现,逼我终于稍稍补上这份无知……)}。

还有第四个反思的方向,在我的数学家过去里进行过:朝着一门既有学科的“彻底”更新。那便是拓扑学里的“温和拓扑”进法——我在《纲领概要》(第5、6节)里对它略有铺陈。此处,一如自中学那些遥远的年头以来的许多次,似乎仍旧唯有我一个人感到那场该做的奠基工作的丰富与紧迫——而这份需要,在我看来,比以往任何时候都明显。我有一份极清晰的感受:温和拓扑视点的展开——照《纲领概要》所唤起的精神——将给拓扑学带来一场其分量堪比概形视点带给代数几何的更新,而这并不要求相当规模的能量投入。此外,我想,这样一门温和拓扑,终将在算术几何的发展中证明是一把珍贵的工具:比如,得以表述并证明“比较定理”——在复数域上有限型的分层概形所系的“投射有限”同伦结构(或更一般地,在该域上一个有限型的分层概形簇所系的同伦结构)与相应的“离散”同伦结构之间,后者以超越的途径定义,并以合宜的假设(尤其等奇异性假设)为模。这个问题,唯有在分层结构的一套精确“拆解理论”的术语里才有意义;而在“超越”拓扑的框架里,依我看,它需要引入“温和”的语境。

\asterisms

\dmark{}回到朋友Pierre~Deligne这个人:在1965至1970年与我数学上贴近接触的那些年里,他有的是机会,全面熟悉我方才以大轮廓巡过的这一整套观念与几何之景\origpage{640}。(温和拓扑的观念除外——如果我没记错,只是从七十年代初起,它们才开始发芽、并叫我着迷。)他对这份宏大纲领的角色是双重的、朝两个相反的方向。一方面,借着现成的$\ell$-adic上同调工具、借着(仍被遮蔽的)Motif理论的观念,他为算术几何纲领的发展作出了显著的贡献。最重要的,无疑是一门混合Hodge系数理论的启动,而尤其,是他关于Weil猜想、及其$\ell$-adic推广的工作。另一方面,撇开他工作直接需要的那些\textbf{工具}与观念(它们的来历,他曾系统地设法让人忘掉)不谈,他竭尽所能,叫其余的一切之自然发展落空:这便是“电锯效应”——我在葬礼反思的进程里已有的是机会谈它,包括(以暗示的方式)在上一条注(n\textsuperscript{o}~136)里。这电锯效应,曾因1981与1982年的几次局部开棺而部分被搅浑——“像几株孱弱的幼芽,重新发出……”——那是迫于眼前需要的骤然推力。(这些应景的开棺,上一条注的结尾处刚刚又提到过。)他还竭尽所能,不断地给人(从不曾明说……)这样一种印象:他所使用的那些观念、想法、技术、成果——他留神隐瞒其来历——应归功于他;而当他不曾慷慨地把它们派给昔日学生或合作者中的另一位时。

算总账——在这份对他那般顽韧地锯断与下葬之物的快速回顾之后——我回到上一条注里占上风的印象:我曾暗示,“东家”、利己的贪欲之插手他的工作,主要限于工作\textbf{主题}的拣选。毕竟,那掘墓人-电锯的姿态,在他的工作里分明可见——例外极少——\textbf{凡}有机会处\textbf{皆然};而我意识到,这些“机会”数不胜数!\textbf{这份掘墓人的综合征}(它无疑与超级阳价值的抬举亲密相连),依我看,\dmark{}对他的工作与著作\origpage{641}起过真正“具有侵夺性”的效应,与他的那些亲阳抉择的效应,毫无公度可言;而且这效应,绝不限于主题的拣选——“东家”把主题交到“匠人-孩子”手边,随后踮着脚尖退下。相反,依我看,整个工作进程中,“东家”几乎从不离开“匠人”:他太担心后者会忘记那些强制性的指令;换句话说,工作本身,常常被一些与发现工作之本性格格不入的\textbf{内在姿态}侵夺——那本性,是冲进未知。这件事,在葬礼反思的进程里曾被强烈地感到过多次;而我在这场关于阴阳的长反思中,曾倾向于把它忘了。

\subsection{天鹅绒中的利爪}
\label{ch18:local:h48}

\subsubsection*{(a) 天鹅绒的脚掌——或微笑}
\tocsubsub{(a) 天鹅绒的脚掌——或微笑}
\label{ch18:local:h49}

\notehead{137}{}
\phantomsection\label{ch18:local:n45}

(12月7日)一个多星期没有继续写注了,除了后勤的工作(包括给前面两条注写子注)。我不得不拔了三颗牙(快到六十的人,就是如此……)——必要而粗暴的一次入侵,害得我近来以略为减档的转速运转。我趁势转向积压的通信。那里,一切似乎已归于秩序……

在前面的四条注里(11月24至28日),我主要试着更贴切地圈定:Deligne与我之间,在气质与数学进路上的亲缘或互补关系——为的是给那“颠倒”定位:那份阴阳角色的颠倒,我曾自信在朋友竭力作出的自我呈现、与对我的呈现之中(至少在双方“数学”人格的层面上)知觉到。一路走来,关于他或我自己的实在的另一些侧面显了形;而在我们两个人之外,还有数学家世界的侧面,或者干脆说,人的世界的侧面。最终,依我看,最惹眼的新东西,是服务的姿态、以及这份姿态在科学世界里消失的种种征兆——我便以给这一阶段反思所起的“主人们与仆人”之名,试着暗示这一点。

\dmark{}回到“定位”某场颠倒这个最初的打算:此刻我有这样的印象\origpage{642},关于我与朋友的真实情势,我已圈定得足够贴近,可以给它下文了。一条不容回避的头一条确认是:那份关于阴阳角色颠倒的出发直觉——它在5月12日“葬礼悼词(1)——或恭维”那场反思的翌日来到我心头——确确实实是正确的。早在11月10日“阴的葬礼(阳葬阴(4))”(n\textsuperscript{o}~124)那条注的反思里,便已清楚:我的朋友竭力给他自己描画一份超阳刚的影像,给我描画一份超阴柔的影像。11月24日那条注“颠倒(3)——或阴葬阳”(n\textsuperscript{o}~133)里提出的问题是:这份呈现,是否确确实实是实在的一场“颠倒”。而“上涨的海……”那条注(n\textsuperscript{o}~122)里显出的“新事实”——即:一如我的朋友,我的数学进法的基调也是阴的、“女性”的——曾一度使人对此生疑。

但最后三条注的反思,驱散了这份疑团。一开始便已清楚:从来,我都被Deligne(一如被我的其他学生与昔日学生)知觉为——至少在意识的层面上——强阳的(也许太阳了……)\fnote{189}{(*)}{况且,时行的价值既然如此,我怀疑一份科学声名能否由一幅(被普遍承认、被普遍接收的)影像来承载,倘若那影像不是一幅“阳”的、乃至超级阳的影像。依我看,唯有在无意识的层面,我数学进法的“女性”本性,才曾被我的朋友兼昔日学生、以及数学公众(至少是与我工作的那类事情多少有接触的公众)所知觉。}。但事情显出:此外,在Deligne与我的数学层面的关系里,在一份强的阴-阴亲缘的底子上,还运作着一份阴-阳的\textbf{互补}(相对于那份扮演“原初”角色的亲缘,可以称作“次生的”)——在那份互补里,扮演“阳”的、阳刚角色的,确确实实是我:凭一份在我身上远比在他身上显豁的“阴中之阳”成分。

我在Deligne身上确认的那份故意——它依我看,从好些方向得到殷勤的回声\fnote{190}{(**)}{我这里想到《导言》里唤起的那“阵阵阴险的轻蔑与含蓄的嘲弄”(见《导言》10,“一份敬意的举动”)。当我看到昔日学生中最显赫的几位亲自定下这个调门时,我无须惊讶。这些年里传到我这里的大量“阵阵”之中,共同的,恰恰是一份居高临下的做派——对着我的数学进法与著作里打有强烈印记的“阴”特征。关于这一点,亦见6月23日那条脚注里的评注,注~n\textsuperscript{o}~96“棺材4——或没有鲜花没有花圈的topos”。}——因此向我显为确确实实一份\textbf{颠倒角色的故意}\dmark{}\origpage{643},更确切地说,\textbf{阴阳角色的颠倒}\fnote{191}{(*)}{这份颠倒角色的故意,头一次在我的反思里出现时,说的是师徒关系里的角色颠倒:我被呈现为我学生的“合作者”,而他本人则扮起平展上同调与$\ell$-adic上同调的\textbf{真正}奠基者与主人。(关于这一点,见“颠倒”与“葬礼悼词(1)——或恭维”两注(n\textsuperscript{o}~68$'$、104)。)耐人寻味的是:在“师——徒”这对“偶子”里,扮演阳的角色的(作为给予者、言说者)、“主动”的,正是师父;扮演阴的角色的(作为接受者、聆听者)、“被动”的,是徒弟。这里又是一次:我昔日学生漂亮地完成的颠倒,可以看作一份阴阳角色的颠倒——取同一方向(阴-阳变成阳-阴),一如构成我葬礼悼词之主要讯息的那场颠倒——那讯息,显形于“阴的葬礼(阳葬阴(4))”那条注里。}。依我看,这是葬礼的又一个重要侧面,叠加到前面已巡过的四个侧面之上(11月13与17日的“回顾(1)(2)”两条注,n\textsuperscript{o}~127、127$'$)。如今该做的,是把这五个亲密相连的侧面,拼成一幅连贯的葬礼总图。

这样一幅总图,若要令人信服,还应当在一份共同的视景里,汇集\textbf{三个相继的“层面”}。第一层面上,只有Deligne:我葬礼的大司仪,师父的非学生、非继承人——那师父被宣告死亡,没有存在的理由,也从未存在过……分明地,他是葬礼仪式的\textbf{中心}人物——除了亡人自己(而他,不过是个亡人,一个沉默的配角)。紧随其后,在第二层面上,是“我那些昔日学生的忙碌的一群,扛着铁锹与绳索”(凭记忆引用“葬礼的行列编成”里对各行列的罗列)。最后,在第三层面上,是(几乎)全体到场的会众:他们前来庆祝我的葬礼(以及四位共同亡人的葬礼——他们在钉得结结实实的“橡木棺材”里立正待命),并为下葬助一臂之力。

\dmark{}在这三个层面之间,似乎统治着一份完美的和谐、一份“\textbf{全体一致的协约}”——一如人们在任何一场照规矩举行的葬礼上\origpage{644}见到的:在满怀虔敬的哀恸的神甫、摆着合宜神色的亡人家属、与广大的会众之间:该起调的地方起调,该静默的地方静默,从不、从不走调。

顺着最后这幅意象说下去:我此刻看到自己被置于这样一个人的境地(不如那位亲爱的亡人自在——他分明置身事外……):面对如此感人的一幕,竟要无礼地猜测神甫、家属与普通信众在那些应景的庄严或悔罪神色背后,究竟怀着怎样的真实想法与动机。

反思已继续了一段时日,其暗中一贯的主导线索,便是那个打算:预备必要的条件,以便把握总图三个“层面”中最切近的那一个——披祭披的神甫的那一层,对不起,我是说我的朋友Pierre~Deligne的那一层。现在我愿把注意放到这一层面上。

我要一上来就说:总图诸侧面(或“联”)之中,在“掘墓人——或全体教众”那条注(n\textsuperscript{o}~97)里被摆到前头的那个——“对一场异见的报复”——在我看来,在我的朋友身上,即便它登场,也只扮演一份极不显眼的角色。我从不曾有过片刻的印象,觉得朋友皮埃尔哪怕丝毫感到被我的“异见”“牵连”。恰恰相反:这份异见是一场天赐的运气——他恐怕从来不敢梦想的运气——使他得以优雅地摆脱一位在场感过强的师父:就在那家机构里,他年方二十五,刚刚获得了数学世界里最令人艳羡(或至少最令人神往)的地位之一。这份异见在随后的月与年里愈演愈烈,依我看(或许不在意识的层面,但那无关紧要),被经验为一份更大的运气:它把一份可观的“遗产”拱手送到他嘴边,没有来自任何方向的抵抗之意向(如他随岁月推移而逐步确认的)\fnote{192}{(*)}{关于这份“遗产”,见“继承人”一注(n\textsuperscript{o}~90),以及“阴,作为仆人(2)——与慷慨”一注(n\textsuperscript{o}~136)的子注(n\textsuperscript{o}~136$_{1}$)。}。这份不期而至的\origpage{645}运气,他可不曾作态\dmark{}抱怨——哪怕在心底里、或在他自己不知不觉中!而且我觉得,等比例地,同样的确认,对我“出发”之前的大多数学生、无论如何对我那五位上同调家的学生中的每一位,想必都成立。他们之中若有人曾暗示过一份不满、一份因我的异见而来的受挫之感——无论是在心底里,还是以或多或少明白的方式\fnote{193}{(*)}{唯一一位曾让我听到这种调门(还带上一丝谴责的口气)的昔日学生,是Verdier,约莫一年之前。而在“生存与生活”的年代,他却似乎同情我的异见。甚至有过一段与他妻子伊冯娜友好合作的情节:记得是(如果我记不错)在Robert~Jaulin(伊冯娜曾是他的学生)倡议的一次巡回展览的组织上,我以“生存派”参与者的身份也参加了进去。}——我倾向于相信:那属于一种\textbf{合理化},为那对着天助地消失的师父的掘墓姿态所作的合理化,而远不是它的\textbf{缘由}(哪怕它是缘由之一)。使我在这份确信中愈发坚定——既就我那些上同调家的学生“一般”而言,就他们那位无可争议的头领Deligne而言——的是:那场即将到来的葬礼的先兆(只要合适的机会出现——而啊,意想不到的奇迹,它出现了!),早在我1970年出发之前便已可见,无论如何,早在那个著名的1965/66年SGA~5研讨班之后便已可见——那研讨班,注定要遭受我所知道的那场屠杀。这肯定不是偶然:如此完美地,他们五位全体\fnote{194}{(**)}{(12月12日)可我该把J.-P.~Jouanolou除外:他终于为研讨班接连编写了三讲,展开了一些观念与技术——那正是他自己的学位论文工作马上直接需要的。}都对那研讨班的命运漠不关心——他们学手艺的地方,连同那十二年间几乎唯有他们有特权认识、使用的漂亮数学一道。关于SGA~5的命运,我在反思中已说得够多,此处无须赘言。关于Deligne,我只提醒一句:在他1970年我出发之前所写的四篇文章里的三篇中,隐藏、或至少尽量遮掩并贬抑我的观念之影响的意图,分明可见——它并没有等我的“异见”。

\asterisms

\dmark{}那么,我的朋友身上那份对我对抗姿态——那份急欲取而代之\origpage{646}、抹而除之的贪心竞争者——的根与特殊本性,究竟是什么?这份姿态,自我与他相遇的最初几年起,便与一份温情的、信赖的同情、与数学层面的一份共融共存。我甚至确信:它自我相遇之日起便以弱音在场,无疑更早;而且,它与其说是被我身上的这个那个特点唤起的,不如说从一开始便是从我注定要在他身边扮演的那份角色中流出的——除非说,是那整个“特点”的总体:它们使我得以在他身边扮演那份角色。这也是他二十年来竭力抹去的角色;那份角色,无疑不出双方所求、纯凭情势之力,含有一个“父性”的侧面。而我心中毫无怀疑:冲突正是系于这一侧面打成结的——一份早已在他身上的冲突,远在他听说有人念出我的名字之前、甚至(无疑)在我们共同的师傅——数学——的名字被念出之前。

这份确信,老实说,不是一场反思的果实,我更不自许去“证明”它。毋宁说,它是在我出发之后的岁月里来到的,我自己也说不清何时如何;我相信,是渐渐的,凭着大大小小的征兆——它们没有一个曾让我驻足,哪怕一瞬——而它们全体,却终于留下一份认识的痕迹:弥散而欠完整,却终究是一份认识,有一天它就在那儿了……我当然可以凭一项费力的工作——把半埋的记忆挖出、逐个探测——深化并落实这份仍然几分难以称量的认识;而很可能,这样的工作会给我留下许多惊奇。可我不觉得自己有兴致去做。这无疑是因为(无论对错)我觉得:这真正说来不是我的工作,而是我朋友的工作——我那里将探测的东西,事关他的,远甚于事关我的。就我而言,我方才表述的这份直觉、或“认识”、或“信念”,已够满足我眼下理解的愿望;我对它毫无保留地信赖。

一如我一生中常常如此,我此刻面对的是一份\dmark{}\textbf{对父亲的对抗}\origpage{647}关系——我在其中扮演替代父亲、被他认作父亲的角色(依我看,与其说是收养他的父亲,不如说是被他认作父亲的人\fnote{195}{(*)}{(12月12日)写这几行时,我有意识:在这样一份角色“非对称”的断言里,该多么谨慎——何况这些角色是在无意识的层面上上演的。我猜,在这个层面上,在真正的数学交流之外,我想必曾在某个时刻,多少进入那被语境备好的“父性”角色。可这份角色,分明在我的生活里、在我与朋友的关系里,分量不能与我对数学的激情相比;它始终是偶发的,而且在我1970年从数学舞台“出发”之后,该不会再有痕迹。相反,我昔日学生对我这个人的眷恋——无论好也罢、(尤其)坏也罢——在其后的整整十五年里不断显形:既在他自己的工作里,也凭那份顶着逆风维持的、与我的持续的个人关系。}),连同我朋友那份颠倒阴阳角色的故意,旋即在我心中与“颠倒(2)——或暧昧的反叛”(n\textsuperscript{o}~132)那条注里唤起的情势联系在一起——我母亲与她父亲的关系,对我而言,是那份情势最极端的原型。然而,那份情势与朋友皮埃尔对我这份关系的差别,一望即知。在他与我的关系里,我从不曾知觉到一丝“反叛”的调门之影,或哪怕一份以些许恶毒、富于攻击性的方式、亮出爪或牙的对抗——哪怕是在一个微笑里。微笑,双方当然都不缺;但那是他一侧的同情之笑(如我所感到的),或有时是无邪的惊讶之笑,又有时几乎是忧伤的笑:当他得以确认(而我终于察觉到其中那一丝内心的得意)某些出其不意、以天鹅绒的脚掌落下的招式,在预定之处命中了。

换一种说法:这份对抗——无论它对着我,还是对着别人(当要经由他们打击那位亡故的师父,而他在他身上又始终好好地在场……)——从来、无一例外,取极阴的形式:那以最精微的温柔的全部表象为乐(并精于此道)、去打击去伤害、乃至去清除去压服的形式。他为数学家的招牌所作的刻意抉择是超级阳的(一如我的抉择也曾是——但没有比他更多成功);而依我看,在关系的层面上——至少对着我、对着那些他认为与我有牵连的人——那基调分明是、彻头彻尾是,超级阴的。(不过,在这件事上我要作唯一的、而且是重要的保留,日后还须回头谈它。)

另一处“一望即知”的差别——皮埃尔与我的关系、与那份“暧昧的反叛”之间:凭我对他家人的所知,我相信皮埃尔的父亲是一位性情温和而谦逊的人,因此绝非会激起一场反叛、再把反叛转嫁到替补父亲身上的那号“形象”\origpage{648}。

\subsubsection*{(b) 颠倒(4)——或夫妻的马戏}
\tocsubsub{(b) 颠倒(4)——或夫妻的马戏}
\label{ch18:local:h50}

\notehead{138}{}
\phantomsection\label{ch18:local:n46}

(12月8日)昨晚结束反思时,我有一种颇不好受的印象:觉得自己越来越看不懂了。睡前,我又逗留片刻,追随那场已写完的反思所引发的联想。我自信看到几点光亮浮现,我想它们将在今天的反思里充当我的灯具。

这些联想中最重要的一个,无疑系于我朋友身上那个“天鹅绒的脚掌”的侧面:以世上最天真的神态、“以最精微的温柔的全部表象”,津津有味地(而且有时入骨而不留情面地)抓挠。这幅意象,是在一场(与先前唤起的一份“反叛”情势所作的)已然触礁的比较的转弯处来到的,当场便向我显为意味深长,显为那份我打算探测的“对抗”的一个要害侧面。而且,回头看,“天真微笑与天鹅绒的脚掌”这幅意象的唤起——它复现着将近二十年贴身经历的精髓——依我看,正是昨天反思里的那个“痛点”、那个意想不到的“亮点”,当时我正在黑暗中摸索。倘若那份摸索与晦暗的印象在更远处仍占上风,那是因为:我太被片刻之前还在脑中的那些念头抓住——那些有待继续、有待安放的念头——以致不曾留心那精微的“咔哒”一声:它在意象显现的那一刻,便已在我身上响起。而其后的半个小时里,顺着系于这幅意象、与那场反思的另一个时刻的几场联想,注意又散了。直到现在,带着一天的回顾,接上被打断的反思之线,我才看到一份先前重读昨天的注时还不曾察觉的视景,正在就位。

倘若我留神追随所有联想中最强、与我亲身经历联系最紧密的那一条——暂且撇开其他更“有结构”、更“智性”的\dmark{}\origpage{649}——得到的是这个。我忽然看见自己回返到那一系列特殊的个案——或作为共同行动者经历、或作为近旁的证人——它们可以由一份单一的印象概括:\textbf{夫妻的马戏}——男与女这对偶子的马戏。这对偶子的马戏:成婚与否、有无孩子、年轻年老或少老颠倒、贫贱中拉着重负爬坡或富足中高车驷马——全然一样,夫妻的马戏并不因此改变。我忽然看见自己回返到它,经由这个马戏中打动我至深的一个侧面(老实说,过了很久,我才看穿其中的门道……):那非常特殊的、十足“天真相”的、“我什么也没说没做”的战术——“天鹅绒的脚掌”的战术——由女人在一个特定的游戏里使出:在那游戏里,始终是她以完美的手法、不动声色地领着;而始终是他跟着(而且常常挨着疼),什么都察觉不到。我所见过的夫妻之中,极少数不是循着这个调子运转的——变化无穷,自不待言:那留给双方即兴的天资,还有各自的脾气与种种情势。而且,就在今天,我还有机会看到一份格外眼花缭乱的演示——不过我谢绝在此为它枝蔓。

对这些马戏游戏的一份多少有色彩、有层次的描述——至少在大轮廓上——或者哪怕只是对其中调门的唤起(正是“她”那一侧的“天鹅绒的脚掌”),正是我刚重读过的11月12日那场反思——“颠倒(1)——或激烈泼辣的妻子”(n\textsuperscript{o}~126)——里缺席的那位要角。分明地,我逆着一阵迟疑推进那场反思,以致它终于摆出一副严峻的“力量与动机”分析的架势——那一天我分明不在状态!那也是在“阴阳之钥”里头一次谈到“阴阳的颠倒”。当时有点迷住我的那个极端个案——它直到昨天还在起作用——是我母亲的个案(重述于11月22日那条注“颠倒(2)——或暧昧的反叛”里,n\textsuperscript{o}~132)。可是,在我那份“四点分析尝试”里,我留神把这三“点”之前的几点圈定为:适用于我有机会多少贴身认识的大多数(即便不是全部)夫妻,而不必然以(暧昧的)“反叛”那份烈性的调门为主导(哪怕它以被遮蔽的形式在场)\dmark{}\origpage{650}。但这并不妨碍还有另一件共同的东西,它那天逃过了我。它只是昨夜才开始显出端倪——在那半小时里,我顺着“在状态”的那场反思的尾流,任思绪游荡。这件重要的共同东西——此前我只在“激烈泼辣的妻子”的极端个案里知觉过它——便是\textbf{阴阳角色颠倒}的精微游戏。

我犹豫:该写这份游戏是那份我方才暗示的权力游戏的“发条”,还是它\textbf{等同于}后者?肯定地,对她(而且常常也对他)构成男性角色之精髓、构成派给男人的角色的,是\textbf{权力的占有}——那占有固然常常是虚的,却无论如何从社会共识中汲取一份实在的元素。或许我曾倾向于低估这份实在元素的力道、男人的\textbf{象征}之力道:作为面对女人的一份\textbf{权威}的代表——尤其,它的力道,作为女人动机里的一个动力要素。我怀疑,对她而言,“做男人”、“做那个男人”,在一切之先,便是\textbf{行使权力}。在利己动机的层面上\fnote{196}{(*)}{阴阳角色的颠倒——在性爱冲动的层面上、在爱的游戏里——别处曾顺带谈到。(尤其见“接纳(阴的苏醒(2))”那条注。)性爱冲动,凭其本性,与自我的游戏无缘——尤其与权力的游戏无缘,即便自我渴望把它变成一件为自身目的服务的工具,并且(至少在某些窄窄的限度之内、以扭曲并残害那份原初冲动为代价)精于把它弄到手。倘若这两种“颠倒”之间——一边是阴与阳两股冲动那份\textbf{既}在情人身上、\textbf{又}在情郎身上的自由游戏,另一边是配偶一方对另一方那份无休止而阴险的权力炫耀的强迫游戏——果真有一层关系,那依我看,这关系多半只能是:两种形态,在任何时刻,都互相排除。},那“角色的颠倒”,想必不多不少,正是\textbf{女人对男人的权力行使}。

鉴于现存的共识,女人这份权力的行使,几乎只能以隐蔽的方式进行。它不在于发号施令,也不在于\dmark{}装模作样地做决定(期待决定被遵行),\origpage{651}而在于\textbf{让人运转}——尤其,把人折腾得团团转,而且始终不留痕迹。这便是那著名的夫妻旋转木马:它转动着,从不停歇!使它保持运转的战术……战术——由母亲传给女儿、由女人或少女传给小女孩、代代相传、无须言语的战术,便是昨天在路的转弯处唤起的那份战术:\textbf{“天鹅绒的脚掌”的战术}。只要留点心,人们便在无限的各式面目之下认出它:从“激烈泼辣的妻子”那极阳的个案——对我来说,以我母亲为化身——到“哀怨的(乃至愁苦的)妻子”那极阴的个案——我曾见它以另一位近亲为化身。

依我看,极少有女人不实行这份亘古的战术、不把它彻底掌握\fnote{197}{(*)}{同样真确的是:极少有男人,当“人家”对他们使出这份战术时,不立刻“入彀”。我自己便在我一生的大部分时间里,不战而降。它真正开始改变,只是随着禅观在我生活中的出现——那年我四十八岁(做好事,永远不嫌晚)。直到今天,我偶尔仍会上当。(说实话,不常,而且从来不会太久……)}。它主要是夫妻马戏里的日常实践,却并不限于此。依我看,女人对女人,它少有施展(或许仅仅因为:让一个女人“转起来”,比让一个男人难)。相反,在某些女人身上,这份战术变成一种第二天性,施于\textbf{一切}男人、或几乎一切——至少,施于那些被她知觉为具有鲜明阳刚品格的男人。

倘若我在这里说“战术”,这其实只表达了一个次要的侧面——恰恰是“战术”的侧面——的一份更要紧的实在:那份积习已深的内在姿态,对着一般“男人”的,或至少对着那个人的——父亲、情郎或丈夫,他在她的生活里扮演一份特权的\textbf{男人}角色,被(社会共识,或被她自己的抉择)授予一份\textbf{权威}。这份姿态,绝非总取支配欲的品格(如“激烈泼辣的妻子”的个案)——至少,不取人们通常理解“支配”一词的意义。它毋宁是一份饥渴——它有时变得吞噬一切:\textbf{不断地对另一方施加影响}\dmark{}、把他“维持在运转中”(言外之意:绕着她这个人运转……)\origpage{652}。为此,常常一切手段都是好的。其中一种行使行动——并由此行使权力的手段,便是\textbf{伤害}:而且有时,伤得越深越好,干脆把人放倒,直至毁掉——身体的或心理的——只要时机合适;而且始终,不着痕迹,“以最精微的温柔的全部表象”。我自己就不止一次被“放翻在地”!也常常,作为共同行动者或证人猝不及防,我被那伤人或毁人的举动之表面无端性惊得透不过气——它伴着一个天真的微笑、或一副心不在焉的神情,却始终不动声色,凭着一份不出错的直觉,抓住时机与地点,打击对方最能被击中的所在——无论那“对方”是父亲还是情郎、丈夫还是孩子,或仅仅一个相识、乃至一个陌生人(只要有机会出手、有机会命中……)。

\subsubsection*{(c) 天真的暴力——或传递}
\tocsubsub{(c) 天真的暴力——或传递}
\label{ch18:local:h51}

\notehead{139}{}
\phantomsection\label{ch18:local:n47}

(12月9日)我在这里触及那个极端、却绝非罕见的个案:\textbf{为暴力而暴力}、暴力与恶毒之中的\textbf{无端性}。这份暴力——无论它打击陌生人,还是最亲近、号称被爱的人——既非女人所专有,也非男人所专有;它既不“阴”,也不“阳”。可我此刻遇见它时所取的那份令人错愕而阴险的\textbf{形式}——戴着一副心不在焉、乃至天真温柔的假面——这份已然变得为我所熟稔的形式,依我看,尤其为女人所专有。这肯定系于那份“父权”的社会共识:它就女人而言,把权威与权力授予男人\dmark{}\fnote{198}{(*)}{况且,这份共识、以及男人在其与女人的关系中的权威,在最近几代人里已大遭侵蚀,如今愈发如此。对这份侵蚀,我将是最后一个抱怨的人!然而,法律与风俗里这份浮皮潦草的改变,似乎并不曾在两性关系的深层发条与“风格”里改变什么——尤其,不曾改变女人对着男人那份彻头彻尾、被精心遮蔽的对抗。这系于一个事实——本注反思的结尾处已强调过:这份对抗的姿态、及其借某份权力游戏(或权力颠倒的游戏)表达自身的手段,与其说是家庭之内种种“客观”条件的结果,远不如说是一份“\textbf{遗产}”之代代\textbf{传递}的结果。}\origpage{653}。这份形式,是她满足一份权力意志的手段——那份意志,纵然被迫(情势使然)走与男人不同的路,在她身上却并不因此少一分强制性、少一分吞噬性——恰恰相反!似乎不能在光天化日之下舒展、被预先注定过一份隐蔽的存在,只会加剧这份饥渴、并使它愈加蔓延,以致在许多个案里,真正“吞噬”她的生活、与她亲近之人的生活。

这份饥渴,远非总能(所幸如此!)达到全方位无端暴力的规模;而它得以施展的各个音区,也并不都落在暴力的调门里。虽然那份含蓄嘲弄的调门最常见——它们给一份蒙着面纱的对抗、或一份隐秘的敌意放风——那些单纯淘气的调门也并不被排除:其中带着一份宽容的亲昵,还透着几分顽皮。而如果说“天鹅绒的脚掌”这份经过考验的战术是女人的特权与首选武器,这份特权却绝非排他的。我不止一次、而且贴身地,见过男人挥舞这件武器\fnote{199}{(*)}{不过,在为我所知的个案里,我记下这份差别:倘若面对亲近或友好之人有一种表面“无端”(我是说,未被挑起)的暴力,那么每一次,那都是一个他对之怀有(哪怕不自觉)一份宿怨或长期敌意的人——那敌意落实为具体的怨怼(即便它们多半未曾成形)。这方面唯一的例外,是我的朋友Pierre~Deligne,在他对我的关系、以及在他归并到我这个人身上、当作属于我“势力范围”的那些人的关系里。于是,那是一份对抗与暴力的姿态(当然,压低了嗓门!),没有“个人”的缘由——我是说:没有缘于他对那些他竭力打击的人所怀的怨怼(真实的或想象的)。而反过来说,这是一种在许多女人身上遇见的行为——而且不仅(如这里)对着亲近的朋友、或单纯的相识乃至陌生人,也对着最亲近者中的某一个:比如情郎或丈夫(当然,而且优先),或兄弟、甚至自己的孩子。},其手法同样纯熟,毫不逊色\fnote{200}{(**)}{况且似乎,这份由无意识掌控的战术,总是从无意识那里继承那份“手法”与那份近乎不出错的稳当——它们在一份完全有意识的行动里,何其罕见。我不记得见过谁使这份战术,而不带着一份大师的手法。}。惹人注意的是:在这些个案的每一个里\dmark{},那据有了这件女人专属武器的男人,都是一个倾向于压抑自己存在中某些阳刚侧面、并(无疑由此)照着\textbf{母亲的模型}自我塑形的人\origpage{654}。

这同一份战术,也常见于孩子们——女孩男孩不分——对着父母、或对着代行父母之职的其他大人所玩的权力游戏之中,而且几乎是通则。这旋即引出另一份联想:在某些国家(过去的或当下的)里,作家或记者的情势——那里盛行着直接或间接的书报检查,使坦诚直白地表达真实思想感情成为不可能或有风险。后一种个案与前几种的主要差别在于:这里,诉诸间接的、蒙着面纱的、有时象征的表达,已不是无意识的作品,而是有意识的思想的作品。缘由无疑在于:当时存在着一份足够广布的共识——支持那些非正统的思想与感情(它们正要不动声色地“传递过去”)——以致当事人不再感到有义务把它们对自己也藏起来,生怕在自己眼中显得是个可憎的异类。唯有在极端酷烈的政治或宗教恐怖的个案里(如中世纪,或斯大林时代的苏联及其卫星国),非正统的意向才被迫(至少在某些人身上)再深潜一层:躲开内心检查官的目光,一如躲开那设立在风俗与警察机器里的检查。

所有这些例子似乎都暗示:“天鹅绒的脚掌”的风格(或“我什么也没说、没想、没要”的风格),或多或少自动地出现在任何一份多少持久的情势里:在那里,一份于我们不利的实力对比,使坦率直接地表达我们的感情、欲望、想法、意图成为不可能、或至少危险——尤其,那些对着被视为对我们施加强制之人(而尤其,那份恰恰自称要拦住我们表达真实感情的那个强制)的敌意或敌视之情\fnote{201}{(*)}{写这几行时,一个念头显现:我方才描述的情势,恰恰正是我们人人无一例外地在童年最初几年所面对的情势。我们无意识的很大一部分(那部分可以称作“垃圾滑槽”,在无意识的层面通常被知觉为一种“垃圾坑”)不是别的,正是我们童年的心灵对周围环境那份压力的回答:那压力强迫我们(这几乎是个生存问题)把身上一切触犯社会检查的东西,作为否弃的表示,深深地埋在自己眼前之外。这份检查,不久便内化为一位内心检查官,他阴沉的在场,保证这场早夭的下葬长治久安。然而,尽管有检查官,那些被规规矩矩埋葬的非正统冲动、认识与感情,却得以表达出来——有时以一种加剧而可畏的效力,以间接的、往往象征的、却依然充分具体的方式。“天鹅绒的脚掌”一栏,便提供一份格外“击中要害”的例子——而且常常,令人错愕……}。

\origpage{655}这并不是上述风格、及其所覆盖的内在姿态显形的\dmark{}唯一个案。常常,这份“实力对比”或多或少是虚的:它与其说对应一份“客观的”实在——顾及那个或那些被知觉为“压迫者”的人的真正姿态(或权力手段)——远不如说对应我们对它的\textbf{观念}(有意识的或无意识的)。这份观念,很少是对一份给定现实作细心的、明智的检视之果实;它几乎总是属于我们在幼年所接收的那一整包形形色色的制约——还须加上某些在我们人生极早的年代便已在我们身上完成的基本抉择。这样,无论在女孩还是在男孩身上,与\textbf{母亲}认同的那份抉择(当然是无意识的),便含着采纳一整套姿态与行为(尤其那些以“天鹅绒的脚掌”风格表达自己的),以及同时又采纳支撑它们的种种观念(多半无意识的,但这无关紧要)——比如关于某份实力对比的观念,及伴随这些观念的对抗反射。在相反的、与\textbf{父亲}认同的个案里——但当父亲本人在自己身上整合了某些典型“女性”的特征(或至少在我们社会中如此被视为的特征)时——可以想见,效应可能与第一种情形全然相仿。

我想在此抵达的点是:在我们当今的社会里——至少在我曾置身其中的那些圈子里——依我看,这份风格(“天鹅绒的脚掌”)、以及我此刻检视的这份“女性”内在姿态,只在很有限的限度内,才是对社会或我们童年周围的特殊情势(乃至我们成年某一时刻的情势)所设立之客观实力对比的一份自发的个人反应;它毋宁是一份“\textbf{遗产}”,得自我们父母中的一位(甚至有时同时得自两位?\dmark{}\origpage{656})——而那一位,又是从他自己的父母之一那里得来的。分明地,这份遗产优先循着\textbf{母系}的线路传递,首先由母亲传给女儿。但我不止一次贴身见过由母亲传给儿子的传递。没有任何东西让我以为,这份传递不能偶尔也由父亲传给儿子、甚至由父亲传给女儿。

\subsubsection*{(d) 奴隶与傀儡——或闸门}
\tocsubsub{(d) 奴隶与傀儡——或闸门}
\label{ch18:local:h52}

\notehead{140}{}
\phantomsection\label{ch18:local:n48}

(12月10日)我想回到几场系于\textbf{无端暴力}主题的联想。那正是昨天反思由之起笔的主题;随后我离开了它,转而检视权力游戏中的“女性”风格(或“天鹅绒的脚掌”),以及它作为对着他人的对抗姿态之表达手段(尤其,对着那些被感到强阳的、或以任何身份处于权威、声望或权力之位的男人)。

一如我昨天所提醒的:(表面)无端的暴力、“取乐”的暴力,并不比女人更属于男人。人人都曾有机会在路的转弯处猝然与它照面:或顶着“最精微的温柔”的面孔,或顶着皮靴踹进肚腹、一梭子机枪扫射的面孔。后一种风格,肯定是“阳”的风格,在眼下这些号称“和平”的时代、在我们这样的文明国度里,总归更罕见。对我们大多数人——教养良好的、在富足国度里处境或多或少优渥的人——这份直言其名的暴力,不属于日常经验,一如另一种——那份压低了嗓门的、天真相的暴力——属于日常。然而,只消翻一翻随便哪家大报的“社会新闻”栏,或听听新闻广播\fnote{202}{(*)}{这些事,说实话,我早已不做了,只满足于经旁人之手得到的偶尔的消息。},便会发觉:“硬”的无端暴力,即便在我们这里,依然随处可见。它还不总发展到\dmark{}顺手割断那位\origpage{657}匿名小老太太的喉咙——只因心血来潮要偷她的家。可当百无聊赖想冒险的年轻人“借走”人家大意地敞开在家门前的汽车时,把它丢在十公里二十公里外的沟里之前,却很少有不好好把它砸个稀烂的。即便在我有幸居住、不必过分担心的宁静乡间,一座最不起眼的农舍或棚屋,空不了多久,便被洗劫一空(这是实用性的),外加被痛痛快快地糟蹋(这是为了取乐)。在我方才唤起的这一切个案里,暴力的无端性显出得格外触目:因为被它打击的那个他(或她)是个陌生人,常常是个从未见过、也永远不会再见的人。

这是一种可以称作“\textbf{匿名}”的暴力。毫无疑问,战争从来便是这类暴力的集体狂欢——那是无偿杀人的机会称王的年代,是一个无名小卒的性命,面对扣动扳机的快感、面对体验自己叫一副无名的孱弱身影在眼前瘫倒的权力,分文不值……

倘若世上有一件事——就我记忆所及——每一次都叫我不知所措、哑口无言,那便是再度直面这份超出领会的暴力:那份只为打击与毁灭之乐而打击、毁灭的暴力。倘若世上有一件事在我们身上打下“恶”的那份不可磨灭之感,那不是死亡,也不是身体所能忍受的苦难,而是这件事。而当这样一份暴力——无论它取硬的面孔还是软的面孔,显得“大”还是“小”——由你所钟爱的一个存在冷不丁地朝你来,它注定击中得又狠又深,注定唤起(或重新唤起……)并向你倾覆一份无名的焦虑。这份焦虑的根扎得最深,当它找到童年、乃至婴幼期那片松软而新鲜的土壤扎下根来。这份焦虑——它是我童年生活、也是我成年生活里“世上守得最好的秘密”——是在我六岁那年,在母亲的手里,在我身上显现的。

直到五十一岁,1980年3月间,我才把这份焦虑在我生活中扎下根的那段情节挖出来。焦虑对我的掌控,早已在此前被大体拆除引信——至少在很大的限度内\dmark{}——随着禅观在我生活中的出现(1976年),它在其中逐步占据\origpage{658}越来越大的位置。我与焦虑的关系中的第三个决定性转折,发生在1982年7、8月间,在对日常生活里焦虑之机制的一份细心检视中。自我童年直到成年,那些制造焦虑的情势,一直是这样的情势:它们在我存在的未知深处,使我重新活过“那超出领会的东西”。也恰恰是那些情势:我又一次面对那份表面无法解释、无法把捉、无法还原的暴力的熟悉征象……这份暴力的骤然闯入,使一阵惶乱无措的焦虑之浪骤然复起、倾覆——随即被制住、被压下。这份本能反应,一如其旧,直到今天,出入无几\fnote{203}{(*)}{(12月14日)更确切的说法是:这份反应保持“一如其旧,出入无几”,\textbf{直到}我1982年7、8月禅观的时刻。从那以后,使我猝不及防的“挑衅”虽多,那份所说的“本能反应”却只现过一次,在一年前。它当时成为一场短短的、几个钟头的“应景”禅观的契机,把情势整个澄清。一旦一份内在的混乱情势被朴素地面对、被承担,那份伴随它来向我们传递我们混乱之讯的焦虑,便消失得无影无踪——只留下一份认识、一份焕新的宁静的痕迹。}。然而,这些年里若有什么改变了,那是由于一份\textbf{反思}的出现——它在焦虑的尾流里,使那曾顶着“超出领会的东西”、“谵妄”的吓人假面显形的东西,变得可懂、乃至显然;而尤其,两年以来,由于一份\textbf{对我自己的注视}的出现:一份对这份焦虑本身的关切而温情的注视——一份威严的反射运动,正想叫我把它对我自己也藏起来。或者换一种说法:我与焦虑的关系,尤其两年以来,已变成一份不再彻头彻尾拒斥、不再驯兽师、不再掘墓人的关系,而愈发是一份\textbf{迎纳}的关系:细心的、温情的,迎纳它带给我的、关于我自己的讯息——既关于我的现在,也关于我的过去、及其在我的现在里的行动。依我看,这是我至今跨出的最后一步,朝着一份对他人、越来越全面的内在\textbf{自主}——而首先,是对我的亲近之人、我的朋友\fnote{204}{(**)}{这“最后一步”,在“接纳”(n\textsuperscript{o}~110)的结尾处已谈过,取一份稍不同的光照:一份相对于\textbf{寻求赞许}或\textbf{寻求确认}之需要的解脱——“那份需要,真正构成了那个‘挂钩’,含蓄而坚不可摧,冲突凭它得以在我们身上‘挂上’,而我们由此……处于他人的依附之下……总而言之,凭它,他‘拿住’我们,并且(不动声色地)随心所欲摆布我们……”。(这一段,分明像是今天写的——可我发誓,一个字也没抄!)
我说不清,前面是否还有其他这样的“步”要跨——它们将给我那份回顾,把眼下这份自主看成仍是相对的、而非完整的(可我却倾向于——也许有点天真地——相信它完整……)。
一份对焦虑的松弛而专注的关系之绽放,确确实实代表一份在对他人的关系里的\textbf{解脱}。因为(如下一段所说),正是他人得以“随心所欲地摆布我们身上那份焦虑的闸门”的可能(尤其借着犒赏与拒斥的交替——由他配量、由他施用、凭他的手法),构成他对我们的主要的权力手段。}。

\dmark{}依我看,正是那份不直言其名的暴力、那份“女性”路数的暴力,最\origpage{659}有力地制造焦虑,远甚于那一拳打在脸上的更耀眼的暴力。那玩弄压低嗓门的暴力的人——并由此也在那些秘密的闸门上作文章:它们在别人身上释放无名的、无面孔的焦虑之波——手里握着一件比一份权威、或一份单纯的强制权力更可畏的武器。而且,带着一副天真的神态、随心所欲、凭兴致去摆弄那些焦虑的闸门,代表着一份权力——无疑更锋利、更可畏,即便它始终隐蔽:胜过一切事实的或法理的、由社会共识设立的权力。这便是女人对男人的“正当报复”,在一个男人自称(或曾自称)支配她的社会里;这也是“他”为那份虚幻的至尊(现在的,或过去的)所付的价钱。倘若她是\textbf{奴隶}(在我们的地界,她越来越不是了),那他在她手里便是\textbf{傀儡}、或相差无几(而他今天,一如既往地是)。

几年来,每当我直面一份无端暴力的情势——无论它施于我、还是施于他人,无论它以粗暴的方式、还是阴险的方式显形——便有一份联想以无可回驳的力量来到:\textbf{对己之轻蔑}——或者毋宁说,我\textbf{看见}这份对他自己的轻蔑,在那个公然或暗自以轻蔑待人的人身上。我毫无怀疑:这在我身上,不是\dmark{}一个简单的按钮机制、一个“哲学的”或“心理的”偏好,我\origpage{660}会乐意时不时拿出来遛遛——比如,作为凭一句有说服力的套话驱除我所谈的那份焦虑的手段,把一张万能的标签麻利地贴在一个吓人的陌生人身上。这只是一份\textbf{认识}:一份本质的、深刻的、(一旦看见便)显然的关系。

这份认识不“清除”任何东西;它只是让我得以给一个未知\textbf{定位}。它绝不是一名哨兵,守在那儿拦焦虑的路、或把焦虑逐出阵地。一份认识——取我所理解的意义——的本性不是这样。一份认识属于一份内在的\textbf{宁静},它有助于给宁静以根基。而那种使我们不停地想拦住“闯入者”、生怕它们搅乱一份拼凑出来的“宁静”的东西,是我们身上的一份躁动。我所谈的宁静不怕闯入者:它使我们迎纳闯入者。与焦虑的新相遇在水面制造的躁动,不打扰这份宁静,反而助成它。

\subsection{暴力——或游戏与刺棒}
\label{ch18:local:h53}

\subsubsection*{(a) 正义者的暴力}
\tocsubsub{(a) 正义者的暴力}
\label{ch18:local:h54}

\notehead{141}{}
\phantomsection\label{ch18:local:n49}

(12月13日)上一条注里,拿“奴隶”与“傀儡”开的那个玩笑,我肯定又找到了得罪所有人的办法,而且(倘若有人读的话……)要被所有人骂个遍!除非那位假想的读者(或女读者)满心欢喜地鼓掌——谁知道呢——深信那幅影像发射得漂亮、适用于全世界,唯独不适用于他自己(她自己);或许至多,还除开那位出言讥讽的作者。不过,凭这份假设,他给了我这个人一份根本配不上的信任。我至多敢承认:这几年来(尤其,自1982年7、8月那场关于焦虑的禅观以来),我开始走出、甚至已经走出了那著名的“马戏”——夫妻的马戏,固然;但也走出那些与它如手足般相似的其他马戏。《收获与播种》第一部分里,甚至有一节预告着这个色调,题为“旋转木马收摊了!”(第~n\textsuperscript{o}~41节,今年3月)。那里说的不是夫妻的马戏,而是一场数学的马戏:像所有人一样,我曾在其中转了半生。可同样真确的是:那节有着许诺性标题的文字几周之后——4月29日——出现了一条注“一脚踏进旋转木马”(n\textsuperscript{o}~72),它的名字似乎要唱另一种调门\dmark{}\origpage{661}!与从前的差别,或许在于:倘若此地彼处,我还会在某个旋转木马上转(而依我看,如今还继续吸引我的,唯有那数学的木马……),那么拽着那些让我原地打转的线的,是我自己(至少是我身上的某个人),而不是任何人;而且那些线,对我而言已不再隐形。

这些保留说完,我可以说:我成年生活的最大一部分(更确切地说,直到禅观发现的时刻),我“入彀”得飞快——一如所有人,仍然——无论在夫妻的旋转木马里(它足足欢快地转了二十年!),还是在别的木马里。我并不为此懊悔,因为我对一切这类木马的认识,首先归功于我自己转过的那些。我转了那么久,是因为学生学得慢——也肯定因为,在不止一个方面,我在其中找得到诱饵。到头来,它们的力量与魅力终于消散了——想必如此……

依我看,在所有这些木马里,我始终是“入彀”的那一个,而从来不是“让人入彀”的那一个。或者换一种说法:我相信,对那著名的“天鹅绒的脚掌”风格,我从不曾有过一丝一毫的癖向——我有时会狠狠地亮爪子;可我相信,从来不曾把爪子泡在一份天鹅绒的绒毛里。这是众多征象中的一个:它们证明,在我之结构的层面、“东家”的层面——即我身上受制约的东西——那基调是强“男性”的,这一次毫无歧义。阴的、“女性”的调门,反而支配着“孩子”的层面、我身上原初的东西的层面——也就是说,也在认识的冲动、创造官能的层面。

我还想就我生活中的“无端暴力”再说几句。三天前那条注里,我从被这份暴力瞄准的人、或至少在他人身上直面它的人的角度(哪怕只作为证人)谈到它,当时我写道:

“倘若世上有一件事——就我记忆所及——每一次都叫我不知所措、哑口无言,那便是再度直面这份超出领会的暴力:那份只为打击与毁灭之乐而打击、毁灭的暴力……”

\dmark{}这几行、及其后的几行,与实在相符——至少与我亲身经历的实在相符\origpage{662},肯定也与像我一样直面过这份暴力的无数男人与女人的实在相符。它们可能给人这样的印象:写下它们的人,自己与这份暴力整个无缘,一生免于这类谵妄。然而全然不是。我记得我一生中某些关系——共四段,其中三段落在我的童年或少年(八岁到十六岁之间):那些关系,打有敌意的印记,却没有任何具体的个人怨怼作根基,而以系统的、毫不容情的嘲弄表达自己,或以拳打脚踢及其他暴行表达。第一次,受害者是一个同学(还在德国):全班的受气包。情形拖了好几年,我记得。接着两个个案发生在战争期间:我(从一座法国集中营出来之后)住在瑞士救援会在利尼翁河畔勒尚邦的儿童之家——“盖斯皮”——1942至1944年间。这一次,“恶人”是我的一个同伴(他的父母,一如我的,该是被当作德国犹太人拘留了)和我们的两位监护员之一:两人都同我一样以德语为母语。他们俩,都多少是一群少年男女——有时毫不容情的一群,我就在其中——的受气包;可我相信,我给他们俩的生活,比全伙任何人给他们过的都更难受。同一个屋檐下的共处、难民营地那份朝不保夕的共同处境、盖世太保搜捕犹太人的恒常威胁,本该在我身上唤起团结与尊重之情——却全然没有。

在这三个个案里,我把那份恶意的箭头对准的人,天性温和、相当胆怯、绝不好斗——我于是把他归入“软”或“懦”一类,而这些品性也就被算作使他显得不体面的特征。在一个被暴力之风、被对个人之轻蔑吹得满目疮痍的时代,而我本人满怀对战争暴力、集中营暴力及其一切随行物的憎恶——我却感到:我对他人施加的那份轻蔑与暴力,完全名正言顺,理由单纯:我乐意把他归入“讨厌”一类(以及其他种种形容词……),此后,一切(或几乎一切)都变得被允许了,甚至变得高度可嘉。我这个自诩头脑“合乎逻辑”而公正的人,\dmark{}当时却不曾看到:我的行为、以及它那份由一份反感所作的辩护(那份反感的真正本性,我当然不曾\origpage{663}想过探测),与三十年代那些道貌岸然的德国人——对着“肮脏的犹太人”——的行为,一模一样(那些事,我在童年曾贴身见过);而且正是这些行为,使得那场席卷当时世界、空前绝后的暴力泛滥成为可能。我当然(循着我父母的脚步)装作远离那份暴力,视之为一场怪异的失序(有时甚至“超出领会”)。我满怀倨傲的居高临下,对着所有那些——士兵或平民——以这种那种方式同意在那些英雄的尸坑及其随行的种种可憎之事中充当主动或被动齿轮的人。而与此同时,以我微末的程度、在我自己有限的行动半径里,我做着所有人做的事……

倘若我试着辨明,这份服务于一份轻蔑与暴力之故意的、如此奇怪的盲目,其缘由何在,得到的是这个。我童年自五岁起自己曾承受的那些暴力——它们从不曾被当作暴力指给我注意——终于在我身上造出一份慢性的紧张状态:它始终无意识,又被一份百炼的意志小心翼翼地控制着。这份紧张、这份没有特定靶子的攻击性之积聚,造成一份攻击性放电的需要。可这份“需要”并非身体性的——在所有这些个案里,凭一项相宜的身体活动来泄愤的机会,从不曾缺——它分明是\textbf{心理}的。肯定有一份积聚起来的宿怨——当然主要无意识,且不曾落实为对某个特定的人(比如说,我的父母之一,或代行其职的某人)的可触怨怼:我本可以把宿怨之情转嫁到那人身上,并给它一份具体的、甚至暴力的表达。我身上想必有一份“空置的”暴力、一份弥散的、游荡的暴力,寻找一个可以卸下的靶子。似乎常常是动物(昆虫、蟾蜍、狗或猫、甚至牛或马……)为这类游荡的、寻找牺牲品的暴力付出代价。我的情形不是这样:我不记得我一生曾折磨过任何动物,无论大小。显然,我需要一个离我更近的替罪羊:一个\textbf{人}!而人找替罪羊,肯定从来不难找到。

\dmark{}我毫无怀疑:我方才写的,确实描述了实在的某一侧面。可我感到,这份描述\origpage{664}仍停在事情的表面:它只圈定某一个“机械论”的侧面,而不曾真正更深地进入那份无意识的亲身经历。眼下,在那份经历的位置上,有一大片“空白”、一片虚空。现在不是越过它的时刻与地方,去更深地探测这片“空白”覆盖着什么、什么东西溶解在这片“虚空”里。是不是那份著名的“对己之轻蔑”?——它三天前那条注里还以那样专断的姿态申张自己,而此刻,既然事关\textbf{我自己},它仿佛无影无踪地消失了?现在是时候了,机不可失,该把这件事弄个水落石出,澄清那份顽固而暧昧的“模糊”——它继续打着我对自己认识的印记,一如当年那份包围着焦虑在我生活中的角色与存在本身的“模糊”。焦虑,曾是我一生“守得最好的秘密”——我是这么觉得。会不会还有另一个秘密,守得更好,自我禅观以来,只在此地彼处、在两三次场合里被擦过?我有这样一种感觉:我仿佛已掌握了一切,足以弄清它的究竟——包括那份熟得不能再熟的兴致的骤然涌入:它告诉我,动身的时机已经成熟!可我感到,我不会在这里做这件事——在这场某种意义上“公开”的、或至少注定要拿去发表的禅观里。这场禅观,至少将在许多别的功效之中,有一功:意外地催熟一桩突然变得切近的问题——它终于被承认为理解我自己的关键,而此前,它不过在一张也许永远望不到头的长长的等候名单上,摆着“百问之一”的架势……

并非全无可能:我将还有机会遇到那三个男人中的一位(其中两位与我年纪相仿)——他们当年曾是我身上一份暴力与攻击性的无辜靶子;或者,至少我将有机会给他们中的某一位写信。能在充分明白事情原委的情况下郑重赔罪,这对我将是一件好事。或许,对他也将是一件好事。可奇怪的是:我没有这样的印象——觉得那三人中任何一个曾真正为此怨恨过我,觉得我的暴力曾在他身上唤起一份专对着我个人的敌意;毋宁说,依我看,他被卷入的整个语境,必定被他经验为一场劫难——一场他根本无从逃脱的劫难——而我这个人,更多地被知觉为这场劫难里众多配角中的一个,而不是那个毫不容情的折磨者\dmark{}\origpage{665}(而我是),被仇恨的。当然,我可能错了,而且可能永远无从知道——正如也可能,我将有幸有一天直面这份我在盲目中播下的业……

我相信,在“盖斯皮”情节之后的年月里,我身上有过一份成熟——尽管(就记忆所及)不曾有过任何对此的反思。总之,其后的岁月里,我身上形成了一些有效的反射:它们会拦住我,不再与全伙人对一个成员的集体暴力为伍。我不相信这类事在我成年生活里重演过,也不相信我曾受试探再去扮演那样的角色——我该感觉得到,那个角色多么虚伪,在那副欢快的、“讲体育精神”的外表底下,多么没有勇气。但这并不妨碍:战后,生活照样大量地在我面前堆起那些浸透蒙面暴力与焦虑的情势,在我身上延续那已然打下了我童年与少年印记的深层紧张。正是在这个语境里,才有第四段关系:它打有偶发的敌意与暴力的印记——我可以称之为“无端”的:不由具体的怨怼奠基或挑起,甚至(我相信)也不由任何可称作“挑衅”的行为挑起。那是与我一个儿子的关系。然而我知道:我对他的眷恋并不少,我“爱”他也并不少于我的其他孩子。可在无意识的某个层面上,我身上想必有一份对他这个人某些侧面的拒斥——恰恰是那些使他比他的兄弟姐妹更温和、更易受伤害、也更难把握的侧面。分明地,他全然不“入画”——比我的其他孩子更不——入我那些我本乐意在孩子身上看到实现的漂亮的超级阳影像;而尤其不入画的是:环绕着他人生最初两年、把他打得很深的那些非常严酷的情势,使他更难与父母结成信赖的关系。总之,在他还与我同住一个屋檐下的时节——直到他十岁上下——我有时会对他施以羞辱性的惩罚,以雷鸣般的嗓音颁布。这些东西,整个沉入了遗忘;连同某种终于浸透了家庭空气的氛围——是两三年前与他的姐妹和两个兄弟的几次对话,才恰逢其时地把这些事情多少\dmark{}重新捞回我的记忆\origpage{666}。或许有一天会来到:他也愿意同我谈这些——他,在我的孩子们之中,也许是那份浸透压低嗓门的焦虑与未被承担之紧张的家庭气氛的最大受害者;或至少,是在父亲手里“挨得最多”的一个——尽管他们每个人都分得了一份很足的“父母的包”。至少我知道——而且我为此高兴:拦着我的任何一个孩子与他父亲保持一份简单而信赖的关系、拦着我们一同谈论并探测那份沉重的过去的,不是一份他们还留着、还要设法藏起来的对我的畏惧。

但在这里,同样,在这些注里深究一份复杂的情势——它牵涉六七个人,一如牵涉我自己——不是地方。对我至关紧要的,首先是无遮无掩地确认:这份表面无端的暴力——它曾那么多次“叫我不知所措、哑口无言”,当我在他人身上遇见它时——曾偶尔在我生活中、在我自己的行为里此地彼处地显形。这份确认,不是怀着什么特殊的“意图”作出的;它不自许去“解释”或“原谅”任何人身上的无端暴力,一如那份暴力也不该被用来解释或原谅我的。并非不可能、甚至很可能:随着反思的深入,两份暴力——他人身上的与我身上的——终将相互照亮。这类事情,多半会自己到来,而且不招自来。我之所以作出这份确认,只因为它就在路上;而(除非甘愿不再真)我不能在这里不提它。

\subsubsection*{(b) 机械与自由}
\tocsubsub{(b) 机械与自由}
\label{ch18:local:h55}

\notehead{142}{}
\phantomsection\label{ch18:local:n50}

(12月14日)昨晚的反思,恰好提醒了我一件人们那么容易忘的事——而尤其(在此)是我,\textbf{我}那么容易忘的:我并不比任何人“更好”;我与所有人裁自同一块布;恰恰像我的某位朋友——我正准备把他架上板凳,置于一份毫不留情的注意的中心……

昨天,我给出了一份关于(表面)“无端”暴力之显形的描述:把它作为一份积聚起来\dmark{}的紧张与攻击性向某个替罪羊的放电——那人\origpage{667}因这样那样的缘由,恰好“对得上号”。这份“机械论”而浮面的描述——它肯定“众所周知”——可能助长一份同样“机械论”的\textbf{态度}:对着这份暴力,在自己身上或在他人身上。于是,暴力被看作一份不可规避的宿命,一份可叹地扎根于心灵结构本身的宿命——我们又能拿它怎么办!这样一种姿态,顶着一副“理性”或“科学”的面孔,依我看不过是\textbf{退位}的一份合理化:在自身与他人之中,对着一份创造性的\textbf{自由}的在场的退位——那份自由,向每个人敞开那份抉择:\textbf{承担}我们所处的情势,而不是被动地顺着那些划定的机制之坡道滑下去——它们时刻准备着接管我们。倘若人们难得动用这份“自由”的抉择为真,那么这份抉择的单纯\textbf{在场}、我们身上那些创造的可能之在场——无论我们用还是不用——已经把事情的本性整个改了。正是凭\textbf{这个},且唯凭这个,那些包含人与人关系、或人与自身、人与周围世界关系的情势,才有一份维度——当涉及的不再是人而是(比如说)计算机,无论多么完善,这份维度便缺席。也正是由此,对我们每个人显出那份\textbf{责任}的特权:对我们行为的、对我们行为之动机的责任。这份责任,绝不因我们常常求那份供我们使用的便利——把自己真实的动机对自己藏起来——而被解除。

回到眼下的个案作例解:倘若我曾扮过好人——同时动用我折磨某位不曾伤害过我的同学的权力——那是因为:在一份表面的“真诚”背后,我曾选择了一种粗劣而惊人的自欺态度,当时便与四十年后回顾时同样扎眼。那确确实实是一份\textbf{抉择}:没有任何东西强迫我做;它等于对我身上积聚的紧张与攻击性闭眼(同时当然还援引漂亮“非暴力”的念头),并把它们“悄悄地”(原文如此!)卸到随手可得的替罪羊身上。这样的暴力——也就是说,几乎全部在人类世界里肆虐的暴力与可憎之事——只有在一个\textbf{条件}下才能发生、其秘密的功能才能完成:这个条件本身严格保持秘密(哪怕它\dmark{}扎眼\origpage{668}之至);也就是说,得让自己给自己“把鱼目当珍珠”,凭一份信服玩一份粗劣的双重游戏,为了事态的需要,把我们最起码的认识官能遮起来。诚然,一直以来的空气在怂恿我们:我们始终看到我们的周围人赶着以他们的共识,批准那些为有他们附和的虚构效力的遁词——无论多么粗劣。而我自己的遁词,在我谈过的那些个案里,确确实实有着周围人的附和或默默的怂恿;不然,我无从把它维持下去、继续我的游戏。

而承担一份情势,不多不少,正是以\textbf{真诚}直面它——取这个词的完整意义;也就是说:不动用供我们使用的便利,凭粗劣的遁词,对它显而易见的来龙去脉闭眼。这也就是说,干脆地:动用我们健全的知觉与判断官能,而不为这样那样事态的需要留神把它们遮起来。事情或许显得奇怪,却同样简单而显然:当我们以这样的姿态——“无邪”的姿态——直面一份情势,它旋即深刻地转变,无论它曾显得多么混乱、多么打了死结。或者更好地说:倘若它果真“打了结”、许久纹丝不动,那是因为我们自己拦着它演化、拦着它按其本性“流淌”;是我们挡着它的自发运动,循着……一如我们自襁褓以来所有围绕我们的人的一致先例。只消\textbf{停止}绷紧,\textbf{停止}阻拦,那些仿佛凝定了的事物便重新动起来;卡住的东西松开了;积聚起来的生硬的紧张,终于得以释放,消解于一份新的、宽广的运动——它终于重现。

这份我们人人都有、又得众人怂恿的“便利”:把鱼目当珍珠,并由此堵住那本为流淌而造的东西——它其实毫不“舒服”!它为我们保有的那份怡然自得的内在不作为,我们付的是一笔高得吓人的价钱——一份内在的僵紧之价,加上为维持这份僵紧与“鱼目=珍珠”的虚构所需要的惊人能量投入。话虽如此,每一时刻,人人按自己的心意行事——这是我们的特权。而在每一\dmark{}时刻,凭我们所\origpage{669}做的,我们播种——为着我们自己,也为着他人。而我们之所\textbf{播}的收成,当下便开始。

\subsubsection*{(c) 贪婪——或蚀本买卖}
\tocsubsub{(c) 贪婪——或蚀本买卖}
\label{ch18:local:h56}

\notehead{143}{}
\phantomsection\label{ch18:local:n51}

或许该是时候回到葬礼的那个“第一层面”了——也就是回到我的大司仪、我的朋友皮埃尔在葬礼中所扮角色的来龙去脉。一周之前,我在“天鹅绒的脚掌——或微笑”那条注里(n\textsuperscript{o}~137,12月7日)已经回到过它,随后又经由那场(连着五条注的)关于“利爪”与“天鹅绒”的枝蔓离开了它。我感到,这场“枝蔓”——一如先前的许多场——不曾白费。

我之所以被引上这场枝蔓,恰恰是因为:我的朋友在扮演那份角色时,也许最惹眼的明显特征,便是那份最纯的“天鹅绒的脚掌”风格的执拗——没有一丝破裂的意向,为一份毫无破绽、从不直言其名的对抗效力\fnote{205}{(*)}{一如我别处已有机会强调的:对抗、或拒斥与嘲弄的故意“从不直言其名”,绝非朋友皮埃尔所专有;而(就我所知)适用于葬礼的全体与会者,无一例外。于是,在这场以嘲弄行事的“阴的葬礼”里,每一位与会者身上的底色(一如这样一场丧仪所应有)本身——是阴的!
关于葬礼这份“隐蔽”的品格,亦见“掘墓人——或全体教众”一注(n\textsuperscript{o}~97)。}。另一个惹眼的事实:在那副听得懂的微笑、和婉的神色之愉快而克制的表象背后,我的朋友身上,对着我本人、或对着他归入“我的”名下的人(在数学工作的层面上),不止一次表达出一份毫无歧义的、表面上无端的\textbf{害人}或\textbf{伤人}的意图。关于这一方向的诸多具体事实,我在葬礼第一部分里已说得够多,此处无须赘言。这确确实实是一份恶意的姿态(似乎严格圈定在科学活动的领域之内),是取一词之强义的“\textbf{暴力}”——即便它严格保持隐蔽:爪子永远泡在精美的、起绒的丝绸里。而这份暴力、这份恶意,带着那份最令人错愕的\textbf{无端性}的全部表象——仿佛它们施展出来,只为害人伤人之乐。

\dmark{}一如每次直面这样一份情势,它显得如此不可思议,以致人们常常\origpage{670}不敢相信自己健全官能的见证\fnote{206}{(*)}{关于这一点,见“中国皇帝的新衣”一注(n\textsuperscript{o}~77$'$)。}。否定这份见证——时下的惯常做法——是不承担一份情势的无数方式之一,并由此把它延续下去;更可取的,肯定是围着事情转一转,寻找或许逃过我们知觉的侧面,它们能提供一条接近的路径,使这份情势得以纳入我们的亲身经历。依我看,一生中从不曾经过这份无缘的恶意姿态的人,必定寥寥——而肯记得自己也曾如此,已是可能的第一步,好\textbf{靠近}那份情势的实况:时下的反射,更巴不得我们把它火速清场。再往深处探测一番也无妨:看看是否有什么暗藏的怨怼,是那份看似无端的暴力之缘由与发条——同样,视情形,也不妨认出某些“怨怼”的赝品,照它们的本相:比如我自己实行过的那一套,说某人是个可憎的家伙、不配任何客气等等。

可在眼下这个个案里,任凭我怎么探测,都看不出任何东西——近看远看——像一份\textbf{怨怼}:我的朋友可能(有理或无理地)怀有的、对着我的怨怼,或对着他择为恶意靶子的任何人之一。他自己,也从不曾有过片刻,暗示过哪怕一星半点朝这个方向的东西;更不必说,对我屡次就他那几桩叫我目瞪口呆的行为所作的探测,他从不曾有一刻承认:他身上可能怀着对任何人的、一丝敌意的姿态。我最终在他身上感到一份暗藏的得意——在我们偶尔见面时:当他端出他那套十足“客观”的好理由,带着那份他特有的、有点好玩的无邪惊讶的神色……总之,我进了一场他随心所欲、凭兴之所至引领的游戏,带着一份我久久未曾察觉的内心满足。(然而,把人折腾得团团转的,他可远不是第一个!)但我终究——迟走出来总比一直陷在里面好——走出了那个木马\fnote{207}{(**)}{那是1981年——即“两个转折”那条注(n\textsuperscript{o}~66)里所谈的“第二个转折”。}!

另一方面,倘若我探测我自己,检阅我与我这位\dmark{}朋友自将近二十年前(1965年)相遇以来的关系\origpage{671},我同样找不到任何痕迹:任何曾在某一时刻成为对我的怨怼之缘由的东西。按事物的常规的、浮面的意义,我可以说:在整个这段时间里,尤其在最初五年贴身接触里,我“只对他好”。可这份确认旋即使我记起另一份,不那么浮面的——我对他的那份\textbf{纵容}:它显形于“置身事外者”与“含混”两条注(n\textsuperscript{o}~67$'$、63$''$)的反思进程里。分明地,这份纵容对他绝不是“一件好事”;同样分明的是:我那年轻而辉煌的学生兼朋友对我的种种姿态,与我自己的种种姿态——尤其与这份纵容——曾在一个紧密的共生中发育。甚至可能:这份纵容,在无意识的某个层面上,曾被我的朋友(不仅被知觉——这反正显然——而且)感受为一份“怨怼”、一份也许太熟稔、被嚼烂了的老套脚本——在他那略带神童色彩的童年时代;而如今,它(哪怕不显眼地)又被端到他面前。他也许曾天真地以为:一脚踏进数学的“上流世界”,一切都将不同于他所认识的——结果不然,还是那套老把戏!(而凭他自己的刻意抉择,今天仍是那套老把戏,而且规模更大,更不必说了……)个中究竟,我恐怕永远无从知道。况且,把它澄清,也不是我的本分——假定我的天线细到足以单凭自己的手段做到。倘若有过“怨怼”,那至多是一份“贴补”的怨怼,它弹了一记指头,把“某样东西”推上路:某一场\textbf{游戏},由一份完全另一个量级的力所驱动;那份力的在场,我久已感到,其本性对我却仍是个谜。在离开葬礼图景的这个“第一层面”之前,我至少想试着掂量那份力的本性。

分明地,有一份\textbf{贪心}:取而代之、排挤、抹除的贪心;还有那份把他人与数学夫人辛劳与爱恋之果\textbf{据为己有}的贪心。可是,我清清楚楚:\textbf{不是}一份单纯的声名、钦慕、荣誉的“贪食症”,甚至不是权力的“贪食症”,才是他在葬礼中所扮角色的深层发条。在我对这份\dmark{}角色的反思进程中,我曾多少次骤然惊觉\origpage{672}:他那份下葬的\textbf{执念},使他把自己也埋了进去!凭他那份例外的天资、凭一份同样例外的时运,他分得了超越师父极远、在整个时代的数学上留下深深印记所需要的一切。只消他放任身上的孩子随意去玩,不用种种指令去烦它、不用这里一道栅栏那里一处禁行去烦它——只消简简单单照看好必要的事,严格地管好后勤——那么,他身上的“东家”(其贪心无疑与任何人半斤八两),既不必推,也不必拉,也不必动用手肘,本来就断然少不了各色可想而知的声名、钦慕、荣誉的标记,外加权力的标记,多到不知道拿它们怎么办——彼时,那位“郁郁寡欢者”却正玩得尽兴,也绝不给“东家”留下多少摆东家架子的闲暇……

分明地,按单纯的“功利”账算,陷进一场葬礼,是一桩很不划算的买卖:它粘在他腿上已十五年以上,若非那碍事的亡人忽然想起掀开棺盖、搅了仪式(恰在其时——正如常理所言——在人们最不提防的时刻!),它本来打算粘他一辈子!(关于那桩不合时宜的事件,对皮埃尔东家未来下注的影响,赌局已开……)或者换一种说法:我的朋友有那份材料(至少凭他的智性手段)、有那份显赫的资格,在数学里做一位\textbf{彼得大帝};可他偏偏选择了做一个小皮埃尔。这看上去千真万确是一桩蚀本买卖——至少,倘若所下的注,确确实实首先是那些虚荣的满足。

\subsubsection*{(d) 两种认识——或对认识的恐惧}
\tocsubsub{(d) 两种认识——或对认识的恐惧}
\label{ch18:local:h57}

\notehead{144}{}
\phantomsection\label{ch18:local:n52}

(12月15日)昨晚反思的收尾处,我身上有过一份轻微的不适,像那种人气十足地端出一份逻辑无懈可击的推理、却撇开一份弥散的感觉的人:总觉得有什么地方不对劲。这份“什么”在我一停笔时便显形了。表述它的一种含混方式是:无意识的“逻辑”——它主宰着我们最关键的抉择——绝不是那份寻常\dmark{}的有意识推理的逻辑,更不是那份“正统”推理的逻辑\origpage{673}。在此个案中,我对青年Deligne在(姑且说)六十年代后半期所握“王牌”的知觉、以及我赋予它们的分量(至少与任何一位消息灵通的理性数学家将赋予的分量同向)——这份知觉与这份分量(我乐意称之为“客观的”)与当事人自己的姿态与感情毫无关系;尤其,与那些关乎他自己能力的姿态与感情无关——而那些,分明是他手中所有王牌里的王牌。

可是我有这样一种印象:至少在意识的层面上、带着谦虚所要求的种种套语,我的朋友早已把那些久已传回他耳中的动听回声——关于他那不寻常的天资——纳为己有。但对我而言毫无怀疑:在更深的层面上——那个一生的主导性抉择被默默做出、无话可说的层面——那份“客观”的事物版本,当时便已(而今天依然)\textbf{失效}。占了它的位置的,是一份阴险的\textbf{怀疑}:任何一份价值的“证明”(或对他人之优越的证明……)都永远拔不了根——一份愈发顽韧的怀疑,因为它永远未被表述。我在朋友身上知觉到它,一如我在其他天资逊色者身上知觉到它:它是同一份。这份怀疑,是一份\textbf{贴身信念}的顽固信使:那份信念同样未被表达,埋得比这怀疑更深:一份对无能的贴身信念——根本的、无可救药的。它也就是那份“对己之轻蔑”——我在《收获与播种》的开头处、在一场仍取“一般”品格的反思的语境里谈到过它\fnote{208}{(*)}{见“他人之无谬与对己之轻蔑”一节(第~n\textsuperscript{o}~4节)。}。一两个月前,它又在一个仍然非个人的语境里、以一副不同的面孔重现:作为一份“裂缝感”\fnote{209}{(**)}{见“一半与全——或裂缝”那条注(n\textsuperscript{o}~112),10月17日。}——那份弥散的感觉:我头一次在我自己身上确认它,是在发现禅观后的第二天。而在葬礼反思的进程中,也有几次,对我朋友身上那份“无能的贴身信念”的骤然而敏锐的知觉,为某一仿佛公然违逆常识的情势投下新的光照……\fnote{210}{(***)}{关于这一点,见“颠倒(3)——或阴葬阳”那条注,其中(及其他)唤起了反思的某些这类“敏感时刻”。}。

\dmark{}我知道,这份贴身信念——在我朋友身上或任何人身上——本身宛如一份\textbf{认识的}\textbf{影子}:对那份确确实实存在的“裂缝”的认识、对一份被承受的“戕害”的认识——而这份戕害,直至今日,仍由当事人自己的默许所批准、所维持。然而,影子并不归还它所由出的那份认识——那份认识,一如一切认识,本身是有益的——它毋宁是它的一个巨大而畸形的漫画、一个稻草人版本。如此歪曲一份认识、使它面目全非的,是一份\textbf{恐惧}:恰恰是怕与这份认识本身接触的恐惧,怕让它从它一直被压抑的深处上浮的恐惧,怕承担它所忠实映照的那份谦卑实在的恐惧。

与这份可畏的认识接触、以一份完全有意识的目光认识那份在深层已知、却被逃避的实在——这才真正意味着:与我们身上那份(人们叫它“力量”也好、“孩子”也好)“自信已亡故一生之久”的东西,全面重新接上联系。因为确确实实是那份力量、唯有那份力量——童年的力量——使我们能够承担对自身那份打了裂缝、被残害、被瘫痪之物的认识。而承担它,也意味着:与那\textbf{另一份认识}重新接上——它先于对我们戕害的认识、且比它更要害:对我们身上那份力量之在场的原初认识——一份既非肌肉之力、亦非头脑之力,却把两者都包含在内的力量。

事情或许显得奇怪:这份失落的认识——关于那份“力量”、那份\textbf{创造权能}在我们身上的在场,作为我们真正本性的显而易见、不可摧毁的一部分——恰恰经由一份\textbf{无能状态}的发现与谦卑接纳而寻回;而那份无能状态,正是凭这份接纳本身得到解决。对一份无能状态的认识,覆盖并遮蔽着那份埋得更深的认识:对我们创造力量的认识。前者宛如一把钥匙,为我们打开后者;两者实际上不可分,如同一份\textbf{同一}认识的正面与反面\fnote{211}{(*)}{在这幅意象里,当然,“正面”是对无能状态的认识——即不真实的状态、“裂缝”的状态;而“反面”,藏得更深,是对我们未分裂的本性、我们的创造权能的认识。这些年里,我屡次得以确认:恰恰是“反面”——两者中埋得最深的那份认识——才是最强恐惧的对象、才招致最激烈的否认。让任何人不安的,并不是那份驯熟的、(或多或少)“有学问”的猴子的熟稔而平淡的状态,而是孩子的那份无邪:他按事物的本相感知它们、直呼其名,他照他所感去做去说,不因与“人家”期待他的样子不同而羞愧。},是同一份\textbf{恐惧}的对象。

\dmark{}当我谈到埋在我们每个人身上的那份“力量”时,这绝不是什么抽象而含混的东西\origpage{675}、什么“哲学家”的口头精致、或半吊子哲学家的心理学的口头精致。正是这份力量,使你能“做数学”(或“做爱”……),像一个孩子呼吸——也就是说,不必小心翼翼地约束自己不偏离前人留下的尾流、不必勤勤恳恳重复他们的手法与配方(或他们的陈词滥调……);也正是这份力量,给你勇气与谦逊——在你自己家里,也在别人家里——把猫叫作猫,不把鱼目当珍珠,哪怕这样做你会顶着最牢固的共识、或顶着你自己身上最积习已深、最运转纯熟的那些机制。\fnote{212}{(*)}{(12月16日)创造力量在每个人身上的作用——更新之力(或“孩子的力量”)——凭它的果实而辨认:无论凭手或心的作品,还是凭日常生活的种种事实,在与他人、与周围生灵与事物的关系中。我得以屡次确认:日常之中的创造性,远比经由“作品”(取常规之义——即凭手或凭心塑成的可触“产品”的创造性)的创造性罕见。
一份持续的创造性在某个人生活中的在场,是与那份在他身上的创造力量的一份持续“接触”的标志——哪怕何等零碎、何等不完美。这属于另一种本性,不同于“天资”的单纯在场、不同于为开发利用它们而作的持续能量投入——后者表达为一份或多或少的可观、或多或小时行的产品,却本身没有创造的德能、更新的德能。
在我的智性求索之中、尤其在我的数学工作里——天资平平(而投入可观)——我觉得,我与那份力量之间的这份“接触”、亦即我对它那份默默而深刻的认识,几乎原封未动。也就是说:出入无几,我在我生活的那个(诚然极零碎的)领域里,凭我(创造)手段的全部“运转”,几乎没有任何损耗、挪用或阻塞——没有惯常的那些“摩擦效应”。其中最常见的一种,是某种怯懦:它如此经常地使我们听不见内在的声音——那声音正叮嘱我们该做什么,当它所教的恰恰是“新”的东西,即把我们领上唯有我们自己踏足的小径。这类抑制,在我的数学关系里几乎缺席(而且依我看,随岁月推移愈发如此);可在我生活的其他领域,它一如在任何人身上,存在着——尤其恰恰在“日常生活”的领域。
回到数学活动:我在我那辉煌的昔日学生身上看到一份某种意义上颠倒的关系。他握有那份从来令我惊叹倾倒的天资——与我的不可同日而语。(诚然,我活得越久,越看清:要在科学或其他领域做出革新之作,这根本不是要害;关于这一点,见“阴,作为仆人(2)——与慷慨”那条注(n\textsuperscript{o}~136)里的反思。)他对数学的投入是可观的,一如我当年的投入;而自他最年幼起,他便享有格外有利的条件:为他的天资之绽放、为构想并打造一部与那些天资相称的作品。二十年之后,我还在等这部作品,而且意犹未尽!在他身上,肯定有某份与创造力量的“接触”——他所做的某些东西之美可为佐证;可这份接触是受扰的、受磨的。我的朋友对他的工作的关系、乃至在他的工作本身之中,是一份冲突的关系:工作随着岁月推移,愈发变成“东家”手里的一个\textbf{工具},用来满足\textbf{他的}饥渴——那些与孩子的认识渴、发现渴格格不入的饥渴。
我怀疑,这样一份冲突的关系,能不经先行承担而得到解决——也就是说,在一切之先:被认出。至少,我一生不曾一次见过这样的事不经此而完成。这正是使我写出下面这句话的缘由:对我们无能的认识,是寻回我们创造权能的全面认识的“钥匙”,并由此,也全面地寻回创造权能本身。在我的数学工作里,问题不曾提出,因为那份工作里不曾有深层的阻塞——一种相当于局部无能、会使我只凭我可能性的小部分“运转”的阻塞。可这个问题,一如对任何人,在日常生活的层面上、在我与他人的关系、与我自己的关系、与我身体及其冲动的层面上,提出过。正是在这个层面上,我一再经验到:对一份阻塞、一份“无能”的认识,确实是那把解放被囚的创造性的\textbf{钥匙}。}

顺带,头一个来到我笔下的例子,颇能说明问题——它肯定足以叫每一个(乃至不那么年轻的)热爱荣耀的年轻研究者心跳。谁不愿意做一门尚在孕育的科学的无畏先驱,并以此在所有教科书里占据显要地位,像开普勒——现代天文学之父!可当要做的(一如开普勒与其他人所做的),是在孤独中、在所有人的冷漠中(更不必说轻蔑或敌意)顽韧地纺自己的线,三十年、哪怕仅仅一年——那时,忽然就没了人影!人们愿意进教科书,总而言之愿意与名人并列;可人们也\textbf{怕}孤单——哪怕一年,甚至哪怕\dmark{}仅仅一天。而那个“认识”到他身上力量之在场的人(为认识它,他从不曾\origpage{677}必须谈起它——对别人、甚至对自己……)——那个人也分明知道:他是\textbf{孤单}的;而孤单,并不叫他不安。至于他会不会进教科书,那是他操心清单上的最后一项——尤其在他工作的时刻。

况且,恰好是这位开普勒,在他的工作本身之中,“顶着科学里最牢固的共识”——而且是竖立了数千年的共识,更不必说了。在他的时代(宗教裁判所还在),这比今天更不便:今天你很可能丢掉饭碗、或找不到饭碗,却不至于冒险上火刑柱。回到开普勒:我不知道在他的日常生活中,对着那些“最牢固的共识”,他如何自处;或许在那些事上,他跟所有人一样规规矩矩。确定的是:今天一如当年、一如从来,肯偏离这些共识一毫的人,也不多。想必还是那套老把戏:\textbf{怕孤单}——一份深层而几乎普世的需要之反面:对赞许的需要、对他人之确认的需要(哪怕世上只有一个肯赞许、肯确认的人)……\fnote{213}{(*)}{我由此经由另一条路径,接上一些已在“禁果”与“孤独的历险”两节(n\textsuperscript{o}~46、47)里、以及顺带在“接纳”那条注(n\textsuperscript{o}~110)里显现过的确认。}

\subsubsection*{(e) 秘密的神经}
\tocsubsub{(e) 秘密的神经}
\label{ch18:local:h58}

\notehead{145}{}
\phantomsection\label{ch18:local:n53}

可我又一次离开了我的话题!我起笔时,是针对这份确认:我昨晚的“推理”驴唇不对马嘴——当时我试图“传递”我身上的这份确信:我的朋友在葬礼中扮演我所知道的那份角色、以我所知的那种方式,其动机不是\textbf{贪婪}(对声名、钦慕、荣誉、权力的贪婪)。诚然,千真万确:拿一份孩子的冲劲换一份\textbf{角色},他做了一桩“蚀本买卖”——甚至按“回报”账、按声名那一侧的账,也是如此。可这什么也证明不了。依我看,这类“亏本账”,在我们主要\dmark{}投入与抉择的(无意识层面的)拣选里,几乎是绝对的通则,而绝不是例外。可即便那推理一文不值,我却毫不怀疑\origpage{678}:我想传递的,恰恰是一份对实在的知觉:\textbf{不是}这份贪心——它千真万确,而且在我朋友的生活里占据了越来越大的、真正吞噬性的份额——\textbf{不是它}构成我的朋友所扮角色的\textbf{神经}:作为我葬礼之操办中那个\textbf{关键}人物。

倘若我试着更贴切地圈定这份极清晰的感觉(至于“确立”它的正当性,则再无所提!),得到的是这个:是那份对抗或恶毒行为中的\textbf{无端性}——那份屡屡叫我目瞪口呆的无端性——与那份万金油的“解释”绝对不“对榫”:贪婪。至少,就声名、钦慕、荣誉而言,甚至就取常义的“权力”而言,我那辉煌的昔日学生兼朋友,在对付他昔日师父时使出那份他深得其妙的“含蓄而配量精巧的轻蔑”——无论在当下还是在更长的期限里,什么也没赢得;或在对付某位地位低于他的研究者、或其现在或过去的工作时,使出同一份轻蔑(或许配量不那么精巧),把那人本来便不如我牢固的自信吓散;或者对付另一个人——那人顶着朋友定调的普遍轻蔑英勇坚持,他便剥夺其坚持之果。倘若说在这最后一个个案里、一如在其他个案里,我的朋友曾作态把别人在孤独中(有时顶着长辈的轻蔑)成熟的果实据为己有,那么这份“收益”(取“\textbf{暂停}”的风格\fnote{214}{(*)}{关于我这位辉煌的朋友兼昔日学生这份据有风格,见“暂停!”(n\textsuperscript{o}~77)与“占有与轻蔑”(n\textsuperscript{o}~59$'$)两条注。}),一想到\textbf{据有者是谁},便寒酸到这种地步,以致那份端出的“解释”自己便化成青烟!

至于我,我分明知道——凭一目了然的认识:这类据有的“神经”,不是\textbf{那份}收益。我从中反而感到\textbf{某份权力的陶醉}——一份比常规意义的权力更精微、或许更醉人的权力:如同某位有分量的科学家惯常行使的那种——坐在各种委员会、理事会、评审团及同类机构里\dmark{},主管一所研究所\origpage{679},或指导一批年轻而辉煌的研究者,或对部长耳语。我所谈的那份“陶醉”,在“乖僻”那条注(n\textsuperscript{o}~76)里显现过(反思中头一次):我在那里骤然面对“一份\textbf{炫弄}之举,一份在如此全面的权力中的陶醉——它甚至容许自己(象征性地……)炫示它那‘乖僻地’剥夺他人的真正本性”。

那是一份炫弄之举:辉煌而张扬,却同时\textbf{隐蔽}、未曾言明,不动声色地滑进那里,甚至还为“反常层”这个奇怪的名称附上一段临时凑出的解释——多么自然啊,三言两语就能向你说明白——此外再列一张小清单,说说“哪些东西本该写进我们那篇谦逊而辉煌的文章”……\fnote{215}{(*)}{见“变戏法的人”那条注(n\textsuperscript{o}~75$'$)。}。

我在这里又认出了那份最纯的“天鹅绒的脚掌”风格——别名“暂停!”的风格;而且,在那份我已在不止一人身上变得熟稔的风格之齐一背后,我也感到那份\textbf{共同的神经}:那份强制性的、吞噬一切的\textbf{渴}:行使一份权力;一份\textbf{特定的权力},取一种特定的方式——猫玩老鼠的权力:当猫施展它的\textbf{大游戏},带着那份完美的优雅(唯独老鼠无从领教其妙处),并且肯定“以最精微的温柔”——或者,也是一位精明的妻子对她的傻大个丈夫的权力……

从我的朋友提出的个案出发,我已被引着谈过那份“风格”及其意义——在形形色色伴侣关系的一般语境里。那是在一周前的反思里,“颠倒(4)——或夫妻的马戏”那条注(n\textsuperscript{o}~138,12月8日)里。正是在那里,“天鹅绒的脚掌”(别名“暂停!”)游戏的“神经”头一次以它应有的清晰显形:作为一份\textbf{权力的游戏}。然而,作为一份性质十分特殊的权力游戏:游戏对行它之人的那份魔力、那份往往吞噬一切的魅力,恰恰在于经由它行使的\textbf{权力的隐蔽品格}——那份“看不见、摸不着”的品格:它使人得以玩弄对方(\textbf{玩弄}他,却从不\textbf{与他一起玩}……),随心所欲地叫他原地打转,始终领舞,而对方笨拙地一次接一次跟着转,对那些凭无形的线使出的、随兴所致的小小招式,报以笨拙的应答\dmark{}\origpage{680}……

终于白纸黑字写下那份多年来晦暗地感到的东西——我从不曾费心对自己明白表述——短短一份努力,把长久弥散的东西凝成言语,便已足够:昨天还向我显为“谜”的东西(即某位朋友身上一份“某种力量”的本性),骤然向我敞开其显然的意义!他身上这份“力量”、或(如我方才所写)这类看似“无法解释”(乃至“超出领会”)的行为的“神经”,我在12月8日那场反思里早已圈定。可尽管那场关键反思的出发点,确实是我那位辉煌朋友的某场“谜”样的游戏,滋养那场反思的,却是\textbf{另一份}亲身经历——比我与他这个人相联系的亲身经历更丰富、更强烈;一份已整个消化(或相差无几)的经历,它向我口授一份已然成形的认识——那份我与朋友皮埃尔零星关系的皮肉经历,当时本无力传给我。

当然,终究要理解的,正是那份经历,并由此把它全面承担;倘若我当时毫无内心保留,便投入那场关于“夫妻旋转木马”的枝蔓,那是因为我分明感到:那个木马,对我的朋友关系,有话要说。对他的念头,继续在背景里在场,像一份低调的底色。

然而,两者的全面“会合”,那天没有完成,随后几天也没有。无疑,时机还没有全然成熟。要使会合毫无保留、毫不费力、带着显然的从容完成,我须先“清理场地”:执拗而不慌不忙地,把最强制性地要求我注意的联想,一一追随过去。我不曾强迫事情;而且我知道,这恰恰是我该做的:照管那份执意呼唤我的东西,而不被一个“话题”、一根(反思的)“线”、乃至一份要收尾的纲领引开。

我这边锄草松土,天地之力在做它们的工。到了晚上,只须来收取熟透的果实:它落进那只张开迎候的手……

\subsubsection*{(f) 激情与饥渴——或逐步加码}
\tocsubsub{(f) 激情与饥渴——或逐步加码}
\label{ch18:local:h59}

\notehead{146}{}
\phantomsection\label{ch18:local:n54}

\dmark{}\origpage{681}(12月17日)依我看,随着前天的反思,发生了一份仿佛解冻的事情:一份一直犹豫不决、有点发懵的理解——面对堆积在我面前的一大堆事实与直觉,乱得像一座不成形的山——像一幅拼图,我不过勉强此地彼地拼上了几块。此刻我有这样的印象:我撞上了那幅有待复原的未知图景的\textbf{那}块“要害”拼块——围绕它,其余的拼块终于将毫不费力地各就各位。无论如何,我毫无怀疑:我确确实实触到了朋友皮埃尔在师父及其(或多或少)信徒们之下葬中所扮角色背后的“神经”,也同时触到了他与我这亡故师父之关系的“神经”。

这份玩弄某份权力的饥渴——不显眼地、带着一副天真无邪的神气,拉动那些无形的线——这份饥渴,必定早在我认识他之前很久便已存在,他本人与所有人都不知情。在我与他相识的最初几年——我出发的情节(1970年)之前——我不曾见它显形;这无疑是因为:在那些强烈学习、精微而有力的思想起飞的年月里,我朋友的能量被整个吸在别处。条件确实是理想的:足以作他那些例外手段的跳板。我出发的情节——先脱离我们共同隶属的那家机构,继而(次年内)脱离数学舞台——不仅是我自己灵性历险的关键转折,肯定也是他的转折。正是这份情节,骤然为他敞开一些权力的手段:前一天他还不敢梦想。首先是把一位前师父从那地方“排挤”出去的权力——那师父在那里占的位置太大,而他此前只求不显眼地与之保持距离\fnote{216}{(*)}{关于这份先求保持距离、继而排挤的心思,见“驱逐”(n\textsuperscript{o}~63)与“兄弟与配偶——或双重签名”(n\textsuperscript{o}~134)两条注,及后者的子注(n\textsuperscript{o}~134$_{1}$),以及“未竟的收成”一节(第~n\textsuperscript{o}~28节)。};继而,当那位师父显然要退出舞台时,是更醉人的权力:叫那所以亡故师父命名的“学派”无影无踪地消失;而这样做,最终是干脆利落地——在每一根主枝上(他自己栖身的那根除外)——砍断一份宏大纲领的绽放:它服务于一份宏大的愿景,而他本人曾长期以之为滋养\fnote{217}{(**)}{甚而,关于一个“学派”的清算与“电锯效应”,见“继承人”“共同继承人……”“……与电锯”(n\textsuperscript{o}~90、91、92)三条注,以及“灵车”行列的头四条棺材注(棺1至4,n\textsuperscript{o}~93--96)。关于那份被下葬的愿景,见“我的孤儿”(n\textsuperscript{o}~46)两条注里(以两道不同光照)给出的两份概观,以及“阴,作为仆人(2)——与慷慨”那条注的子注~n\textsuperscript{o}~136$_{1}$。
人们会注意到:正文里“而这样做……”一语(“……干脆利落地砍断一份宏大纲领的绽放……”)并不贴切。清算一个学派,是“干脆利落”地“锯断”一整批主枝的\textbf{第一}记“电锯”,但不是最后一记(所引的那些棺材注,n\textsuperscript{o}~93--96,即为明证)。}。

我朋友生活中这场大转折的意义,向我显为他这个人身上两股支配性力量的相互霸权之中的一场颠倒——依我看,它们凌驾于其余一切之上:数学的激情,与权力游戏的“饥渴”(“以天鹅绒的脚掌”)。前一份力量,在本质上取“冲动”的品格\fnote{218}{(*)}{数学的激情“取冲动的品格”、是“孩子”(别名“匠人”)的表达,这并不妨碍(同一段里有力地提醒)它也或多或少强有力地被“东家”的“饥渴”所占据——这属于“匠人”与“东家”关系的共同命运(我本人也毫不例外,一如任何人)。},后一份取利己的、习得的品格。转折之前,是认识的冲动支配着我朋友的生活(就我所知),而权力的饥渴或多或少沉睡,处于空置状态。在短短几年间一段令人眩晕的社会攀升之后\fnote{219}{(**)}{关于这一点,见“高升”那条注(n\textsuperscript{o}~63$'$)。},在一个骤然出现、逼人作严苛抉择的情势里,取胜的(我相信是断然取胜,且毫无交战之意)是权力及其秘密陶醉的诱惑,压过了认识的激情。后者并未退出舞台,却从此成了饥渴的封臣与谦卑的仆婢,成了后者手中的一个\textbf{工具}。激情(别名“匠人”)在饥渴——别名“东家”——的醋眼里操持它的活计;后者寸步不离。既然匠人有好工具(并非全被没收)、有一双好手,即便这样被辖制,他仍旧磕磕绊绊地维持着一份产量、维持着这家的名声。但已不复是当年:当年匠人(颇带几分顽童气)整日整日自得其乐,而东家远在天边,一季才来转一遭!

\dmark{}其后的演进,依我看,更多是数量性的,而非质量性的。那是东家某一份\textbf{战术}\origpage{683}的渐进演进——风格始终如一,而东家-匠人的关系纹丝不动。这位东家禀性谨慎,只在有赢的把握处才肯冒险。为此,须对地形有把握——或者,对“(几乎)全体大会”的默认附和有把握,从亡故师父那些昔日学生的更小圈子算起。顶风冒雨维持的那份个人关系之演进,是“地形知识”之演进的忠实映像。权力与轻蔑之游戏的胆量,有一份渐进的\textbf{加码},历时十二年,在1981年以反常层会议的壮举达于顶点:在普遍的亢奋中,一切克制(甚至一切谨慎)都欢天喜地地被抛进水里\fnote{220}{(*)}{关于“反常层会议”,见第七行列“会议——或Mebkhout的层与乖僻”,注~n\textsuperscript{o}~75--80。}。这样,竟须十二年,我的朋友才使自己确信:地形已如此相宜,任何谨慎都属多余:招招皆赢!时机分明成熟了:终于可以把那件秘密武器公之于众——\textbf{Motif}:次年便换了个爹被开棺出土\fnote{221}{(**)}{关于Motif的开棺,见“一个梦的回忆——或Motif的诞生”与“下葬——或新父亲”两条注(n\textsuperscript{o}~51、52)。}。

我没有兴致在此追溯这份十二年加码的逐步脚印——尽管我手段样样在手,完全可以做。那会是一份编年史家的工作——那类工作,我在葬礼第一部分那场未曾预料的“侦查”(或“中国皇帝的新衣”)里已经做得够多。这场加码的诸“步”,依我看,是“\textbf{试探}”——我的朋友朝一位沉默的大会掷出的一记记探针——每一次都得到同一回答:尽管放手干!将近十五年,大会是他的无声盟友与他的保人;而他,在不知不觉中、大概也不曾操心,是大会一件驯服的工具\fnote{222}{(***)}{见“掘墓人——或全体教众”一注(n\textsuperscript{o}~97)。}。

\subsubsection*{(g) 心软的爸爸}
\tocsubsub{(g) 心软的爸爸}
\label{ch18:local:h60}

\notehead{147}{}
\phantomsection\label{ch18:local:n55}

\dmark{}\origpage{684}我不知道我朋友这份饥渴,是否也施于我之外的其他人——施于那些嗅得出我的“气味”的更年轻数学家。这方面没有回声传到我这儿。可对我分明的是:正是经由他与我的关系、借着一份在科学世界里诚然罕见的情势,他身上那份在阴影里苟活的癖向,一夜之间变成一份吞噬性的饥渴。在我出发的情节里,当他一脸严肃地向我解释:他已把一生毫无保留地献给了数学\fnote{223}{(*)}{关于这段情节,见“兄弟与配偶——或双重签名”那条注(n\textsuperscript{o}~134)。}——他无疑“相信”他所说的话;而我本人,尽管有点发懵,也不曾想到怀疑他的话。可倘若我的耳朵更灵,或者更好地说,倘若我当时已有那份成熟,去倾听并信赖一份“更灵的耳朵”——它确确实实存在于我身上、一如存在于每个人——我本会知道:他关于他自己所说的话,前一天或许还是真的,可那一天已不是真的。那是一份高尚的缘由,为一桩可疑的行为而设:那桩行为的意义尽管昭然,他与我当时都没有那份质朴去直视它。那一天攫住他生活缰绳的、至今不曾松手的,是\textbf{别的东西},而不是那样的激情。

因此,是我这个人、或毋宁说他与我之关系中的某种东西,(借着合适的机会)成了触发器:触发那份支配他生活的力量之严苛的变质、触发他投入数学的意义与方向的严苛的变质。此刻该记起葬礼那些著名的“联”或“侧面”了——它们在11月13日的反思(“回顾(1)——或一幅画的三联”那条注,n\textsuperscript{o}~127)及紧随其后的注(“回顾(2)——或死结”,n\textsuperscript{o}~127$'$)里被摆到聚光灯下;此后,它们在路上多少有些被淡忘了。十天前,在“天鹅绒的脚掌,或微笑”那条注(n\textsuperscript{o}~137,12月7日)里,我曾多少作态记起它们。在那里,我尤其重新接上了一份直觉:那份“过继的父亲”的永世角色——我注定要在我年轻朋友身边扮演的;而依我看,它在他身上保存至今、至今仍在起作用。借着这场反思,我再次毫无保留地表达一份确信——它想必是在逝去的六七年里逐步\dmark{}成形、成体的(甚至可能更早\origpage{685}):“正是绕着这一侧面(他对我这个人的把握中的父性侧面),冲突打了结——一份早已在他身上的冲突,远在他听说有人念出我的名字之前……”(这便是那个著名的“超级父亲”之联;而“超级母亲”之联,至少眼下,仍在冥府。)

况且,只隔一页,那份著名的“微笑与天鹅绒的脚掌”风格便头一次匆匆亮相,成了注意的对象。系于它的联想,在随后几天里,似乎先把我引离我朋友这个人,也引离那份“父性”的隐蔽侧面——他在自己生活中派给我的角色。直到今天之前,这个侧面再没被谈过——人不可能同时想一切事,更不可能同时谈一切事!可是,就“想”而言,依我看,在某个地方、在那份模糊却依然在场、依然起作用的背景里,那个父性侧面的念头必定在场:它必定像一份有效而低调的兴奋剂,推动着那场关于“天鹅绒中的利爪”风格的长长枝蔓。毕竟(事后我才对自己说清楚;可它当时必已作为一份弥散却强有力的动机在场……),“父亲”的形象与那份著名的风格,绝非无缘——恰恰相反。甚至可以说:小女孩(或小男孩,也无妨)一生中见到的第一个被这份风格精微而利落地(虽不总是温柔地)引领的人,不是别人,正是爸爸!

而只要那无邪的小丫头(或小捣蛋)把这份风格、这份手艺收为己有——它几乎与此同时学会说话一般,变成一份第二天性,或相差无几——那么头一个试验品与受益人,毫无疑问,正是那位傻大个爸爸!通常,当我见过这份游戏上演时,其中还叠着一份宿怨的暗恨,外加一份嘲弄的故意。当然,在大多数家庭里,对着父亲的宿怨缘由并不缺——还不算那些由贤惠的妻子巧妙地暗示(乃至凭空炮制)出来的。可在我朋友身上,我从不曾有过片刻感到这样一份宿怨或暗恨的调门。当我见他“为取乐”而伤人害人时,那是\textbf{真的}(如我\dmark{}所感到的)\textbf{只为取乐}\origpage{686};我相信,不是为那份他施加的苦难或羞辱本身之乐,而毋宁是那份暗藏的陶醉:凭他的兴致、在他那手已臻化境的特殊风格里,行使一份\textbf{权力}——无疑更醉人、更够味,凭那份带“\textbf{乖僻}”“\textbf{禁忌}”意味的作料(为\textbf{取乐}而伤人、或叫人受苦)——而那,唯独他能做得起:精微地、不动声色地,此外一切照旧,想玩多久玩多久,管够……\fnote{224}{(*)}{作为周详的例解,尤其见“乖僻”那条注(n\textsuperscript{o}~76)。}

\subsubsection*{神经中的神经——或侏儒与巨人}
\tocsubsub{神经中的神经——或侏儒与巨人}
\label{ch18:local:h61}

\notehead{148}{}
\phantomsection\label{ch18:local:n56}

(12月18日)随着昨晚的反思,我感到葬礼图景的那个“第一层面”——以朋友皮埃尔与我的关系为中心——继续走出未被理解的迷雾、走出那份浑然未明的感受。好久以来,我便面对这样一项任务:在这个第一层面里,安插某个“超级父亲”之联;而我虽不曾真正对自己言明,那份联却似乎不很乐意体面地就位。倘若说我有过一个学生,始终让我感到与我“自在”至极、毫不紧张——就我记忆所及,任何时刻皆然——那正是他!诚然,我对我们最初几次见面的记忆已不多,也无从断言:当时他身上是否已有那份紧张——往往难以察觉、却实实在在——它在我们头一次接近一位被(以这样那样的名义)授予权威或声望的人、对其怀有特殊期待时出现。至少很可能:这样的紧张曾在场,而我对它的注意,不会多于对任何一位我偶然结识的年轻研究者。确定的是:倘若初接触时有过紧张,它很快便消失得无影无踪。借昨晚显出的意象说:他与我在一处,如同一个孩子(或曾经的孩子)与一位他从来不曾需要惧怕、也很少拒绝过他什么的心软爸爸在一处。

昨晚停笔后,我又把那份情势想了一遍。此刻向我显出:我朋友与我的关系,运作在两个截然分明的层面\dmark{}上,\origpage{687}而且(似乎)相互之间没有沟通。其一,大概在我们相遇之后的几个星期、几个月内便已确立,是个人关系的层面——也就是“心软爸爸”的层面:好得不能再好,毫不唬人,自己也有几分顽童气,包括在他的工作里;以致他对我的态度里有一层——我几乎想说\textbf{母性}的——色调:我已有一两次机会提到过:那是人们给一个毛手毛脚、有几分闹腾、尤其天真得没人比得上的孩子所给的那一种。而且,说得更进一步:在工作层面上、客观地说,他确实没有什么可被唬住的理由。当然,数学里我懂得许多他不懂的东西(他花几年便轻巧学会了),尤其,我有一份数学的经验,那还是他所缺的。可他有那份消化的敏捷、那份眼光的锐利——在纠缠混乱的情势里迅速辨明方向——常叫我叹为观止,而那是我所缺的。倘若我自己偶尔叫同事们吃惊,那主要凭我工作中那份不寻常的\textbf{快手},它主要归因(我相信)于我对数学工作的某种进法。可我的年轻而辉煌的朋友,断然没有理由为此唬住:他自己的快手,只要他肯动笔(那事他并不讨厌),分明比我的还要有效。

我朋友与我的关系的这个层面——“心软爸爸”的层面——依我看,涵盖了他对我的意识影像的全体,也无意识影像的很大一部分。依我看,正是这份影像,经由多半自童年起便已铺设的路径,唤起一份回应性的反射性欲望:那份著名的“天鹅绒中的利爪”的游戏——一份恰恰要求人对伙伴完全“自在”、完全“有把握”、并由此也自信的游戏\fnote{225}{(*)}{(12月29日)这份断言,只在表面上被某些个案(不包括我的朋友)驳倒:在那里,“领舞者”似乎(至少初看)被他所玩弄的人唬住、乃至慑服。可那是一份\textbf{架势},为事态的需要而设——而演员本人当然是头一个受骗的(我是说在意识的层面):这对赋予那份架势某种非即兴的“真实”外貌,必不可少!我所知道的这类游戏最极端的个案,是我母亲对我父亲。关于这一点,见“颠倒(1)——或激烈泼辣的妻子”与“颠倒(2)——或暧昧的反叛”两条注(n\textsuperscript{o}~126、132)。}。这是那份全然自信的层面:它安身于对一份情势的贴身认识,又进一步由经验印证——那份经验,被有意识与无意识的知觉与评判官能\dmark{}\origpage{688}以完全一致的方式解读。游戏本身是隐蔽的、对当事人自己也是无意识的(至少我这样推断);可那份自信之感、及作为其根基的对实在的知觉,属于意识的、理性的、“客观”的领域。

而另一个层面,则整个是无意识的(至少是我的印象)、不受控制的、也不可控的:一份非理性的本性,它似乎公然违逆、并嘲弄对“客观”实在(我方才唤起的)的一切有理有据的认识。在这个层面上,本来意义的个人关系——系于对他人的一份多少写实的知觉——消失了。我在其中显为一个巨人:强大而被暗暗嫉妒;而我的朋友自感侏儒,被他自己那份无可救药的微不足道之信念压垮,同时被一份疯狂的欲望吞噬:不是要他自己成为巨人——他以不可改变的境况生为侏儒——而是以这样那样的方式把自己\textbf{抬到}那个高度、至少\textbf{装出}巨人的样子;或者,更隐秘、更阴险地——那份疯狂的欲望:\textbf{做那个巨人本身},或至少,\textbf{冒充是他}。我在这份欲望里辨出又一层细节:它宛如那份欲望——在场于接近表面的层次的那份——在更深层的回声;表层那份欲望,恰恰在“以天鹅绒的脚掌”这份游戏里找到象征的满足,而它正是游戏的神经与发条:\textbf{角色颠倒}的欲望。在较浅的层次,涉及的是阴阳角色的颠倒:被支配-支配者、客体-主体。可这份关系在此并不合用:因为巨人绝无支配侏儒的操心——他满足于做巨人,并由此,在不知不觉、也不曾操心之中,做一份永恒而灼人的\textbf{挑战}:对着那个自感被一份无可救药的侏儒境况压垮的人……对方感到自己被圈在这份傲然的无知里,他把它感受为一份默契的轻蔑、一份冒犯。他心火烧灼,要颠倒的正是这份关系:他自己显为巨人,把巨人贬入微不足道——以\textbf{遗忘}致微不足道,不然便以\textbf{嘲弄}致微不足道,作为对那份他自感身处其中的无知与轻蔑的正当回报。

我方才说过,两个层面——“心软爸爸”与“巨人”——“似乎没有相互的沟通”。细细想过之后,眼下我倒觉得\origpage{689}:两者之间确有一份沟通——哪怕只凭那份颠倒的欲望:一个层面上的欲望,此刻显为对另一个层面所确认的相似欲望的一份“回声”。初看,我曾觉得:那个“侏儒-巨人”的更深层面上的角色颠倒,不属于阴阳型的角色颠倒。真确的是:这份颠倒确实不属于被支配-支配者的类型。可是,再想一层,毫无怀疑:巨人身上所体现的\textbf{价值},是阳的、超级阳的价值;而侏儒显为阴的非价值的化身——我是说,按我朋友的意识形态抉择的措辞;那些抉择,与我们关系最初年月里我自己的抉择,相去不远\fnote{226}{(*)}{阳的或超级阳价值之拣选上的这份一致,一直延续到我出发的1970年。其后的年月里,我有意识层面的价值系统,“翻转”向阴的、超级阴的抉择——见“阳扮阴——或主人的角色”那条注(n\textsuperscript{o}~118)。}。

这份断言将变得清楚,一旦我在“侏儒与巨人”的影像与实在之间架起一道桥——至少,显豁这份影像在我朋友与我关系的史前史与历史中的起源。关于“史前史”,几乎无须说明:这样一个有意识或无意识的影像,唯有借着那份埋藏极深的“对己之轻蔑”方能诞生——我在反思中已屡次唤起它;或者更好地说:这样一份影像,不是别的,正是这份轻蔑的一份可触的、多少具体的\textbf{落实}。甚至或许可以说:这份“秘密的信念”守候着一份能给它充当支撑的情势,并同时催生那幅表达它的稻草人影像。我相信:在心灵的万事万物之中——无论埋藏多深——都活着一份力,催促它表达自己,往往以象征的方式。这份表达,固然常常自身仍是无意识的,却并不因此少一分作用——恰恰相反——在日常生活的可见言行层面上。

回到这一次是我朋友与我这人关系的\textbf{历史}:它肯定也始于我们相遇之前。他大概是在1960年前后、在布鲁塞尔最初接触数学界的时刻,听说过我的——也就是说,比我们相遇\dmark{}\origpage{690}早四五年,当时他十六七岁\fnote{227}{(*)}{(12月29日)这条编年信息,我取自Pierre~Deligne1975年为(比利时)“国家科学研究基金会”“五年一度奖”撰写的那份两页的“传记小记”(Rue~d'Egmont~5,~1050~Bruxelles)。我打算在一条后续的注里回到这份小记:谈Deligne去年10月的来访。正是在那次来访中,我由他得知这份小记的存在;经我请求,他后来惠寄给我。也正是在这份小记里,我找到了我朋友身上某个影像的具体形态“侏儒与巨人”——一份弥散的观念,已在葬礼反思的进程中逐步析出。它开始显形于“下葬”那条注(n\textsuperscript{o}~61),并在“驱逐”“死结”“颠倒”“屠杀”“……与电锯”“葬礼悼词(2)——或力量与光环”各条注的反思中逐步明确。唯有借着本注,这份知觉才开始“落位”于葬礼“第一层面”的一份连贯总图。
(1985年3月)关于Deligne的传记小记,见“信条——或假中之真”那条注(n\textsuperscript{o}~166)。}。这肯定不是偶然:他请求教他数学家手艺的、正是我而非别人;或至少,请求教他那份即将成为他著作之中心主题与中心工具的东西(即代数几何)。在我们相遇之前,我以什么面貌向他显现(至少作为数学家),想必只能是我的招牌影像:它把我做成一份英雄式的、声名显赫的化身——数学世界里通行的那些主导价值的化身;而彼时,他自己还是个普通的学生,刚从中学毕业。他对我怀有的这份影像——它恰恰也是我乐于示人的那份——况且还不是一幅单纯的“埃皮纳勒”版画,专供渴望荣耀的中学生们做梦之用。它由可触的实在筑成;而他无疑颇有直觉,早在那些年头,便在与年富力强、十足在局内的数学家的接触中,嗅出它的气味。自1965年起,他更是比任何人都有资格,亲自来掂量我的斤两。我当时在他身上感到一份着迷:为一份向他敞开的景——它在我身上诞生、成熟于已逝的十年之中,并继续在他眼前铺展、发育。当时对我毫无怀疑:这些他“仿佛从来就认识”般纳为己有的景,将公开为他效力,充当灵感与工具,去发展更宏大的景、一部与他的手段相称的作品。事情全然不是如此——只是借着这份将近二十年之后、对一场葬礼的长程禅观的光照,我才隐约看清:那份对我所传递之物的精微而热烈的知觉,如何\textbf{同时}被用来充实、支撑一幅\textbf{稻草人影像}\origpage{691}——一份畸形的影像:它凭第一手的、无可辩驳的实在的要素而加固;一份足以\textbf{瘫痪}人的影像,一如它所表达的那份“贴身信念”。他对我在传递之物中的那份“伟大”与深度的知觉之敏锐——那份他独力(且毫不费力)整体纳为己有的知觉——正是这份敏锐与这份鲜活,曾构成他的力量,却旋即反过来指向他自己:把那幅畸形的影像,衬得愈发触目、愈发专断。

三天前,我曾自信触到了我的朋友将近十五年来所扮角色的“神经”——而当时确凿无疑:我刚触到一处要害:那份吞噬一切的\textbf{饥渴},渴着某一份\textbf{游戏}——一份精微的权力游戏;它同时,是一份特定角色颠倒之欲望的象征性的、昙花一现的满足……随着今天的反思,沉入更深的层,我觉得此刻触到了\textbf{神经中的神经}、那份更隐秘的\textbf{刺棒}:它不停地唤起并滋养这份饥渴。因为在“心软爸爸”的层面,固然有着机会与全副自由,安安全全地玩这份游戏:潇洒而精微地领舞,稳操每一招的胜券。可无疑,一旦没有刺棒,那轻易机会的魅力便钝了。而一如我昨天还确认的:面对心软的爸爸,没有那份咽下去的怨怼的刺棒、没有那份暗藏的宿怨——这正是人们叫它“心软”爸爸的缘由!这份缺席的刺棒,我方才忽然触到了:顺着种种联想,仿佛凭一份早已备好的认识的口授,我被引去描述那个“另一层面”——那份“不受控制也不可控”的层面:侏儒与巨人并肩居于其中。

最初那份直觉尚且含混的印象——两个层面之间没有相互沟通——顿时消散,让位于一份理解:它由“神经中的神经”与“刺棒”这双幅意象同时表达、同时引发。这次改用“层”的措辞——一些浮面,一些深层——我愿以第三幅意象再说一遍:深层滋养、维持着浮面的运动;深层是它们的根基,牢牢锚定在“我”的结构之中。没有这份根基,表面的躁动很快便会消散,终于让位给别的东西……

\subsection{另一个自己}
\label{ch18:local:h62}

\subsubsection*{(a) 缓期中的宿怨——或事物的回转(2)}
\tocsubsub{(a) 缓期中的宿怨——或事物的回转(2)}
\label{ch18:local:h63}

\notehead{149}{}
\phantomsection\label{ch18:local:n57}

\origpage{692}\dmark{}(12月20日)自五天前的反思以来——尤其自那天第二条注“秘密的神经”(n\textsuperscript{o}~145)里继续的反思以来——我感到:对葬礼图景那份著名“第一层面”的工作,骤然换了气象。在那场反思之前,我处于一个有点尴尬的位置:面对一幅拼图,却觉得对它不甚了了。自4月起,我便一直努力把拼块一一收拢、细心编目。倒不是缺拼块——不是的,我反倒觉得太多了!总之,拼块肯定足够拼出一幅图:也许只是局部,却是站得住的一幅。我掷上桌子的最后一块拼图,是那份“颠倒”(阴阳的)——它自“阴阳之钥”的发端起便被存而备用(作为一场我应许自己要回头谈的“联想”),终于在11月10日那条注“阴的葬礼(阳葬阴(4))”(n\textsuperscript{o}~124)里,以预料不到的强力闯了进来。其后的三十五天——直到五天前——大体用来把已出土的拼块翻来覆去,顺着最强制地要求我注意的种种联想\fnote{228}{(*)}{作为整场阴阳反思出发点的那个“拼块”(自10月初起),直到十四天后——11月24日,在“颠倒(3)——或阴葬阳”那条注(n\textsuperscript{o}~133)里——才重新登场、得到显豁。}。我本以为,这样做下去,那些拼块终会自行拼合,终于显出那幅未知的图景。事情全然不是如此。恰恰相反:它们继续互相顶牛,一如十张各不相同的报纸碎片被乱扔在那里,却要我来拼!我开始盘算:到头来,我是否会被迫作一份拼块的最终清点、再作一份关于其拼合的问号清单,就此收场……

五天前,情形变了:翻来覆去地摆弄那些著名的拼块、摩挲、嗅闻之后,终于有什么“咔哒”一声对上了——其中一块(那份藏在某个\textbf{风格}背后的\textbf{饥渴})骤然被承认为“要害”。我当即有一份\textbf{质变}的印象:一份此前一直缺席的\textbf{视景},已然\dmark{}开始就位\origpage{693}、要从那块拼块出发组织起来。仅仅两天后,我在下一条注(“激情与饥渴——或逐步加码”,n\textsuperscript{o}~146)里接续反思时,正是以这样的措辞表白自己的。而我的预感当天便已开始得到印证:随着“\textbf{心软爸爸}”那块拼块的现身——人们简直要说,它是应“要害拼块”之召而来的,为的恰恰是严丝合缝地嵌入!

那份“\textbf{超级父亲}”拼块——它一直在场(从《收获与播种》第一部分继承而来,并自“阴阳之钥”的发端起便被重拾\fnote{229}{(*)}{见“敌对的父亲(1)(2)”两节(第~n\textsuperscript{o}~29、30节),以及“超级父亲(阳葬阴(2))”那条注(n\textsuperscript{o}~108)。})——此刻似乎转眼成了盈亏账上的弃子,仿佛它只是不小心在那里迷了路。在“心软”新拼块\fnote{230}{(**)}{“新”拼块这个说法,或许不完全站得住。但至少,它是一块此前逃过编目的拼块——正因为它太显而易见了!}的新鲜印象之下,我倾向于忘掉:那份著名的超级父亲(这一回,一点儿也不“心软”)确确实实与朋友皮埃尔和我的关系有涉,即便它不曾领着那台戏(差得远了……)。可我终究在下一场工作里想起来了——必然如此:恰恰在那一刻,我正准备对自己解释:为什么拼图里那块永世的拼块,其实与此毫无干系!总之,它是“心软”拼块的“恰恰反面”——后者方才以那样的从容自行就位。可不然:细看之下,这块据称与游戏无关的拼块——它的轮廓一直最模糊——骤然显出了形状:“取得”了那幅力量意象的形状(道出这幅意象的,不是别人,正是朋友皮埃尔本人\fnote{231}{(***)}{个中细节,见上一条注“神经中的神经——或侏儒与巨人”(n\textsuperscript{o}~148)的最后一条脚注。}):\textbf{侏儒与巨人}。起初,见它以如此鲜明的轮廓重现,我料想它会与那对已就位的“要害双拼块”(由心软爸爸、与那份强制性的“让他转起来”的饥渴——一个电话这边、一个小动作那边……——构成)“不相沟通”。可它偏偏显形为“神经中的神经”、一份更要害的拼块,与已就位的拼图部分毫无摩擦、毫不脱胶地拼合!

这块拼块,顶着它的旧名“超级父亲”,其实已被擦过许多次\dmark{}\origpage{694},甚至被拿在手里、像其他拼块一样翻来覆去,甚至(我现在记起)被宣布为关键拼块、“图景的心脏”云云;可是——或许由于它不曾由一幅醒目的意象来化身(那意象由当事人自己提供),而尤其无疑由于它那荒谬的、有违常理的本性——按当下通行、举世公认的粗浅“常识”来看,它甚至全然滑稽——我对这块该死的拼块尴尬而近乎羞愧;它在我手里发烫:永远不会有人(包括那个还执拗地活在我身上的某位“我自己”……)肯拿它当回事!还不如悄悄打包收起,拿些更拿得出手的拼块来“玩”!

我方才谈“关键拼块”、“图景的心脏”云云——就那块变成了“侏儒与巨人”的拼块而言——我想的当然是“对己之轻蔑”的侧面,而非“超级父亲”的侧面。眼下,对这块刺棒拼块、或“神经中的神经”作“超级父亲”的指称,还为时过早,也无从论证。我是说:至少初看,那位著名的无面孔、双手硕大无朋的巨人,并不显得有多少父性。倘若非得给它一个名字,配得上它的,似乎是“超人”或“超级雄”,而不是“超级父亲”。所以,细细算来,后者眼下确确实实还“记在账上”——至少眼下如此,一如“超级母亲”那块拼块(或“联”):那份我日后同样要回头谈的。

眼下,最要紧的依我看,是试着以我的朋友这个人身上的阴阳动力学,给图景中已就位的部分定位——连同那份“秘密的神经”与更隐秘的“神经中的神经”。这方面,我手头有三件裸的事实。其中两件,由那份阴-阴的“双重签名”表达\fnote{232}{(*)}{“双重签名”的观念,是随着“兄弟与配偶——或双重签名”那条注(n\textsuperscript{o}~134)进入反思的。}:朋友皮埃尔基调属“阴”,无论在可称作“习得人格”的层面——它主要在与他人的关系的调门里表达——还是在“先天人格”或冲动的层面——它主要(至少对我这样一个外在观察者而言)经由自发的、免于“东家”插手的工作风格表达。第一件事实关乎习得人格、关乎“我”的结构(或以更形象的措辞,“东家”的\textbf{脑袋}),它似乎表明:这份结构是在童年、在人生最初几年里完成的,凭与一个“阴”性模型的认同。这在先验上并不排除:模型曾是父亲——倘若父亲本人(如我看来确系如此\dmark{}\origpage{695})便有一份基调属阴的“习得人格”。但另一方面,我朋友身上那份癖向:渴着一份权力游戏——它在我们这些地界(即便不是处处、事事)典型地(若不说排他地)是“女性”的,而更确切地说,它是妻子惯于与丈夫玩的那份游戏中的\textbf{头号}游戏——这份癖向使我推测:认同是与母亲这个人完成的;正是从她那里,他“继承”了这份饥渴(或一份对此类饥渴的癖向);也正是从她那里,他接过了那份合宜的“风格”(或“战术”):“天鹅绒中的利爪”。

有可能,父亲既是一位心软丈夫,又是一位心软父亲;而我的朋友,久已有的是机会拿他做自己第一个“试验品”,在他身上磨爪子(和绒毛!)。但也有可能:我朋友身上这份癖向或禀赋,一直闲置未用,直到与我相遇之后——因为头一个现成的靶子、即他的父亲,缺少足够鲜明的阳特征,不足以“\textbf{挑起}”这份饥渴、又不足以让那份对付倔家伙的考验过的战术\textbf{得手}。老实说,凡我记忆所及、落在我与朋友相识最初年头里的印象,没有一份能暗示:他对那份游戏已然熟稔,甚至已然实践过。无论如何,即便带着回顾,我在他与我的关系里、在他与别人的关系里,也辨不出一丝痕迹——比如稍微带点“被惯坏的孩子”味儿的做派。所以我反倒倾向于认为:他身上这份癖向当时还潜伏着;它得以发育、并取得我所知的那份对他生活与工作的掌控,是在我1970年“亡故”之后(他时年二十六),借着一份格外诱人的情势。

这里该回顾的“第三件事实”,是我的朋友对一份与通行价值相合的价值系统的抉择:对“阳刚”(或阳)价值的抉择。而这份抉择,在逝去的十五年里,依我看,在他身上愈发转向“超级阳”。在他的个案里,这份抉择有一道一望即知的矛盾:一面采纳\textbf{阳}的“官方”价值,一面却在大多数要害特征上,照一个\textbf{阴}的模型自我塑形\fnote{233}{(*)}{这是一类矛盾,尤其多见于女人;我自己的生活则与它无缘。}。而这并不是说:这份价值的抉择纯属“赝品”、只是为情势所需而竖起的假幌子、只在心灵的边缘层次\dmark{}流通\origpage{696}。侏儒与巨人的力量意象,从深层起作用;倘若对阳的抬举不曾也在这些层次里被内化,这份意象便会失去意义,连同它所引发的那份强制性的颠倒之饥渴。毫无怀疑:这份矛盾,必定给那份“裂缝”“阴险无能”的“贴身信念”追加一份活力——而与此同时(或许只缺童年里一个可供模仿的合宜“模型”),他自知(在心底)根本\textbf{不同于}他“\textbf{理应}是”的那个人!

倘若我的朋友——这在我看来是说得通的——不曾在他父亲身上找到那些照当时周围的共识\textbf{本该}在的、而他本可纳为己有的特征,这必定在他身上唤起一份弥散的\textbf{宿怨}:一份无法攀住任何具体怨怼的宿怨——对着一位心软爸爸,他唯一的缺点竟是太“软”!这份宿怨,苦无一个可以攀挂的“挂钩”,想必始终“空置”、\textbf{待命}一份相宜的靶子——恰恰是这样一份靶子:首先,(由语境)扮出父性面孔;而且,担纲此角的\textbf{资质}昭然,凭那份不容置疑、甚至耀眼乃至过分的特征的在场——正是他“原装”父亲身上所缺的那些特征。也恰恰是这些特征,使新来的“父亲”成为理想的\textbf{靶子}:在一场此地早已备好、只待相宜伙伴的“游戏”里——伙伴别名“替补父亲”,别名(我们终于来到这里!)“超级父亲”!

忽然,我觉得自己回到了一片非常熟悉的地界——直到此刻我才认出它。这是一片我做过二十年囚徒的地界:在我一生唯一的一场婚姻里(我的五个孩子中有三个出自这场婚姻)。在上面几行里,我完全不是存心(毋宁说像一个小心翼翼在黑暗中摸索、认识周围事物的人),也依次描述了那个曾是我妻子的人对其父亲、继而对我的关系中的要害力量。我无从说清,这份认识(或毋宁说那份无可辩驳的直觉)——关于这两股力量在她身上的无声而执拗的在场、及它们的相互关系——是何时、如何来到的。有一天我懂了,却从不曾费心稍加反思:那份无情地支配着我妻子与我的关系的力量——自我们成婚的最初日子起——正是被那份对我宿怨所驱动:怨我不曾在一段无依无靠的童年岁月里,守在她身边,做一个\textbf{别的、真正的}父亲……

\dmark{}诚然,我知道:我朋友的童年毫无“无依无靠”可言;而他发育成形、并为我所认识的那个人格——六十年代直至如今——与我前妻的人格相去甚远\origpage{697}。但这并不妨碍:在那些明显的差异之外,我看到——在正从阴影中浮现的图景那部分里——一份惊人的相似,与另一幅我熟之又熟的“图景”相似。这份相似,显于对父亲关系的本性(系于一份阳特征匮乏的父亲气质),以及它对一份成年关系的回响:在两人身上,这份成年关系都支配了其一生,都成为两人身上冲突力量之瞄准的焦点\fnote{234}{(*)}{(1985年2月19日)朋友皮埃尔与我的关系、与我妻子(自婚后最初日子起)与我的关系之间,确有一份惊人的亲缘。这份亲缘更超出对单一我的关系:两人都终于发育出一份癖向,把某些与我以情感相系的存在(一者是我的孩子,另一者是我的学生)变成\textbf{器械}——经由他们来打击我。}。

差一点,我便会略而不提第三份“相似”——它在我自己的生活里却不无分量:在这两份关系里,担纲的\textbf{主角},每一次都\textbf{不是别人,正是我}。而在这两个个案里,把我指派给那份我注定要扮演的“超级父亲”角色的,是(除一份不成熟之外)还有一样东西:它自我童年起,或许便比世上任何东西都更亲近于我——也是我在其中投入最无度的东西:一副天生超常的“阳刚”身量……

就这样,我以一道比八个月前更不同、更透彻的光照,重新寻回那份“事物的回转”之感\fnote{235}{(**)}{见“事物的回归(或一脚踩进菜盘)”那条注(n\textsuperscript{o}~73)。}——今天一如当年,带着一份难以置信的惊异的调门(事情似乎落得太“巧”,巧得不像是真的!)。而且,这一回也如此——只是调门比当年那阵骤然的大笑更收敛——有一份滑稽的知觉,为这些无情的“回转”添上幽默那份更柔和的音符。

\subsubsection*{(b) 天真与冲突——或绊脚石}
\tocsubsub{(b) 天真与冲突——或绊脚石}
\label{ch18:local:h64}

\notehead{150}{}
\phantomsection\label{ch18:local:n58}

(12月22日)昨天我又没找到时间在注上工作,只除了细心重读并修改前天的注。最近这些天,我的精力被通信\dmark{}\origpage{698}等任务岔开;我又有点按捺不住急切(这可不是新鲜事!):急着重新单独面对自己,推进已经开始的反思。在《收获与播种》的这第三部分——以眼下这场反思“阴阳之钥”为中心——写作分明更慢了:阴阳的动力学,是深入葬礼意义的那根恒常的引线。倘若不留神上好闹钟,以便在工作约三小时后安排一次休息(活动活动筋骨,或提醒自己夜已深、该停下了),一整夜会像一瞬般过去!那三个钟头,每一次都是这样过去的:我只觉得刚刚开始(或刚刚接上),只打出两三页可怜的纸——有时只有一两页——刚刚够巡完某场看似平淡、我本想一口气跨过去的联想……

按每小时或每天的页数计,有一种极慢的推进之感——对这份印象的自然反应,既然一份滚烫的实质正在我鼻子跟前拽着我向前,本该加倍再加倍地赶——一如直到最近几年我还惯做的那样。可我知道:这正是该避开的陷阱——发现工作中那份出奇的“轻捷”的陷阱\fnote{236}{(*)}{见“陷阱——或省力与衰竭”那条注(n\textsuperscript{o}~99)。}:只消轻轻“推”一下,便确信自己在切实前进——也许慢,但稳;像一个人双手扶住一辆钢犁的扶手,前面由一双强健而沉稳的公牛牵引:他慢慢地、稳稳地犁开一条路,一垄又一垄,穿过一片紧密的、有时倔强的、却又柔韧的土地——它顺从那把闪亮的犁铧:犁铧不慌不忙、精微地把它翻开、穿透、掀转,成大片大片冒着热气的黑土,把一份强烈而蠕动的地下生命带到光天化日之下。节拍或许慢,田地却广,每一垄犁沟似乎才动了动那片尚待开垦的荒地。可一天终了,一垄又一垄,田已犁遍;犁田人满意地收工:对他而言,这一天不曾虚度。他的辛劳与他的爱,是他的种子;他劳作中的快乐、每一垄尽头与长日终了的那份满足,便是他的收成与他的报偿。

\asterisms

\dmark{}\origpage{699}随着前天的反思——也许是在《收获与播种》的写作中头一次——我有这样的印象:我踏上了那片不确定的地界:尚未被直接知觉或感到、而依然(或许将始终)\textbf{属于假设}的地界。苦于没有能在看似半明半暗之夜里看东西的眼睛,我摸索着开出一条犹疑的路,毫无把握它是否“对”。当路分岔,我确实不曾掷硬币;我凭直觉与常识,为自己指出最说得通的继续方向——却毫无任何念头:它将把我引向何处。我这样循着、或这样画出的路,看上去处处“贴合”我所知的事实;这是个好兆头。但这并不排除:尤其在那些事实稀薄之处,另一条全然不同的路同样会“贴合”——前提也许是再把某条尚属生涩的事实挖掘一番……然后,在路的转弯处、在我自己的惊异之中,我忽然发现自己来到“一片非常熟悉的地界”:当年我曾长久而艰难地跋涉过、终于认识并离开的地界。一份几分钟前还显得晦暗的情势——被“想必”与“或许”的疑云包裹——骤然被另一份情势的光照亮:而后者,是被懂得的。当我追问那份朋友与我的关系中的冲突、在我与在他身上的遥远根源时,它们似乎经由一份骤然瞥见的深刻相似而显露——那份关系与另一份关系之间的相似:另一份曾以全然不同的分量,压了我漫长二十年的一生。

这份相似的显形,力量如此之大——我承认——那份迟疑、无把握、摸索之感旋即消散,让位于一份有把握、确信之感。当我在反思的结尾谈到那份感觉(“难以置信的惊异”):事情“落得太巧,巧得不像是真的”,那份感觉是对另一份的回答——另一份在底色里说:事情“落得太巧,巧得不\textbf{能}不是真的”!而后一份感觉——就我手头的事实现状而言,肯定仓促而毫无根据——其间并未调整过来:它始终作为底色在场,无论我愿不愿意。肯定地,若非借助某些我终于懂得、终于承担的经验——尤其,借我婚姻生活的漫长经验——那份“空置状态的宿怨”(总而言之,一份“缓期中的宿怨”)的念头,几乎无从来到;而恰恰是这同一个念头,又是那个“路的转弯处”:几瞬之间,它使我重新驶出那片婚姻经验的“非常熟悉的地界”。

\dmark{}\origpage{700}当然,可以说:一份无意识的故意,把我引到一个早已被指定的地点——那个地点,或许正说明一些关于我、关于这份故意的东西,而对他人身上的动机,则说明不了什么。同样可能:一份被承担的经验,使我得以把握他人身上的一份实在——否则,它会整个是个谜,怪我没有足够灵敏的“天线”(也怪我手头没有关于我朋友童年、他父母各自人格的可触事实)。

依我看,我离完成葬礼“第一层面的图景”的速写(信手写来!)已经很近。为拼上我手中剩下的最后几块拼图,必要时我会动用前一条注的反思中显出的领会要素(无论它们多么属于假设)。这倒也是一份办法:检验它们与我所知的其余事实的融贯。

前天的反思里,是拼图的“超级父亲”那块,显出了形状与轮廓。我起初把它认作“侏儒与巨人”那块——有点匆忙;在那块里,巨人毋宁显为一个尺寸骇人的“超人”,而不是“父亲”、或“超级父亲”。可后者那块,终于在同一场反思里重新显现——这一回,作为一份“缓期中的宿怨”的靶子:一份正在寻找靶子的宿怨——仿佛那位“超级父亲”恰是\textbf{被}这份宿怨所召、应这份召唤而来,兑现一份弥散的期待。倘若果真如此,可以说:倘若超级父亲(借这场合借得我的身量与我的面貌——它们分明是度身定做的)不曾出现在我朋友的生活里,那就得发明一个!至少,对那个我曾做其丈夫的人——而且我曾是“那支被盼了整整年轻一生的箭靶”——这份说法,对我已毫无假设可言。

于是,超级父亲显为“侏儒与巨人”那块拼块里那位“无面孔而双手硕大无朋的巨人”的“有面孔的一侧”。“侏儒”看巨人,想必主要看背影:巨人正忙着做他那著名的“力量展示”(10月5日“超级父亲”那条注(n\textsuperscript{o}~108)里谈到的)。于是,“超级父亲”那块拼块终于就位:与“侏儒与巨人”拼块的“巨人”一侧拼合。

至于“侏儒”那一侧,它的轮廓也在前天的反思里显得更清楚:它在此接上10月17日那条注“一半与全——或裂缝”(n\textsuperscript{o}~112)的反思。还是仿佛——常常如此——那份对“阴”、“女性”特征的永世拒斥(为了让位给“阳”、“男性”的特征),使我朋友“根本不同于他‘理应是’的那个人”,而他本人却照一个以“阴”为主导的模型自我塑形\dmark{}\origpage{701}。

这里须要强调:在整场已往的反思里,我从不曾想过、也不曾想过暗示:我朋友这个人打有以阴为主导的\textbf{失衡}的印记,即:在他习得人格的阳的、阳刚特征一侧,有一份匮乏、一个“空洞”。关于这一点,我提醒:至少在我认识他的最初几年,从他的这个人身上透出的那份印象,恰恰是一份\textbf{平衡}、一份和谐的印象——它使他对我、对我以为对所有当时认识他的人,都如此可爱。这份印象与另一份系得很紧——我在别处谈过\fnote{237}{(*)}{关于这一点,见“孩子”那条注(n\textsuperscript{o}~60),第五行列“我的朋友皮埃尔”内。}:他似乎保有了几分孩子的新鲜与天真——在他把握事物(尤其数学)的方式里;而且依我看,也在他对人的方式里。这份平衡、这份“新鲜”或“天真”,对我而言无可怀疑:它们是\textbf{事实},谈不上要把它们变没了。它们在我朋友身上表达为一份精微的感受性,而且时机一到,便以有层次而不含糊的方式,表达被知觉与被看见的东西。有一份坚定,正如有一份温柔。温柔在岁月中消隐了,只留下一副毡垫般柔化而空洞的甲壳——那消失了的温柔的甲壳;而坚定变成了封闭与坚硬,藏在一层借来的、矫饰的半音调式的门面背后。一份精微的阴阳平衡,在岁月中(无疑无人察觉)转变成了那永世的阳失衡——正是它,只是风格不同,曾支配我自己的生活,自我童年起。这是他的抉择;而抉择可以变——牌局从来不曾定死!无论如何,在我朋友的生活里,我从不曾知道有任何一段时期,打有阴失衡、疲软、放任或不一贯的印记;我也不认为曾有过。

这一切使得至少很可能:那个在他童年充当“模型”的人——那人肯定有强烈打上印记的阴特征——并不因此缺阳特征去与它们平衡。倘若(如我倾向于相信)那个人是他的母亲,那么我推想:她有着足够鲜明的阳特征(尤其对着父亲身上那些无疑较不鲜明的特征),以致显得是“最佳选择”——充当一个男孩的“男性”模型;而与此\dmark{}同时\origpage{702},凭这样的选择,助成一份和谐气质的绽放。

至此,一切似乎都尽善尽美,在最好的世界里,在一个(或许)没有任何不和搅扰的和睦家庭里。一切本都会尽善尽美——倘若没有那一小块绊脚石,取一份无声共识的形态、貌似再平淡不过:男孩理应像父亲,而不是像母亲……

\subsubsection*{(c) 天赐的情势——或封神}
\tocsubsub{(c) 天赐的情势——或封神}
\label{ch18:local:h65}

\notehead{151}{}
\phantomsection\label{ch18:local:n59}

(12月23日)依我看,要把葬礼第一层面的“拼图”拼完,只剩最后一块该就位。那是我在11月11日那条注“超级妈妈,还是超级爸爸?”(n\textsuperscript{o}~125)里称作“超级母亲”的那块。这份“超级”式的命名,起初由我葬礼悼词里那份以超级形容语一笔一笔画成的我的“肖像”所启发\fnote{238}{(*)}{见“葬礼悼词(1)(2)”(n\textsuperscript{o}~104、105)两条注,以及“阴的葬礼(阳葬阴(4))”(n\textsuperscript{o}~124)。}。肯定,一份对称的反射也起过作用:既然“超级父亲”已经在空气里,且不止一重!可是细细想过,我给那幅方才显现的影像起的名字,并不十分贴切。那幅超级阴的影像所唤起的,毫无“母性”的意味。倘若它与另一幅影像有对称关系,那是与“超人”的对称——钢的肌肉、IBM软件的头脑——而不是与“超级父亲”的对称。所以,在此情形,毋宁说该是“超女”或“超级妞”:垂到肚脐乃至更下的沉重乳脯(更不必说到膝盖……),臀部亦然,足以叫赫拉克勒斯做梦——至于头脑,且不必提……大抵是这样的调门。语言的不足,想必也多少逼了我的手:著名的“超人”(它本身也是晚近的发明——大力神赫拉克勒斯的一个现代版,分明已被时势超越)没有现成的“阴性”对等物。将就着,“超级妞”也罢,聊胜于无……

得承认:我拖着这块起错了名的拼块,拖了将近一个半月,不曾真正动它分毫,除了此地彼处以备忘的形式重提一句,作为一种许诺:总会去处理它,不过往后。最终,它实在不曾给我多少灵感;而这或许正系于那个贴不紧的名字\dmark{}\origpage{703}。毕竟,在我直到今天还在的数学世界里的所有朋友、(昔日)学生与其他同事之中,我会很难找出哪怕一位:对着他,我曾多少扮演过一个“母性”的角色;或者,我曾有印象,他觉得我担着这样一份角色。甚至那些对我而言我扮演的角色偏于“阴”、偏于接受的人——而不是“阳”的施教者、传达者、传递者的角色——也必定极少:粗粗一望,我几乎看不到谁(1952、53年我写论文的年头之后)——除了Serre,勉强……倘若我试着回忆我与别的数学家打交道时通常(甚至不妨说恒常)的心态,那主要是:我总有崭新的“地毯”要“铺出去”(借我当时流行的说法),更不必说那些(也出自我手)不那么新、却(依我看)几乎没用过、而在我看来为一家数学之宅的体面所不可或缺的“地毯”——在数学的某个我熟悉的街区里。换句话说,在我与我的数学“同行”的关系里——即便我们一同几乎只谈数学(在这点上,我该比我的任何同事朋友都糟!),我习得气质里的阳优势(或毋宁说超级阳的失衡),便重新主导一切,一如在任何其他关系里。甚至或许更强,鉴于我对数学的无度投入——那份投入,取利己的品格(无须说破),况且,恰恰由我积习已久的超级阳抉择所驱动!

正是这些显而易见的侧面——它们在我与其他数学家的关系的每一步显形——必定把它们那\textbf{另一份}事实遮蔽了,对同事们一如对我自己:我的数学工作风格、我对数学的进法,是阴占强势主导、“女性”的。依我看,正是这份特质——它在科学世界里显然颇为例外——也使这份风格如此\textbf{可辨认}、与任何别的数学家如此\textbf{不同}。这份风格确实“与众不同”,这一点,自我发表数学以来——至少自我的论文工作(1953年)以来——已由无数的回声传回我这里。况且,这份风格不曾少引起抵抗——我乐意称之为“发自脏腑”的抵抗:我是说,那些抵抗,依我看来当时(而且今天依我看依然)不能用可以称作“客观”或“理性”的“缘由”来辩解。这使我记起:我的论文(我在其中尤其引入了核型空间),投给《美国数学会回忆录》,曾被第一审稿人否决——一位在该领域工作过的、名声端正的数学家,他把我的工作看得多多少少有些“花架子”。多亏Dieudonné的有力斡旋,我的论文才不顾审稿人的否定意见得以发表\origpage{704}。几年前我得知:它属于过去二三十年数学文献里被引用最多的百篇文章之列\fnote{239}{(*)}{或许,此处我的记忆在骗我:说的可能是泛函分析里被引用最多的百篇(或二十篇?)文章。}。我推想,倘若我们前面还有二十年、三十年数学,同样的话对SGA~4也将成立——作为几何拓扑里topos视点的基本参考文献之一;而SGA~4,曾被我的辉煌朋友兼昔日学生Pierre~Deligne判为“读不通”(还有同一水质的其他考语\fnote{240}{(**)}{见“白板”那条注(n\textsuperscript{o}~67)。})。我知道(他自己也知道):它是我花时间最多、照料最极端的数学文本之一——尤其把所有关于site与topos、及范畴预备的部分,彻底重写、再彻底重写。这份格外照料的缘由,是:我分明感到,这是我正在奠基的“算术几何”的一块真正的奠基石——那份奠基,我已做了十年\fnote{241}{(***)}{这肯定也是那份缘由:Deligne如此执意要败坏这份文本的名声,以致他有时竟忘了自己偏爱的半音调式风格,毫无节制地来贬它!关于这一点,见前一条脚注已引的“白板”那条注。}。我也知道:当我做这份工作时,我久已(不想自夸)握有那份写数学的大师手法:既\textbf{清晰}——主导观念如一条无所不在的引线始终摆在前头——又\textbf{便览},便于查考\fnote{242}{(****)}{况且,正是借着熟悉SGA~4已誊清的部分(1965年,他刚到我的SGA~5研讨班),并亲手编写其中某些讲义(凭我的手稿笔记),这位Deligne在与我的接触中,学会了编写数学文本的手艺——尤其,把一份繁密而复杂的实质清晰地呈现出来。}。倘若我或许错了——提前四五十年的光阴,写(并让人写)一份周详的参考著作——那么,那些曾经成熟的年代(六十年代)骤然不再成熟,这笔账依我看记不到我头上!

\dmark{}这些与Deligne的联想,把我带回我出发之后的时期:同方向的回声,不止一次传回我这里——“如一阵阵阴险的轻蔑与含蓄的嘲弄”\origpage{705}。这层\textbf{嘲弄}的调门,是我出发之前那些“对我的工作风格的发自脏腑的抵抗”的征兆里所没有的——我方才唤起过。我在其中辨不出任何敌对的意图、或对我这个人多少恶毒的意图。甚至在Bourbaki之内,我也有机会唤起过这类征象\fnote{243}{(*)}{尤其见那条(无名的)注~n\textsuperscript{o}~5,《收获与播种》第一部分内。}——至少(如果我没记错)直到1957年前后:那一年,我关于Riemann--Roch--Hirzebruch--Grothendieck公式的工作,驱散了可能仍有的对我作为数学家之“扎实”的怀疑。1957至1970年(我“出发”之年)之间,我不记得知觉到对我工作风格的抵抗——除了偶尔在Serre身上\fnote{244}{(**)}{关于这一点,见“兄弟与配偶——或双重签名”那条注(n\textsuperscript{o}~134)。}——却从不带一丝敌意的调门:那宁可是一份皮肉的恼火反应。相反,我有这样的印象:我的朋友们有时感到不堪重负,因为我走得太快,而他们不愿把时间整个花在跟上我的全部著作——随着我把一部部砖头寄给他们,或(凭信件或当面)讲我正在炮制的东西。

我相信我已懂得那份对我的风格的“发自脏腑的抵抗”的本性——我方才唤起过。它的缘由,依我看,独立于那场后来才有的葬礼(那份抵抗,在其中终究扮演了一个重要角色)。这份抵抗,不是别的,正是对着一种“女性”\textbf{进法风格}——对着一门科学(在此即数学)——的(发自脏腑的)\textbf{反应}。这样的反应,在一个科学世界里,是常见的、“合乎事物本性的”——那个世界,与我们当今社会的任何其他局部小天地相比,不多让地被\textbf{阳刚价值}、及随之而来的感情、姿态、反应(尤其顾虑与拒斥的反应)所浸透。对那份特定的工作风格的抵抗反应——一份以“女性”为底色的创造性进法的化身——不过流自今日与最近几十年科学人的惯常制约——至少,流自我从来认识的那个科学世界。

一如任何出自制约的反应,这份反应实则毫无“理性”可言;而在它显形的那个人身上,有\dmark{}相当大的抵抗:仅只想到去检视它的意义,都有抵抗\origpage{706}。它被强烈地感到是\textbf{自身的辩白}——有点像道貌岸然的圈子里对“变态”的厌恶,或我们这里对“外来佬”的厌恶。可是,在我关切的个案里,我没有在这份反应本身之中感到一丝(有意识的或无意识的)对着我的敌意调门,而毋宁是一份\textbf{保留}、一份不利的成见——\textbf{单对着我的工作}。只是从那一时刻起——当我分明凭着我的风格(或不管我的风格,那都一样!)做出了此前人们不曾做出来的东西(而且事后也做不出真正别样的东西)——只是那时,这些保留才被收起,或许还带着几分不情愿……

无论如何,倘若在某些人身上,这些保留仍以默契、无意识的形式存留,那我太埋头于我的工作与我的任务,不曾知觉到它们。

老实说,依我看,至少很难想象:这样一份“发自脏腑的反应”,会如变戏法般消失,只因为某位先生证明了几样此前无人证出的定理。在有意识与无意识的接纳与拒斥的故意们做与散的那个层面上,两件事(“这样的工作方式不该被允许”与“某先生证明了某某定理”)真正毫无关联!

人们会说:既然如此,事情在我退出数学舞台之后起了变化,乃属正常——一旦我不再在场,去堵那些人的嘴:他们面对我的风格挑挑拣拣、撇嘴挑剔,却做不到用他们自己的风格做到同样的事。可这份“解释”瘸着腿走:因为它不顾那份嘲弄、那份压低嗓门的恶意的调门——它们从前不曾有。而且,凡我所知,没有任何东西能让我设想:1957至1970年间,我曾有时间把自己弄得叫我那些同行大会如此讨厌,以致一份宿怨或报复的动机会在我走后起作用。与我离开的那个世界的众多朋友,我保持过热忱的、有时深情的关系;而且(如我在别处所说)我不记得与任何一位数学家同事,在1970年之前有过一段敌意的关系。

然而,大会确实有过一份\textbf{后来的}对我的怨怼——一份集体“宿怨”的缘由,无论如何,一份集体的“报复”之举\dmark{}\origpage{707}:它虽保持默契,却“毫无破绽地有效”。这份“对一场异见的报复”的侧面,我在5月24日那条注“掘墓人——或全体教众”(n\textsuperscript{o}~97)里探测过。在那条注里,我把这些报复——对着我和那些不慎抬出我的人——的某个调门搁在了一边:恰恰是嘲弄的调门,它超出单纯的“不予受理”。而每当我感到那一“阵”,被点名的靶子,始终是某一份\textbf{风格}。换一种说法:使这份风格区别于一切其他的那份特质——它的“阴”或“女性”的本性——恰是天赐的情势:集体的无意识赶着抓住它,来洗刷一场异见所造成的冒犯,在以\textbf{排斥}施行的报复之外,再添上\textbf{嘲弄}的维度——一种嘲弄:它号称经由某一份风格,指认\textbf{无能}那份无可辩驳的征象。

而如今,随着“无能”这个词,某个不曾说出的东西终于被说出了:变得显然的是——这同一份“天赐的情势”,叠加到我“亡故”的情势之上,成了我的朋友兼昔日学生兼昔日继承人Pierre~Deligne的一次绝妙契机,使那份角色颠倒变得可触、可信、\textbf{被当真}——那个自感立于“\textbf{巨人}”跟前的“\textbf{侏儒}”那份疯狂而显然无望的欲望!(“骑在巨人的肩膀上”——借他简历里那句结语\fnote{245}{(*)}{关于这一点,见“神经中的神经——或侏儒与巨人”那条注(n\textsuperscript{o}~148)的最后一条脚注。})——从此,在众人眼里,当“巨人”的是他;而他将把那个“纯皮货”的巨人——千真万确!——指给全体大会嘲弄:一个吹牛大王、一个徒劳挥臂的“侏儒”——而那巨人毕竟曾(并且终究仍旧……)是“一份永恒而灼人的挑战,对着那个自感被无可救药的侏儒境况压垮的人……”

这场在“侏儒”与“巨人”角色分配中的惊人颠倒——在他与那个他者(那个被感到是一份\textbf{挑战}、必须不惜一切取而代之的人)之间——这场颠倒同时便是“\textbf{女性}”与“\textbf{男性}”角色的颠倒。分明地,正是作为\textbf{女性}的化身(臃肿、松弛、没有轮廓——从不曾被明白说出、却被热切否认),那个曾(并且终究仍旧……)是巨人的人,被交给\dmark{}人群(而首先交给变戏法的人本人……),当作可怜的侏儒、嘲弄的对象;也分明地,正是作为\textbf{男性}气概的英雄而典范的化身,那个曾是侏儒的人(而他,不管怎样、在他心底最深处“知道”自己是侏儒、且凭不可改变的境况仍是侏儒……)发现自己成了钢手的巨人,受着同一群赶来喝倒彩他者的人群的欢呼\origpage{708}。

这场颠倒,尽管是象征性的,分明与那个“私人”性质的“颠倒”毫无公度:后者凭一份考验过的战术(号称“天鹅绒的脚掌”)在一个狭小而无足轻重的“面对面”圈子里完成;一份可爱的小木马,他感到自己握着让“他者”运转、打转的线……侏儒让巨人转起来,行;可他终究还是、不可救药地还是侏儒!而侏儒的封神——他变成了巨人、还骑得更高,并把那个他骑在其上的人交给众人的嘲弄——这份封神,在大庭广众之下铺开:面对人山人海、欢天喜地的人群,他们来欢呼一位“侏儒”亡人的葬礼悼词,把它当作一场漂亮而令人陶醉的葬礼仪式分明的高潮。

\subsubsection*{(d) 否认(1)——或提醒}
\tocsubsub{(d) 否认(1)——或提醒}
\label{ch18:local:h66}

\notehead{152}{}
\phantomsection\label{ch18:local:n60}

(12月24日)随着昨天的反思,我有这样的印象:“拼装”葬礼图景的这个第一层面,已大体完成——至少,凭我手头现有的拼图“拼块”,我能做到的程度。当然,在葬礼反思的这第二部分(《收获与播种》的第三部分)里,我的用意已不再是收集事实材料(在“侦查”部分、I至X行列的进程中,我收集得尽够了),而是达到对葬礼\textbf{发条}的理解:经由众多当事人每一个身上那份秘密的动机(多半无意识)\fnote{246}{(*)}{(12月31日)这份“用意”,若按字面来取——鉴于“众多当事人”为数之众(哪怕只有十个!)——当然完全够不着。除朋友皮埃尔之外,我至多能形成一份总体的想法:在一个“集体无意识”里,多少圈定一些“动机”与“意图”——而那集体无意识,充其量不过近似地涵盖某个特定“当事人”的那些。}。这些动机,首要地,出自当事人与我这个人(作为“亡人”)关系的本性;或者更确切地说,出自我的所是之于他意味着什么——无论系于我从数学舞台的离开与否、系于环绕它的情况与否。

\dmark{}\origpage{709}这个“第一层面”,除我自己之外,便是那位在我葬礼上扮演“披祭披的神甫”、或“葬礼大司仪”角色的头号人物。他也是——在我出发之前的数学世界里曾为朋友或学生的人们之中——与我系得最紧的:凭一份格外强力的数学亲缘;也是唯一一个在我出发之后仍与我保持个人关系的人:那份关系延续至今。正是为这一切缘由,关于他,我掌握着一份材料:它的丰富,与葬礼与会者中其他任何人为我所知者,毫无公度可言。最后,在我认识的所有数学家之中\fnote{247}{(*)}{甚至在我认识的所有人之中——除两个例外。},他无疑也是远远最重的一个:那份他在自己生活中派给我的角色,压得最沉——分明远比通常派给师父的角色的分量更沉,哪怕那是人们倾注了全部身心的艺术(一如我自己曾倾注的那样)。这方面,我大约十年前便终于有所察觉:他派给我的那份角色,也漫溢到他的数学激情上(以及终于取而代之的东西上)。我身上这份知觉——多年来一直弥散——在葬礼反思的进程中,直到昨天,大大精确、充实起来。

依我看,随着昨天的反思,与那个以朋友皮埃尔和我的关系为中心的第一层面一道,终于就位、拼合的,还有“第三层面”:“(几乎)全体大会”——欢天喜地赶来,以它殷勤的首肯参与葬礼与下葬。如我昨天所写:5月24日那条注“掘墓人——或大会全体”的反思里析出的图景所缺的,是那份\textbf{嘲弄}的调门——它被放进对那个被当作亡人、“外人”、“圈外人”对待者的排斥之中。这份嘲弄的意义——它早在(11月10日)“阴的葬礼(阳葬阴(4))”那条注里便已清晰——昨天被重提并重新定位:这是对着被(在未成形的层面)感到为“女性”之物的嘲弄;而那份“女性”,经由(同样未曾成形的)一份对“女性”与“无能”的混同,成了“发自脏腑”的拒斥反应的对象:唯有男人——以他得胜的阳刚——才被认定是“力量”、创造力量的承载者。我也强调过:这类出自制约的出于本能的混同,全然不认常识与道理——当它所引发的观念与意象\dmark{}被带着那样强的信念与显然之力感到时,它们通常便被当作自身的辩白\origpage{710}。

然而有一个侧面,它随着“阴的葬礼”那条注的结语骤然以“闪光”显形,却还不曾被重拾。那条注里结束反思的几行如下:

“这不再是一个人的葬礼,也不是一部著作的葬礼,甚至不是一场不可容忍的异见的葬礼——而是‘数学中的阴柔’的葬礼——而或许更深地,在每一位为葬礼悼词鼓掌的众多与会者身上:是那被否认的、活在他自身之中的女人的葬礼。”

依我看,如今想来:这个侧面,在朋友皮埃尔本人的个案里,也曾或多或少被略而不提——尽管关于他,我第一手的事实并不缺!倘若这个侧面曾多少在场、或许被一位细心的读者感到,那必定是在字里行间:当时注意主要被“阴阳颠倒”侧面的种种角度吸住——(那个侧面,至少初看,似乎是我朋友在葬礼中的人格与特殊角色所专有的)。这份遗漏提醒我:日后(几天之内)我还须谈谈我朋友最后那次来访——10月10日至22日(10月21日那条注里提及,当时我许诺“几天之后”回头谈)。依我看,那将是最合宜的时刻,去检视“颠倒”的最后一个(?)角度——伴随着我朋友\textbf{这个人本身}之中原初阴阳平衡的颠倒。那又是一份\textbf{下葬}:下葬他身上某些原初的阴特征——在一份后起的、来占地盘的阳特征的鞭笞之下。我由此又以一道新的、更深的视景,站在那份惊人确认的跟前——它已不止一次强加于我\fnote{248}{(*)}{这份“确认”,头一次显形于“下葬”那条注(n\textsuperscript{o}~61)的反思里。}:以为自己下葬的是曾经的师父(而师父终究始终是一个朋友),可实际上,他亲手埋的,不是别人,正是\textbf{他自己}!

所以,倘若我再度回到“第三层面”或“底色层面”——那份“大会”别名“数学共同体”——方才引用的那几行便暗示:我在朋友皮埃尔的个案里感受得如此强烈的东西,或许对“每一位为葬礼悼词鼓掌的众多与会者”同样为真。依我看,正是这个侧面\dmark{}\origpage{711},还须稍加检视,我才会感到完全满意,才把葬礼图景的“底色层面”(连同前景)当作(暂时?)完成。

(12月25日)昨天,我借着平安夜,给自己放了个假,让自己痛快大干了一场:在注前一直待到凌晨三点过(下不为例!)。诚然,整个白天散给了别的任务;(重读完前一天的注)夜里只剩几个小时,倘若我想当天继续的话。一如常常:我终究甚至没能碰到我坐到白纸前时心里想着的任何东西!取而代之,我对葬礼“图景”的进度作了一点盘点,并突出了“第一层面”与“底色层面”中一个仍模糊的侧面:“\textbf{被否认的女人的下葬}”——她活在我的葬礼每一位与会者身上。

清楚的是:在这段引文里,“下葬”一词用作意象,指一份\textbf{否认}与\textbf{压制}的行为(或按通行的术语,“压抑”)。要谈否认、压制某样东西(在此,某样“活”在自己身上的东西),首先须确认:这份“东西”确确实实在场、“活着”(哪怕可怜地活着)。这里说的是每个存在者——无论女还是男——身上的“女人”:即其这个人身上由“女性”、“阴”品格的特征、品质、冲动或力量构成的那个“坡面”。奇怪的是,这个简单而要害的事实——在每个存在者身上,无论女还是男,同时活着“女人”\textbf{和}“男人”——至今普遍被忽视。我自己,也只在八年前——我四十七岁那年——才懂得它\fnote{249}{(*)}{关于这一点,见“接纳(阴的苏醒(2))”那条注(n\textsuperscript{o}~110)。}。

当然,“精神分析家”们早已“知道”它、谈论它。想必书很多,人人也多少听说过,一如我自己当年听说过。甚至,“人人”都乐于承认:这里头必有些真东西,只要说这话的是公认的内行、只要有为它写的书,等等。可是,“听说过”与“乐于承认……”,甚至读过一本、\dmark{}乃至十本\origpage{712}这方面的书,甚至(我斗胆说)自己写过一本、乃至几本——这本身并不意味着我们“知道”这件事;至少,并不意味着我们以一种比单纯死记现成套语更强、尤其更有用的方式知道它:诸如“弗洛伊德(或Jung,或老子……)说过……”。这类套语构成一份文化行囊、一张“有教养”者的名片——“懂”这“懂”那,甚至有时(配上文凭)是这或那的专家;就此,甚至可以承认它们有某种“效用”;确定的是:人人对自己这份在四面八方、在学校里、书里、“有趣的谈话”里攒起来的行囊,都宝贝得很,顶着逆风带着它,像一份亮闪闪而累赘的战利品,直到咽气。倘若我方才曾不敬地暗示:这份贵重的行囊“无用”,我的意思是:它对一件谁也不操心、甚至人人避之如瘟疫的事无用——即:自我的学习。或者换一种说法:这份行囊,对\textbf{承担自己的生活}无用——也就是说,对消化并同化自己亲身经历的实质无用,并由此,成熟、更新……

倘若我得以用几句话概括我这场关于阴阳的长反思的要害内容,那便是借着“提醒”这份“简单而要害的事实”——我方才恰恰提醒过。倘若有一位读者一路跟到这里,却还不曾以他自己的亲身经历感到这个事实——他身上有“女人”,哪怕他是男人;她身上有“男人”,哪怕她是女人——那么,他做了这番徒劳的努力来“跟着我”,却把时间花在给一份想必已然沉重的行囊,又压上一份新的分量:一份贴着“收获与播种”标签的分量。而倘若他是男人——即便他不属于那些葬礼与会者:在读我之前,他既不知情也不曾起疑——大可以打赌:他也在日复一日、在不知觉中,“下葬一个活在他自身之中的被否认的女人”(一如我自己当年也曾如此,且在我一生大部分时间里)。

一个男人“下葬”活在他身上的女人,一如一个女人“下葬”活在她身上的男人,有一千零一种方式\fnote{250}{(*)}{同样的话,也适用于“下葬活在他身上的男人”的男人、或“下葬活在她身上的女人”的女人:这类姿态,远不像人们想象的那么罕见。}——也就是说:去否认它、压制它。\dmark{}\origpage{713}“下葬”活在自己身上的东西,最常见的途径之一,是当那份东西在他人身上显形时,作出拒斥它的姿态或行为。这份拒斥,恰恰便是昨天在一个具体个案里我谈到的“发自脏腑的反应”。拒斥反应的发自脏腑的强力,其实\textbf{不是}(如我昨天似乎暗示的)因为被拒于他人的那份东西,顶撞了一套得到我们全盘而毫无保留之认同的“价值”总体。自知“强”的人,不会被“弱”的景象冒犯。反应的活力,恰恰相反,来自这样的事实:那份在他人身上被看见、“没有存在理由”的东西,\textbf{把我们自己牵连进去}。它像一份阴险的\textbf{提醒}——旋即被否认——提醒某件关乎我们的事:那件我们心底\textbf{知道}、却想对自己也对别人瞒住的事;一份提醒,由此便取一份无声而可畏的质问的调门。在这样的语境里,一份善意的宽容姿态——对着他人身上那份表面“毛病”——会显为一份危险的共谋自供,须不惜一切避开。而凭一份拒斥的姿态,相反,我们毫不含糊地与对方划清界限;总而言之,我们交出令人信服的保证(而首先,交给我们自己身上的内心检查官):我们自己白璧无瑕,我们始终合乎规矩、“道貌岸然”。拒斥反应,在是一份对某些\textbf{价值规范}的无条件的\textbf{效忠行为}(那些规范分判可敬与不可容忍)的同时,也是一份\textbf{象征性的下葬行为}:凭着它,我们身上那份“没有存在理由”的东西,被赶忙“归档”为一件“\textbf{不存在}”的东西。\textbf{至少不在我们身上}!

在这幅图景里,拒斥所取的形式——其变化无穷——依我看无关紧要。可以是那份被冒犯的拒斥,带着愤怒或厌恶的一切迹象;也可以是经由反讽、或“配量精巧”的轻蔑的拒斥。它可以以清晰不含糊的言辞表达,也可以仅仅被暗示:凭影射的或双关的话语,甚至无须言语,凭一个放在恰切之处的会意的微笑(或不笑……)。拒斥可以全然有意识,也可以只滞留在刚刚浮入眼帘的朦胧里,或躲进那份目光永不穿透的完全的阴影。

拒斥反应的强度,也变化无穷\dmark{}\origpage{714}:视那份“牵连”被感到是无足轻重,还是确实可畏。或许引发最强反应的“牵连”,是那些直接触及\textbf{性}的。这份极端的敏感,最近几代人里稍稍缓和了。可我确认:像爱欲冲动中所谓的“同性恋”与“自慰”(或说得好听些,“自恋”)的侧面这样普遍的东西,今天一如当年,仍引发强而有力的拒斥反应。至少,当你直面它的场合,不是一场关于罗马时代风俗、或深层心理学的“有趣谈话”,而是在你的日常生活里。甚至面对面,人们也很少谈起自己这个人身上这些性冲动侧面的显现(它们一般被感到是几处有点碍眼的“越轨”——至少这么说)。

在我这里关切的个案里,我出发脱离数学舞台之前所面对的拒斥反应,当然不曾有一份与我方才唤起的相当的力量。诚然,这份拒斥的对象——即“女性”的存在与做法,而人们本该“在男人之间”——确实有一份“性”的意味:取这个词更宽的意义,而不只是围着“屁股”及其余打转的言行的意味。我不怀疑:这份意味通常被感到——在无意识的层面\fnote{251}{(*)}{关于这一点,尤其见“阴的葬礼(阳葬阴(4))”那条注(n\textsuperscript{o}~124)。}。可它足够含蓄、足够间接,排除了任何粗暴的反应:超出对我作为数学家之“严肃”、“扎实”的一份单纯“保留”的反应。还要加上:我的“毛病”所处的领域——一份纯粹智性活动的领域——有助于给它一份相对平淡的外表:远(你上这儿来找什么呢……)远离男人-女人扭着腰肢跳肚皮舞那份令人不安、有伤风化的联想!但这并不妨碍:自我最初接触数学界(1948年)之后,还须将近十年,我的风格所引发的那些保留——在一个善好的小天地之内——才终于消失:至少,从我的视野里消失。可我出发之后,情形又变了:那份对我怀着善意、友情与尊重的氛围,骤然变了(在其后的六年里,我自己还不曾察觉)——被那个小天地所感到的一份“异见”、一份否认所变。

\asterisms

\dmark{}\origpage{715}老实说,我不敢肯定:这份氛围之变,是否真如我方才所说那般“骤然”。或者更好地说:我确认,我手头几乎没有事实,能使我对这份变化——1970年我出发之后如何发生、而1976年骤然(这一回,该用这个词了)摆到我面前——\textbf{如何}发生,形成任何想法\fnote{252}{(*)}{我提醒:那是借着我为使Yves~Ladegaillerie的论文得以发表所作的徒劳努力。这段情节,在“谁也拦不住进步!”与“棺材2——或锯成段的切割”两条注(n\textsuperscript{o}~50、94)里谈到。}。诚然,这整段时间里,我与那个我离开了的世界几乎毫无接触:本可让我感到某种“温度”及其演变的接触。对我清楚的是:在这份演变里,所有那些曾是我学生的人的小组、及其无可争议的头领Pierre~Deligne的姿态,扮演了一个决定性的角色。葬礼得以举行、引发葬礼的氛围得以确立,唯有凭一份“\textbf{全体一致的协约}”\fnote{253}{(**)}{关于“全体一致的协约”这份确认在反思中的头一次出现,见同名的(大写!)那条注(n\textsuperscript{o}~74)。}、一份毫无破绽的协约:它同时涵盖这场葬礼的“三个层面”:“继承人”(别名葬礼大司仪)、“共同继承人”或“近亲”的小组——由其他十一位“昔日的学生”构成——以及最后,“大会”(或许终究不是“全体”——日后还须回头谈……)。这份完美的协约如何设立、确立,对我仍属未知,或许将永远未知。眼下,我不觉得有兴致去探测它;我怀疑会有别人替我去探测(恰恰相反!)。

这提醒我:写上一条注“天赐的情势——或封神”时,有个问题擦过我心头:\textbf{二者之中},究竟是“大会”还是“披祭披的神甫”,代表了葬礼中运作的\textbf{主导}力量,而另一个多少算是它的“工具”\fnote{254}{(***)}{我提醒:在5月的反思里,在“掘墓人——或全体教众”那条注里,我曾意识到:我的朋友曾是一份毫无破绽的\textbf{集体意志}的“\textbf{工具}”。接下来的几行并不真正驳倒这份直觉,而是补充它:让“大会”与“披祭披的神甫”之间关系里的某种对称性保持开放。}。我当时没有驻足,甚至不确定这问题有没有意义——它分明像那道著名的先有鸡还是先有蛋!确定的是:两者之中任何一个——“神甫”也好、“大会”也好——都无法缺对方的助力来操办葬礼\dmark{}\origpage{716}。

可是另一个问题,依我看意义更清楚:两者之中,谁在这份工程里投入得更强。诚然,“大会”不是一个人,谈“它的”投入不很恰当。可同样真确的是:对我而言,这个人格化的实体,取得具体的形貌——凭十来个或二十个我熟识的\textbf{人}:与其中每一位,在一二十年间甚至更久,我保持着持续的、友好的关系。所以,当我谈大会的“投入”时,我具体想到的,是那些旧日朋友之中参与我葬礼的所有人的投入之“总和”。如此限定之后,依我看,这个问题便不再有什么修辞的成分。

对这个问题来到我面前的回答——不带一丝犹豫或怀疑的调门——是:在“继承人”的投入与大会的投入之间,\textbf{毫无公度可言};一如在一场普通的葬礼里毫无公度——何况遗产在继承人眼里越大(而大会里没有任何人能从中为自己捞到什么)、把他系于亡人的纽带(吸引的或冲突的)越强、在他的生活里扮演的角色越要害,便越甚。倘若这样的情势里有疑问,那多半只能来自亡故者近亲之中“共同继承人”的在场。(所以,这里说的是“第二层面”,而不是由大会的主体构成的“背景”。)在我关切的个案里,这些“近亲”与共同继承人之中,唯一一个其在我葬礼中所占的份额、可与主继承人Pierre~Deligne所占的相比者,依我看是Jean-Louis~Verdier:他扮演葬礼的\textbf{副司仪}。这份称呼并非凭空:在葬礼的进程中,我不止一次确确实实看到两位一同主持,配合得天衣无缝!可是,一如我在别处已写过的:撇开Jean-Louis~Verdier的某些公开行为不论,自我与他失去联系以来,关于他我知道得很少:太少,无疑,不足以对那份他与我的关系的来龙去脉、或他与那位显赫的“保护人”兼朋友的关系,形成一份多少周详的想法。

\subsubsection*{注~153}

\notehead{153}{}
\phantomsection\label{ch18:local:n61}

(12月26日)昨天的反思里,我试着精确化那份11月10日“闪光”般显出的直觉:在我葬礼的“每一位众多与会者”身上,我的葬礼代表着“活在他自身之中的被否认的女人”的象征性下葬\dmark{}\origpage{717}。当我在这里一次又一次说与会者的“每一位”时,这是一个稍嫌脱口而出的说法,或许最好不必完全按字面来取。至少,我深信:这份直觉,确确实实对每一位这样的人生效(而这样的人肯定很多):他们身上多多少少发生着那份对我特殊数学风格的“发自脏腑的拒斥反应”——那份反应,正是三天来我注意的中心。

另一方面,清楚的是:这样一份反应\textbf{不在}我的朋友皮埃尔身上,或至少,在我出发之前的五年里,那里非但没有它的痕迹,恰恰相反。正是我的数学进法风格与他自己的风格那份深刻的\textbf{亲缘},才造成了这些年里那份完美的交流;它也是我们在数学层面那份罕见的亲缘之缘由——那份亲缘,他与许多别人想必都曾感到,一如我自己曾感到。也正是这份亲缘,肯定是我这个数学家的人格、我的著作对他施加的那份\textbf{着迷}之缘由:不仅在那些年头(那时它“以正面”表达自己),也在其后的年月直至今天(那时它主要“以负面”表达,却同样雄辩\fnote{255}{(*)}{或者毋宁说:这份着迷,起初想必是“以正面”行的那份力(即\textbf{认同}于那个被感到相似的人之力)——在建立那份与我的暧昧而冲突的认同关系的那两股力之中。})。我不怀疑:倘若在那些最初的年头里,他对我把握数学事物的工作与进法,有过哪怕最轻微的保留、最轻微的不适,我不会觉察不到。

诚然,从那些年头起,我的朋友便竭尽所能,对外抹去那份在我身边为他所备的角色——哪怕只作为那个教过他、传给他一些有分量的东西的人,作为他从中取得工作所需的重要观念的人——更不必说,也抹去那份亲缘、乃至着迷的关系。我出发之后,便有对我这个人之否认的逐步加码:不仅凭沉默,也凭一份对我工作风格、对我引入的许多观念与想法的做作的轻蔑。我所知道的这类做作的\textbf{头一}份痕迹,定在1977年——借着“SGA~4$\frac{1}{2}$工程”\fnote{256}{(*)}{关于这一点,尤其见“两个转折”与“白板”两条注(n\textsuperscript{o}~66、67)。}\origpage{718}。

我没有试着亦步亦趋地追随这份加码的进程,也不觉得自己有兴致去追(一如昨天就一个切近的问题所说的)。

这份对一份与他自己近亲的进法风格、对一份他自己的著作由之而出的著作的否认,分明近于一份\textbf{对自己的否认}。方才,想着这份对我风格与著作的否认时(而我主要还沉浸在我1970年出发之前五年贴身数学接触的印象里),我本倾向于把它看轻,只给它一份某种意义上“战术”的意义:作为一份格外诱人的\textbf{手段},用来取而代之、满足对抗的冲动,借着抓住一份“天赐的情势”的运气。这确实便是三天前那条注“天赐的情势——或封神”(第~n\textsuperscript{o}~151)的调门。而我刚刚记起的这件事——我出发之前的年月里,没有对他自己风格或对我的拒斥姿态的痕迹——也正朝着这个方向,而不朝着昨天检视的情势的方向:即一份对“活在自己身上的女人”的否认(哪怕部分地经由对数学的某种进法)——一份\textbf{先于}葬礼之操办而有的否认。

但这并不妨碍:谁若拣了这样的手段,不管愿不愿意,\textbf{就得付它的价}。这份对某一份风格的“做作的轻蔑”,要奏效,便不但得对着别人演,也得、而且尤其\textbf{对着他自己}演。可是,人们无法当着别人、也当着自己,否认一份同时也是他自己骨子里的“风格”,\textbf{却照常实行着它},仿佛什么也没发生。这份对着别人的“战术性否认”,凭事物的逻辑,要经过一份对自己的否认、一份对自己一部分的\textbf{压制}——在此,即对他那份凭创造力量的原初本性而有的数学进法风格的压制。

这份确认,在此不是作为对一份事实的直接知觉之效而来到。它是一段简短反思的终点:动用已知的事实,从中得出常识的“结论”。我学会了提防这类结论(尤其在数学之外!),并只信那些事后由其他事实印证的\dmark{}\origpage{719}。可我在此恰好便利地记起:我曾被引着——以我对Deligne著作所知的措辞——确认:在他的著作里,找不到痕迹:他身上某些癖向(“阴”性的)——它们在我出发之前的年头里分明可见,我也在自己身上认得——在他的“之后”的著作里消失了。关于这一点,我在一个月前(11月26日与28日)的两条注“阴,作为仆人——与新主人”与“阴,作为仆人(2)——与慷慨”\fnote{257}{(*)}{即第~n\textsuperscript{o}~135、136两条注。还须加上第二条引注的子注(n\textsuperscript{o}~136$_{1}$)。}里表白得相当周详。这些事里最重要的,或许是一份谦逊:它让人看见(并描述,不怕显得傻气)那些极简单、极平淡的东西——此前无人肯对它们稍加注意。我自己献给数学的最好的东西\fnote{258}{(**)}{关于这一点,见前一条脚注所引的子注~n\textsuperscript{o}~136$_{1}$。},恰恰正是这一水质的。倘若我曾否认我天性的这份癖向——它可不是人人都喜欢的……——那么,我的著作、以及我那位最辉煌学生的著作的精要,都不会被写出来。这份癖向(或“倾向”)与另一份系得紧密,不然它的效应会极其有限。那又是一份谦逊、与“服务”的姿态:当要做的是认识并以精微的方式、面面俱到地描述那份被众人轻蔑的新东西时,不觉得自己的时间太贵,不肯为它花上十页,如果有必要(而不是满足于两行:东西在这儿——随你们怎么办!),甚至一万页;为它花上一整天(对一个并不缺别的苍蝇可拍的人而言……),甚至一生,如果有必要。

当我说到有待发现的“新的世界”时——调门或许有点倨傲——我说的不是别的,正是\textbf{这个}:看见并接收那份显得微末的东西,怀上它、滋养它九个月或九年,需要多久便多久,必要时在孤独中——好看见一件强健而活的东西发育、绽放:它本身便为生育与受孕而造。

倘若这份癖向——人们可以称它“母性”的——今日成了嘲弄的对象,那是“为了”那些被感到“阳刚”的姿态的利益:它们只容许一种可能的数学进法:“肌肉”的进法,把“肚肠”排除在外。“真正的数学”,又号称“硬数学”(或“硬邦邦的数学”),与那些(不怎么可口的)“软数学”(或“绵软的数学”——不说“瘫软的”,呕!)相对——\dmark{}那是十页五十页排得密密麻麻的证明,证明那些竞赛型定理(困难尽人皆知,不然便不算游戏!),万弹齐发——把一切“众所周知”的理论与观念、把左右可得的全部事实,统统打出去。\origpage{720}至于“木柴”,只消它在那儿;它在那儿就是为了这个!而那些耐心开荒的人——那些播种、栽种、施肥、剪枝的人,贯穿季季年年,好叫这些挺拔的大树之林生长铺展——它们如今待在那里如此合宜(那里曾是密不透风的莽丛),以致人们会以为它们自创世以来便在(大约充当背景布景,充当“万能木柴”的储备……)——这些人,只配下些大河文章(甚至大河书、大河书的大河丛书,倘若找得到疯到肯印的出版商),还外加读不通——他们是“软数学”的落伍者,不说“瘫软的”——不过,人再阳刚,礼貌总还是要讲的……

随着这段漂亮的咏叹,我忽然觉得回到了这场关于阴阳的长反思的出发点——10月初头一条注“肌肉与肚肠(阳葬阴(1))”(n\textsuperscript{o}~106)。确确实实是同一场下葬:踏着阅兵的步伐、吹着军号,下葬那份“女性”的东西——由铁臂,别名钢脑,别名超人的男性轻蔑所掩埋。这场下葬当然不只在数学的小天地里;它的分量超出一切个案——尽管个案恰好能多少贴身地让人嗅一嗅它的气味。而这份气味,确确实实是葬礼教给我的主要功课之一:在其中,我扮着未老先亡的亡人。

当我把注意的场子收得更紧,专注我朋友皮埃尔所扮演的特殊角色时,我在葬礼里看到又一层意义。我在其中辨出的,又一次是一场\textbf{颠倒}。一如我昨天预告过——没想到这么快便回头——这里,不再是一份系于他人的\textbf{关系}之中(真实的或虚拟的)的颠倒,而是一场发生在他\textbf{这个人本身}之中的颠倒。它不被为其自身的功德所求(或许作为一份“疯狂的欲望”的对象……),也不再限于纯然的象征(一如在一场漂亮的变戏法的收场处:那个自感“侏儒”的人,并不因此少觉得自己是侏儒,仿佛他并不曾刚说服自己已变成“巨人”……)。这是一场颠倒——我不说不可逆——但至少全然\textbf{实在}。它出发自一份和谐的平衡状态:“女性”与“男性”的创造冲动\dmark{}\origpage{721}的平衡,带一份女性的主导音符。它抵达一份战争与压制的状态:种种\textbf{姿态}与\textbf{架势}(利己的,一如一切姿态与架势)——打着“阳刚”的旗号——执拗地压制着那份\textbf{创造力量}:它被嘲弄、“象征性地”下葬,化作一具怪诞而瘫软的模拟像,长着“超级雌性”的面孔。

用不那么有层次的、但更形象、或许更触目的措辞:一个“\textbf{女性}”的存在——精微而强健、柔韧、\textbf{活着}——经一场永不停歇的变戏法,变形为一个“\textbf{阳刚}”的存在:不可摧毁、僵硬、而且\textbf{死}了。

\subsubsection*{(f) 布景——或“第二天性”}
\tocsubsub{(f) 布景——或“第二天性”}
\label{ch18:local:h67}

\notehead{154}{}
\phantomsection\label{ch18:local:n62}

(1985年1月1日)五天过去了,被种种营生占去。年尾正是写积压了几周几月的信的时机,更不必说几张贺年卡——回那些圣诞节前后收到的。还须做的是:堆肥——两个多月前收进的粪肥,加上园子里与开荒下来的植物废料、或从市政垃圾场运回的,要垒起几座肥堆,好叫开春时园子里有现成的好土。因地有坡度,还得再垒一级台地,靠着已有的那一座——那一座是给生活垃圾“日复一日”堆肥用的。这一切之下,我几乎没找到时间在注上工作,除了后勤。我极其细心地重读——此改彼修——了“主人们与仆人”部分以来(即自11月24日那条注“颠倒(3)——或阴葬阳”(n\textsuperscript{o}~133)起)的整场反思,补上为最近十五天的注预备好的脚注。主要是要有份可送去打字的手稿;但撇开实际问题不论,这次重读是有用的:重新取得近四五周反思的总览。一如在一场长程的数学反思中:当反思的“时刻”日复一日地被一份高度聚焦的注意的强光柱照着,反思的“线”、及其在数周乃至数月里走过的曲折轨迹,便倾向于在路上失落、没入半明半暗的混沌之中。我说不清:这在一切长程\dmark{}的研究工作\origpage{722}里是不是普遍现象,或者它系于我生活中那份“过去之下葬”的系统性机制——我已有机会暗示过它\fnote{259}{(*)}{这份机制,是在我童年发生的那场“翻转”的时刻启动的——我把它定在1936年夏(我九岁那年)。这一“我”结构化中的关键情节,在“超级父亲(阳葬阴(2))”那条注(n\textsuperscript{o}~108)里、及其子注~n\textsuperscript{o}~108$^1$里被暗示。}。无论如何,在一场长反思的日日夜夜、乃至数月之中,我身上有与它先前诸阶段的失联,表现为工作中一份渐增的不适。这份不适,终于由一场对刚完成的工作或多或少的深入回顾所解决:由它,那份逐渐松弛的联系重新接上。我观察到:这些“回顾性停留”在我的工作中扮演一个重要角色。每一次,我都重新启航,帆上新风,卸下了那份“不适”——它曾向我标示:对我进行中的工作在时间中的\textbf{连续}之整体知觉的一份渐进失落。在我的数学工作里,这样的回头并不罕见——甚至不妨说是通则:它引我把已做的工作彻底重想,并在一份新的视景里,既看已做的工作,也看将做的工作\fnote{260}{(**)}{关于“偶发的回顾”在长程工作中的角色,类似的反思,亦见“回顾(1)——或一幅画的三联”那条注(n\textsuperscript{o}~127)的第二部分,尤其其中所系的脚注。}。

但无论是一份数学工作,还是一场对我生活的禅观,我所谈的那份“不适”,始终是一份仍不完美的领会的征象:不仅(理所当然)是对尚待做的工作的领会,也同样是对已逝工作中所做之事的领会。这份不完美,实际上绝不归结为对反思各阶段及其时序的记性失灵(在一份数学反思里,这些是相对次要的方面:那里注意的对象是一份数学情势,它本身与研究者的心理特质、与这场检视的波折无缘)。它依我看,更是一份\textbf{太一}之缺的征象、一份不够充分的整合:把反思诸阶段之果实的那一整批部分领会,整合起来。这些部分领会,同样不完美、甚至属于假设,只要它们不曾被整合进一幅总景——在那幅总景里,它们相互照亮\dmark{}\origpage{723}。再用一次\textbf{拼图}的意象:对一份未知实质的探查,类似于拼一幅拼图的工作——但拼块不是预先给定的,须在工作进程中一一发现。更有甚者:每一块出土的拼块,起初只以一份含混而近似的形式显形,甚至相对于那个尚且未知的“正确”形状而言,被粗粗地歪曲着。反思的“局部”工作,在于把拼块一一辨出,并试着多少猜出每一块的轮廓:主要凭对所检拼块内部融贯的掂量、或凭对它与其他被预感到相邻的拼块的融贯掂量。可是,这些拼块的每一块,唯有当它们被拼进那幅它们由之而出的、尚且未知的总图,才显出其真正的本性、其精确而最终的形式。我所谈的那份“不适”,便是向我发出信号的:面对一堆被完美定位、却以一份或多或少的无形堆积着的拼块,该是把它们终于拼起来的时候了——或者,倘若(或多或少的局部)拼装已有,那是它仍太过零碎、或它歪着,须整个重来。要找到\textbf{那个}对的拼装,我撞见拼块的时序,无疑常常是次要的。可是,把拼块一一重新拿在手里(顺手便按那个顺序),怀着那份知道它们必定拼合、并等着它们各自落位的心境——无疑是工作中不可或缺的一步:为的是终于看见它们果真拼合。

上一条注(六天前)的“结语”,试着以言语圈定我心中一份强烈的印象:一场\textbf{变形}——它或许在朋友皮埃尔身上、在我离开数学舞台之后的十五年里,岁月流转,完成了。这些年里,我曾在此地彼处知觉过它的零散征象,它们有时叫我目瞪口呆;却(就我记忆所及)从不曾有任何时刻停下来,去形成一份对发生之事的\textbf{总}印象。得说:尽管我嗅出某阵“风”、以及我朋友在其中扮演的一份特殊角色(尤其借Motif的下葬——我对此只是朦朦胧胧地有所察觉\fnote{261}{(*)}{(2月30日)关于这份感觉——它长期保持着未成形、弥散的状态(直到去年4月19日发现“葬礼的全部辉煌”为止)——我尤其指出《收获与播种》第一部分里若干偶然的暗示(写于去年2、3月):关于\textbf{Motif}观念的命运,尤其见《导言》4(“一场追寻显然之物的旅程”)与“做梦人”一节(第~n\textsuperscript{o}~6节)。这份感觉的表述,在第一部分末节“一段过去的重量”(第~n\textsuperscript{o}~30节)的最后几页里大大清晰起来,自那个段落“我几乎可以把那封‘致……的信’(读作:Daniel~Quillen)视为……”起——它构成反思中的骤然转折。由这一天的反思最后阶段催生的头几条“注”——尤其月末写的双注“我的孤儿”与“对一份遗产的拒绝——或一个矛盾的价格”(n\textsuperscript{o}~50、51)——在相当程度上,给先前以弥散状态感到的东西“做了个盘点”:关于我的数学著作之待遇、关于对着它与我这个人吹来的那股“时风”。关于这份“弥散之感”就Motif所取的特定形态的一份描述,见“墓”(n\textsuperscript{o}~71)与紧随其后的“一脚踏进旋转木马”(n\textsuperscript{o}~72)两条注。}),对那场盛大的下葬——我的朋友\dmark{}\origpage{724}正手法老到地编排着的、对我这个人、对我全部著作的下葬——我远谈不上起疑。正是这场下葬在过去一年里的逐步发现,终于成了一记够强的\textbf{冲击}:它推动了我身上的一份惰性,并激励我终于“落定”在一份仿佛没入遥远过去之雾的情势上。所以,10月我朋友的来访,我也是以一份与当年会面时那份有点儿“例行公事”的心态全然不同的心态接待的:一份愕然惊诧的注意。正是在那次来访中,出现了这份印象、或毋宁说这份骤然知觉——它想必早已在场甚久,而我一直乐于无视:那份“变形”的知觉——恰恰是我在上一条注里经由另一条路径重新撞上的那个。倘若我又一次找回这份印象——这一次,经由对我朋友数学著作之所知——那肯定不是凭天大的巧合,而是被近两个月来与他这个人的直接接触所教给我的东西所引导。这份变形印象的确凿之力——它归结于一个“阳刚”的、不可摧毁的、僵硬的、而且\textbf{死}了的存在——当然不会作为一场比较、拼合事实(或其他性质的局部印象)的反思的终点而来,而只能凭一份直接的亲身经历而来:它一直未被说出。而这份经历,就在此刻,依然未说\fnote{262}{(*)}{(1985年2月30日)此刻它依然未说,尽管我终于写出了我朋友来访的记述——在“完成的义务——或真相的瞬间”那条注(n\textsuperscript{o}~163)里。}。

在上一条注里,我写道:这场“颠倒”(在我朋友这个人本身之中)、或这场“变形”(借用出现在“结语”里的说法\dmark{}\origpage{725}),并非“为其自身的功德所求”——还加了括注:“(或许作为一份‘疯狂的欲望’的对象……)”(即那份颠倒的欲望:“神经中的神经——或侏儒与巨人”那条注里谈到的)。可是,次日重读反思诸注时,我不再那么有把握了——也拿不准:我那份\textbf{把}我在葬礼中辨出的这两场“颠倒”\textbf{对立起来}的故意,是否真有根据。毕竟,在侏儒与巨人的意象里,“巨人”体现(如我不止一次强调的)着“阳刚”的价值,“侏儒”则被“阴柔”的非价值压垮。而即便这幅意象\textbf{坐落}在我朋友这个人之外——它被贴在他的关系上,对着另一个人(在此即我)——它在我朋友这个人之外,却毫无“客观”的存在:它反倒是——那场正发生在他\textbf{自己}身上的冲突实在——贴着外面的投影(贴在他的与某人的关系上)。换一种说法:侏儒与巨人的意象,显为那场真实\textbf{冲突}的象征性\textbf{布景}——那场冲突发生于比意象所居更深的层;而它不多不少,正是他这人身上“阴”“阳”两个“坡面”之间那场永世的\textbf{冲突}。

一份内在冲突的这样一份\textbf{外化}——它必须严格保持隐蔽——正属于无意识用来“清场”的少数尽人皆用的手法之列:尽可能“排去”那场原初的真实冲突,用一个显得更“可接受”、或至少不那么扰人的冲突取而代之。在此,所选的那幅避雷针意象本身仍然无意识(至少我推断如此);甚至,我倾向于相信:它被圈定在无意识中相当深、却比真实冲突之认识更接近浮面的层。(后者不多不少,正是“两面认识”的“正面”——“两种认识——或对认识的恐惧”那条注(n\textsuperscript{o}~144)里谈到的。)

这便暗示:上一条注括注里重提的那份“疯狂的欲望”——“\textbf{做那个巨人本身},或至少,\textbf{冒充是他}”——那份欲望,不多不少,是关于一场他自己身上“变形”之欲望的“外化的”\textbf{转写},就着侏儒与巨人的避雷针意象;一场变形——即使不真实,至少表面——在其中:他存在中被感到不可容忍的那份主导——那份“阴”调门的主导(被感到“软”、可鄙)——将被“颠倒”、变形为“阳”或“阳刚”调门的主导(被感到“英雄”、唯一值得艳羡)\dmark{}\origpage{726}。这两份欲望,远非在内在本质上有哪怕些微的对立,如今在我眼中反而是不可分割的:一份是另一份的影子、其\textbf{象征而可触的表达}。至于那场“变形”——我朋友来访时终于知觉到的(迟做总比不做好!)——它此刻显为那份“疯狂的”而强制性的欲望之\textbf{实现}或兑现;兑现,不是凭一份天赐的恩典的干预,而是作为“东家”那份执拗意志的长效效应:他要“校正瞄准”,照一些借来的特征把自己\textbf{重塑},并把这些特征强加给匠人-孩子(自然,此类“东家”式的操作,从不咨询匠人)。

上一条注里,我强调过这场“颠倒”(或这场“变形”)的\textbf{实在}品格。此刻我更清楚地辨出这份“实在”的本性与限度。它是一个\textbf{架势}的实在:竭力照一个被感到为理想目标的模型自我塑形。对模型的选择——即所采取的架势类型——无疑远早于我们的相遇。但依我看,在我们相遇时及其后的岁月里,投入这份架势的能量与耗散仍微乎其微。我朋友身上的这份投入之规模的骤然、严苛的变化,我相信,是由我的离去所创造的“非同寻常的”契机:先是从我的机构离去(我的朋友一夜之间必定在自己眼中显为偷偷地\textbf{取代了他的“对手”}),旋即又是我从数学舞台离去。第二个更重要的实在方面是:凭一份无度的投入,这份架势终于确确实实变成“\textbf{第二天性}”。正是在我们最近的会面中,我知觉到了这份“第二天性”。它压着一份巨大的惰性——一如我自己的情形也曾如此。这并不妨碍:在我这里,一场更新遂曾发生;而我身上发生了,并不削减我朋友身上那股反对他自身更新的惰性。

在他身上渐渐确立的这份“新”实在,并没有“解决”他身上的冲突——正如一个邻国占领一个国家,并不“解决”冲突。毋宁说,我朋友身上的冲突被“冻住”在某个“力量对比”之中;而很可能,它将如此直至生命的尽头\dmark{}\origpage{727}。无疑可以说:自我的结构——即行为机制——确确实实改变了,有时以触目的方式。然而,这类由“东家”的意志强加的变故,不改变原初的本性——匠人-孩子的创造力量的本性。它们只类似那些强加给匠人的刑具:匠人得想方设法,照样干活——在“东家”多疑的目光之下;当东家不干脆把工具从他手里夺走、好教匠人知道该怎么做的时候!

这并不妨碍:企业照转、照赚,而东家,大体上,满意。氛围确实够糟;可像大多数东家一样,他皮厚,不为所动,只要回报还好就行。

\subsubsection*{(g) 另一个自己——或认同与冲突}
\tocsubsub{(g) 另一个自己——或认同与冲突}
\label{ch18:local:h68}

\notehead{155}{}
\phantomsection\label{ch18:local:n63}

(1月2日)自12月24日那条注“否认(1)——或提醒”(n\textsuperscript{o}~152)以来,已有一周多了:我一直觉得,葬礼图景的前景层面,我差不多做完了。结果不然——已经连续三次,我得回到这个那个看着不太清楚的点上:大概只差三句话,好给最后一个ｉ点上最后一点。可每一次,这个“最后一点”都占了我整整一晚的工夫,而原来那个看着“不太清楚”的东西,其实一直相当晦暗;回头给它找它自己的光照,绝非奢侈。我预感今天也不会例外:我打算回到那个(最后一个?)点——它曾在下一条注“否认(2)——或变形”(n\textsuperscript{o}~153)里被一笔带过。那是我扮演“过继父亲”角色的关系独具的侧面之一:我朋友对我的“暧昧”\textbf{认同}。这一侧面,在引注的一条脚注里以三四行唤起。当夜它不再被谈;可次日,重读前一天的注,我感到必须回头谈它。昨晚接续反思时,我本打算从这里接下去;可最终,是前一场反思悬着的另一个“最后一点”占住了我,直到深夜。

在《收获与播种》中我屡次被引着记录:在某某朋友或学生的关系里,扮演收养的父亲或被收养的父亲\dmark{}\origpage{728}(的侧面)——每一次,都是借着这份关系中冲突性特征的显形。于是,无任何故意,我注意与强调的中心,总是这份“父性”意味的关系中的\textbf{冲突的}侧面。我分明感到:在这样的关系里,总有一份或多或少的\textbf{对父亲的认同}成分;唯有一点保留:这份认同有时取“负面”的形式——凭与一份被否认的父亲影像的“负面”(或反面)认同\fnote{263}{(*)}{尤其,这曾是我三个儿子与我的关系的情形——他们,这一回,全然不是“被收养的”,更不是“收养者”……}。这份认识留在背景里,不以可见的方式介入反思,却带来了它的一份贡献:对这样那样关系的一份弥散的把握、一幅尚模糊的轮廓之形成。我记得只有一次——我以为——以一般措辞、顺着认同的方向,在“敌对的父亲(1)”一节(n\textsuperscript{o}~29)的结尾表白过:

“……这便是那同一个对父亲的冲突原型之重现:那位父亲,既受钦慕又遭惧怕、既被爱又被恨——那个须要与其交锋、战胜、取代、乃至羞辱的\textbf{男人}……却也是那个人们暗中想成为的人:想夺他的力、据为己有——一个\textbf{另一个自己},被畏惧、被仇恨、被逃离……”

无须多言:在这些行间——写于一场“对我数学家往昔的回顾”——倘若有一个具体的个案在动笔时指引着我的笔,那便是与我的“隐蔽继承人”、那位不说自己名字的昔日学生的关系:Pierre~Deligne——尽管在写的那一刻,我对那场由他亲手编排的盛大葬礼,在意识的层面至少,毫无察觉!此刻复现这几行——写于九个多月前——我已惊于:它们仿佛预示、并(可谓)召唤着侏儒与巨人的意象——那意象仿佛正是为给这幅直觉以可触的形式而成形、落实的。然而,对我毫无怀疑:这幅意象绝非在我这个编年史作者兼探索者身上成形,而分明在我朋友自己身上;它也正是从非他莫属的某人那里取得\fnote{264}{(**)}{关于这一点,见“神经中的神经——或侏儒与巨人”那条注(n\textsuperscript{o}~148)的最后一条脚注。}!

冲突性的认同,清楚地显现在“那个人们暗中想成为的人”\dmark{}\origpage{729}这几个字里;更强、毫无含糊:“一个\textbf{另一个自己}”。在侏儒与巨人的意象里——它12月18日来到我的笔端(在那条注“神经中的神经——或侏儒与巨人”,n\textsuperscript{o}~148)——谈到的是那份“疯狂的欲望:\textbf{做那个巨人本身},或至少,\textbf{冒充是他}”——这几行似乎正是对方才引用的“那个人们暗中想成为的人”的回答。可这一次,我在这里停住(一天的本分一天尽!):离“另一个自己”还差一步——后者九个月前作为天经地义的事信手而来!诚然,这一次,既然是一份“看物件”的工作、在一个十足具体的个案里,便须迥然不同地小心、审慎——不同于那种语境:在那里,人们(不动声色!)作态抛出一份一般品性的断言,不涉及任何特定的人……

但想想:诚然,对那个渴望\textbf{象征性}满足的无意识而言,这确确实实只差一小步——它大可以凭自己制造的意象付这部账——从那份“疯狂的欲望”(而且显然有一股相当大的力)想成为这个那个,到与那所望之物的\textbf{认同行为}。要使这份认同——无论多么无意识——多少可信,并使它所带来的一切满足可以被带着最少的安全感来赏味,还须:它有某些(在此)被认同之人身上的“客观”相像特征的保票。我推想,在我关切的个案里——我朋友与我的关系——第一份“客观特征”、足以助成一份相似之感与一份认同行为的,是渗透在我们共同情人——数学——的进法之间的那股强亲缘。这想必便是那份“以正面”行之的力、那份“\textbf{认同于那个被感到相似的人}”的力——今天反思开头引的脚注里顺带谈到的。

可是,一如我在关于朋友与我的关系反思中已屡次提示的:自这份关系的最初几年起,他不曾漏看我自童年起所扮角色中的“超级阳”失衡侧面——那早已成为我的“第二天性”。我无从说:在意识的知觉层面,我朋友是否清楚地分辨了我这个人身上这两个截然不同的侧面。(我倾向于怀疑。)无论如何,我企业里“东家”的超级阳侧面,想必在他身上唤起了\dmark{}\origpage{730}两类截然不同的反应。一类——也是直到最近几个月我唯一知觉到的、也是他身上唯一有意识(我推断)的——偶尔以一份略带忧伤的惋惜姿态表达:我曾唤起过的;那份姿态从未离开友好或温情的调门。另一类反应,细看之下,自身也是“暧昧的”:由两份方向表面相反的成分组成。一份,“正面的”,朝着对我这个人的毫无保留的\textbf{抬举}:作为“英雄”价值、“超出常人的”价值的化身;那些普遍公认的价值固然,人们在少年时代像呼吸空气一样吸收它们,可他童年的直接近旁大概没有给他提供任何多少令人鼓舞的“模型”。这份成分——一如前面谈过的那份(全然另一性质的)\textbf{亲缘之感}——朝着对我的\textbf{认同},不带任何对抗因素。而那对抗因素,则进入另一份成分,或更好地说,进入这份我方才描述其“正面”的认同的另一面(或“\textbf{反面}”),它对我始终更费解。想必正是在这里,那份“父性”角色——我朋友凭我符合某个号称体现这类价值的理想“轮廓”而派给我的角色——扮演一个关键角色。当我试着摸黑探测——凭我手中极少、极细碎的材料——这份对一个“过继父亲”(特征十足的“超级父亲”! )的认同之强烈对抗内容的深层缘由时,两周前,我撞上了一份说得通的、但仍是假设的“脚本”:在12月30日那条注“缓期中的宿怨——或事物的回转(2)”里。

不是在这里回头谈那份脚本。依我看更有趣的,是把“侏儒与巨人”的意象(它刚在前天那条注里出现),放回这份我朋友对我的冲突性认同的光照下来看。于是一下子显出:意象中的两位主角——侏儒与巨人——\textbf{都不过是他自己};或者毋宁说,\textbf{他自己身上两个截然不同的侧面}。“侏儒”代表被我朋友感到为\textbf{原初}、“不变”的存在侧面:扎根于他的童年,远及他的记忆、无疑更远……它也是被感到为他这个人平淡、无足轻重、乃至可鄙的侧面。它是\textbf{被否认的}侧面,并由此,也是被感到“无可救药”、“不堪重负”的侧面:他存在的\textbf{耻辱的}、可鄙的\textbf{一极}。“巨人”相反,代表那个令人眩晕的\textbf{理想}:人们绝望于永远达不到\dmark{}\origpage{731};至多,有望些许相像——不惜自欺欺人、也欺他人,凭一切可用的手段。其中一个手段,便是取代那个显为这份理想之可羡化身的人,凭一切想得到的方式“证明”自己优于对手。至于巨人本身,此刻显为与对手和父亲有别:他是\textbf{被拔高、被吹大的一极:自我的理想、英雄的一极}。“东家”的至高犒赏,便是这一切:足以滋养那份错觉——人们\textbf{就是}这理想的一极、那份贪得无厌地自我膨胀的精神的投影。可正是这份对犒赏的饥渴,泄露出几分不安——“一份深深埋藏的怀疑”——它告诉我们:当事人“在心底最深处,并不上这些重要性与‘价值’的虚假征象的当”……\fnote{265}{(*)}{引号内的引文,出自“他人之无谬与对己之轻蔑”一节(第~n\textsuperscript{o}~4节)。}

在心灵的较浮面层次,这些“虚假征象”\fnote{266}{(**)}{这些征象尽管“虚假”,却常常终于铸成一份固若金汤的“第二天性”——“不可摧毁”(借用“否认(2)——或变形”那条注(n\textsuperscript{o}~153)里结语的说法)!}却属于方才所谈的“(或多或少)客观特征”之列:它们号称使一份对理想模型的认同行为“可信”(无论那模型保持“巨人”的非个人形式——活在他自身中的无面孔者——还是取敌对之父、对手的熟悉面容)。

\subsubsection*{(h) 敌对的兄弟——或传递(2)}
\tocsubsub{(h) 敌对的兄弟——或传递(2)}
\label{ch18:local:h69}

\notehead{156}{}
\phantomsection\label{ch18:local:n64}

(1月3日)昨天下午,趁等朋友登门的个把钟头空档,我翻了翻C.G.~Jung的自传——一位朋友凑巧刚带来。读到的寥寥几页,便强烈地钩住了我。这是我头一次手里拿着一篇Jung的文本;至今,我对他只有一份极模糊的概念:一个弗洛伊德的叛逆弟子——(凭零零星星传回我的回声)他懂得把奥秘那明暗微妙的运动,重新引回主人的笔直小径里。到此为止,相差无几。此刻,我感到一个像你像我一样活生生的人;而且,他不浪费时间绕弯子;尤其:一个直奔真正问题的人——那些他凭自己的光感到要害的问题,并且不满足于(当偶发的疑问古老如世界时)饱学之士那现成的答案\dmark{}\origpage{732}。

“传记”一面(那是要发表的)当然特别令我有兴味:因为我正写着的这些注,颇有几分像一部传记——而且与Jung的精神十分相近:外部事件始终从属于内在历险,它既是内在历险的揭示者,又是偶发的刺激者。令我震惊的是:Jung直到83岁才写一部自传(或更确切说,才为一部传记贡献他的一份),而且,尤其:在他一生的任何更早时刻,他从不曾费心深入地检视他自己的童年。在我看来,对弗洛伊德的弟子们,这本该是理所当然:要熟悉无意识之路,头几件事——若非头一件事——便是探一探他们自己身上的路!我甚至毫不怀疑:一份号称的“对无意识的认识”,若限于大学课程里学到的东西(哪怕由弗洛伊德本人这样一位显赫的大师讲授)、外加对若干“临床病例”的分析,则仍是一份未经整合的知识、一份零碎的、“死的”知识——一份本身不提供、也不促进对自己、对他人、对世界之理解的知识。

但同样真切:对自己这个人的勘探,凭其本性,不能列入一份制度化的“课程”之中——正如(且说一个“病人”)一份被扰动的心理平衡,就其根本而言,不能凭任何一位“学家”的干预为果实——他不过是施展万金油式的技术而已。那份“被扰动的平衡”,绝不限于神经崩溃或神经官能症那个社会不可接受的阶段;它几乎人人身上都查得出(程度多半\textbf{更深}而非更浅)。心理学家们自己(或民族学家、社会学家、其他“学家”),无论何种派别,也并不比旁人更豁免!而一份对被扰动平衡的真正恢复,绝非一份介入第三方“医疗行为”的本性。它是\textbf{当事人自己}的行动、非他人——\textbf{一份爱的行动}——他既可做、也可不做。它不是心理机制那无情运转的结果(无论有没有那位心理机制专家的干预),而是取这个词完整意义的一个行动:一份\textbf{创造},一场\textbf{重生}。

在写完上面那句断然武断的话之前——关于那份“号称的\dmark{}\origpage{733}‘对无意识的认识’”——我便意识到:语境会使人显得多么狂妄。对Jung的作品一无所知(前面刚谈过他),我竟仿佛把他、连他的“号称的”无意识认识,一起毫不客气地打发掉——就因为他显然不曾(在83岁之前)费心探一探他自己那无意识生长的土壤。可我倒猜:读他的传记,将会显出:即便不曾致力于这样一场“勘探”,Jung想必另有\textbf{别}的接触他自己无意识的路径(那些路想必自己也长期无意识)——由此,那被指摘断言的先决条件,想必不适用于他。

有另一桩全然不同的事,翻看词汇表时叫我错愕。在“四元性”条目下(按:这是法文版),Jung坚持数字四的“总体化”品格。直到十年前,我对数字的哲学或“神秘”运用还十分抵触——任何这种方向的思辨或谈论,在我看来都是胡话、孩子气、“变戏法”(德语叫Hokuspokus,指廉价的魔术)。我后来学到的关于《易经》(或“变化之书”)的那一点,使我没那么武断了。昨天,我把归给数字四的“宇宙”品格,与写“阴阳之钥”时自发形成的那份分组联系了起来:一般四或八条注合成一“包”,归在一个共同标题下。第一个“包”只有一条注;诚然——可(我在写完第六组“阴阳数学”时曾带着满意记下:那组有七条注,而不是八条)把它并入其后一个显得它最自然归属的组,又凑成八条(7+1=8),便又是四的倍数。这个“样式”一直延续到如今:最近完成的组,是第十组“暴力——或游戏与刺棒”(156$_{1}$)。得说:自第七组(“阴阳的颠倒”)起,我便任这份不召自来显出的“样式”引导——却从不曾给它寻、给它设一个“意义”,除了一份形式上的“数学”规整之感,被感到和谐。

这让我记起我写过的唯一另一篇堪称“宇宙”主题的文字——依旧以生活中的、创造行为中的阴阳动力学为中心\fnote{267}{(*)}{即《乱伦颂》——在注~n\textsuperscript{o}~43里谈到过(那里引的是“非古鲁的古鲁——或三足马”一节,第~n\textsuperscript{o}~45节),而主要在“行为”那条注(n\textsuperscript{o}~113)里,原书pp.507--509。亦见“事物的动力学(阴阳的和谐)”那条注(n\textsuperscript{o}~111)的开头。}\dmark{}\origpage{734}。这篇文字,表面上无任何初衷、也肯定从未费力,便按一份严格的数字编排归了组。我忘了它是怎样:方才一看(好奇者恒好奇!),显出:七节“判词”,每节四“诗章”。于是,这又是一份按四分组。诚然,判词的数目是七,不是四的倍数——于是按Jung的标准,“总体性”的品格对整部作品并不成立\fnote{268}{(**)}{这部拟议中的作品(顶着那个挑衅性的名字“乱伦颂”)本应含三部分(“天真”、“冲突(或坠落)”、“解脱(或寻回的童年)”),唯有第一部分做完了。这里说的就是它。},而只对构成它的七节“判词”成立。可这里我又有得脱身:那部著名的“诗集”,还配有一个天赐的“尾声”(更不必说一章没完没了的“序言”——我庆幸放下了),于是又有7+1=8——得救了!

该回到昨天的反思,从我中断的地方续起。我试过以我朋友对我的认同,理解他身上的侏儒与巨人意象。显出:“侏儒”与“巨人”代表(或“\textbf{布景}”着——借用昨天前那条注的说法)我朋友这个人身上的两个“\textbf{极端}”(我是说:“东家”设立为“极端”甚或更准确说是“两极端”的东西):一个“耻辱而可鄙的极端”,另一个“理想的、英雄的极端”。老实说,借一份重音与光照的差别,我在这里与前一天对同一力量意象找到的解释相接了:前天那条注“布景——或‘第二天性’”(n\textsuperscript{o}~154)。那里谈的是对那份由东家、“我”在存在的“阴”“阳”两个坡面之间设立的冲突的“布景”。这份对原初冲突的表述——以两个“坡面”的措辞——会对应一份未被歪曲的冲突认识;而我深信:这份认识,确确实实存在于(绝非不可触及的)心灵深层。昨天那份以两个“极端”的措辞的表述,代表的是一份\textbf{被歪曲的}冲突视景——被东家的一份故意歪曲:他抬举其中一个“坡面”\dmark{}\origpage{735},把它做成一个理想的、英雄的“极端”,贬抑另一个,把它做成另一个与前者相对立的极端:一个耻辱的、可鄙的极端。我猜,这份居间的意象,居于较浅的层、居间的层,或许与那幅外化的意象、那幅更接近意识表面、并与表层部分交叠的“侏儒与巨人的布景”部分共居\fnote{269}{(*)}{这份关于侏儒与巨人意象的推想,当然,系于那幅意象在其结语里如此明白的表达——在Pierre~Deligne亲自所写的那份传记小记里(暗示已于“神经中的神经——或侏儒与巨人”那条注(n\textsuperscript{o}~148)的最后一条脚注)。}。最后,在这些表层里,我提醒:统治着“心软爸爸”那份田园诗般的意象——边上有点瘫软的:一个恭敬而体贴入微的儿子,天鹅绒显露无遗,爪子隐在天鹅绒之面……

相较于前天的反思,昨天的反思更像是在给它增添细微层次,并由此稍稍磨锐其轮廓——却还不曾带来什么本质的新东西。诚然,当因时辰太晚而止笔时,我毫无已经抵达我所投入方向的尽头之感——“暧昧认同”的方向。事后回想,我意识到:由于一份积习——把自己“看作阳”——对我而言似乎天经地义:当涉及对我这个人的认同时,它只可能关乎我的阳特征。在此,在侏儒与巨人的舞台意象里,到目前为止,我是在“\textbf{巨人}”里认出自己的——虽然以一份歪曲的形式,却仍清晰可辨。然而,凭借我朋友身上的“颠倒”综合征之效,我被执拗地呈为“\textbf{侏儒}”\fnote{270}{(**)}{这个“侏儒”本身,不是别的,正是“超级妞”的隐喻——一个“假巨人”的面孔,臃肿、瘫软的形态……(1985年2月)}——这份混同(带着明显恶意的意图)被我当即拒绝:凭一份本能反射,它普遍而强力:直面一份嘲弄的意志——它拿我身上(在此即阴的)十分真实的特征当靶子,同时把同样十分真实的互补特征(后者享有抬举的共识)压回沉默——这样的情势,在我身上唤起那永世的反应:若非否认被指摘的特征,至少在默契中把它们缩小,并抬出那些恰恰被不公正地隐匿的特征,仿佛要去\textbf{对抗}它们。

凭这份“发自脏腑的”反应,我确确实实进入了冲突的轮舞——恰如我所该做的!它向我标示那个永世的“\textbf{挂钩}”:它在哪儿可以拿住我,把我拽进轮舞\dmark{}\origpage{736}。我自己的实在视景,也随之被歪曲了——作为对一份挑衅性歪曲的回答。所以,我昨天从牙缝里(或从打字机的键上)写道:

“‘第一份客观特征’——足以助成一份相似之感与一份认同行为——是渗透在我们共同情人——数学——的进法之间的那股强亲缘”。

这是白写的:我写它时,竟乐意忘了:这份“强亲缘”,由一份\textbf{阴}的、“女性”的\emph{进法}构成——在对事物的发现与认识中;而这恰恰是我作为与他“相似者”显形的一面:正因它,我与他一样,也显为“侏儒”:那是他秘而不宣、易受伤害、羞耻的一侧——待相宜时机一到,他便拉它上场,取而代之、“颠倒”。这份“天赐的情势”\fnote{271}{(*)}{见同名那条注,第~n\textsuperscript{o}~151。},即我认识冲动中的阴占主导,不仅是一份可疑朋友手中的\textbf{武器}——它首先也是、并首要地是,他对我认同的一份“客观根基”;这一回,不是对\textbf{父亲}的认同,而是对一个\textbf{兄长}、乃至一位“\textbf{姐姐}”的认同。

当我在这里用“客观”一词,是说要表达:这一回,涉及一份“认同”——它扎根的,不是“东家”想(或怕……)成为这个那个的虚构,而是一份\textbf{实在}:深刻、可触、无可置疑——双方原初本性之间那份\textbf{亲缘}的实在。无论如何,这份亲缘绝不会逃过他、一如逃过我;而我不怀疑:在某个深层面,这份亲缘的\textbf{意义}同样被知觉。而且我至少猜测——虽无十足的把握——这份知觉确确实实充当了他对我认同的材料。于是,这份认同是在\textbf{两个层面}上完成的:一方面是“理想”的层面,我在其中成为某类\textbf{价值}的化身——那些他但愿自己是其典范化身(哪怕仅在表面上,而模型显得遥不可及、理应着实实现理想)的价值;另一方面是“真实”的层面:认同凭一份被正确知觉的\textbf{亲缘之实}而确立——但那是些被公认\textbf{要命}的、可怜的共同特征上的亲缘\dmark{}\origpage{737}\fnote{272}{(*)}{这两个“层面”,于是对应着我这个人身上\emph{两种}彼此对立的“原型”:父亲(别名“巨人”)的原型,与兄弟、乃至姐妹(别名“侏儒”)的原型。后者也出现在“心软爸爸”的意象里——那意象由血肉之躯的父亲“如其所是”(唉!)所暗示,而非“如其所应是”……}。

此刻该记起:在我们相遇的当口、乃至其后十多年间,我身上肆虐着这份对我“女性”特征的压制——与我最终在我朋友身上确认的那份相同。带着回顾,依我看:我们相遇时,我朋友身上的这份压制已有几分程度,但主要尚属潜伏;无论如何,比我身上的弱得多。如我不止一次强调的:我这个人长期打有超级阳失衡的印记,而他透出的则是一个和谐平衡的印象。自那时起,他身上与我身上,有过\textbf{反向的演进}:在我朋友那边,从一个阴阳平衡的状态走向强的阳失衡;在我这边,从强的阳失衡走向一个(相对)的阴阳平衡。

随即浮现的念头是:我的朋友——或许凭这份对我的双重认同——(带着约莫三十年的时间差!)追随了我自己从八岁起所追随的那场演进:一份原初平衡的退化。有可能:一份对“阳刚”价值的适度高估——以“阴柔”价值为代价——在我的接触、或我所属圈子的接触之下,变成一份一竿子到底的高估。可是,一如我在别处强调过的:那场由他编排的葬礼的“神经”(或“活力”),以及在他自身变形(它同时是——在“东家”之手——他身上的孩子的下葬……)中的“神经”——这份神经,决计不能仅仅寓居在对这样那样、或极端或(甚至)疯狂的价值体系的单纯采纳里。同样的话,也适用于我这个人认同中的“神经”、以及这份认同在我朋友生活中扮演的过大角色。毫无疑问:是同一份“力”在运作,而它的根,深深扎入他的童年\fnote{273}{(**)}{关于朝着这个方向的更精确的直觉,尤见“缓期中的宿怨——或事物的回转(2)”那条注(n\textsuperscript{o}~149)。}。

这里又有另一个奇怪的念头浮现。仿佛是我一生四十年里拖过的那副最重的担子——那场“男性”对我身上“女性”的压制,它同时也近似于“大东家”对我身上孩子的压制——这份担子,被我的朋友“\textbf{接了过去}”\dmark{}\origpage{738}:恰在那他本可显得与此类担子无缘的时刻。那大约是在我的价值系统朝阴的方向\emph{翻转}的时刻——那场演进,预示了十五年后的“与孩子的重逢”;彼时,我骤然感到卸下了一副巨大的重负\fnote{274}{(*)}{关于价值系统这场“翻转”,在“阳扮阴——或主人的角色”那条注(n\textsuperscript{o}~118)里谈到;关于“重逢”,在同名那条注(n\textsuperscript{o}~109)里。}。随即便浮起与印度教\textbf{业}的念头的联想。对我清楚:在这八年里,我卸去了自童年起拖着的那份业的一个实质部分。我本会以为(而且现在我仍倾向以为):这份卸除,不是“以牺牲”任何人完成的;它不仅对我是善的,而且“对全世界都善”。我甚至可以说,我\textbf{分明知道}正是如此——即便竟显出:另一个人选择了(甚至,另一个人的命运注定要)把这份业接过去。诚然,这份我所卸下的业,我并不视之为一个“恶”。它曾是我一场\textbf{成熟}的滋养实质——那场成熟,正摆在我面前。我知道:我吃了、以它为滋养,是好的,对我对一切人都好;而一份认识,是在无知的滋养母体里成形的\fnote{275}{(**)}{朝着同一方向的反思,见“循环”那条注(n\textsuperscript{o}~116$'$)的结尾,尤其其中最后一段。}。我本觉得:这份实质或这份业,一旦转化为认识,不留残渣,它消失了。老实说,我不知道印度教或佛教传统在这个问题上教什么——是否有一份“业之守恒”的律法(类似于物质守恒),全然不受食入、消化、同化那些创造生命的过程之影响。

出于得体的顾虑,我刚刚在这些“生命过程”中,把\textbf{排泄}偷藏了过去。可它(一如整个机体的死亡)是一份回收所吸收之物的关键过程:让“死”的有机质,回到那永无穷尽的转化循环——活着有机质的循环——凭它,生命永恒地重生自死亡\fnote{276}{(***)}{关于生与死的循环,亦见“行为”那条注(n\textsuperscript{o}~113)。}。

\subsubsection*{注~156$_{1}$}
\notehead{156$_{1}$}{}
\phantomsection\label{ch18:local:n65}

(2月30日)这份“样式”,终于被最后的第十二组打破:它竟不幸地有\dmark{}\origpage{739}\textbf{六}条注,把构成“阴阳之钥”的注之总数推到62。我原预计这组“冲突与发现”会有\textbf{八}条注——那正与“总体性”的标准相符,并将把注的总数推到64=8×8=4×4×4——也是《易经》的卦数!我的预期没有实现,我很遗憾;可我也不肯为凑数而“作弊”,把关于Pierre~Deligne来访我处的那两条注塞进“阴阳之钥”——它们的天然位置,依我看,在“葬礼仪式”的续篇里,置于“阴阳之钥”\textbf{之后}。

尽管,我对第十二组仍留着一份不满之感:它是“阴阳之钥”十二部分中唯一不留下\textbf{统一}的灵感与用意之印象的。这份统一的缺失,依我看,并非系于“冲突与发现”这一主题本身,而系于反思进程中异己事件(有时扰人的)的闯入。

(3月7日)昨夜重读我归入一条注(n\textsuperscript{o}~162)、题为“信念与认识——或传递”的1月14日反思时,我感到了对这名字的不满\fnote{277}{(*)}{它也是“阴阳之钥”的最后一条注。}。一方面,这“主”题与副题乍看并不相合——事实上,它们分别对应反思中第一、第三“乐章”:彼此表面毫无关联:前者是对一份认识之绽放过程的描述(以一份骤然的\textbf{信念}的形式),后者是对业之无穷链条与“传递”、一代传一代、一个人传另一个人的唤起。何况,那份最贴身个人化的内容、那份对我自己“要害”的内容——它构成反思“第二乐章”的实质(而且它正是那道“桥”,从第一乐章通向第三乐章)——这份关键内容,竟没有出现在所选的名字里。(当然,对我毫无怀疑:这回偷梁换柱,绝非纯属偶然……)既然三个主题各有其重要,而我一时看不到任何把三者都唤起的“佳”名或双名,我终于\dmark{}明白\origpage{740}:最好把这条注一分为三,各取一个提示性的名字:“信念与认识”、“最烫的铁——或转折”、“无穷的链——或传递(2)”(n\textsuperscript{o}~162、162$'$、162$''$)。

我是事后才骤然明白:凭这场操作——可谓由反思的实质本身所口授——那份我拖了将近两个月的“美学”不满,也同时得到了解决:正当“阴阳之钥”这第十二个、也是最后一个部分(我先前唤它作“冲突与发现”)执拗着不肯(以自然的方式,自不待言)让自己补全为一组\textbf{八}条注,而只肯容纳已经写就的那六条。而我因没有向那份轻易的诱惑低头——“作弊”一把,把两条“信手拈来”、位置本在别处的注“贴”到“阴阳之钥”的末尾——得到了我的回报!这“阴阳之钥”的最后一部分(它终将取名“恶之谜——或冲突与发现”),同时也呈现出一副漂亮的对称结构:两个注包(各含三条注)围绕中心主题而聚拢,分居于两条“离题注”——一条关于Fujii~Guruji,一条关于我的僧友们——的两侧。

\subsection{冲突与发现——或恶之谜}
\label{ch18:local:h70}

\subsubsection*{(a) 无恨亦无容赦}
\tocsubsub{(a) 无恨亦无容赦}
\label{ch18:local:h71}

\notehead{157}{}
\phantomsection\label{ch18:local:n66}

(1月4日)在昨天与前天的反思里,我着力求取的,是与我的朋友认同于我这个人这一实在的接触,并由此辨明它的分量与含义。这是一件我做起来仍像在幽暗中摸索的事——不说在漆黑里的话。或者倒不如说:我的眼睛依旧闭着,我的眼睑对一束我依旧无力知觉的光,是不透明的。不管怎样,我不记得在我与我朋友的关系的任何时刻,曾“感到”或“看见”过这份认同;一如我不曾“感到”或“看见”过他对我怀有的敌对意向。然而我\textbf{知道},毫无怀疑的余地,凭一束丰富而相互印证的事实:这份对我这个人的认同,与这份如影随形的敌对,是\textbf{实在}——如同一个生而失明的人“会知道”太阳、日光、颜色、明与暗存在,纵然他从未见过它们。他知道它们,却不\textbf{认识}这些东西。或者说,倘若他毕竟有一种极为弥散的认识——经由一种或许更精微的触觉(或经由一个“回忆”——它并非只植根于他一个人的生命,而是植根于无数代有视觉的存在者\dmark{}\origpage{741}的生命,那些先他而在的人)——这份认识也仍是间接而微弱的,如同一副温暖而洪亮的人声,借一道遥远而不确定的回声,传到我们这里。

这两天的工夫,做得仍像一个权宜之计,像一份缺席的直接知觉的替身。凡是我所理解的“禅观”工夫,多多少少都是如此。这工夫时时在逆着一份\textbf{惰性}的水流\textbf{推进}——逆着眼睑那铅一般沉沉的惰性!诚然,在眼睛全然睁开、全然警觉的那些时刻,无须什么禅观、什么工夫:只消看,便见到。既然这样的时刻稀少,与其抱臂坐等它们到来,我宁可抢在前头,不计较这工夫是否笨拙而“缓慢”。它再慢,甚至有时比往常更慢——却从不曾原地踏步,也从不曾兜圈子。只要有工夫——我是说真正的工夫,由一个真正的愿望驱动——就有进展:某件事在被做成,在成形,在转变,此时不知不觉,彼时有目共睹……而有时,在一场笨拙而执拗的推进的尽头——那推进在一片无形无廓的幽暗里绵延数小时、数日,乃至数月、或许数年——奇迹发生了:盲人\textbf{看见}了!而被看见的,并非一道倏忽即逝、仿佛从不曾存在过的幻象,只留下一缕微弱的回忆痕迹。这是一份诞生于这些幽暗劳作的\textbf{认识},一份新的认识,亲切如在,恰似我们所爱之物的滋味。

我在前天的反思里写道:若说九个月前,在写“敌对的父亲(1)”那条注(我刚刚引过它)的末几行时,有一件个案,它的想法曾“引着我的笔”——那便是我朋友皮埃尔与我关系的个案。然而,当时与我更贴近的其他“个案”,想必也在这份反思的背景深处,呈于我的心头。当我在那里谈到一位“既受敬慕又受畏惧、既被爱又被恨的父亲”,继而又谈到“另一个自己,被畏惧、被恨、被躲避……”时,“畏惧”“恨”,乃至多半连“躲避”这个词,都\textbf{不}适用于朋友皮埃尔与我个人的关系。无论是凭直接知觉——哪怕何其倏忽轻微——还是凭对我所知的确凿事实的印证,我从来不曾得到过丝毫迹象,指向我的朋友对我怀有一份畏惧,或一份仇恨、哪怕仅仅一份敌意。反倒是真的\dmark{}\origpage{742},一如我不止一次有机会强调的。而恰恰是这个情况,使那份无缝无隙、表面上无缘无故的敌对显得如此令人错愕:它在已往十五年间一路渐强,顶着“暂停!”的风格——别名“天鹅绒的脚掌”\fnote{278}{(*)}{见“暂停!”与“天鹅绒的脚掌——或微笑”两条注(n\textsuperscript{o}s~77、137),以及紧随后者之后的那些注——它们构成“阴阳之钥”的“天鹅绒中的利爪”部分。}——最终抵达一种安然放肆的调门:(只要守得某些规矩)确信全然免于惩处……

这一份令人错愕、难解其故的演进,随即与另一份同样“令人错愕”、“难解其故”的演进(这一回,这可是委婉的说法!)联结在一起:同样在大约十五年间,那份与曾是我妻子的那位之间的伴侣关系里持续恶化的演进,连带波及我们所组建的家庭。由于没有任何迹象向我示警,说我的妻子对我怀有仇恨或长期敌意的意向,我花了十年,才在这段关系的无情恶化中(其间我最充沛的精力都被数学占去——它扮演着那著名的鸵鸟沙堆的角色……)终于确认:在我继续爱着的那个人身上,存在着一份执拗、神秘而毫不容情的毁灭意愿——它借着那些我疼爱的人,施于我身。那是1967年,早于我离开家宅五年,也早于这场冲突为我解开十年——我曾把它感受为我一生背负过的最重的担子。凭一份久已承担下来的关系所给出的纵深,我只能确认那至今于我仍是奥秘的东西:一份永不餍足的毁灭意愿,与此同时,却是\textbf{没有仇恨}、甚至没有敌意——面对那些成人或孩子,只要时机凑巧,便遭无情打击。

(在适当比例下)这是同一个奥秘,正如我现在在我朋友对我的关系中所面对的,带着这份区别:那份“执拗的毁灭意愿……借着那些我珍爱的人施于我身”,严格地局限在数学家世界的层面上,而且它的工具和人质,不是我“肉身所生”的孩子,而是象征性地充当这一角色的存在:那些多少“承我的名字”的学生及同类。在这两种情形中,我不仅察觉不出仇恨亦无敌意,更有甚者,对我还怀有同情、甚至常常是毋庸置疑的深情厚谊。

这些并不是我仅在他人身上直面一份伤人意愿——甚至一份摧毁的意愿\dmark{}\origpage{743}(在最强烈的意义上\fnote{279}{(*)}{所谓“最强烈的意义”,我指的不是为使痛苦而使痛苦,不是毁掉某样为对方所珍爱的有限之物,而是对他人施以心灵上(乃至身体上)毁灭的意愿;是(在可能之际)面对“超出理解的东西”种下一份洗刷不掉、摧垮一切的绝望的意愿。在“反常层会议”那副光鲜可亲的门面背后,我似乎在它最耀眼的两位活动家身上,重新发现了这个极端的维度……})——的仅有情境:在其中,我察觉不出丝毫仇恨或敌意的痕迹。对我一生打下最深烙印的那一例,发生在1933年,我生命的第六个年头,主角是我的母亲——就在那一年,我们构成的那个家——我的父母、我的姐姐和我——被永远地摧毁了\fnote{280}{(**)}{关于这段往事,见“超级父亲”那条注,n\textsuperscript{o}~108.}。

我近距离经历过的这类情境——无恨而欲毁灭,或欲伤人至深、深至所能,而我从中察觉不出丝毫敌意——彼此显得非常不同。我很怀疑能为它们找到一个共同的“解释”,或至少在主角们遥远的身世里找到一个共同的特征,足以暗示一条深的因果纽带\fnote{281}{(***)}{然而,一种恶毒的、深埋于内的对己之轻蔑,想必为所有这些情境所共有。或许,这样一种恶毒(当它不曾凭一个恩典的行动、凭一场深刻的内在转变而解消时——也就是说,只要它还不曾被“承担”)总要寻得出口,借种种毁灭性的行动、借一份毁灭的意愿表达出来;当它无法在他人身上找到目标时,便反身扑向自己。在不止一个男男女女身上,乃至在亲近的人们中间,我多次得以确认一份毁灭意愿的同时作用:既指向自身,又指向某个从亲近者中选出的外部目标(母亲、父亲、配偶或孩子……)。(1985年2月)亦见“无缘之暴力的缘由”那条注(n\textsuperscript{o}~159)里的反思——它比本注的反思晚三天,而本注显然为它作了铺垫。}。或许比一份解释更要紧、无论如何更在先的一件事,是先得\textbf{确认}这样一桩东西的存在:\textbf{无恨而毁灭的意愿}。我在这里接上了“无端的暴力”的主题——此前已从另一个侧面触及过它\fnote{282}{(****)}{见“天真的暴力”那条注,n\textsuperscript{o}~139.}。这里涉及的,是对一个亲近的存在、或一个被视作“朋友”的人施行的无端的暴力(有时是毁灭性的)。这样一种暴力——它难得自道其名\dmark{}\origpage{744}——在日常生活中单单\textbf{存在},就是每个人生活里的一个重要事实——人类生活中的重要事实之一。确认这一事实——顶着那些根深蒂固、时时撺掇我们把它藏掖过去的机制——是承担它的第一步。这一步,任何理论、任何推理、任何“进路”,都无法替我们省去。

我不知有朝一日我是否会\textbf{领会}这一事实;我觉得,领会它,也是“领会冲突”的一部分。对我清楚的是:这样一种领会,不可能出自一份“理论”,也不可能出自一份“经验”(单凭经验之力)。它不是某种累积(“知识”的或“经验”的)的“总和”,一如它不属于单纯智性的层面,甚至不属于单纯“智力”的层面\fnote{283}{(*)}{(3月5日)至少我知道:这样一种领会,唯有借对这份暴力\textbf{在我自身之中的}领会,才会来到我这里。}。我不敢肯定我认识什么人——哪怕只闻其名——在他身上活着这样一种领会。但我觉得:一个人,在一份无可否认、有千副面孔的实在面前,经过百次千次的闪避之后,终于抵达对这一事实的单单\textbf{确认}——谦卑地,无怨苦亦无反叛,无屈从亦无愤慨——如同确认一个可畏的\textbf{奥秘}:或许,它的意义逃过他的把握,他却预感到它的广度与深度;一个令他着迷、向他发问,却不再使他惊惧、不再使他惶惑的奥秘——这个人,就不曾白白活过。

\subsubsection*{(b) 领会与更新}
\tocsubsub{(b) 领会与更新}
\label{ch18:local:h72}

\notehead{158}{}
\phantomsection\label{ch18:local:n67}

(1月5日)并非预先谋定,昨天那份反思的结尾声调,依然全然是葬礼悼词的调门——只是这一回,由逝者本人亲口宣读(或咏唱)。人从来是自己侍候自己最好!

昨天,我再度直面“冲突之谜”最令人错愕的面向之一:一份无恨、亦无表面缘由的毁灭意愿,在暗处执拗而不懈地施于一个亲近的人,或这样的亲近者、朋友们。这样一份意愿,有时会终于失控加速,演成一场四面八方撒开的毁灭狂热:凡显出脆弱之相的,都成了受欢迎的靶子。这像一股按捺不住、专事倒行逆施的“行动”贪食症;当旁观的人(乃至刚吃了这份妨害的人)颇有幽默感\dmark{}\origpage{745},而那位“演员-提线人”支配他人的能耐又颇为有限时,它的重复性(像小丑把戏),连同它在牵线艺术上炉火纯青的绝技,效果可以滑稽之至。而当这些马戏游戏的妨害承受者中间有孩子在场——哪怕它们只在转喻的意义上“见血”——情形就更其严肃、更其后果攸关;或当那个被毁灭之渴攫住的人——无论男女——对某些同类握有可观乃至全凭己意的权力时。历史给我们传下了一些暴君的名字,他们被这样一种不加区分的毁灭疯狂攫住,把自家领地变成一座巨大的藏骸所。人们会想到伊凡雷帝,或斯大林,或中国某位皇帝(名字与朝代我都忘了)——他最终被走投无路、手持棍棒木桩的臣民们打倒了\fnote{284}{(*)}{这位皇帝惧怕民众暴动,曾禁止人民使用一切金属器物(如刀、叉之类)——凡可用作武器的——唯每村留一把刀,用坚固的铁链拴在一处公共场所。

三位人物有一个共同点:除这份毁灭之渴外,他们同样被\textbf{恐惧}攫住——怕遭暗杀,而在这恐惧之外,无疑还有对自身那无可避免的\textbf{死亡}的恐惧——可他们正四下播种着死亡。这个巧合决非偶然。我还注意到:斯大林(三人中唯一我有稍许翔实材料的一位)进入政治生涯时,恰恰是一位大行家——行在牵线的艺术,行在拿他人的虚荣与贪欲做杠杆来摆布人。他最早习得的风格,似乎是“天鹅绒的脚掌”,直到哪一天,连掩藏利爪都成了多余。

我未把我那位(前)同胞希特勒列入上述例子里,并非由于我对他怀有什么特别的同情,而是因为在他身上,我察觉不到刚才谈到的那种“四面八方”的毁灭癖。轻蔑、继而是毁灭的目标,是那些被指认为“他者”、“外人”的人:先是“犹太人”(还有纳粹行话津津乐道的共产党人与其他“犹太-布尔什维克”),继而是“亚洲人”和其他非雅利安的外来杂种。在希特勒治下,不沾犹太血统的正派德国人过得安逸至极——至少直到盟军头几轮大空袭、战局对他们真正急转直下之前。}。毫无疑问,在我们自己的地界也曾有过类似的情形,规模或许小些,而“历史”对它们则更其缄默……

昨天我写道——不带丝毫假谦虚——我不领会我刚确认过的那个“事实”:无恨的毁灭之渴。这绝不意味着我对此毫无想法,恰恰相反。我手里有的,明显不止一些单纯的“想法”,而正是几份强韧的直觉。它们诞生于、生长于我生命的沃土——那土壤富于冲突,那些冲突有时仿佛把它摧垮,如一场场无尽的风暴,在一片凝然不动的冬日大景里发作,毫不容情地拔走该拔走的\dmark{}\origpage{746}\fnote{285}{(*)}{这意象刚随笔势落纸,我便发觉它只是部分贴切——几乎带着一股“陈词滥调”的余味!在这余味上略一驻足,我找回了心里那个久已有之的蓄意图谋:“以阳观看我的一生”:运动、箭与风暴……

还来不及驻足,却分明觉出这意象有毛病(可它偏偏就是来到我笔下的那个意象,没办法!),我便在正文里“纠正弹道”,顺势接上“静候着的沉睡大地”——这可就是阴了!这正是把“假和音”(或“不协和音”)“解决”掉的那个和音。有一个意象,在许多方面都比风暴的意象——“拔走该拔走的”——更贴切,而且调门恰是更阴的:那虫子,啃噬“该啃噬的”——而它终究塌落——但一切都化作静候着的大地的腹中之物,而春天回来时……(接续不变!)。}。而一切,都为静候中的沉睡大地化作腹中之物。春天回来时,在僵卧于地、一动不动的巨树残干的凹处,浓烈的生命在攒动;到次年春天(若不是当年),那里已见青草与野花绽放。

这些“强韧的直觉”,我相信,全都关乎冲突的“成分”。其中某些成分,我稍稍谈过,又再谈过;头一个便是“对己之轻蔑”,以及它与我们原初存在的某些本质面向与力量的压制之间的关联——诸如“阴”“阳”这两个“坡面”,其中一个常遭否认。我也曾多次有机会谈到虚荣:它像一张名片,是我们之中冲突在场的所有征象里最普遍、最显眼的一个;在我看来,它是同一枚奖章的“正面”,其“反面”便是对己之轻蔑。还有对他人的轻蔑:对己之轻蔑向外的投射,同时又是它的掩体,或者更确切说,是它的宣泄口与驱魔仪。对他人的轻蔑,说到底无非是对其存在的刻意无视——把他当作一个有感受、与我们同等分享这个世界的存在。无端的暴力,唯有在这样的轻蔑的土壤上才能发芽、增殖。还有认识的恐惧,对实在的恐惧——这份恐惧的神经中枢,那个“黑点”,一份焦虑旋涡的震中——一遇风吹草动便待引爆,是对认识自己的恐惧:怕认出自己那些姿态与遁词,哪怕是最粗劣的;也怕认出我们之内那份创造的力——它日复一日被我们用同样的姿态与遁词拒斥、掩埋。

在我的生活里,恐惧现身于六岁,当时(照我看)还没有任何虚荣。虚荣想必是后来才出现的,在那场“翻转”发生的时刻(我猜想)——约在八岁\dmark{}\origpage{747}\fnote{286}{(*)}{关于这场“翻转”,见“超级父亲”那条注(n\textsuperscript{o}~108)。}。而恐惧也是第一个消失的,不留痕迹,随着一份既仁厚又不敬的好奇心的现身——它固然好奇,却绝不被那些装神弄鬼、阴森可怖的大场面机关布景震慑,诸如“黑点”公司一类。相反,自认识的恐惧消失以来的八年间,虚荣的机制一直原位未动,看不出变化。改变的,只是这些机制对我生活的把持:在一份警醒的好奇心在场的时刻,它们被拆除了引信——那份好奇心可不吃糊弄这一套!

我手里握着的,是冲突成分的一整个谱系——凭第一手、不带分毫怀疑,我知道:这些确确实实是成分,而且是本质性的成分。而多少年来,我也同样一切在手,只等我高兴的时刻,把这些成分“组装”起来:借我在自己身上、在他人身上得以观察的东西,细心阐明它们相邻与相依的关联。这是一件几天、几个星期——依我看甚至用不着几个月——的工夫,想必会非常有教益、非常有用。而我至今不曾费心去做,把优先让给了更直接切身的其他方向;这无疑是因为我心知肚明:对冲突的“领会”,不可能来自这样一种成分的“组装”——用一般性的措辞、而我这个人缺席其中的组装(除非我作为“例子”之一);也不可能凭把一定数目的单质并排放在一起、“组装”甚至混合——单质,即构成某一化合物的一些“成分”——就把这化合物复原出来。要让“复原”得以发生,首先得有一场“化学反应”——某种东西,以全然另一种方式把诸成分置于接触与作用之中,凭一种全然另类的力——那是单纯的“组装”或混合所不能为。

对生活之事的领会,亦复如是。单凭智力,勉强能辨认出“冲突”这样一桩东西的种种成分;而在成分已知的情况下,借助于有关的事实(第一手或第二手所知),它至少能把它们组装得貌似有理,甚至“正确”。这样的工夫自有用处:偶或在某一场冲突的情境里认得路径,从中得出一份多少精确的“病因学”——但这仍不是一份“对冲突的领会”。我倒是会说:朝这样一种领会迈进了一步,是在那一天——我的\textbf{与冲突的关系}得以转变的那一天\dmark{}\origpage{748}。我在这里说“我与冲突的关系”——这首先、自然,指的是我自己这个人身上的冲突,以及(由此出发)偶尔使我与这个或那个人对立的冲突;而最后,是我在日常生活中看到的、在或近或疏的人们身上运作的冲突——它常常借他们彼此之间的冲突表达出来。

在已往的八年里,确实有过这样一场朝着领会冲突的推进,也就是说:也是我与冲突的关系的一场转变——或不如说,接连的转变。我曾有机会提到其中两三段经历\fnote{287}{(*)}{关于这一点,尤见“接纳(阴的苏醒(2))”与“奴隶与傀儡——或闸门”两条注,n\textsuperscript{o}s~110、140.}。或许,对冲突的圆满领会,等值于对冲突之存在的圆满接纳——无论它在哪里,以何种方式显现\fnote{288}{(**)}{这样一份“圆满接纳”的意义,可能引起无数的误会。它与一份默许是全然不同的性质。它不排除\textbf{拒绝}——干脆而毫不含糊的拒绝——它把拒绝包含在内。关于这一点,见“伴侣——或‘恶’之谜”那条注(n\textsuperscript{o}~117)里的反思。}。显而易见,我离此尚远!而或许,对冲突的圆满领会,还意味着冲突在我自己这个人身上的全然消解。我离此更远!

然而关于那股力量的性质——它从一堆组装好的成分里,骤然激起一份更新此人的领会——我相信我还知道一件事。恰恰是这股力量,不“属于智力的层面”。我怀疑有任何智性工夫——譬如读书,且不管那些书多么博学、深刻或崇高——能丝毫促动它的显现。当它喷涌而出时,只在静默之中,只在接触我们这个人、我们经历里最切身之处;那是一样东西,任何书、任何人——哪怕是基督或佛陀——都永远无法向我们揭示。

当我说“最切身之处”,这并不意味着那是我们无法谈论的东西——对自己,或对他人——而且有时,谈一谈是有益的。但哪怕以天使之音、以先知之音来说,被说出的也不是那桩东西本身。这桩东西——早已认识,或许却深埋着;与它的接触,能骤然激起一份新的认识——这桩东西,天使不认识,先知不认识,乃至最亲近、最被深爱之人也不认识,唯独\textbf{你}认识。

回到冲突,回到“无恨的毁灭”\dmark{}\origpage{749}——在我看来,它是冲突最坚硬的“硬核”,最抗拒领会,也就是说:也最抗拒\textbf{接纳}。我还相信我知道,在摆在我面前、要更深入其中的下一步里,那桩“最切身”的东西是什么——我首先得重新找到与它的接触;在此处,它扮演的,正是那个出名的、被如此执拗地避开的“黑点”的角色!那是“无端的暴力”、对他人之轻蔑(或许还有“无恨的毁灭”)种种情境里的亲身经历——在其中,主角\textbf{是我}:行暴力的是我,从轻蔑中得计的是我。唯有接触这份实在——否则永无可能——我才有机会就那个出名的“对己之轻蔑”弄个水落石出,终于\textbf{看见}——在一切“毫无疑问”、一切“或许”之外——它是否确确实实是恶的深根,而不只是在“除我之外的所有人”身上!

\subsubsection*{(c) 无缘之暴力的缘由}
\tocsubsub{(c) 无缘之暴力的缘由}
\label{ch18:local:h73}

\notehead{159}{}
\phantomsection\label{ch18:local:n68}

(1月7日)前两条注里的反思,围绕着这桩怪东西的存在的奥秘打转:一份毁灭的意愿(或伤人的、羞辱人的、损害人的意愿),却在全然无恨、无敌意的情况下。促成这份反思的,是我朋友皮埃尔对我的关系,它随即唤起我的前妻对我的关系的联想。在整个关于葬礼的反思里,我不止一次被引着去确认,或去记起:在这两例以及别的情形里,在我这个人身上起作用的,正是某些特质——我自八岁起在自己身上培养的那些“超级阳刚”的特质——它们充当了这类敌对冲动的刺激物与“吸引子”。如果我没有弄错,这一点头一次谈到,是在10月5日那条注“超级父亲(阳葬阴(2))”(n\textsuperscript{o}~108)里。这层关联,在紧接着的10月9日那条注“重逢(阴的苏醒(1))”(n\textsuperscript{o}~109)里被重新接起。

在这条注里,我回到我一生中头一次知觉到这层关联的时刻。那是1976年10月18日,恰是与我心内孩子重逢的那一天,在见证这个日子——我成年生命里所有日子中最要紧的一个——的那几条注的末几行里。在那几行(已转录于所引的那条注)里,我谈到我曾爱过的三个女人的“隐秘的恨与怨怼”,其中一位,在那时还是我的妻子(尽管我已五年不曾与她同住)。凭回顾,我觉得:在我当时想及的三个个案里,每一个的这份“隐秘的恨”的印象\dmark{}\origpage{750},严格说来,并不对应实在——我是说,并不对应于任何时候我曾有的、对这样一份恨的直接知觉\fnote{289}{(*)}{(3月6日)写下这几行后,我记起:在我婚姻生活的岁月里,毕竟有过两段插曲——第一段几天,第二段几分钟——其间我觉得自己像被两束仇恨的强光袭中,它们从当时还是我妻子的人的眼中迸射出来。

第一次,我妻子正经受人们(委婉地)所谓的一场“神经抑郁”,时值我们婚姻的第五个年头(1962年)。这段插曲深深烙进了夫妻的生活与家庭的气氛。这也是我一生中——在所有我有清醒记忆的时刻里——活得最惨烈的时刻,给我留下的烙印也最深(而它也确实被指望如此)。

除非内里有一副异常稳固的根基(当年我远没有这份根基,那时还谈不上成熟),以我们为靶子的仇恨——当它来自所爱、所亲之人时更是如此——对我们的心灵有一种毁灭性的效果,只要它在我们心中唤起一种同样性质的、指向我们自身的毁灭性仇恨。我们身上似乎总有什么东西,非要不惜代价给“超出理解的东西”找出一个“意义”来,哪怕这个“意义”无非是一场自我谴责、一次无可挽回的自我拒斥:既然我们被恨(何况这恨的“缘由”我们全然无从把握……),那就是我们可恨……

我之所以被这段插曲创击得如此之深——在其后的六七年间,它始终像一把达摩克利斯之剑,悬在我的生命之上——想必是它与一场我童年的创伤经历发生了猛烈的共振。那经历已从清醒记忆中消失;可正因如此,每当我骤然直面一份无从解释的恶意或仇恨,它反而愈发作用——它们同样骤然、同样无从解释,一如那股在我五岁上袭中我的毁灭意愿,而它当年恰恰来自那样一个人:无论我在记忆里回溯多远,她都曾是宇宙那安宁而可靠的中心。

关于我有时成为其靶子的那份恶意或仇恨,我在生活里终于学到的重要一课是:我却根本不是它真正而直接的\textbf{缘由}(纵然我这人身上的某些方面——我既不否认也不推拒——有助于把它引到我身上)。可这份认识,此后多年仍嫌浮于表皮,不足以拆除那深植于我体内、每当我直面一份表面“无缘”的恶意或暴力便启动的机制。要拆除它,首先得让我溯到它的根,循着那些被遗忘的、浸满焦虑的日日夜夜寻踪而进——那时,我的母亲骤然间、神秘而不由分说地,成了一个陌生人,一个敌意的、可畏的陌生人……}。我曾知觉到、并有大把机会承受其效果的,是一份毁灭的意愿,或一份使人痛苦、伤人的意愿——既持久,又表面无从解释,无缘无故——我曾把它\textbf{解释}为一份恨的征象,一份“隐秘的”恨,因为它从不曾宣之于口。此外我还相信:对其中两位,正是在所引的那几行里——自我认识她们\dmark{}\origpage{751}以来头一次——我确认了这份在我看来是“隐秘的恨”的东西。在我当时所处的那一点上,我不可能不犯下我刚才指出的这层混淆。但这层混淆,丝毫不减损作出这份确认的重要性——在这份确认里,我把自己卷入得那样深,一如那些与我至近相连的女人们。

至于“怨怼”——它与“隐秘的恨”一并被谈到——从那一刻起我就感到:倘若我身上的一股“超级阳”的“某种力量”把这三个女人各自的怨怼引到了我这个人身上,那也是为着一些怨情——可我对它们毫无责任——为着一些创痛与损害,受于“早在她们认识我的存在之前、在一段缺爱童年的仓皇岁月里”。这份知觉,经年累月沉淀下来,如一段浓烈经历的果实;对我去年12月20日那条注“缓期中的宿怨——或事物的回转(2)”(n\textsuperscript{o}~149)里的反思,想必起过隐形向导的作用:在那条注里,出现了这样一份直觉:同一个怨怼移位的过程——一份原初的怨怼、一份“虚位以待的宿怨”的移位——想必也发生在我朋友皮埃尔身上,约在我们相遇的时刻,甚或更早。我所知的那些事实,至少使这份直觉显得可信。

然而与我前妻的情形相比,与重逢之后那场反思里谈到的另外两例相比,有一处重要的不同。的确,我全然不觉得我朋友的童年有多少“仓皇”或“缺爱”。这份不同,在我看来,显现在我朋友对我的那份敌对的调门里:它从不曾抵达那种激烈的音高——那样的音高,在那三段关系里,于我是何其熟悉。同样,在我朋友与我的关系里,敌对迹象的显现,起初极为隐微而零星;甚至在我1970年离开之后,又过了八年,这份敌对才直接而毫无疑义地表达出来,指向我这个人本身\fnote{290}{}{关于这一点,见“两个转折”那条注,n\textsuperscript{o}~66.}。这一切,恰与一份原初“怨怼”的存在相呼应:它始终弥散、无从称量,没有一个坚硬的“内核”与之相应——那本该是一份受辱或受害之感(哪怕藏于清醒目光之外),被感受为或许无可挽回……

在前一条注里,回想起那份毁伤他人的意愿\dmark{}\origpage{752}——它施于人身,却在\textbf{无恨亦无敌意}的情况下——一个表面的矛盾浮上心头(而且相当执拗),我本想立刻回头谈它。它是这样的。在我注意中心的两个个案里——涉及我从前的学生(和我推定的数学“继承人”)与曾是我妻子的那位——谈的确实是一份无意识的“宿怨”,他们把它转到了我这个人身上。可“宿怨”或“怨怼”这个观念本身,似乎与“敌意”或“不和”的观念连在一起:人们会想说:宿怨(或怨怼)是滋养一份敌意的可能方式之一——而且是最常见的方式之一。这份断言,在一种可称作“直接的”宿怨的情形里,肯定站得住:一份“真”宿怨,由针对当事人的一份怨情(真实的或想象的)所驱动,由当事人施加于我们的一份委屈或损害所驱动。但在我所关切的情形里,涉及的不是这样的宿怨,而是一份间接的宿怨——不妨说“凭代理”的宿怨:它从一副原初的潜在靶子上,挪往一副“认养的”或替代的靶子——前者出于这样那样的缘故而不配\fnote{291}{(*)}{这样的“缘故”有很多,它们常常使那个(有意或无意)造成损害、施加伤害的人,反倒“不配”充当一份怨怼或敌意——乃至一份仇恨或毁灭意愿——的靶子,尽管那意愿确确实实由他引发。最常见的或许(尤其当对方是母亲或父亲,或是一位凭其身份、社会地位而被视作不可触犯的人)是那道栏障:不敢违逆一份权威的禁忌——早已内化于心的禁忌。这类栏障力量极大。(对我而言,近十五年来,它们倾向于日渐消隐……)

反过来,也可能有这种情形:当事人“不配其分”,受不起一份与所受委屈相称的宿怨——显得太无足轻重、太躲闪或太怯懦,当不起那份本该名正言顺归于他的角色。

最后,我也可以设想:在某些情形里,所受的损害太过飘忽、太过精微(说到底,按当事人早已内化的通行共识,甚至“不存在”),除了生成一份弥散的宿怨之外,生成不出别的什么;那份宿怨无力“凝缩”,无法在一段本身调门温和、没有明显棱角的关系里成形、发力。这无疑只是前一情形的一个变体,出现在“缓期中的宿怨——或事物的回转(2)”那条注(n\textsuperscript{o}~149)的反思里。},后者则看来与事态的需要“对得上口径”。值得注意的是:这样一份“移位的宿怨”(这回可是实至名归的措辞!)——它作为那份执拗的力在运作,藏在种种态度、行为与行动背后,那些东西的性质,让人会说它们是被一份仇恨\dmark{}\origpage{753}、或一份“无缘的”敌意所驱动——这样一份“宿怨”,却\textbf{不带丝毫仇恨或敌意之感}!况且,正是“无端的暴力”——最强烈意义下的,即我此刻正在考察的这份——这两个\textbf{方面}的相互结合,使它如此令人错愕,如一桩真正“超出理解”的东西\fnote{292}{(*)}{关于这种“超出理解”的暴力(“德文作~unfassbar”),见“奴隶与傀儡——或闸门”那条注(n\textsuperscript{o}~140)。我在这里谈“最强烈意义”下的无端暴力时,当下只把它称作“超出理解的东西”而不另加限定;我那时心中所想的确切含义,在接下来的说明里被圈定:凭对那两个在它之中相互结合的“方面”的阐明。}:这份暴力全然没有理性的、可触的“缘由”,无论在吃这份妨害的人那一侧(他并未以伤人或损人的态度、行为或行动挑起它),还是在行使它的人那一侧(他并未被他对靶子怀有的——“有理或没理的”——仇恨或敌意之感所驱动)。

或许,在我所关切的这些情形里(人直面一份显得“无缘”、未被挑起的暴力),仇恨或敌意是否存在这个问题在这里是相对次要的;可以肯定的是,一如在我这里曾有的那样,在承受这份暴力的人的经历里,从被施的暴力变得有意识的时刻起,必会浮起一份“隐秘的恨”或“敌意”的印象——出自施加它的那一方。可这份印象,全然不是知觉的效果(那知觉仿佛被魔杖一点、骤然现身),而是生吞活剥的\textbf{等同}的效果:暴力=仇恨(或敌意)\fnote{293}{(**)}{(3月6日)不过在某些情形里,对一份确确实实在场的仇恨,完全可能有知觉,哪怕它根本不曾被挑起。(关于这一点,见本注上面另一条同日的脚注。)此时涉及的仇恨,除非在例外的情形,始终圈禁在无意识的深层,而且在那里还处于一种“虚位”状态,没有指定的靶子;可它恰恰是暗中驱动种种暴力行动的力量(多以阴险的形式)——而那些行动,却确确实实、毫不动摇地指向同一个选定的靶子……}。

但在我看来要紧得多的一件事,是确认:不仅“凭代理的宿怨”这样一桩东西存在着——它表面何其反常、何其悖理、何其违背那些最根深蒂固的“常识”反射:从它的“原初靶子”\dmark{}\origpage{754}(或诸原初靶子)挪往一副“替代靶子”(一副几乎纯为图方便的靶子!);而且要确认,更进一层:这是最常见的机制之一,在每个街角都会遇到,无论是在自己这个人身上(那恰是人们最不会想到去找它的地方……),还是在亲近者、朋友们身上。我甚至有这份印象:这一机制具有普遍的本性,它属于人心那些基本机制之列,属于构成“在实在面前逃避”综合征的那几种万能机制:拒绝去认识实在,害怕承担它。

说得更确切些,我有这份印象:今天,我把手指点在了所有“无端的暴力”情境——无一例外——共有的那个机关上。这份印象是带着一份骤然信念的力量浮现的,当时我正着手考察(三段之前)那个“表面的矛盾”。我那时感到:我一生积攒的、围绕这份“超出理解”的暴力那最疼痛点的、一大堆零碎而不相属的印象,骤然排起了秩序,霎时获得了一直欠缺的一份视角——这份视角出人意料地在那里现身——在一次反思收尾的转弯处,而我正准备给最后一个~i~点上最后一个点……

\subsubsection*{(d) Nichidatsu~Fujii~Guruji——或太阳与其行星}
\tocsubsub{(d) Nichidatsu~Fujii~Guruji——或太阳与其行星}
\label{ch18:local:h74}

\notehead{160}{}
\phantomsection\label{ch18:local:n69}

(1月8日)一周以来,一股少见的寒潮——零下十五度或更低,而当风从“旺图山”吹来(这名字名副其实——“旺图”(Ventoux)里就带着“风”),想必还要更冷。据说这股寒潮在世界各地肆虐(据一位听新闻的人说),而在南方,自1956年那场有名的寒冬与晚春以来,还不曾见过。我在德国的童年经历过这样的冷,但那时有雪,雪护着土地,给空气、给万物添上一层温柔。眼下这场无雪的冷,地表冻得像一整块冰。几天工夫,园子就被剃了个精光——春天来时,我们播下种下的还会剩些什么,我不知道。残存的韭葱、芹菜、莙荙菜、野苣、甜菜、刺菜蓟的叶子,都像冰做的叶子,冻菜一般。人们赶着一天天尽量多收,随收随吃,赶在解冻之前——一解冻,一切都得进堆肥。而昨天\dmark{}\origpage{755},厨房里的进水管冻住了,幸而楼下旧车房里还有自来水,那边不那么受寒。今天一位朋友带着便携燃气喷灯来了,把水弄通了。往后得让水细细地流着,免得它立刻又冻住。所幸饭厅里有一只好柴炉,我把活儿挪到了那里。坐在炉边,实在舒坦。我烧葡萄藤茬取暖,每天拿斧子劈,满满一葡萄筐还冒尖——天就是这么冷。当风整个下午刮个不停,光是顶着风劈上一刻钟、二十分钟的柴,就足以让人把手冻坏。更别提留在外面的车,再打不着火了——据说汽车这东西,顶不住大冷,防冻液也一样。还是那位好心的朋友,刚才把车发动了起来;可明天它还走得动吗——去校读那位我交付了活儿的打字员的打件?总之,冬天一股寒潮——有时是夏天一股热浪,或随便什么时候来一场好好的小病——就足以提醒我们几桩存在的实情;而当一切如意运转,人们容易把它们忘掉……

三个月来,不知不觉,我的工作节奏挪向了夜间的时辰。我工作到夜里两三点,睡到十一点、中午。就这天时,倘若我在床上由着自己一次,十二个小时睡下去毫不费事——反过来,一旦上了工,我就不再上床!目前我试着保持一份合理的平衡。作息的错位,我不太惊慌,只要觉还睡得好,不至于一连几个小时躺在床上睡不着、那台思考的机器还转个不停。即便眼下园里没什么活,每天总还有足够多样的事情可做,包括备柴火,还有这里那里一点小体操。我有一种满意的生计平衡之感:其中,发现的工作并不摆出吞没其余一切的架势,却也绝不受亏待。自9月22日重新上工以来,我平均每天得花上五六个钟头。不算多,可“收成”似乎不比从前差多少。“砍伐量”(每月约一百页)与写《收获与播种》前两部分时大致相当,相去无几。而从质的一面看,我毫无怀疑:最深的正是这第三部——它教我认识自己、认识他人,教得最多。

\asterisms

\dmark{}\origpage{756}Na mu myo ho ren ge kyo !

正当我写完这一小段回顾——关于冬天的严酷,关于我生计平衡的演变——接到了日本山妙法寺团体一位僧友的电话,告知他们尊崇的“preceptor”\fnote{294}{(*)}{“Preceptor”,英语词,与“teacher”约略相当,指“师父”、施教者。日本山妙法寺(Nihonzan Myohoji)是该团体日本名称的音译,意为“日本传道团”。这是一个佛教“传道”团体,主要宗旨是和平主义。详情见下文。}Nichidatsu~Fujii与世长辞——他更广为人知的名字是Fujii~Guruji,亲近的人则唤他“Osshosama”。我巴黎的朋友刚从东京的电话里得知消息,我猜想,Fujii~Guruji就在今天逝世\fnote{295}{(**)}{事后证实,他当时刚去世不过几个小时。消息传得真快!}。去年8月6日,他刚满一百岁;身体已衰弱,头脑却极佳。

奇巧的是,8月6日这一天,还是另两桩重要事件的纪念日,一桩具有历史的分量,一桩于我个人切身。它是广岛原子弹(1945年8月6日)的纪念日——日本人把它纪念作“广岛日”(Hiroshima day)。(正因如此,Fujii~Guruji的寿庆反倒多在7月末举行,好把8月6日前后的日子留给和平与反核的集会。)另一桩:我的父亲生于1890年8月6日,比Fujii~Guruji的诞生——一日不差——整整晚六年。

继Claude~Chevalley之后,Nichidatsu~Fujii是在《收获与播种》写作期间辞世的第二位、在我生命中起过不可忽视作用的人。面对他的离世(它其实算不得意外),我尤其庆幸:直到去年,我与他还有过一回满含暖意的通信。我曾受邀出席老师父的百岁寿庆典礼——它将在东京以罕见的排场举行。(人们甚至赶工印出一本关于他本人的证言小书,备作这一场合的献礼。)那给了我一个机会,像几乎每年一样,写几句预致的贺词,并为不能出席7月30日的典礼致歉——写这封信时,我自己还多多少少卧床。(诚然,我也不怎么喜欢大型公开典礼,可我当时觉得,信里不必提这个。无论如何,我一味固执地谢绝出席任何“大场合”\fnote{296}{(*)}{这类“大场合”中主要的,是“和平塔”(Shanti Stoupas)的落成典礼。这类塔——为世界和平而设的静修之所——的建造,溯自佛教世界一条极古老的传统(由印度的阿育王开启),也是Fujii~Guruji毕生的主要关切之一。他推动了世界各地大量和平塔的建造,其中三座在欧洲,一座在美国。},想必让不止一位僧友失望和伤心了\dmark{}\origpage{757},而他们邀请我,从不厌倦。)我还在信里添了几句话,谈病有益的一面:它不由分说,强使我们从诸务中“脱开”,把身体索要的还给身体。Fujii~Guruji自己,过去这一年也大半卧床;以他那倾于行动的性情、那份少有的精力,想来于他是沉重的。我已有七年多不曾收到Fujii~Guruji的私人音讯,却惊讶地收到他的一封来信——他卧床时亲自口授的。信(我刚刚重读了一遍)写于1984年7月13日。信中满纸体贴:他挂念我的健康,难过不能派人前来照料我。他也谈自己的健康,谈他以怎样的心绪承受这份被强加的无为。信以这样几句作结——那腔调非常“日本式”,须拿一颗(大大的!)盐粒来听——而它们向我显出的,或许比全信的其余部分更多:他的精气神一如既往\fnote{297}{(**)}{信是以日语口授的(那是古鲁吉唯一会说的语言),随即直接译成英文。所引几行的法译:“诚然,我是一个十分衰老、衰朽无用的人,即使我还能恢复正常的健康。可我仍愿意活下去,看看这世界怎样转下去。”}:

\begin{quote}
Indeed I am a very old decrepit man of no use even if I may get back to normal life. Yet still,
I would like to live and see how the world turns。
\end{quote}

随后,他果然还看着这世界转了近六个月……

我与日本山妙法寺团体的渊源,上溯到1974年。在这里,这些多插曲、方方面面都有的关系,连勾勒一个轮廓也谈不上——那得用一整卷书。它们是“生存与生活”那段插曲\fnote{298}{(***)}{《愚妄与更新》(《收获与播种》第一部分)里,多次提到这段插曲。“生存与生活”(起初就叫“生存”,再无其他)是一个团体的名字,起初以和平主义为宗旨,继而也兼生态关怀,1970年7月(在蒙特利尔大学一次“Summer School”的边缘)诞生于一群科学家(尤其是数学家)中间。它很快朝“文化革命”的方向演进,同时把听众扩展到科学界之外。它的主要行动手段,是同名的(半定期)公报;历任主持依次为Claude~Chevalley、我本人、Pierre~Samuel、Denis~Geygi(四人皆数学家)——还不算一个英文版,由Gordon~Edwards勉力维持(一位年轻的加拿大数学家,我在蒙特利尔结识他;他是团体与公报的少数发起人之一)。

第一期公报全部出自我的笔下(天真而满怀信念!),印了一千份,在尼斯国际数学家大会(1970年)上散发——那届大会(照例每四年一届)聚了数千名数学家。我本指望会有大规模的响应——记得响应一共两三份。我尤其感到同事们中间一片难堪!谈到科学家与军事机器的协作——它们已渗透到科学生活的各个角落——我几乎是一脚踩进了最肥腻的菜碟……正是在科学的“上流大世界”里,我感到最大的难堪——从那边传来的同情回应,最后只落到Chevalley与Samuel两人。倒是在我别处称作科学界“沼泽”的地方,我们的行动找到了某种共鸣。公报最后印到约一万五千份——顺带说,后勤是一桩疯狂的活儿,发行全靠手工。Didier~Savard那些多汁的插画,想必为我们这份小报的相对成功立了不少功。在我与Samuel先后离开之后,它终于滑向一个小小的左翼团体,行话尖锐、论断不容置辩,公报也终于寿终正寝。该理解的、该说的,在1968年那阵亢奋余温尚近的时候,已被理解、被说了。此后再把一张唱片无休无止地转下去、倒回来,已没有什么意思……}最丰厚的“落果”之列——那段插曲\dmark{}\origpage{758},随我1970年至1972年底的离开而来。1972或73年,某份(或某些?)日本报刊谈到过这个团体、同名的公报,也谈到我的“离开数学”与我的“轨迹”。那些“科学批判”的一面、对军事机器的揭发,或许还有“文明批判”的一面,想必多少“透”进了某篇文章,引起日本山妙法寺一位僧侣的注意。他把它讲给了别人,尤其是同城(鹿儿岛)一位更年轻的僧侣——后者在他的影响下出家,颇有几分“弟子”的气象。1974年春,这便是该团体头一位登上“西方”海岸的传道僧——确切地说,到了巴黎\fnote{299}{(*)}{他确确实实向我保证过:在佛教历史上,他是第一位远赴西方的佛教传道僧——这份材料是否可靠,我可不敢担保!再说,出家当传道士,对佛教是否真是一大“进步”,也未可定。打一开始,日本山妙法寺团体的这个面向便在我心中引起保留,而岁月只是让这份保留一再获得证实。}。几个星期后,他来找我,而且\dmark{}\origpage{759}不经通报,来到我当时住的那个偏僻村子——离蒙彼利埃五十来公里。自那个难忘的五月天起——正午的日头底下,我看见一个装束古怪的人,一路敲着鼓唱行,正朝着(没错儿……)我独自劳作的那个园子走来——自那一天起,我有了这份荣幸与喜悦:看到古鲁吉的众多信徒与支持者\fnote{300}{(*)}{恰恰是其中一位,有此“荣幸”:作为“身份不合规的外国人”,成了法国判例里某部“1949年法令”中某条相当不可思议的条款头一次被照字面适用的由头。而我也有幸进了轻罪法庭,罪名是“无偿收留并藏匿”了这么一个法外之人。关于这段插曲,见“我的告别——或异乡人们”那节(n\textsuperscript{o}~24)。},络绎不绝地来到我家中做客。与他们的接触,给了我很多。1976年11月初,我甚至有了这份殊荣与欢喜:在我朴拙的居所接待Fujii~Guruji本人——时年92岁——同行的是七或八位僧侣、尼众与弟子组成的一行人。前一年,我已见过他:在团体位于巴黎第十八区的殿堂的隆重落成典礼上。在例行的客套话之外,那时有过强烈的接触,有过即刻的相知。而这次数日逗留寒舍——一个更贴身、更个人的情境——自然给了我一份丰富得多的领会:既关于Fujii~Guruji其人,也关于他与团体的关系——他正是那团体的头颅与灵魂。

有意思的是,Fujii~Guruji的这次来访,与那场我生命中的关键转折——它完成于同年10月15日至18日之间,别处已经谈到\fnote{301}{(**)}{见“欲与禅观”那节(n\textsuperscript{o}~36)与“重逢(阴的苏醒(1))”那条注(n\textsuperscript{o}~109)。}——相隔不过两周。紧随那些危机与更新之日的几个星期,是我一生最激荡的几个星期:每一天都带来它意外的收成——内在的事件与发现。说实话,这次来访——几个星期前就谋划、筹备停当,由整整一个僧侣与尼众的团体簇拥着他们尊崇的师父——来到那里,像一段奇异的插剧,像当时吞没我全部存在的那场冒险中的一次岔神。正是对客人们的敬重——尤其对亲临寒舍的Fujii~Guruji的敬重——让我毕竟得以在这几天里,保有这场合所要求的从容。一如我常有的情形,唯有置身事件的中心,我才明白:它绝非什么“插剧”或“岔神”,而\dmark{}\origpage{760}正是我正在经历的那场冒险的一部分。顶着它那一副十足的“东方故事”的门面——极尽的体贴,不寻常的魅力——这个所谓的“插剧”,让我面对的是与我相似的男人和女人们,面对那些我向来认识的男人和女人们——只是处在没那么异国情调、外表没那么不寻常的情境里。正是感到了这份亲缘,我也才在我的客人们身上感到了朋友与兄弟,而不是直直走出《一千零一夜》的人物——对不止一个目瞪口呆的村民,想必正是后一种情形。而Fujii~Guruji本人——他的“亲近者们”还守着尊崇师父所要求的那份得体的距离,他却同我说话如此家常——我感到他非常、非常遥远(对我,一如对他的亲近者),却又同时切近,仿佛他是我的父亲,或一位宽厚仁慈的兄长。

而一如在一位父亲或兄长身上——哪怕最宽厚的——并不少见的那样,他对我怀有一份期待,对此他也并不掩饰;一份陪伴他的人们共有的期待——他们个个都是我的客人。而我也知道,我无法回应它。我的冒险与Fujii~Guruji的冒险相连——凭一些我难以辨明的纽带,或许比我能看到的更深;也与他的弟子们的冒险相连——他们闭着眼睛追随他。但这场冒险,既不属于我这位显赫而宽厚的客人,一如它不属于我的父亲——他于我同样显赫而宽厚,极其切近却又不同:另一个人,另一个命运。

要让他们“明白过来”:在一桩属于他们的事业里——我并不觉得它属于我——我不会是他们中的一员,这并非易事。照人们想必向Fujii~Guruji和他的信众们描画的那幅我的图景,这该是他们最想不到的事——个人层面的关系越是水乳交融,就越想不到:无论团体整体还是它的各个成员与我之间,都近乎一场真正的蜜月。也正是在这次来访期间,某些源于我所受教育的、积年的抗拒消散了:我加入了客人们,与他们一道唱他们的曼怛罗,以鼓伴奏:“Na mu myo ho ren ge kyo !”

这曼怛罗是他们宗教修行的根基,是他们的阿尔法与欧米伽。他们唱它,多伴着祈祷鼓,晨一小时,晚一小时。依日本先知日莲的教诲,这击鼓而歌,本身即是至善,给唱它的人、也给他周遭带来和平。所以,对我\dmark{}\origpage{761}的日本朋友们,这歌唱便是人们通常所说的“祈祷”。他们赋予它的意义——与日莲相合,也与他们直接的“preceptor”Fujii~Guruji相合——是一份\textbf{尊崇的行动}:对着向之咏唱的那个人,并透过这个人,对着宇宙间一切活着的存在——就其(依《法华经》)被许诺终将成佛、成为圆满智慧的化身而言。这七个音节,也用作对任何他人的问候,乃至对任何想要问候的存在,带着那层对他人身上本质神圣之物的尊崇之意。它们也用作饭前的谢恩。说实话,我觉得几乎没有哪一种场合——无论是惊诧、感怀还是凝神的时刻——对一位日莲的追随者是不宜说出那句神圣之言的。至于我,虽不分享僧友们的宗教信仰\fnote{302}{(*)}{我不觉得自己属于任何一种特定的宗教。因父母所受的教育,我是个无神论者(带一点反宗教的色彩),直到十四岁。我的博物学老师作过一场精彩的讲述,谈地球上生命演进的历史——它让我懂得,毫无分毫怀疑地:宇宙之中,有一种创造性的智慧在运作。这份理解,当时还停留在单纯智性的层面;它在我此后的成熟中扩展、精纯起来——这成熟,在我1970年离开数学舞台之后仍在继续。},却满心欢喜地加入他们——每当时机来临——行~Odaimoku,击鼓而歌他们所谓的“祈祷”。正是怀着对他们的纪念,怀着对他们师父Nichidatsu~Fujii~Guruji的一份深情尊崇的行动,我也把“祈祷”纳入了我的日常生活:每天两顿主餐之前各唱一遍——至少当我自家居家、或身在朋友家中、或与那些我知道不会因此不安的人同席时\fnote{303}{(**)}{我尤其避免在每周去学院吃的那顿饭上唱祈祷——与几个学生或同事同桌——我怕他们中会有人从中感到一种强制,觉得我借年长者或“东家”的地位把它强加于人。}。这是我从Fujii~Guruji和他那些我认识的弟子们那里受惠的贵重之物之一——他们给了我他们的情谊,对我或近或远参与他们传道事业的推拒,从不厌倦。

在日本,有数百万日莲宗的佛教徒,分属众多面貌迥异的宗派。日本山妙法寺团体按人数是最小的之一:几百位僧侣、尼众与活跃的同情者。然而它在日本内外都颇知名,在所有传统宗教团体之中,它以一份毫不含糊的政治担当自我区分\dmark{}\origpage{762}——其主旨是为和平而斗争,是反军国主义的行动,尤其反核。越战年代,它是(如果我没记错)唯一明确站在美国人对立面的佛教团体,并为反对美军基地留驻日本而斗争(那些基地,正是越战延续下去的后勤支柱)。近年,Fujii~Guruji还与美国印第安人解放运动的领袖们——“美国印第安人运动”(AIM)——有过密切接触。日本山妙法寺的僧侣们参加过美洲印第安人组织的进军,还不算世界各地的其他和平进军。印第安领袖们显然被Fujii~Guruji不寻常的人格吸引、折服。这位精力不屈不挠、年近百岁的老人,竟是另一种宗教信仰的大传道者——这个事实似乎丝毫不使他们为难。恰恰相反:尊者师父那一竿子到底的“反美”倾向里的宗教维度,无疑连同他的高龄,是他们把古鲁吉当作自家人迎接的原因之一——当作一位极受敬重的父亲或祖父,一个人们在他身上认出自己的人\fnote{304}{(*)}{为给出一个印象,说明联结印第安领袖们与古鲁吉其人的那份信任与尊崇,我在这里指出:在围绕“太阳舞”举行的那场盛大的年度入会节庆上,也有古鲁吉的弟子僧侣们参加,从日出到日落,擂着大祈祷鼓,合着那扣人心弦的节奏:\emph{Na mu myo ho ren ge kyo}!这些大鼓,由整段的树干凿成,绷着牛皮,声威不凡,而(我猜想)一连十二个钟头怕也难以忍受。(在巴黎殿堂的落成典礼上,我领教过两个钟头——经验是有结论的……)不管怎样,Robert~Jaulin(他与僧侣们一道,是受邀参加节庆的少数非印第安人之一)向我转述:印第安人以斯多葛式的坚忍,从头到尾承受了祖父古鲁吉的圣鼓——那祖父的鼓声,会是入会礼的诸多试炼之一……}。

无疑,这份宗教维度之于我也起了同样的作用——它使Fujii~Guruji于我更切近,尽管我并不自认任何一种界说分明的宗教信仰。若问在他身上最吸引我、最触动我的是什么,我看到好几种东西。最显眼的,是一份内在的\textbf{喜悦}。这份喜悦似乎自发地源于他这个人身上的一份\textbf{统一},或者或许不如说,源于一份\textbf{对他自己的忠实}。人们感到这个人是幸福的,因为他一生都毫不犹豫地做了他感到自己该做的事。他于我并非没有矛盾,却没有暧昧。他的某些行动或不行动的意义,我无从把握\dmark{}\origpage{763},但对这个人全然的正直,却不曾有半点疑问掠过我心。倘若确是如此,这并非凭对那些经他人转述所知的他的分析。只消见过他一次,便知道这是一个不知暧昧为何物的人,一个与自己深切相合的人。美国印第安人运动的领袖们必是感到了这一点,才在中间给了他他们所给的位置。他对那些自认是他一伙的人的超凡感召,无疑也正系于此:那些男男女女,意识形态与哲学倾向从彻底纯粹的马克思列宁主义,一直铺到一家连锁大百货公司总裁那种体面的随波逐流。把他们联起来的,不是对某部《法华经》的崇奉——他们中或许无人曾自大到去读它\fnote{305}{(*)}{古鲁吉不止一位弟子曾让我明白:摆出读《法华经》的样子,会被视作自大——纵然它已有日译本存在。唯有一个精神极为深邃的人——如他的师父Fujii~Guruji本人——才有能力、才配得上读这部神圣的文本;它远远超出俗人的智力。显然,这些男男女女的信仰,直接系于Fujii~Guruji其人,而非某个或多或少被神化的历史人物,如佛陀,或圆满的菩萨、先知日莲。};也不是某篇源出巴利语的祈祷文——经汉译辗转还原为日文,其教义正是崇奉这部《法华经》。把他们联起来的(或者该说:曾把他们联起来的?),是一个\textbf{人}——他对他们施加着一种感召;而他并不曾刻意寻求这份感召,一如太阳不曾寻求它的行星。

我还看到:这个人是\textbf{孤独的},而孤独并不压在他身上。那是他的自然处境,或许从来如此。这份孤独,与这份正直、这份与自己深切相合,在我看来,是同一桩东西的诸多不同面向。这同一桩东西还有一个面向:\textbf{力}——一种没有暴力的力,它也不在意自己是不是、或看起来“强”。这仍是太阳的力:太阳只消是它自己,它周遭便自会生出那个力的场,生出那些行星循行的轨道。

无疑,这也是我在《收获与播种》里不止一次谈到的那份“我们之中的\textbf{力}”——区别只在:在这个人身上,它全然显露,凡接近他的人都感觉得到;在另一个人身上,它埋得或深或浅,有时深到让人以为它并不存在。而如果我某些僧友看起来在自身之中否认它,他们所信奉崇敬的那部《法华经》,连同他们日复一日唱着的祈祷本身\dmark{}\origpage{764},却分明宣告:一种这样的力,活在造化中一切活着的东西里——如他们一样,也如他们尊崇的师父~Osshosama 本人一样,被许诺了佛陀的命运。

\subsubsection*{(e) 祈祷与冲突}
\tocsubsub{(e) 祈祷与冲突}
\label{ch18:local:h75}

\notehead{161}{}
\phantomsection\label{ch18:local:n70}

(1月13日)\fnote{306}{(*)}{(1月23日)本注的前半部分,是顶着强烈的抗拒写下的——抗拒在于要提及干扰我工作的种种风波。这些风波显出几分可笑的模样,单是提它们,就有点像自请领教鞭、任人责打!另一方面,这些风波——“真能把人锯断的”——在我的工作里已变得如此刺耳、如此侵扰,尤其这一两周,若装作无事一般讳而不提,反倒会是一种作假、一种见证里的不真实。此外,十天之后,我在“Jung——或‘恶’与‘善’的循环”那条注里,重新谈了我的这些倒霉事。
(3月7日)后面这条注——关于C.G.~Jung自传的一整套“读书注”的头一条——最终被划归《收获与播种》的末一部分,即由这部自传所引发的那部分反思。}又是整整四天,我不曾有工作的闲暇与平静——我是说,继续写注的闲暇。缘由主要在于:要把《收获与播种》这第三部分誊打成清稿,我遇到的困难简直不可思议。交人打字这类活,我做了三十多年,从未见过这样的事。显然,手头这份性质极为切身、甚至可谓私密的文章,在负责打字的人身上触发了(想必是无意识的)力量可观的反应,每一次都奔着对托付给他们的工作的真真切切的破坏而去。几个月间,同一个剧本一连上演了三次,只差些变奏,三位接续的秘书,她们彼此却不曾串通\fnote{307}{(**)}{那些对我怀有善意的人,在这里有得可说:该给我扣上一顶“受迫害妄想”的帽子了——搬家师傅公会之后,又轮到打字秘书们的公会总动员,一心跟我过不去!关于前科,见“屠杀”那条注(单看注的名字,关于我这个人,它已经说得够多了……)p.~538,谈的是我朋友Ionel~Bucur的搬家……}!而这第三回,又添上一段卑污的插曲:秘书J太太摆出要用那份交在她手里的、颇为异常的稿子作筹码,敲诈一笔赎金。这是一位资深的老秘书,行业经验十分老到\dmark{}\origpage{765}。头打的十一页无可挑剔,几乎没有一个错字——为的是显一显她的本事;而单单随后的十五页里,就跳行了十一行——把文本糟蹋到这般地步,我也少见!为取回我的稿子与打件,赎金的数目(除已打好文字议定的报酬之外)我没有去问——我毫无意愿鼓励这类行径。这意味着,我多半将不得不诉诸司法途径。

所幸我手头还留着稿子的一份草稿,需要时用得上。但这路数的马戏——尤其当它带上了重复性——真能把人“锯断”。当我设想我这块沉思而自传的拙砖会掀起怎样的困难与对抗时,我当然不曾想到:头一拨麻烦竟会从这一头袭来——从打字秘书们的公会(而非我可敬的数学家同行们的公会),而且是消耗战的路数!眼下,把这份文稿(一旦取回)再交到第四位秘书手里,我已不太热心:我没有任何依据,能预见她会比她将接替的那些人多一分怜惜。而自己来做秘书的活计,得投入整整一个月的时间——我绝对无意提供。

或许我将不得不放弃把这第三部分誊成清稿,直接以草稿的形式交予出版者。(排版付印的排字工们,我倒不预期会有同一路数的麻烦!)这尤其意味着:我将放弃把第三部分收入《收获与播种》的限印预版——它应由我的大学USTL经办,在同事与朋友间以私人名义散发。或者,倘若我终于找到一位把活做正派的秘书,也可以以后再印。这一部分(三部之中肯定最“难”的一部),我只寄给那些确实想收到它、并在已收到前两部的人之中明确提出请求的人。我确实巴不得把前两部印出来寄出去(第三部分,我倒不那么急)。这两部分的打字早已完工数月,是由USTL的秘书们经手(而且毫无问题)的。它们本可以早早就印出来,若不是我想在里头收入涵盖《收获与播种》全三部的总目录——而三个多月以来,我相信\dmark{}\origpage{766}自己已处在给这没完没了的第三部分收尾的关口。这里,我打算给自己宽限到本月底完成;不然,就先操办前两部的印行(《愚妄与更新》,与《葬礼I,或中国皇帝的新衣》),而不把第三部分(《葬礼II,或阴阳之钥》)完整而确定的目录收入其中。

而在这些不快的插曲之后,如今我得千方百计找回一份被打断的反思的线索。

1月9日,Fujii~Guruji在他第一百零一个年头上辞世,这给了我一个机会:借他这个人,唤起我生活的一个此前未曾触及的面向。既然无缘再见古鲁吉于灵床,无缘在他的亲近者们中间守一场灵,我便在他死后的那一夜,独自守了一夜的灵,直到清晨,记下这桩事件引发的一部分追忆与思绪。事后我想:借着这个时机,我也该试着说说:与Fujii~Guruji、与他的那些我曾过从甚密的弟子们的相遇,给了我什么。

五天前的那几条注里,我已谈到 Na mu myo ho ren ge kyo 之唱——它多年来已进入我的生活,是一桩善益。还有那份情谊:由Fujii~Guruji本人、也由他好几位弟子——年轻的与不那么年轻的——所给予的。正是这份情谊,无疑,给了我自他们那里领受的歌唱以价值与美——那歌唱本身,便是对造化中一切活物的一份尊崇与情谊的行动,连他们的人、连我的人,也在其内。

同样,我与日本山妙法寺僧侣与尼众的接触,是我头一遭、也是仅有的一遭,与这样一些男男女女的密切接触:他们主要的、乃至全部的倾注,都投向宗教动机的事务(一如长年间我自己的倾注,投向数学发现的劳作)。这让我有机会确认:一如别处,在超出某种因共同志业(所谓宗教的)而生的亲缘、超出对同一位强势而动人的人格的效忠之外,气质的差异、制约的差异、乃至深层\textbf{抉择}的差异,依旧同样醒目,在人与人的关系里依旧同样起作用。换个说法:某些人力图依某种宗教理想把自己\textbf{塑造}成什么(这里是“菩萨”的理想——佛陀教诲的不倦传扬者),这份努力所抵达的\dmark{}\origpage{767},是或多或少敏感应激的\textbf{姿态},而非一个\textbf{内在转变}的进程、一份成熟。另一方面,采纳一个“信经”(无论多么崇高)并全力投入一项所谓“宗教”的活动,对于惯常那些唯我机制的游戏,似乎并无实质影响。在修道院、女修院、寺院与各种教门的宗教团体里,冲突一如世界的任何角落一样在场。而宗教志业,又常常被当作诸多手段中的一件,用来把冲突打发掉——说服自己:凭信经之力,冲突已然消失。

同样真的还有:在不同的场合,我某位僧侣客人身上,有一种内在的平和与喜悦在放光——我感觉得到,凡接近他们的人都感觉得到;对他们自己、对一切人,都有益。显然,这样一种和谐、圆满、深切相合的状态,与一切“成为这个那个”的努力格格不入——那是一种“无须用力”的状态,一种全然自然的状态。

在我感到过这样一份光辉的四位僧侣身上,我有这份印象:那是他们惯常的状态,已历多年,乃至数十年。Fujii~Guruji本人,尤其如此。对我的另外两位朋友,我在别的场合也见过他们,与任何人一样纠结、一样撕裂。仿佛我认识他们时的那份和谐状态,连同作为其征象之一的那份对事物的自发领会,已化为一纸空文——仿佛它们不曾在他们身上留下任何痕迹。然而我深信:确有一份不可摧毁的“痕迹”,比记忆里记下的一道标记更深——一份其本性为\textbf{认识}的痕迹。一如任何人,这些朋友随时可以顾及那份在他们的存在的创造时刻存入他们之内的认识,让它起作用、让它结果;一如他们也可以无视它、埋葬它,总而言之“装傻”。说到底,这是世间最平常不过的事……

又有一个念头来到心头:这样一种全然自然的状态、这份与自己深切相合,连同伴随它的光辉,却\textbf{并非}那么寻常的东西。相当触目的是:在我得以接待到家里的那一小群僧侣中间——或数日,或数周——竟有这么多的人,我在他们身上找到了这种内在和谐的状态、这种十足意义上的\textbf{力}——谦卑与坚忍、温润与犀利在其中合一的那种力。这莫非到头来确确实实是一个信经、或表达信经的那份“祈祷”的作用?这祈祷——\dmark{}\origpage{768}显见得单凭它自己造不出一个恩典状态——却或许倾向于\textbf{促成}这样一个状态的显现,以及它日复一日的更新?毕竟,全副身心投入地唱一支美的歌,这本身就多少是一份“恩典状态”——而一支歌(或一篇祈祷)单是它的美,就已促使我们“全副身心投入”。

同样真的还有:最美的歌,当我们心不在焉地反复唱它时,依然不起作用——错在我们不曾向它打开自己。或者更确切说:我们这样反复唱的,\textbf{并不是}我们以为在唱的那支歌,我们的灵魂并不以它为滋养——纸做的玫瑰、塑料做的玫瑰不是玫瑰,蜜蜂也不会来采蜜。

\subsubsection*{(f) 信念与认识}
\tocsubsub{(f) 信念与认识}
\label{ch18:local:h76}

\notehead{162}{}
\phantomsection\label{ch18:local:n71}

(1月14日)一周前结束那份反思时,我有这份感觉:我把手指点在了一桩重要的东西上。当夜,我想以一句简峻如刻石的话,把这件“东西”表达在给一条注取的名字里:“无缘之暴力的缘由”(注~n\textsuperscript{o}~159)。我也知道:这道骤然领会的闪电,绝非一份反思的收梢,甚至不是它的终点——那份反思,一个多月以来\fnote{308}{(*)}{确切说,自12月7日那条注“天鹅绒的脚掌——或微笑”(n\textsuperscript{o}~137)起。}正围绕“无缘的暴力”或“无端的暴力”之谜打转。恰恰相反:这份骤然现身的新“视角”,宁可说更像一个新的起点。“移位”的机制——为久远岁月里受下的委屈与损害而生的宿怨或怨怼,从那些被感受为不可企及或“禁忌”的真正的责任者,移往一副\textbf{可以接受的}靶子——这个机制,我起初只是在生活中这样那样孤立的个案里零零星星地认出它,心里默默认它是无意识里一种古怪而飘忽的畸变;如今它终被认作“人心基本机制”之一。同时,它又显为“无缘的暴力”无数令人不安的显形的责任者:无论夫妻之间、情侣之间、父母与儿女之间肆虐的那一种,还是那“匿名”的暴力——它在战争或社会大动荡的年代抵达它的顶点。

我不知这些关联是否早已进入心理学或精神病学的入门常识\dmark{}\origpage{769}(且假定有这样一门“科学”),也不知我在这里就它们所说的,会不会显得像“精神分析的业余爱好者”的凭空幻构。既然我的主旨不是提交一篇心理学的博士论文,甚至不是为某个新旧理论挺枪出阵,而是借我这个人在场所处的种种情境去理解我的生活,那么,我偶尔把手指点在其上的这桩东西、或我看见骤然在此处彼处打开的那些“视角”,它们的“地位”如何,于我无关紧要。我很清楚:无论如何,只要我想理解哪怕最小的一桩东西,我就省不得一份个人的反思——无论在数学里,还是在生活里,在与我的生活这样那样相连的那些生活里。而越是当所要理解的东西乍看便像要挑战理性,当我看到此界内外人人像躲瘟疫一样,拿一堆让人安心的套话把它躲开——情形就更是如此。(而且我觉得:心理学业内人士,一旦自己的这个人直接成为问题,也并不比所有其他人更能例外。)

我很清楚:那道在“给最后一个~i~点上最后一个点”的转弯处现身的“骤然信念”——即“我刚把手指点在了所有无端暴力情境共有的那个机关上”——丝毫免不了我的这桩任务:把这份刚进入意识目光之域的新直觉、这份还远未脱开方才从雾里浮出的东西那片弥散光晕的新直觉,逐件地、针脚细密地加以检视。恰恰相反:这正是头一件该做的工夫;在其中,我已看到一大堆新问题冒头——有专属于这样那样个案的,也有一般的。如果说这份生吞活剥的“信念”里有任何确定性,或者更确切说,有一枚可靠认识的核,那么它告诉我的,绝不是:我刚给这份信念取的表述是“真的”、“对的”,或许没有重要的保留或修改;而倒是:我确确实实刚把手指点在了一桩\textbf{新的}(对我而言)、\textbf{本质性}的事实上,一份看待暴力的新\textbf{视角}确确实实刚刚确立\fnote{309}{(*)}{写下这几行时,一个比较强加于我:与我在1968年孟买会议上提出的关于代数闭链的“标准猜想”相比。它们当时(今天依然)在我眼中,与奇点消解一道,构成代数几何里最灼手的问题之一。提出这些猜想时,我分明感到:一份“新的视角……刚刚确立”——这一次,是在代数闭链之上,在它们与Hodge理论、与Weil猜想的关系之上。最触动我的是:我看到一道通往Weil猜想的进法正在显现——它将是“纯粹几何的”,我的意思是,无须(至少表面上)经由一道上同调理论的迂回。

一如我在别处强调过的(在“肌肉与肚肠”那条注的子注~n\textsuperscript{o}~106$^1$里),这份“新视角”的实在与它的分量,全然独立于这个问题(它仍悬于未来的雾中):这道猜想将被证实为真,还是为假。一道猜想,于我不是一场(或赢或输的)\textbf{赌注},而是\textbf{一记探针}——而无论答案如何,我们从中出来的,只能是“赢家”,我的意思是:带着一份更新的认识。(与“错误与发现”那节(n\textsuperscript{o}~2)里的反思比较。)假定这道猜想被证为假,凭直觉我已看到它的两三个“不那么乐观”的变式——它们随即把它精细化,其中最弱的一个,实质上等值于:一个域上半单Motif的“合理”理论的存在。

清理出这些变式,对一个多少在内行里的人,是一两个下午的练习(或许也是一场通往未知的漫长旅行的起点……)。清理出最初的表述(一如往常,受Serre一个想法的启发,发表于他的文章《Weil猜想的Kähler类似》)——那不是一道练习题,而确确实实是\textbf{一份发现};或者(借用Zoghman~Mebkhout信中的说法——它引于“一次教学的失败——或创造与自负”那条注,n\textsuperscript{o}~44$'$)\textbf{一份创造}。而当Zoghman怯生生地贸然说“我的学生们不太知道创造是什么”——那还是客气的说法;或者不如说,我会讲:他们曾知道,却早已忘了个干净——他们全副身心,都在给一辆灵车推轮子……}。至于这桩新事实、这份新视角的精确而细致的含义,它确切的分量,或许还有它未曾预料的延伸与回响——只消我向它投入必要的工夫,它们自会浮现。这份刚现身的“认识”特别告诉我的是:做这样一件工夫的时机成熟了——更深入一步,去领会暴力,无论如何去领会“无端的暴力”的时机成熟了;我献给这桩任务的每一小时、每一天,只要走到方才现身之物的尽头,都会使我更深地进入这份领会。我不记得:这样一种感觉——一桩新的、本质性的东西正在现身(纵然它还弥散而近似)的感觉,连同能够更深入领会这桩东西的那份贴身信念——曾经骗过我。如果说在我的研究里有一个可靠的向导,替我把投资“安放”到这个或那个方向,那便是这份\textbf{新}东西现身的感觉,与这份贴身的信念——它告诉我,什么时候时机成熟,可以更深入那个被瞥见的“新东西”,去认识它\fnote{310}{(*)}{与“孩子与海——或信与疑”那条注(n\textsuperscript{o}~103)比较。}。

\dmark{}\origpage{771}这并不意味着:每当时机成熟,该向这个方向纵身、该去认识这样一些东西了,我就确确实实纵身而去!早在把全部精力投入数学的岁月,这已不可能:渐渐地,我发觉自己先是十块铁、继而一百块铁同时在火里!\fnote{311}{(*)}{见“炉中百铁,或:干熬无济于事!”那条注(n\textsuperscript{o}~32)。}在禅观里——也就是说,在对我自己的发现里——亦复如是。在意识工夫的层面上,我们很遗憾,一次只能做一件事(而这已不算坏了,只要肯下工夫把它做好……)。对“一百块在火里的铁”之中的\textbf{一块}所下的工夫,的确可以循无意识的神秘途径,惠及所有其余,或至少其中若干——它可以让它们“热起来”,使我们一旦转向它们,它们便更乐于承接意识关注的铁砧上落下的锤击。但还得一开始就知道,从一百块里挑出“对的”那一块——对它的锻打,将同样推进对其他铁块的工夫——那些正与它一同热起来的铁块。

\subsubsection*{(g) 最烫的铁——或转折}
\tocsubsub{(g) 最烫的铁——或转折}
\label{ch18:local:h77}

\notehead{162$'$}{}
\phantomsection\label{ch18:local:n72}

在关于葬礼的反思里,我遇到了不少“铁”——它们要求我在它们上面下工夫,冷热不一,视情形而定。我觉得,它们都在工夫里热了起来,有的多些,有的少些。这些“铁”里的头一块,是对己之轻蔑的问题——在我自己这个人的情形里:它最初是顺带着、为良知起见而提出的,在《收获与播种》最初那个胚胎的边上\fnote{312}{(**)}{见那条注(n\textsuperscript{o}~2)——它指涉1983年6月的一节“他人之无谬与对己之轻蔑”(n\textsuperscript{o}~4)。}。它一直温而不热,直到12月13日(一个月零一天前)的反思——在“正义者的暴力——或发泄”那条注(n\textsuperscript{o}~141)里。我想,这是我一生中头一回,把一场反思——无论多么粗略——献给生活中的那几桩个案:我自己施行并让人承受“无缘的暴力”、“超出理解的暴力”的那些个案。近年来我不时想到它们,但总是顺带一想,不作停留,尤其:不曾为之写下一场反思。

然而,那份不自道其名的暴力,曾深深烙进我的生活——那是一桩关键性的东西,乃至是万般之中那桩最关键的东西:为了理解我的生活,以及一般而言的“生活”、人类的生活,我必须尽可能深地\dmark{}\origpage{772}理解它。可事情确是如此——这只要肯想一想便显而易见——却一直被遮蔽着。它最终浮现出来,像出于偶然,在12月13日之前那些日子的反思的边上——反思延续于合起来题为“天鹅绒中的利爪”(n\textsuperscript{o}s~133--136)的四条注。正是在这些注里,“暴力”头一次在《收获与播种》中被道出其名,并成为一份关注的对象。它一直处于关注的中心,直到如今,或至少直到1月7日(一周前)那条注“无缘之暴力的缘由”。

这个引人瞩目的题目,会让人以为最后这条注是那场持续了整整一月的暴力反思的一种登顶。而它确实也是主要果实之一。然而我心知肚明:倘若骤然有这份新视角现身,有这份贴身信念之感——关于某道骤然瞥见的关联——那是因为\textbf{我这个人自己}也直接卷入了方才现身的东西之中,在那“一生积攒的、一大堆零碎而不相属的印象”之中。所有这些印象里最后、也最新鲜的那一份——当时感到它确实很“零碎”、很不充分——恰恰要回溯到12月13日那场关于\textbf{我自身之中的暴力}的反思。这场反思,在浮泛的读者眼里,或许只是关于葬礼的追查里诸多插笔中的一笔;而于我,此刻,凭回顾,它却显为关于我自己的反思里一个神经中枢般的时刻、一个关键的转折(至少在潜能之中)。而且就在当天,我便感到自己终于起步,迈出了头一步,走进一个此前一直被我避开的方向——它将笔直地领我走向我这个人之内冲突的心脏。那块“温吞的铁”——它作为存照被搁在那里,已历十个月——骤然间烧得通红:只消我停步,鼓风吹打,它便烧成白热,向我显出一份形状、一道讯息。而今天,依然如此。

但显然,此处不是锻打这块铁的地方。在《收获与播种》进程中现身的所有铁块里,它于我确是最烫的一块;紧随其后,是那道与它密不可分、随“无缘之暴力的缘由”一同现身的铁。倘若孩子背上没有驮着一位太成年的东家——他执拗地铆在长线事务与它们强加的“优先事项”上——那么此刻\dmark{}\origpage{773},我必定早已纵身扑向这个方向——它领我走向我自身之内、他人之内冲突的心脏——无须自己掂量!可一如它的名字所示:下单的多半是东家,而不是孩子;决定投资的,也是东家。“恶之谜”因此要等一个更合宜的时刻——东家度假的时候(那是千载难逢的),或者他不再被尖端的“优先事项”压得过重的时候——比如终于写完《收获与播种》!

\subsubsection*{(h) 无穷的链——或传递(3)}
\tocsubsub{(h) 无穷的链——或传递(3)}
\label{ch18:local:h78}

\notehead{162$''$}{}
\phantomsection\label{ch18:local:n73}

但在回到葬礼之前,我想至少记下一周前那场反思所引发的联想之一——一条或许不如其他的显豁、因而倘若此刻不记下便有消散无痕之虞的联想。它与印度教的业之观念相连,与“敌对的兄弟——或传递”那条注(n\textsuperscript{o}~156)里出现的联想同向:朝着一份微弱的直觉——某种“业之守恒律”。

一个人身上那份原初的弥散宿怨——它此后转为表面上“无缘”的攻击性与暴力冲动——并非生于虚无。它是对真真切切受下的深层侵害的回应,尤其是童年最早期受下的那些侵害。诚然,可以认为:这些侵害中有许多,属压抑性质,算不上严格意义上的“暴力行动”——即出于伤人或损人的意图,尤其在父母对孩子。同样真的是:这样一种意图(几乎总是无意识的),其出现的情形,却远比通行共识所承认的多。但在一份业的创造或传递的视角里,当一份“暴力”确确实实发生——施加“一桩恶”,造成“一桩损害”——那么\textbf{意图}或\textbf{动机}(显明的或隐秘的)的问题,或许是次要的。这话我不敢说定。不管怎样,在大多数情形里,一双浮泛的眼睛会有这样的幻觉:所受的那桩“恶”化为一纸空文,它被承受下来,而一经承受,便“消失”了,不留痕迹。而事实是:那些把自身的焦虑、把自身“无法做自己”的无力撒进孩子心里的父母,最终从这同一群孩子手上,直接收到\dmark{}\origpage{774}他们当年播下的东西——这并不那么常见;或者至少,人们得到的印象是:他们从中收获的只是微不足道的一小部分!换个说法:在他们于孩子身上唤起的那份弥散宿怨里,只有一小撮会凝成一份“硬”宿怨、指向他们自己——对此他们叫苦连天,俨然它是天底下最黑的忘恩负义——仿佛这不证自明!可这份积攒的宿怨或“业”的其余部分,并不因此湮灭。它有效地找到用武之地,以一副或许显得无从解释的方式:凭那“移位”的机制,把宿怨挪往临时的靶子;有时是飘忽无定的靶子,有时也是特意搭配、专任其职、可以说备受娇宠的靶子——被孵育一生!

在寻常的时日里,业这桩强烈的工夫——如一个深植于人们生活中的脓疮——在幽暗中完成;而人人以无视它为己任,只肯把它看作此处彼处偶发的“越轨”,相对于那被视为正常与得体的东西。唯有在非常的时日——当战争或苦难肆虐(或在非常的处所,如监狱与疯人院)——这份地下的工夫才迸发出来,自由自在地铺展在大白天里:轻蔑与杀人癫狂的一场恣意的火光冲天,由大言煌煌的旗帜所激励——旗帜飘扬在“英雄的”藏骸所上空,飘扬在一座座赤裸而寒冷的城池上空……

